\RequirePackage{fix-cm}
\documentclass[a4paper,11pt]{ltjsarticle}
\usepackage{amsmath,amssymb,amsthm,mathtools,bm}
\usepackage[margin=25mm]{geometry}
\usepackage{booktabs,longtable,array}
\usepackage[unicode,colorlinks=true,linkcolor=blue,citecolor=blue,urlcolor=blue]{hyperref}
% Use the Japanese fonts bundled with TeX Live; no system Noto fonts required.
\usepackage[haranoaji]{luatexja-preset}
\numberwithin{equation}{section}
\theoremstyle{plain}
\newtheorem{theorem}{定理}[section]
\newtheorem{proposition}[theorem]{命題}
\newtheorem{lemma}[theorem]{補題}
\newtheorem{corollary}[theorem]{系}
\theoremstyle{definition}
\newtheorem{definition}[theorem]{定義}
\newtheorem{assumption}[theorem]{仮定}
\newtheorem{example}[theorem]{例}
\newtheorem{conjecture}[theorem]{予想}
\theoremstyle{remark}
\newtheorem{remark}[theorem]{注意}
\renewcommand{\proofname}{証明}
\newcommand{\C}{\mathbb C}
\newcommand{\A}{\mathbb A}
\newcommand{\Gm}{\mathbb G_m}
\newcommand{\OO}{\mathcal O}
\newcommand{\PP}{\mathbb P}
\newcommand{\Gr}{\operatorname{Gr}}
\newcommand{\Spec}{\operatorname{Spec}}
\newcommand{\Spf}{\operatorname{Spf}}
\newcommand{\Hom}{\operatorname{Hom}}
\newcommand{\End}{\operatorname{End}}
\newcommand{\Ext}{\operatorname{Ext}}
\newcommand{\Aut}{\operatorname{Aut}}
\newcommand{\Isom}{\operatorname{Isom}}
\newcommand{\QGr}{\operatorname{QGr}}
\newcommand{\qgr}{\operatorname{qgr}}
\newcommand{\Qcoh}{\operatorname{Qcoh}}
\newcommand{\coh}{\operatorname{coh}}
\newcommand{\Def}{\operatorname{Def}}
\newcommand{\Sym}{\operatorname{Sym}}
\newcommand{\Perf}{\operatorname{Perf}}
\newcommand{\rank}{\operatorname{rank}}
\newcommand{\HH}{\operatorname{HH}}
\newcommand{\Version}{0.9}
\title{次数1・2の非可換del Pezzo多様体\texorpdfstring{\\}{ }構造対象付き変形とモジュライの局所構造}
\author{}
\date{2026年11月20日\quad v\Version\quad 公表用研究稿}
\hypersetup{pdftitle={次数1・2の非可換del Pezzo多様体：構造対象付き変形とモジュライの局所構造},pdfauthor={}}
\begin{document}
\maketitle
\begin{abstract}
滑らかな複素del Pezzo多様体 $X$，$n=\dim X\ge2$，
$-K_X=(n-1)H$，$H^n\in\{1,2\}$ の非可換変形を扱う。
構造対象 $\OO_X$ だけを標識し，全ての標識を保つ完全Abel同値を射とする。
既知の変形・量子化の比較と相対Poisson構造定理を組み合わせ，
特殊ファイバーを固定した全Abel Artin変形群oidを計算する。
次数1の形式環は正則であり，次数2では滑らかな形式基底と
$\Gr(2,n+1)$ のPl\"ucker錐の頂点完備化の積となる。
各可換点の形式分枝は一つで，枠を忘れた形式商は
$[\operatorname{Spf}R_X/\operatorname{Aut}(X)]$ である。
有限Ore被覆と対角のcompact持上げにより，全Artin変形の飽和性，豊富性，
smooth proper性と相対Serre公式 $(-(n-1))[n]$ を証明する。
さらに全構造対象付き同値の認識，反標準section環とVeronese復元を調べる。

有限パラメータについては，両次数・全次元で独立な全Poisson方向を含む
有限表示ambient族を明示的な主開集合上に構成する。
同じ代数の中心を中央多項式係数上で持ち上げ，方向空間上で貼り合わせることにより，
全可換係数の零切断全体を含む有限型開近傍上で
普遍的に正則な中心元と次数別平坦な中心超曲面族を得る。
次数2ではGrassmann多様体上のコホモロジー消滅と錐への降下を使う。
次数1では4次までの座標不変な中心補正を厳密計算で与え，
残る高次の貼合せ障害の消滅と有限線型方程式によって代数化する。
この一つの有限型中心商族のArtin基底変更が，
可換sliceを選んだ後の全framed Abel変形を実現することを証明する。

さらに，活性最高次曲線の特異・可約・非被約な退化を含む相対切断環を構成する。
次数別のKrull交叉と，方向頂点の完成および頂点外の忠実平坦被覆を組み合わせ，
この全方向族の強Noether性と積代数のNoether性を証明する。
同じambientの無捩れな相対CY同型から，相対tails・Serre公式と有限compact生成窓を得る。
対角の有限部分解消により，滑らかな全可換模型の零切断を含む有限型開近傍で，
同じ族が実際のsmooth properな非可換del Pezzo族となる。
全モジュライのatlasと全既約成分数，追加開条件の明示公式は未解決として区別する。
\end{abstract}
\tableofcontents

\section{問題と主結果}
\label{sec:introduction}
次数1・2のdel Pezzo多様体のPoisson構造は、一般の次元でも単純な相対錐により記述できる。
本稿の目標は、この記述を構造対象付き非可換射影多様体の変形とモジュライへ結び付けることである。
基本偏極や豊富列を固定した表示の分類と、構造対象 $\OO$ のみを保つ全Abel同値による分類を区別する。
比較の方法は非可換三次三様体の研究\cite{CubicModuli}を踏まえるが、
重み付き生成元を持つ本件では次数一生成や二次包絡をそのまま仮定しない。

\subsection{中心となる形式局所結果}
$b_{n,d}$ を通常のdel Pezzoモジュライの次元とし、
\begin{align}
 b_{n,1}&=\binom{n+3}{4}+\binom{n+5}{6}-n^2,
 &b_{n,2}&=\binom{n+4}{4}-(n+1)^2
 \label{eq:base-dimensions}
\end{align}
とおく。特殊ファイバーを固定し、全Abel変形の射を自然同型で割ったとき、
$\OO$-marked変形の同型類関手を表す環は非標準的に
\begin{equation}
 R_X\simeq
 \begin{cases}
 \C[[t_1,\ldots,t_{b_{n,1}},D_1,\ldots,D_n]],& d=1,\\[2pt]
 \C[[t_1,\ldots,t_{b_{n,2}},p_{ij}\ (0\le i<j\le n)]]/(q_{ijk\ell}),& d=2,
 \end{cases}
 \label{eq:main-formal-ring}
\end{equation}
である。ここで
$q_{ijk\ell}=p_{ij}p_{k\ell}-p_{ik}p_{j\ell}+p_{i\ell}p_{jk}$
($i<j<k<\ell$) である。証明は第\ref{sec:formal-abel}節に置く。
この環は特殊ファイバーの枠を固定した変形を表す。
第\ref{sec:formal-groupoid}節では微小自己同値の消滅を証明し、
枠を忘れた形式群oidを $[\Spf R_X/\Aut(X)]$ と同定する。
$\OO$ の標識を保つ自然自己同型は自明である。
定理\ref{thm:all-artin-are-nc-dp}により、同じ表示が
$\mathfrak M^{\OO}_{n,d}$ とその重み付き超曲面部分問題のArtin完成を与える。

次数1では全一次変形が延長する。次数2では
一次Poisson方向 $Q\in\wedge^2\C^{n+1}$ の二次障害は $Q\wedge Q$ であり、
通常の多様体の変形方向を加えても相殺できない。
これらは全ての通常の局所Artin係数環を対象とする主張であり、
固定した一つの形式量子化族だけの接空間計算ではない。

\subsection{先行研究と本稿の範囲}
Lowenのalgebroidによる全Abel変形の比較\cite[Theorems 3.14, 3.26]{Lowen2008}と、
Yekutieliのtwisted quantization\cite[Corollaries 7.5, 11.2]{Yekutieli2015}を用いる。
後者のPoisson変形では、可換積自体も変化する。
De Deken--Lowen\cite[Theorem 4.5, Proposition 4.6]{DDL2011}は
豊富列の持上げとsection $\mathbb Z$-代数による全Artin変形の表示を与える。
これらの比較原理そのものは既知の結果である。

Lowen--Van den BerghのAbel変形とinjective線型圏の比較
\cite[Theorem~8.17, Appendix~B]{LowenVdBDeformation}を用い、
$HH^1(X)=0$ から全Abel同値の微小剛性を得る。
compact性の持上げにはDe Deken--Lowen
\cite[Propositions~5.8--5.9]{DDL2011}を用いる。
一般的なdeformed schemeのcompact生成は
Lowen--Van den Bergh\cite[Proposition~5.1]{LowenVdBCompact2013}も扱うが、
本稿では重み付き中心商の有限Ore被覆を構成し、任意基底変更、飽和性、
proper性と基底変更を検証し、compact持上げの論法を対角dg加群へ適用してsmooth性を得る。
ambientの無捩れな双対性はShoikhetの鎖formality
\cite[Theorem~2.2.7]{Shoikhet2003}を冪零係数上で有限twistすることで得る。
unimodularityとVan den Bergh双対性の関係自体は
Dolgushev\cite[Theorem~3]{DolgushevModular2009}の既知の原理である。
相対Serre公式には別途、smooth proper圏の対角bimoduleの剛性を用いる。

本稿固有の入力は、\cite[定理7.3, 定理8.8]{PoissonDP}の
任意基底変換と両立する相対Jacobi錐と全族の商表示であり、
式\eqref{eq:main-formal-ring}はそれらを既知の変形比較へ適用した帰結である。
曲面でのPoisson変形の接空間次元10・9と高次障害空間の消滅は既知\cite[Theorem 3.9]{HongXu2011}であり、
この数値を新発見とは扱わない。
一般次元の本稿の帰結と復元定理について、優先権の包括的確認は未了である。

具体的な非可換del Pezzo曲面の構成として、Etingof--Ginzburgは三変数の
Calabi--Yau代数とその中心商を研究し、filteredな族の半普遍性を示した
\cite[Theorems 3.4.4--3.4.5]{EG2010}。
Rainsは楕円代数のfiltered変形と中心拡張を詳しく扱う\cite{Rains2022}。
これらは本稿の重み $(1,2,3)$、$(1,1,2)$ の活性な三変数を具体的に代数化するための先行例であるが、
任意の $n-1$ 個の中心Casimir変数に依存する全族の平坦性や、
全 $\OO$ 保存同値による分類が自動的に従うわけではない。

Abdelgadir--Okawa--Uedaの三次曲面のコンパクトモジュライ\cite{AOU2024}は、
例外集合に付随するquiverの関係式を標識したGIT問題である。
本稿では構造対象一つを残して豊富列の選択を忘れるため、同値関係が異なる。
Okawa--Uedaの球面helixからのAS正則代数の構成\cite{OkawaUeda2020}も、
将来の全 $\mathbb Z$-代数型の表示を考える際の基礎となる。
これらの先行研究を、一般次元の全ncモジュライの代数性・成分数の既知の解答とは扱わない。

以下では、全Abel Artin変形、許容表示のモジュライ問題、
Serre・豊富性条件を満たす部分問題を定義上区別し、可換点のArtin完成での一致を証明する。
全非可換モジュライの有限型atlasや成分数は、形式局所環の計算だけからは従わない。

\subsection{形式族と同値関係についての結果}
第\ref{sec:ncdp-anticanonical-reconstruction}節では、相対Serre双対、
反標準列の豊富性、section復元を明示して、全 $\OO$ 保存同値を
反標準軌道環の同型として認識する。
第\ref{sec:ncdp-veronese-recovery}節では、二つ以上の中心一次変数と普遍中心条件の下で、
任意の $r$ に対するVeronese同型の持上げが $\mu_r$-torsorとなることを示す。
次数1・2の非零rank2模型には共に $n-1$ 個の中心一次変数があるため、
$r=n-1$ の反標準環に適用するための次元上の障害はない。
ただしその中心条件は可換点を含むArtin族には成立しない。

第\ref{sec:equivariant-quantization}節では、標準形に対する一つの量子化を固定し、
monicな中心商の分裂によって任意の冪零係数と基底変更に対する平坦性を証明する。
正規形の座標変換に関する同変性と、次数スカラー作用の整合的な恒等化から、
Poissonモジュライの零切断の形式完備化上に構造対象付きAbel圏の族が得られる。
定理\ref{thm:concrete-quantization-all-artin}は、この具体的量子化自体が
各可換点の全Abel Artin変形の同型類に自然な全単射を与えることを示す。
初commutatorによる接写像の同型と、同じ完備局所環の自己写像に対する
逆写像補題を組み合わせるため、次数2の特異なPl\"ucker錐にも適用できる。
従って全Artin変形は通常の次数付き中心超曲面で表示でき、
第\ref{subsec:nilpotent-weighted-as}節の相対augmentation Ext条件も満たす。
第\ref{sec:artin-ample-perfect}節は有限Ore被覆から
section saturation、有向ACM、正周期の豊富性とsmooth proper性を証明する。
第\ref{sec:relative-serre-artin}節で相対Serre公式を証明し、
全Artin変形が定義\ref{def:nc-dp-stack}を満たすことを示す。
また命題\ref{prop:artin-full-graded-recognition}はArtin近傍で全 $\OO$ 保存同値を
通常の次数付き環同型を次数スカラーで割ったものとして認識する。
この局所結論は普遍中心条件を使わず、可換点でも成立する。
一方、大域モジュライの代数性と成分数には有限型族と一般nc点の解析が残る。

第\ref{sec:finite-auxiliary-quantization}節では，特殊なPoisson方向について有限型の平坦Abel族を構成する。
補助次数が真に下がるJacobian括弧では，グラフ量子化とCasimir持上げが各入力に対して有限和になる。
別に，分離型potentialの関係式に対する唯一の重なりを解消し，
次数1では $h^6$ まで，次数2ではWeyl対称化の $h^2$ 補正で書ける厳密な中心元を与える
（定理\ref{thm:explicit-finite-pbw-casimir}）。
これは形式級数の打切りではなく，任意の係数環・任意の有限 $h$ で成立する公式である。
potentialの対角複体とその自己双対性によりambientは相対CYであり，
定理\ref{thm:finite-parameter-universal-center}は $h\ne0$ における普遍中心を計算する。
従って二つ以上の中心一次変数を用いるVeronese持上げの代数的仮定が，有限非零パラメータで実現する。

この構成にはEtingof--Ginzburgの低次数potential変形
\cite[Proposition~7.3.1(ii), Corollary~7.4.1(ii)]{EG2010}と共通するfiltered量子化の機構がある。
対角複体はGinzburg\cite[Proposition~5.1.9, Theorem~5.3.1]{Ginzburg2007}の構成に基づき，
ここでは中央多項式係数と任意基底変更を直接扱う。
これらの一般原理や量子化の存在を新規性として主張しない。
全次元への適用，明示式と普遍中心の整理についても包括的な優先権調査は未完了である。
第\ref{sec:finite-tails}節では，分離型ambientの相対CY同型と
中央正則元のchange of ringsから，$r=n-1$ として
\[
 R\Gamma_{B_+}(B)\simeq B^\vee(r)[-n-1]
\]
を任意のnoetherian係数基底上で証明する。
その結果，飽和性，有向ACM，$\chi$，全シフトの豊富性，compact生成と相対proper性を得る。
局所化bimoduleの双対を計算することにより，perfect圏の相対Serre関手は
$(-r)[n]$ となる。これは体上のAS--Gorenstein双対性
\cite[Theorem~4.2, Corollary~4.3]{YZ1997}と整合するが，
非正則係数基底への主張は本文で直接証明する。
有限個の次数の余核を用い，正周期の豊富性を滑らかな可換点の開近傍まで拡張する。
perfect圏のsmoothnessとcoherent圏の有限ホモロジー次元は，
第\ref{sec:finite-smoothness}節で別の対角の議論により得る。

第\ref{sec:elliptic-degree-two-finite}節は，負の補助次数では扱えなかった
次数2の一般楕円型最高次項を扱う。
Belmans--Raedschelders\cite[Example~3.11]{BelmansRaedschelders2020}に記された
三次type-A代数の正則領域を使い，$h=0$ を含む有限型開集合でambientを構成する。
固定topでのlower potentialのPBW性と中心元の存在は
Etingof--Ginzburg\cite[Proposition~7.3.1(ii), Theorem~3.3.2(ii)]{EG2010}の既知の結果である。
本稿では普遍係数環での平坦性，低次係数に依存しない有限個の分母，
任意基底変更で正則な中心商を証明して，全次元へ適用する。
最高次中心元にはDe Laet\cite[\S5.2, Theorem~5.3]{DeLaet2017}の公式を用い，
自由代数内の恒等式と可換極限を別に検算する。
一般の量子化やSklyanin中心元の存在自体は新規性の主張ではない。
第\ref{sec:elliptic-relative}節では，二重対角線に退化するpoint curveを解消し，
最高次中心商を滑らかな相対種数1曲線のtwisted homogeneous coordinate ringと同定する。
Keeler\cite[Theorems~7.13--7.14]{Keeler2003}の相対Noether性を任意のnoetherian基底変更後に適用し，
中心拡大とfiltered持上げから，全lower係数・全次元の強Noether性を得る。
相対CYと有限augmentation解消を合わせ，飽和性，ACM，proper性，Serre公式と
有限compact生成元を一般楕円型族へ拡張する。
既知の $\sigma$-ampleness原理そのものは新規結果ではなく，
ここでの有限開近傍・可換極限・普遍係数との比較が本稿の構成である。
さらに曲線の積上の評価イデアルと相対Fujita消滅を使い，
$B^{\mathrm{op}}\otimes B$ のNoether性を独立に証明する。

第\ref{sec:elliptic-degree-one-finite}節では，重み $(1,2,3)$ の
Weierstrass型potentialを，有限忠実平坦被覆と重みを保つ変数変換によって
Stephensonの楕円代数へ移す。
同氏のTheorem~1.1\cite{Stephenson1997}による11個の例外点を避ける条件は，
本稿の係数で $(1+\alpha h^4)(5+\alpha h^4)\ne0$ と具体化できる。
正則代数の分類自体は先行結果である。
本稿では，この有限開条件から普遍係数上の相対CY性と全lower項のPBW性を得て，
活性次数9以下・23未知数の中心方程式を解く。
定理\ref{thm:degree-one-finite-central-family}は，分母が最高次係数だけに依存する
普遍的に正則な中心商を与え，系\ref{cor:degree-one-elliptic-general-dimension}により
任意の $n\ge2$ へ拡張する。
さらに，有限被覆上でpoint curveを解消し，無限遠の切断 $P$ による
次数1線束 $\mathcal O(P)$ の相対THCRと最高次中心商を同定する。
定理\ref{thm:degree-one-relative-thcr}と
系\ref{cor:degree-one-product-noetherian}により，
次数1でも任意noetherian基底変更後の強Noether性と積Noether性を得る。
定理\ref{thm:degree-one-relative-tails}は相対tailsと長さ $n+5$ のcompact生成窓を与える。

\subsection{有限型開集合上の非可換del Pezzo族}
第\ref{sec:finite-smoothness}節の定理\ref{thm:finite-family-smooth-open}は，
両次数の分離型族と一般楕円型族のそれぞれについて，
滑らかな可換模型を与える点 $s$ の有限型Zariski開近傍 $V$ を構成し，
\[
 \operatorname{Perf}(\QGr B_V)\quad\text{が相対smooth properである}
\]
ことを示す。任意の許されるnoetherian基底変更にも両立する。
証明では，二重次数付き対角の自由解消を $q>2n$ で切り，
相対Serre双対により対角からこの有限複体へのHomを係数上perfectにする。
可換ファイバーの $X\times X$ の有限ホモロジー次元と $HH^1(X)=0$ によって
対角への写像の右逆を有限開近傍へ持ち上げ，対角をperfect対象の直和因子として得る。
一般のproper dg代数に対する未確認の開性定理は使わない。
反標準豊富性の開集合と交わらせると，
系\ref{cor:finite-family-ncdp-open}により実際の
$\mathfrak M^{\OO,\mathrm{hyp}}_{n,d}$ の族となり，
幾何学的ファイバーのcoherent圏のホモロジー次元は $n$ となる。
開集合は有限な局所化による存在証明であり，具体的な判別式はまだ計算していない。
命題\ref{prop:one-parameter-atlas-obstruction}は，これらの共通パラメータ $h$ による中心商族の
Poisson Kodaira--Spencer像が高々1次元であり，可換点でsmooth atlasにはならないことを示す。
全ncモジュライのatlasには独立なPoisson方向の同時中心商族が必要であり，
これらの構成の包括的な優先権確認も残る。

\subsection{全方向の有限中心商とその全Artin普遍性}
第\ref{sec:simultaneous-poisson}節では，次数1の $n$ 個の線型方向，
次数2のPl\"ucker錐の全座標を直接関係式へ入れる。
次数1では，Weierstrass最高次代数のAS条件に楕円判別式は不要であることを示し，
方向 $D$ に対する $(1+A_4(D))(5+A_4(D))\ne0$ を明示する。
次数2ではMori--Ueyama\cite[Theorem~3.4]{MoriUeyama2019}の三次CY判定から，
四次式の対称四重線型化と交代形式 $P$ による作用素
$T_{U,P}(ab)=\widetilde U(P^\flat a,P^\flat b,-,-)$ を用いた
$g_2(U,P)=\det(1+4T_{U,P})\ne0$ という十分条件を得る。
$g_2(U,0)=1$ であり，最高次四次式に重根があっても適用できる。
これにより，両次数で可換係数の零切断全体を含む有限主開集合上の
普遍的に次数別平坦なambient族を得る
（定理\ref{thm:simultaneous-degree-one-ambient}，\ref{thm:simultaneous-poisson-ambient}）。
この同じambient族に，第\ref{sec:relative-center-lift}節は中央多項式係数上の
規格化形式中心元を構成する。Etingof--Ginzburgの形式中心持上げ
\cite[Proposition~8.1.1(i)]{EG2010}を踏まえ，相対formalityのgauge同型によって
中心元を元の代数へ戻す。中央差のKoszul分解から絶対CY性へ移り，
Dolgushev\cite[Theorem~3, Proposition~1]{DolgushevModular2009}の絶対定理を適用する。
不変体積の重み零成分と相対多項式Poincar\'eホモトピーを使うため，
中央変数の分母や活性最高次式の楕円判別式は必要ない。
有限中心方程式の余核を使えば，零切断全体を含む有限主開へ代数化できる。

方向間の貼合せは両次数で異なる。
次数2では $\Gr(2,n+1)$ 上の $\operatorname{Sym}^4K(m)$ の $H^1$ が
全 $m\ge0$ で消滅する。
Borel--Weil--Bott\cite[Theorem~2.3]{GuoLuSharpe2017}とproperな方向解消から，
全四次係数の零切断を含む錐上へ中心を降下する。
退化ファイバーで中心が増える場合も，普遍係数環上の核を先に扱うことで含められる。
次数1では $\mathbb P^{n-1}$ 上の $\operatorname{Sym}^6K(m)$ の消滅は
$m\ge5$ で用い，それ以前の障害を
式\eqref{eq:degree-one-fourth-center-jet}の普遍的な4次中心jetによって越える。
従って定理\ref{thm:simultaneous-degree-two-central-family}と
\ref{thm:simultaneous-degree-one-central-family}は，全 $n\ge2$ で
有限かつ任意基底変更後も正則な中心超曲面族を与える。
追加の主開条件は存在を証明したもので，その最小行列式は未計算である。

定理\ref{thm:simultaneous-central-all-artin}では，この同じ有限型族を
可換係数の横断sliceへ制限すると，全framed Abel Artin変形を表すことを示す。
De Deken--Lowen\cite[Proposition~4.4]{DDL2011}でArtin基底上のAbel平坦性を得て，
初交換子による全接空間の同型と，式\eqref{eq:main-formal-ring}の
全完備局所環の比較から証明する。
有限基底上の圏論的性質は，この形式比較とは独立に次のように証明する。

\subsection{全方向族のNoether性と相対幾何}
第\ref{sec:degree-two-gorenstein-curve}節では，
全二元四次式のpoint curveを解消し，$\mathbb P^1$ の有限平坦二重被覆を得る。
可約曲線と非被約な二重 $\mathbb P^1$ も含み，後者のnilradicalは
$\mathcal O_{\mathbb P^1}(-2)$ である。
第\ref{sec:degree-one-integral-curve}節では，三次根による尺度変更を取り除き，
階数2の有限被覆上で節点・尖点を含む整なGorenstein種数1曲線を得る。
これらの曲線のコホモロジーと成分ごとの次数を使い，
Keeler\cite[Theorems~1.3, 7.14; Erratum~Theorem~E2.5]{Keeler2003}の
既知のample filter・Noether性・相対Fujita消滅を適用する。
THCRのtensor積を多重斉次切断環として扱う研究には
Keeler\cite{KeelerMultiple2003}の先行研究がある。
同文献は書誌とabstractを確認した段階で，本文の定理番号と一般基底への仮定は未確認である。
従ってここではその定理を入力とせず，確認済みの相対Fujita消滅から積の議論を直接記述する。
最高次中心商のTHCR同定，二重切断環，中央拡大とPBW持上げにより，
全topを許す方向chartで強・積Noether性を得る。

方向解消は頂点で平坦でないため，そのまま忠実平坦降下には使えない。
第\ref{sec:noether-descent}節では，任意の可換Noether環 $R$ と
方向イデアル $J=(a_1,\ldots,a_s)$ に対する忠実平坦射
\[
 R\longrightarrow\widehat R_J\times\prod_iR[a_i^{-1}]
\]
を用いる。完成側では，次数別有限性とKrull交叉により，
可換な零方向商からNoether性が持ち上がる。
頂点外では方向chartの相対曲線を使う。
定理\ref{thm:simultaneous-strong-product-noether}は，
全可換係数の零切断を含む一つの有限主開上で，
任意Noether基底変更後の $A,B,A^{\mathrm{op}}\otimes A,B^{\mathrm{op}}\otimes B$ の
左右Noether性を与える。体ファイバーだけのNoether性に基づく結論ではない。

第\ref{sec:simultaneous-relative-geometry}節では，有限な次数零Hom複体を用いて
全ファイバーのprojective性から相対有限対角解消を構成する。
最低次数の双対化線束と左右作用を調べ，同じambientの無捩れな相対CY同型を得る。
従って相対tailsは飽和性・ACM・$\chi$・豊富性・Serre公式
$(-(n-1))[n]$ と長さ $n+5$（次数1），$n+3$（次数2）のcompact生成窓を持つ。
定理\ref{thm:simultaneous-smooth-ncdp-open}により，
滑らかな全可換模型の零切断を含む有限型開近傍で，
この同じ全方向族が実際のsmooth properなnc del Pezzo族となる。
全モジュライの被覆と全 $\mathcal O$ 保存同値の一般点での認識はなお残る。
同時中心・退化曲線・頂点での降下・相対幾何のこの組合せについて，
包括的な優先権確認は未了である。

% !TeX root = ../nc-del-pezzo.tex
\section{次数1・2の構造対象付きモジュライ問題}
\label{sec:definitions}

\subsection{重み付き模型と相対Poisson錐}
基礎体を $\C$ とし、$n\ge2$ を固定する。基本偏極を $H$ と書く。
可換な滑らか模型とその切断環のHilbert級数は次の表で与えられる。
\begin{center}
\begin{tabular}{ccc}
\toprule
次数 & 方程式と生成元の重み & Hilbert級数 $h_{n,d}(t)$\\
\midrule
$1$ & $z^2=y^3+A_4(x)y+A_6(x)$、$(1^n,2,3)$
    & $\displaystyle\frac{1-t^6}{(1-t)^n(1-t^2)(1-t^3)}$\\[5pt]
$2$ & $z^2=f_4(x)$、$(1^{n+1},2)$
    & $\displaystyle\frac{1-t^4}{(1-t)^{n+1}(1-t^2)}$\\
\bottomrule
\end{tabular}
\end{center}
$h_{n,d}(j)=[t^j]h_{n,d}(t)$ とし、$j<0$ では $h_{n,d}(j)=0$ とする。
変数の個数はどちらも $n+2$ であり、重みの和から関係の次数を引くと $n-1$ となる。

\begin{lemma}[滑らかな標準形と通常の射影スキーム]
\label{lem:weighted-free-locus}
上の標準形が滑らかなdel Pezzo族 $X\to T$ を与えるとき、
その重み付き射影超曲面はambientの安定化群が非自明な部分と交わらない。
従ってその重み付き商スタックは $X$ 自身であり、切断環 $B$ に対して
$\QGr B\simeq\operatorname{QCoh}(X)$ となる。
\end{lemma}
\begin{proof}
次数2では $x_i$ のどれかが非零なら重み1により安定化群は自明であり、
全てが零なら唯一の候補点 $[0:\cdots:0:z]$ は $z^2=f_4(0)=0$ を満たさない。
次数1でも $x_i\ne0$ であれば同様である。全て $x_i=0$ なら $z^2=y^3$ であり、
射影点であるため $y,z$ はともに非零である。安定化群は
$\mu_2\cap\mu_3=1$ となる。この議論は幾何学的ファイバーで成り立ち、
位数が可逆な対角化可能群の作用であることから相対的にも自由な部分に入る。
次数付き加群のtorsion商はこの自由な重み付き商上の準連接層圏である。
\end{proof}

標準式が滑らかな多様体を定める係数空間の開集合を $U_d$ とする。
以下の相対錐と商表示は \cite[定理7.3、定理8.8]{PoissonDP} を入力とする：
\[
 C_{n,1}=\A^n,\qquad
 C_{n,2}=\operatorname{Spec}\C[p_{ij}]/(q_{ijk\ell}),\qquad
 \mathcal P_{n,d}=[(U_d\times C_{n,d})/G_d].
\]
ここで $\mathcal P_{n,d}$ は相対Poisson構造付きの滑らかなdel Pezzo族のスタックであり、
零構造も含む。標準形の変換群とその中央部分群は
\begin{align*}
 \widetilde G_1&=\mathrm{GL}_n\times\Gm,
 &Z_1&=\{(\lambda I_n,\lambda)\},\\
 \widetilde G_2&=\mathrm{GL}_{n+1}\times\Gm,
 &Z_2&=\{(\lambda I_{n+1},\lambda^2)\},
 &G_d&=\widetilde G_d/Z_d.
\end{align*}
作用は
\begin{align*}
 (A,t)\colon (A_4,A_6,D)&\longmapsto
 \bigl(t^4A_4(A^{-1}x),t^6A_6(A^{-1}x),t^{-1}AD\bigr),\\
 (A,b)\colon (f_4,P)&\longmapsto
 \bigl(b^2f_4(A^{-1}x),b^{-1}(\wedge^2 A)P\bigr).
\end{align*}
これは偏極の自明化のスカラーを除いた作用であり、Poisson構造自体を
任意の非零スカラーで同一視する操作ではない。

\subsection{豊富列の存在と標識の選択}
モジュライ問題は局所noetherianな複素スキームのfppf site上で考える。
まず各affine基底 $T=\Spec R$ 上で、$R$-flatな局所noetherian Grothendieck圏
$\mathcal Q$ と $R$-flatな対象 $\OO\in\mathcal Q$ を取る。
基底変更は $R'$-作用付き対象の圏 $\mathcal Q_{R'}$ により定め、
構造対象の基底変更は $R'\otimes_R\OO$ とする。
以下の表示条件はfppf局所的な存在条件とし、表示そのものは標識に含めない。

\begin{definition}[豊富列型の許容表示]
\label{def:ample-model}
$(n,d)$ 型の許容表示とは、$E_0=\OO$ を満たす相対豊富列
$(E_i)_{i\in\mathbb Z}$ であって、
\[
 b_{ij}=\Hom(E_j,E_i)\quad(i\le j),\qquad b_{ij}=0\quad(i>j)
\]
が正・連結・局所有限なnoetherian $\mathbb Z$-代数をなし、
$b_{ij}$ が階数 $h_{n,d}(j-i)$ の有限局所自由加群で、これらのHomと合成が
任意の基底変更と可換し、標準関手
$\mathcal Q\simeq\operatorname{Qmod}b$ が対象 $E_i$ と両立するものをいう。
これらのnoetherian性、局所有限性、豊富性、復元と
$\mathcal Q$ の局所noetherian性は、許される任意の局所noetherian基底変更後にも要求する。
相対豊富性には、有限表示対象への十分大きい添字からの有限和による全射と、
全射に対する十分大きい添字でのHomの全射性を要求する。
\end{definition}

このような表示を局所的に持つ $(\mathcal Q,\OO)$ の射は、
\emph{全ての} $R$-線型完全同値 $F:\mathcal Q\simeq\mathcal Q'$ と
$\alpha:F(\OO)\simeq\OO'$ の組を、$\alpha$ を保つ自然同型で割ったものとする。
選んだ豊富列、次数シフト、生成元、またはAS包絡を保つことは要求しない。
このfibered groupoidのfppf stack化を
$\mathfrak A^{\OO}_{n,d}$ と書く。必要なら固定した宇宙で小さい模型を取る。
stack化は降下データを有効にする操作であり、有限型atlasの存在を意味しない。

\begin{lemma}[標識と自然自己同型]
\label{lem:marking-removes-two-aut}
許容表示を持つ圏では
$\operatorname{Nat}(\operatorname{Id}_{\mathcal Q},\operatorname{Id}_{\mathcal Q})=R$ である。
従って $\OO$ の標識を保つ同値の自然自己同型は恒等だけである。
\end{lemma}
\begin{proof}
自然変換は $\End(E_i)=R$ 上でスカラー $\lambda_i$ として作用する。
$\Hom(E_{i+1},E_i)$ は正階数局所自由なので自然性から
$\lambda_i=\lambda_{i+1}$ となる。豊富列からの全射と、
各対象がnoetherian部分対象のfiltered colimitであることを用いて、
圏全体で同じスカラーとなる。標識を保つなら $E_0=\OO$ 上で1である。
\end{proof}

\begin{definition}[非可換del Pezzo多様体のモジュライ問題]
\label{def:nc-dp-stack}
$\mathfrak M^{\OO}_{n,d}$ は、$\mathfrak A^{\OO}_{n,d}$ の対象で、
導来圏 $D(\mathcal Q)$ がcompact生成され、
$\Perf(\mathcal Q/R):=D(\mathcal Q)^c$ とその標準dg enhancementについて
次を満たすものからなる
fullな部分問題をstack化したものとする。
\begin{enumerate}
\item 幾何学的ファイバーのcoherent圏はHom有限でホモロジー次元 $n$ を持ち、
      その有界導来圏は滑らかで固有なdg圏をなす。
\item 相対perfect部分圏は豊富列の対象を含み、その相対RHomはperfectで基底変更と可換し、
      相対Serre双対が存在する。$\mathcal Q$ の完全自己同値 $\mathsf H$ が存在し、
      その導来関手の相対perfect部分圏への制限が $\mathsf S^{-1}[n]$ と同型である。
      $(\OO,\mathsf H)$ は相対的に豊富である。
\item 基底変更は上のcompact部分圏を保ち、相対Serre双対と両立する。
\item $r=n-1$ とする。上記と同じ許容豊富列を、
      $\mathsf H(E_i)\simeq E_{i-r}$ を整合的に満たすようにfppf局所的に取れる。
      さらに $i\le j$ に対して $R\Hom(E_j,E_i)$ は
      次数零に集中する階数 $h_{n,d}(j-i)$ の局所自由複体である。
\end{enumerate}
これらの性質は許される基底変更後にも要求する。
これらの構造の選択を追加の標識にはせず、その存在と保存を要求する。
標準dg enhancementには複体のdg圏を非輪状複体で局所化したものを用いる。
完全Grothendieck圏同値は全余極限とcompact対象を保ち、複体への項別適用から
このenhancementの同値を誘導する。Serre構造は双対の表現として一意なので、
これらの条件は全 $\OO$ 保存同値に対して不変である。
最後の条件を\emph{Serre周期 $n-1$ と有向ACM条件}と呼ぶ。
周期の同型、豊富列、基本的な次数1自己同値を標識に加えていない。
これは有限長のfull例外集合の存在を要求する条件でもない。
相対的な基底線束の捻りはfppf局所で自明化してこの条件を検査する。
\end{definition}

さらに、$\mathfrak M^{\OO}_{n,d}$ の対象のうち、
fppf局所的に $\mathcal Q\simeq\QGr(A/(\Omega))$ と書ける部分問題を考える。
$A$ は各次数が有限局所自由な連結次数付き代数で、幾何学的ファイバー上で
生成重み $(1^n,2,3)$ または $(1^{n+1},2)$、大域次元 $n+2$、
Gorensteinパラメータ $n+5$ または $n+3$ のAS正則代数とし、
$\Omega$ は次数6または4の正規元で、任意の許される基底変更後にも左右正則とし、
商の各次数は平坦でHilbert級数を $h_{n,d}$ とする。
この\emph{重み付き超曲面による部分問題}を $\mathfrak M^{\OO,\mathrm{hyp}}_{n,d}$ と書く。
中心的な $\Omega$ はその特別な場合である。AS包絡の選択は忘れ、射はやはり全 $\OO$ 保存同値とする。

\begin{remark}[定義と証明すべき性質]
以上でモジュライ問題は定まるが、代数スタック性、有限型性、
重み付き超曲面による部分問題との一致は定義からは従わない。
一方、Serre周期と有向ACM条件により、$\mathfrak M^{\OO}_{n,d}$ では
命題\ref{prop:ncdp-intrinsic-degree}を使って次数の内在性を証明できる。
広い $\mathfrak A^{\OO}_{n,d}$ について同じ結論は主張しない。
第\ref{sec:formal-abel}節で全Abel Artin変形の同型類を計算し、
第\ref{sec:formal-groupoid}節で全同値を含む群oidを同定する。
compact性・相対Serre条件・反標準豊富性の持上げを証明した後、
定理\ref{thm:all-artin-are-nc-dp}でそれを $\mathfrak M^{\OO}_{n,d}$ と
$\mathfrak M^{\OO,\mathrm{hyp}}_{n,d}$ のArtin完成に同定する。
非可換点を含む零切断近傍の全点での全Artin比較は，
定理\ref{thm:all-points-abel-artin-versality}で別に証明する。
定理\ref{thm:zero-neighborhood-open-full-moduli}はこの近傍を全問題の開部分として実現するが，
一般の有限型基底上で全モジュライと超曲面部分問題が一致するとは主張しない。
この区別が、次数付き表示だけを分類する商と、構造対象のみを標識する分類を比較する際に必要である。
\end{remark}

% !TeX root = ../nc-del-pezzo.tex
\section{可換点の全 Abel Artin 変形と形式局所環}
\label{sec:formal-abel}

本節では滑らかな複素 del Pezzo manifold $X$ を固定する．
$n=\dim X\geq2$，$-K_X\simeq(n-1)H$，$H^n=d\in\{1,2\}$ とする．
全ての試験環は剰余体 $\mathbb C$ を持つ通常の可換局所 Artin
$\mathbb C$-代数である．全 Abel 変形を扱う形式的な結果と，
有限型基底上の非可換モジュライの構成とは区別する．

\subsection{構造対象と豊富列の持上げ}

\begin{definition}[framed な構造対象付き Abel 変形]
\label{def:framed-o-abel}
$\mathfrak D^{\mathcal O}_X(R)$ の対象は，
Lowen--Van den Bergh の意味で平坦な $R$-線型 Grothendieck 圏
$\mathcal D$，その $\mathbb C$-対象の部分圏への同値
$\iota:\operatorname{Qcoh}(X)\xrightarrow{\sim}\mathcal D_{\mathbb C}$，
および $\mathcal O_X$ の $R$-平坦な持上げ $\mathcal O_R$ の組とする．
射は $R$-線型完全同値 $F$ と同型
$\alpha:F(\mathcal O_R)\xrightarrow{\sim}\mathcal O'_R$ の組で，
$F_{\mathbb C}\iota\simeq\iota'$ を満たすものの，
$\alpha$ と両立する自然同型類である．
この還元後の自然同型自体を追加の frame として保持しない．
その同型類集合を $\mathrm{Def}^{\mathcal O}_X(R)$ と書く．
基底変更には Grothendieck 圏の係数変更を用いる．
\end{definition}

\begin{lemma}[有向な線対象列の消滅]
\label{lem:directed-line-vanishing}
任意の整数 $a\leq b$ と $q>0$ に対して
\[
 \operatorname{Ext}^q_X(\mathcal O_X(aH),\mathcal O_X(bH))=0.
\]
特に各 $\mathcal O_X(aH)$ は例外対象である．
\end{lemma}
\begin{proof}
左辺は $H^q(X,\mathcal O_X((b-a)H))$ である．
$(b-a)H=K_X+(b-a+n-1)H$ と書くと，
$n\geq2$ より右辺第二項は ample である．Kodaira 消滅を適用する．
$H^0(X,\mathcal O_X)=\mathbb C$ だから例外性も従う．
\end{proof}

\begin{proposition}[全 Artin 変形の section 表示と標識忘却]
\label{prop:directed-section-deformation}
$E_i=\mathcal O_X(-iH)$ とし，
$a_{ij}=H^0(X,\mathcal O_X((j-i)H))$ $(i\leq j)$，
$a_{ij}=0$ $(i>j)$ を section $\mathbb Z$-代数とする．
任意の平坦 Abel $R$-変形 $\mathcal D$ に対して，
各 $E_i$ の平坦持上げ $E_{i,R}$ が同型を除いて一意に存在する．
さらに
\[
 b_{ij}=\operatorname{Hom}_{\mathcal D}(E_{j,R},E_{i,R})
 \quad(i\leq j),\qquad b_{ij}=0\quad(i>j)
\]
は $a$ の平坦線型変形を与え，
$\mathcal D\simeq\operatorname{Qmod}(b)$ である．
各 $b_{ij}$ は有限自由で，階数は $a_{ij}$ の次元に等しい．
$b$ は正，連結，局所有限かつ noetherian である．
自然同型を法にした定義\ref{def:framed-o-abel}の $1$-群oidでは，
$\mathcal O_R$ を忘れる関手は全 Abel 変形群oidとの同値である．
\end{proposition}
\begin{proof}
$H$ は ample なので $(\mathcal O_X(iH))_{i\in\mathbb Z}$ は
$\mathbb Z$-generating sequence である．
補題\ref{lem:directed-line-vanishing}は
De Deken--Lowen \cite[Theorem~4.5, pp.~495--496]{DDL2011}
の有向側の $\operatorname{Ext}^{1,2}$ 消滅仮定を満たす．
同定理から，線型変形と Abel 変形の比較，
平坦な線対象の持上げ，および Hom の平坦性を得る．
同論文 Proposition~4.6 により局所有限性と noetherian 性が持ち上がる．
$R$ は局所 Artin 環なので有限平坦な Hom は有限自由であり，
その階数は特殊ファイバーによって決まる．

特に $\operatorname{End}(\mathcal O_R)$ は階数1の有限自由加群である．
単位写像 $R\to\operatorname{End}(\mathcal O_R)$ は還元後同型だから同型となる．
従って任意の変形同値 $F$ に対して
$F(\mathcal O_R)\simeq\mathcal O'_R$ が存在し，
二つのそのような同型の差は $R^\times$ の元である．
この差は $F$ 全体のスカラー自然自己同型で補正できる．
同様に，二つの同値間の自然同型はスカラーを補正して
指定された $\alpha$ と両立させられる．
従って標識忘却は本質的全射かつ完全忠実である．
\end{proof}

\begin{remark}[曲面と重み付き生成の注意]
この適用では $a\leq b$ の消滅だけを用いた．
曲面の負の捻りに対する $H^2$ の消滅は主張しない．
また，次数1の切断環は重み $(1^n,2,3)$，関係次数6，
次数2では重み $(1^{n+1},2)$，関係次数4を持つ．
したがってこれらを通常の次数一生成環と扱ってはならない．
上の比較は有限生成 $\mathbb Z$-代数の tails topology によるものであり，
この比較だけから通常の graded 周期化や重み付きAS包絡の存在は結論しない。
定理\ref{thm:concrete-quantization-all-artin}と
系\ref{cor:all-artin-weighted-hypersurface-presentation}では、
具体量子化と全形式環の比較をさらに使い、通常の次数付き中心商表示の存在まで証明する。
非切断の $2$-群oidではスカラー自然変換が残るため，
標識忘却についても同じ結論をそのまま適用しない．
\end{remark}

\subsection{全変形の量子化比較}

$\mathrm{PoisDef}_X(R)$ を，$\mathcal O_X$ の平坦な可換代数層の変形と，
特殊ファイバーで零になる $R$-線型 Poisson 括弧の組の，
還元を保つ同型による同値類集合とする．
可換積自体も変化を許す．

\begin{proposition}[既知の量子化比較の適用]
\label{prop:abel-poisson-formal-comparison}
全ての局所 Artin $\mathbb C$-代数の射に自然な全単射
\[
 \mathrm{PoisDef}_X(R)\xrightarrow{\sim}
 \mathrm{Def}^{\mathcal O}_X(R)
\]
が存在する．この比較は一次の Poisson 括弧を保つ．
\end{proposition}
\begin{proof}
$X$ の有限アフィン被覆とその有限交叉を取り，
各被覆開に含まれるアフィン開全体を考える．
これは $\mathcal O_X$ に関する acyclic basis であり，
準連接層の制限は局所化を与える．
Lowen \cite[Theorems~3.14,3.26]{Lowen2008} により，
全 Abel 変形は $\mathcal O_X$ の algebroid 変形で記述される．
補題\ref{lem:directed-line-vanishing}から
$H^1(X,\mathcal O_X)=H^2(X,\mathcal O_X)=0$ である．
従って Yekutieli \cite[Corollary~7.5]{Yekutieli2015} により，
twisted な代数層変形の同値類は普通の結合代数層変形の gauge 同値類と一致する．
同論文 Corollary~11.2 は，この結合代数層変形と
Poisson 代数層変形の間に，係数環の射に自然で一次括弧を保つ全単射を与える．
最後に命題\ref{prop:directed-section-deformation}の標識忘却を用いる．
\end{proof}

\begin{remark}[比較定理の適用範囲]
使用した Yekutieli の結果は可換積も変化させる Corollary~11.2 である．
固定した $R\otimes\mathcal O_X$ 上の星積を扱う Corollary~11.3 は
$H^1(T_X)=H^1(\mathcal D_X)=0$ を追加仮定するため用いない．
また，高い $H^q(X,\bigwedge^pT_X)$ の一括した消滅は必要ない．
上の全単射は既知の一般定理の適用であり，その量子化原理自体を本研究の成果とはしない．
\end{remark}

\begin{lemma}[Poisson 代数層変形の幾何的解釈]
\label{lem:poisson-geometric-artin}
$\mathrm{PoisDef}_X(R)$ は，$X$ を特殊ファイバーとして固定した
滑らかな次数 $d$ の del Pezzo $R$-族と，特殊ファイバー上零の相対 Poisson 構造の組の
同型類関手に等しい．
\end{lemma}
\begin{proof}
平坦な可換代数層変形を考える．滑らかなアフィン代数は形式的に滑らかなので，
アフィン開上ではその可換 Artin 変形を自明化できる．
これらは重なり上の代数同型によって貼り合い，
滑らかな有限表示 $R$-スキーム $X_R$ を与える．
特殊ファイバーは $X$ であり，冪零厚化の下で固有性は保たれる．

小拡大の核を $I$ とすると，$H$ の持上げの障害は
$I\otimes H^2(X,\mathcal O_X)$，同型類の差は
$I\otimes H^1(X,\mathcal O_X)$ に属するので消滅する．
従って $H$ は同型を除いて一意に線束 $H_R$ へ持ち上がる．
同じ一意性により $\omega_{X_R/R}^{-1}\simeq H_R^{\otimes(n-1)}$ であり，
ample 性は冪零厚化に不変だから $H_R$ は相対 ample である．
次数も特殊ファイバーから変わらない．
滑らかさの下で二重導分と相対二重ベクトルは一致し，
Jacobi 恒等式は $[\pi,\pi]=0$ に等しい．逆構成は構造層を取ればよい．
\end{proof}

\subsection{形式環と唯一の形式分枝}

\begin{theorem}[可換点の全 Abel 変形の形式局所環]
\label{thm:abel-formal-ring}
\[
 b_{n,1}=\binom{n+3}{4}+\binom{n+5}{6}-n^2,
 \qquad
 b_{n,2}=\binom{n+4}{4}-(n+1)^2
\]
とおく．定義\ref{def:framed-o-abel}の同型類関手は，非標準的に
\begin{align*}
 d=1:\quad R_X&\simeq
 \mathbb C[[t_1,\ldots,t_{b_{n,1}},D_1,\ldots,D_n]],\\
 d=2:\quad R_X&\simeq
 \frac{\mathbb C[[t_1,\ldots,t_{b_{n,2}},p_{ij}\ (0\leq i<j\leq n)]]}
 {(p_{ij}p_{kl}-p_{ik}p_{jl}+p_{il}p_{jk}\mid i<j<k<l)}
\end{align*}
で prorepresent される．有限群 $\operatorname{Aut}(X)$ による商は取っていない．

次数1の環は次元 $b_{n,1}+n$ の正則局所環である．
次数2の環は正規かつ Cohen--Macaulay な整域であり，
Krull 次元は $b_{n,2}+2n-1$，接空間次元は
$b_{n,2}+\binom{n+1}{2}$ である．
$n=2$ では正則で，$n\geq3$ では可換点で特異である．
各場合とも可換点の形式近傍はちょうど一つの既約分枝を持つ．
\end{theorem}
\begin{proof}
命題\ref{prop:abel-poisson-formal-comparison}と
補題\ref{lem:poisson-geometric-artin}により，
Poisson 族の零構造点における形式関手を計算すればよい．
次数1・2の相対構造定理と全族の商表示
\cite[定理7.3，定理8.8]{PoissonDP}を用いる．
通常の del Pezzo スタックは滑らかで次元 $b_{n,d}$ であり，
$H^0(X,T_X)=0$ なので，特殊ファイバーとの同一視を保つ通常の変形関手は
形式円板 $\operatorname{Spf}\mathbb C[[t_1,\ldots,t_{b_{n,d}}]]$ となる．
例えば標準形の係数空間で座標変換群の軌道に横断する局所 slice を取ればよい．

この形式基底上で基本線束と標準座標を局所的に選ぶと，
任意基底変換と両立する相対 Poisson スキームは
次数1で $\mathbb A^n$，次数2で $\operatorname{Gr}(2,n+1)$ の
Pl\"ucker アフィン錐の積族となる．
特殊ファイバーの Poisson 構造を零とするため，
パラメータ $D_i$ または $p_{ij}$ は $\mathfrak m_R$ に値を取る．
原点で完備化すれば表示した環を得る．

次数2の錐の次元は $2n-1$ であり，
Pl\"ucker 方程式は全て二次なので原点の接空間次元は
$\binom{n+1}{2}$ である．$n=2$ では方程式がなく，
$n\geq3$ では接空間次元と次元の差は
$(n-1)(n-2)/2>0$ となる．
Pl\"ucker 座標環の正規性と Cohen--Macaulay 性\cite{Svanes1973,DeConciniProcesi1981}は積と完備化で保たれる．
ここで有限型 $\mathbb C$-代数の優秀性を用いれば，
頂点の正規局所環の完備化も正規局所環であり，整域となる．
従って既約分枝は一つである．
\end{proof}

\begin{corollary}[一次方向と二次障害]
\label{cor:abel-quadratic-obstruction}
次数1では同型類関手 $\mathrm{Def}^{\mathcal O}_X$ は形式的に滑らかである．
次数2の一次方向を $(v,Q)$ と書く．
ここで $v\in H^1(X,T_X)$，
$Q\in H^0(X,\bigwedge^2T_X)\simeq\bigwedge^2\mathbb C^{n+1}$ である．
この方向が $\mathbb C[\varepsilon]/(\varepsilon^3)$ 上へ持ち上がるための
必要十分条件は $Q\wedge Q=0$ である．
従って通常の多様体の一次変形を加えても，非零の Pl\"ucker 二次障害は相殺できない．
条件を満たす方向は，全ての $\mathbb C[t]/(t^{r+1})$ 上の
両立する Abel 変形同型類へ延長できる．
\end{corollary}
\begin{proof}
定理\ref{thm:abel-formal-ring}を用いる．
次数2で $P=\varepsilon Q+\varepsilon^2Q_2$ と書くと
$P\wedge P=\varepsilon^2Q\wedge Q$ であり，
式には通常の変形パラメータ $t_i$ は現れない．
$Q\wedge Q=0$ なら $P=tQ$ と任意の形式円板上の直線を選べば，
全次数の両立する点を得る．
\end{proof}

\begin{remark}[形式的な結果と大域モジュライ]
本節で決定したのは，全 Abel Artin 変形の同型類関手である．
重み付き AS 包絡などを追加条件として課した大域非可換モジュライ関手を用いる場合，
全 Artin 変形がその条件を満たすことを示すまでは，
その関手の完備化を $\operatorname{Spf}R_X$ と同一視しない．
また，形式分枝の一意性は，全非可換モジュライの既約性，
可換 locus に接しない成分の不存在，有限型 atlas，
形式量子化の代数化，内在的な Poisson モジュライ射，
または Kuznetsov 成分の比較を意味しない．
それらは別の研究課題である．
\end{remark}
\par

% !TeX root = ../nc-del-pezzo.tex
% Conditional reconstruction results for nc-del-pezzo.
% Citation keys to provide: MoriUeyama2021, StacksProject, CubicModuliSource.
\section{反標準軌道環と構造対象を保つ全同値}
\label{sec:ncdp-anticanonical-reconstruction}

本節の圏論的結果は，明示した相対復元条件の下での命題である．
重み付き代数の Hilbert 級数や体上の AS--Gorenstein 性だけから，
その条件が任意の基底上で成立すると主張しない．
反標準自己同値を用いる発想は，三次三様体の偶数列による復元を
一般の指数へ拡張するものである．

\subsection{体上の反標準復元}

$\mathcal C$ を $\mathbb C$-線型 Abel 圏とし，$D^b(\mathcal C)$ は
Hom 有限で Serre 関手 $\mathsf S$ を持つとする．
$\mathsf H:=\mathsf S^{-1}[n]$ が heart を保つとき，
その heart 上の制限も $\mathsf H$ と書く．これは canonical bimodule の
逆に対応する自己同値である
\cite[Definition 3.4 and Remark 3.5]{MoriUeyama2021}．
対象 $\mathcal O\in\mathcal C$ に対し
\[
 R(\mathcal C,\mathcal O)
 :=\bigoplus_{m\ge0}\operatorname{Hom}_{\mathcal C}
       (\mathcal O,\mathsf H^m\mathcal O)
\]
と置く．$a\in R_i$，$b\in R_j$ の積は
$ab=\mathsf H^j(a)\circ b$ とする．
この規約は右加群の section 環と整合する
\cite[Definition 2.1(5)]{MoriUeyama2021}．

\begin{proposition}[反標準軌道環による全 $\mathcal O$ 保存同値]
\label{prop:ncdp-anticanonical-hom}
$(\mathcal C,\mathcal O)$ と $(\mathcal C',\mathcal O')$ は上記の条件を満たし，
さらに各々について次を仮定する．
\begin{enumerate}
\item $(\mathcal O,\mathsf H)$ は豊富な対であり，
      $\operatorname{End}(\mathcal O)=\mathbb C$，$R_1\ne0$ である．
\item 負の反標準次数の切断は消滅する：
      $\operatorname{Hom}(\mathcal O,\mathsf H^m\mathcal O)=0$ for $m<0$．
\end{enumerate}
このとき $R=R(\mathcal C,\mathcal O)$，
$R'=R(\mathcal C',\mathcal O')$ に対して自然な全単射
\[
 \operatorname{Isom}_{\mathrm{gr}}(R,R')/\mathbb C^\times
 \ \simeq\
 \pi_0\operatorname{Eq}_{\mathcal O}(\mathcal C,\mathcal C')
\]
がある．左辺の作用は次数 $m$ 上で $c^m$ を掛ける次数スカラーであり，
右辺は全ての完全線型同値 $F$ と同型
$\alpha:F(\mathcal O)\simeq\mathcal O'$ の組を，
$\alpha$ を保つ自然同型で割った集合である．
標識を保つ自然自己同型は自明であり，これらの同定は合成と両立する．
\end{proposition}
\begin{proof}
$E_i=\mathsf H^{-i}\mathcal O$ と置く．各 $E_i$ の自己準同型はスカラーであり，
$\operatorname{Hom}(E_{i+1},E_i)=R_1\ne0$ である．
恒等関手の自然変換は $E_i$ 上のスカラー $\lambda_i$ により表され，
隣接 Hom の自然性から全ての $\lambda_i$ は等しい．
豊富列から各対象への全射を用いると，圏全体でも同じスカラーになる．
従って $\operatorname{Nat}(\operatorname{Id},\operatorname{Id})=\mathbb C$ であり，
$\mathcal O$ 上で恒等な自然自己同型は自明である．

完全 Abel 圏同値は有界導来圏の同値を誘導する．
双対同型を備えた Serre 関手の Yoneda 一意性
\cite[Lemma 57.3.5, Tag 0FY3]{StacksProject} により
$F\mathsf H\simeq\mathsf H'F$ を得る．
その整合的な正負の反復と $\alpha$ は，
$F(E_i)\simeq E'_i$ と，すべての
\[
 \operatorname{Hom}(E_j,E_i)=R_{j-i}
\]
に対する同型を与える．反復の整合性により，これは一つの次数付き環同型
$\phi:R\simeq R'$ から来ており，すべての添字での合成を保つ．
比較同型 $F\mathsf H\simeq\mathsf H'F$ の選択を変えたときの差は，
準逆を合成すれば恒等関手の自然自己同型となる．
従って差は一つの非零スカラーであり，$m$ 回の反復でその $m$ 乗が現れる．
このため $\phi$ の次数スカラーを法とした類は選択によらない．

豊富対による復元
\cite[Theorem 2.6(2)]{MoriUeyama2021} は
$\mathcal C\simeq\operatorname{tails}R$ を与える．
この同値は $\mathcal O$ と $\mathsf H$ をそれぞれ
$\pi R$ と次数シフトに移す．
したがって任意の次数付き同型 $R\simeq R'$ は構造対象付き Abel 圏同値を誘導し，
上述の section $\mathbb Z$-代数全体の比較により，元の $F$ もこの同値へ自然同型となる．

核を調べるため，$\phi\in\operatorname{Aut}_{\mathrm{gr}}R$ の構造移送が
標識付き恒等関手へ自然同型であるとする．
$E_i$ 上のその自然同型を $\lambda_i\in\mathbb C^\times$，$\lambda_0=1$ と書く．
隣接 Hom の自然性は $\phi|_{R_1}=c\operatorname{Id}$ と
$\lambda_{i+1}/\lambda_i=c$ を与える．
さらに全ての $\operatorname{Hom}(E_{i+m},E_i)=R_m$ に自然性を適用すると，
$\phi|_{R_m}=c^m\operatorname{Id}$ となる．
ここでは $R$ が次数一で生成されるとは仮定していない．
逆に次数スカラーは，各次数に適切なスカラーを掛ける自然同型により
標識付き恒等関手を誘導する．これで核と一意性を得る．
\end{proof}

\begin{remark}[通常の Veronese と一致する条件]
\label{rem:ncdp-serre-twist}
$\mathcal C=\operatorname{tails}B$，$\mathcal O=\pi B$ とし，
section 写像が飽和していて
$\mathsf S\simeq(-r)[n]$，$r=n-1$ ならば
$R(\mathcal C,\mathcal O)\simeq B^{(r)}$ となる．
AS--Gorenstein パラメータが $r$ という条件だけでは，この結論は出ない．
一般の balanced dualizing complex は
$B_\nu(-r)[n+1]$ であり，Nakayama 自己同型 $\nu$ の捻りを持ち得る
\cite[\S1.2, after Definition 1.3]{MoriUeyama2021}．
その場合の反標準軌道環には対応する Zhang 型の捻りが現れ，
通常の $B^{(r)}$ と無条件に同一視できない．
また，重み付き代数では全シフト列の豊富性から
$r$ 倍シフトだけの豊富性を推論してはならない．
\end{remark}

\begin{proposition}[相対版と基底線束]
\label{prop:ncdp-relative-anticanonical-hom}
上の圏の族について，相対 Hom の局所自由性と任意の基底変更，
section $\mathbb Z$-代数による復元，固定模型間の同値と自然同型の fppf 降下を仮定する．
全ての完全線型同値が保存する標準の相対 perfect 部分圏上で，
相対 RHom の perfect 性，基底変更，Serre 双対も仮定する．
さらに反標準列が相対的に豊富であり，その隣接 Hom は正の階数のベクトル束であるとする．
このとき命題\ref{prop:ncdp-anticanonical-hom}は fppf 同型層の同型
\[
 \underline{\operatorname{Isom}}_{\mathrm{gr}}(R^{\mathrm{can}},(R')^{\mathrm{can}})
       /\mathbb G_m
 \simeq
 \underline{\operatorname{Eq}}_{\mathcal O}(\mathcal C,\mathcal C')
\]
へ拡張する．特に局所的な graded 模型で
$\mathsf S\simeq(-r)\otimes L[n]$ と書ける場合は
\[
 R^{\mathrm{can}}_m\simeq B_{rm}\otimes L^{-m}.
\]
商は fppf 層商であり，大域切断の集合の商ではない．
\end{proposition}
\begin{proof}
完全線型同値を複体へ項ごとに適用し，相対 perfect 部分圏の保存と
相対 Serre 双対の表現の一意性を使う．
相対 Hom とその合成の比較は前命題の証明と同じである．
隣接 Hom の正階数局所自由性によりスカラーの比較は基底上で忠実であり，
恒等関手の自然自己変換は基底の一つのスカラーに限られる．
相対 section 復元を用いて得られた局所同値と自然同型を降下する．
$L$ の自明化を変更すると反標準軌道環の次数 $m$ には
その逆の $m$ 乗が現れるので，表示した線束の捻りが必要である．
次数スカラーを除いた同型層はこれらの変更と両立する．
\end{proof}

滑らかな可換 del Pezzo $n$-fold では，$\mathsf H$ は
$-\otimes\omega_X^{-1}$ であり，反標準束の ample 性から上の豊富性を得る．
したがってこの場合の反標準軌道環は通常の反標準環である．
非可換族への適用では，相対 Serre 同型，saturation，反標準列の復元を別々に確認する必要がある．

\section{二つ以上の中心一次変数による任意指数の Veronese 復元}
\label{sec:ncdp-veronese-recovery}

以下の定理は $B$ の次数一生成を必要としないので，重み付きの環にも適用できる．
ただし中心の条件は強い条件であり，可換点や冪零な量子化パラメータを含む全族には課さない．

\begin{theorem}[任意指数の相対 Veronese 持上げ]
\label{thm:ncdp-veronese-lift}
$T$ を複素スキーム，$r\ge2$ とする．
$B,B'$ は非負次数付きの単位的 $\mathcal O_T$-代数で，
$B_0=B'_0=\mathcal O_T$ とする．
階数 $k\ge2$ のベクトル束 $E\subset B_1$，$E'\subset B'_1$ があり，
$P=\operatorname{Sym}E\hookrightarrow B$ と
$P'=\operatorname{Sym}E'\hookrightarrow B'$ は中心的な埋込みとする．
次を仮定する．
\begin{enumerate}
\item $B$ は $P$ 上平坦であり，$B'$ は $P'$ 上平坦である．
\item 任意の基底変更後に $Z(B)_r=P_r$，$Z(B')_r=P'_r$ が成り立つ．
\end{enumerate}
このとき $R=B^{(r)}$，$R'=(B')^{(r)}$ に対する制限射
\[
 \underline{\operatorname{Isom}}_{\mathrm{gr},T}(B,B')
 \longrightarrow
 \underline{\operatorname{Isom}}_{\mathrm{gr},T}(R,R')
\]
は $\mu_{r,T}$-torsor である．核の作用は次数 $m$ 上で
$\zeta^m$ を掛ける次数スカラーである．
従って次数スカラーを法とした fppf 層には自然な同型
\[
 \underline{\operatorname{Isom}}_{\mathrm{gr},T}(B,B')/\mathbb G_m
 \simeq
 \underline{\operatorname{Isom}}_{\mathrm{gr},T}(R,R')/\mathbb G_m
\]
がある．右辺の次数スカラーへの写像は $\lambda\mapsto\lambda^r$ である．
\end{theorem}
\begin{proof}
全ての議論は任意の基底変更と両立するので，束を自明化した基底の affine 開集合上で証明する．
$u,v$ を $E$ の基底の最初の二元に取る．
$P$ 上の平坦性により $u,v$ は $B$ 上の非零因子である．
まず $z\in Z(B)_1$ に対し
$zu^{r-1},zv^{r-1}\in Z(B)_r=P_r$ であるから，
$z\in P[u^{-1}]\cap P[v^{-1}]=P$ を得る．
従って $Z(B)_1=P_1$ であり，$P$ は $B$ の次数一中心が生成する
内在的な部分代数である．特に任意の $B\simeq B'$ は $P\simeq P'$ を誘導する．
$z\in Z(R)_1$ と $b\in B_m$ に対し，$m+a\equiv0\pmod r$ となる
$a\ge0$ を選ぶと $u^ab\in R$ なので
$u^a[z,b]=[z,u^ab]=0$ となる．$u^a$ を消去すると
$Z(R)_1=Z(B)_r=P_r$ を得る．
従って $R$ の次数一中心が生成する部分代数は，内在的に $P^{(r)}$ である．
同じ議論は $R'$ にも適用できる．

$\phi:R\simeq R'$ は $\theta:P^{(r)}\simeq(P')^{(r)}$ を誘導する．
まず $\theta$ を $P\simeq P'$ へ持ち上げる．
Proj を取ると偏極 $\mathcal O(r)$ 付きの射影空間束の同型を得る．
各幾何学的ファイバー上でこの同型は $\mathcal O(1)$ の類を保つ．
$\mathcal O(1)$ の差は基底の線束であり，基底変更に両立する
$H^0(\mathcal O)=\mathbb C$，$H^{>0}(\mathcal O)=0$ と評価写像からこれを確認できる．
その基底線束を局所的に自明化すれば，次数一切断により線型同型
$E\simeq E'$ が得られる．その $r$ 次対称冪と $\theta$ の差は一つの単元である．
この単元の $r$ 乗根を fppf 局所的に取れば $P$ の持上げを得る．
二つの持上げの差は射影空間上で恒等な線型写像なのでスカラーであり，
その $r$ 乗が $1$ である．従って持上げ層は $\mu_r$-torsor である．

この持上げを $\theta_1:P\simeq P'$ と書き，$u'=\theta_1(u)$，$v'=\theta_1(v)$ と置く．
局所化には自然な表示
\[
 B[u^{-1}]
 \simeq
 R[(u^r)^{-1}][t,t^{-1}]/(t^r-u^r),\qquad t\longmapsto u
\]
がある．右辺の $t$ は中心的で内部次数 $1$ とする．
実際，$b\in B_m$ に対して $m+a$ が $r$ の倍数なら
$b=(bu^a)u^{-a}$ と書けるので全射である．
$1,u,\ldots,u^{r-1}$ を係数 $R[(u^r)^{-1}]$ 上で考えると，
それぞれ異なる内部次数の剰余類を持つので直和であり，単射性も従う．
したがって $\phi$ と $u\mapsto u'$ は局所化の同型を一意に定める．
$v$ についても同様であり，重なりでは両者は一致する．
例えば $vu^{r-1}\in P_r$ を用いると，$u$ 側の延長も $v$ を $v'$ へ送ることが分かる．

任意の係数環上で Laurent 単項式を比較すると
\[
 0\longrightarrow P\longrightarrow
 P[u^{-1}]\oplus P[v^{-1}]
 \longrightarrow P[(uv)^{-1}]
\]
は最初の二つの非零項で完全である．
これを平坦な $P$-加群 $B$ とテンソルすることにより
\[
 B=B[u^{-1}]\cap B[v^{-1}]
 \quad\text{in }B[(uv)^{-1}]
\]
を得る．$B'$ にも同じ表示があるため，二つの局所同型は
$B\simeq B'$ を与える．これは $\theta_1$ を固定すると一意である．

持上げを $\zeta\theta_1$，$\zeta^r=1$ に変更すると，
$b=(bu^a)u^{-a}$ の像には $\zeta^{-a}=\zeta^m$ が掛かる．
従って持上げの差は正確に次数スカラー $\varepsilon_\zeta$ である．
最後に $\mathbb G_m\to\mathbb G_m$，$\lambda\mapsto\lambda^r$ は fppf 全射で，
その核は $\mu_r$ なので，商層の同型を得る．
\end{proof}

\begin{corollary}[一般次元への適用の形]
\label{cor:ncdp-veronese-dp}
次数1・2の非零 Poisson 構造の Casimir 模型には，いずれも $n-1$ 個の中心一次座標がある．
その量子化が前定理の平坦性・普遍中心条件と
命題\ref{prop:ncdp-relative-anticanonical-hom}の相対圏論的条件を満たし，
さらに section 写像が飽和していて，Nakayama の捻りを持たない相対 Serre 同型
$\mathsf S\simeq(-r)\otimes L[n]$ があるとする．このとき基底線束 $L$ を
fppf 局所的に自明化し，次数スカラーを法とした同定を降下すれば，
$r=n-1$ に対して
\[
 \underline{\operatorname{Eq}}_{\mathcal O}(\operatorname{QGr}B,
       \operatorname{QGr}B')
 \simeq
 \underline{\operatorname{Isom}}_{\mathrm{gr}}(B,B')/\mathbb G_m
\]
となる．$n\ge3$ では $n-1\ge2$ より前定理を使える．
$n=2$ では $r=1$ なので Veronese への移行そのものが恒等であり，持上げは不要である．
\end{corollary}

\begin{remark}[中心条件の確認と適用範囲]
$B$ の生成元の次数が有限個の $w$ にあり，その生成元束を $V_w$ とする．
$B_r/P_r$ と $B_{r+w}$ が有限局所自由なら，有限な交換子写像
\[
 B_r/P_r\longrightarrow
 \bigoplus_w\mathcal Hom(V_w,B_{r+w})
\]
のファイバー単射性は開条件である．この開上では局所分裂単射となるので，
$Z(B)_r=P_r$ が任意の基底変更後にも成立する．
ただしこの開の非空性は個々の量子化族で別途証明する必要がある．
特に $\hbar$ が冪零な量子化では $\hbar$-torsion による中心元が生じ得るため，
形式族の一般ファイバーでの中心の計算を，零切断の全 Artin 近傍へ適用しない．
\end{remark}

\section{反標準 Euler 多項式による次数の識別}
\label{sec:ncdp-intrinsic-degree}

以下は，全 $\mathcal O$ 保存同値を取った後にも次数1と次数2を識別するための
条件付きの不変量である．同一次数内の既約性を主張するものではない．

\begin{proposition}[次数の圏論的不変性]
\label{prop:ncdp-intrinsic-degree}
$n\ge2$，$r=n-1$ とし，$D^b(\mathcal C)$ は Hom 有限かつ有限な
Euler 和を持ち，Serre 関手 $\mathsf S$ を持つとする．
$\mathsf H=\mathsf S^{-1}[n]$ と置き，これはheart $\mathcal C$ を保つとする．
$\mathcal O\in\mathcal C$ と $d\in\{1,2\}$ に対し，
次の数値条件とACM 型消滅を仮定する：全ての $m\ge0$ に対し
\[
 \dim\operatorname{Hom}(\mathcal O,\mathsf H^m\mathcal O)=h_{n,d}(rm),
 \qquad
 \operatorname{Ext}^i(\mathcal O,\mathsf H^m\mathcal O)=0\quad(i>0).
\]
ここで $h_{n,d}(q)$ は次の級数の $t^q$ の係数である：
\[
 H_{n,1}(t)=\frac{1-t^6}{(1-t)^n(1-t^2)(1-t^3)},\qquad
 H_{n,2}(t)=\frac{1-t^4}{(1-t)^{n+1}(1-t^2)}
\]
この数値条件は通常の次数付きsection代数の存在を要求しない．このとき
\[
 p_{\mathcal C,\mathcal O}(m)
 :=\chi(\mathcal O,\mathsf H^m\mathcal O)
\]
は $m\ge0$ 上で次数 $n$ の多項式であり，前進差分 $\Delta p(m)=p(m+1)-p(m)$ に対し
\[
 \Delta^n p_{\mathcal C,\mathcal O}(m)=d\,r^n.
\]
従って，$n$ を固定したとき $d$ は全ての $\mathcal O$ 保存完全線型同値の不変量である．

さらに相対的な族で，上記の条件が各幾何学的ファイバー上で成り立ち，
$0\le j\le n$ に対する $R\mathcal Hom(\mathcal O,\mathsf H^j\mathcal O)$ が
基底上 perfect かつ基底変更に両立すると仮定する．
このとき $d$ は基底上局所定数である．
\end{proposition}
\begin{proof}
因数分解 $1+t^3=(1+t)(1-t+t^2)$ により
\[
 H_{n,1}(t)=\frac{1-t+t^2}{(1-t)^{n+1}},\qquad
 H_{n,2}(t)=\frac{1+t^2}{(1-t)^{n+1}}
\]
を得る．$q\ge0$ に対して，その係数はそれぞれ
\[
 \binom{q+n}{n}-\binom{q+n-1}{n}+\binom{q+n-2}{n},
 \qquad
 \binom{q+n}{n}+\binom{q+n-2}{n}
\]
である．ここで二項係数を上段の多項式として解釈する．
$n\ge2$ なので $q=0,1$ に現れる負の係数添字の項も零となり，
この式は全ての $q\ge0$ に有効である．
両者は次数 $n$，最高次係数 $d/n!$ を持つ．
仮定した消滅により $p(m)=h_{n,d}(rm)$ であり，その最高次係数は
$d r^n/n!$ である．第 $n$ 差分は $n!$ 倍の最高次係数なので所望の式を得る．

完全 Abel 同値は導来同値を誘導し，Serre 関手の一意性により
$\mathsf H$ と整合する．したがって $\mathcal O$ を保つ同値は
全ての上記 Euler 数を保つ．
相対的には perfect 複体の Euler 階数が局所定数なので，
\[
 d=\frac{1}{r^n}\sum_{j=0}^n(-1)^{n-j}\binom nj
      \chi(\mathcal O,\mathsf H^j\mathcal O)
\]
も局所定数である．
\end{proof}

滑らかな可換 del Pezzo $n$-fold では，反標準線束が ample であり，
Kodaira 消滅を $\omega_X\otimes\omega_X^{-(m+1)}$ に適用すると
$H^i(X,\omega_X^{-m})=0$ for $i>0$，$m\ge0$ を得る．
標準の重み付き section 環の飽和性と合わせ，上の条件を満たす．
非可換模型では，Hilbert 級数が一致することだけでこの ACM 型条件を省略してはならない．
相対 perfect 条件を満たす全圏の問題では，次数1部分と次数2部分は開かつ閉に分かれるが，
これは各部分の既約成分の個数を計算したことにはならない．

\begin{corollary}[定義したモジュライ問題での次数の開閉分解]
\label{cor:defined-moduli-degree-separation}
定義\ref{def:nc-dp-stack}の $\mathfrak M^{\mathcal O}_{n,1}$ と
$\mathfrak M^{\mathcal O}_{n,2}$ は全 $\mathcal O$ 保存同値で混同されない。
両方の型を許すfullなモジュライ問題の中で、これらは開かつ閉に分かれる。
滑らかな可換del Pezzo族はこの定義を満たす。
\end{corollary}
\begin{proof}
Serre周期により $\mathsf H^m\mathcal O\simeq E_{-rm}$ である。
$i=-rm\le j=0$ に有向ACM条件を適用すると
\[
 R\operatorname{Hom}(\mathcal O,\mathsf H^m\mathcal O)
 \simeq b_{-rm,0}[0]
\]
を得る。これは階数 $h_{n,d}(rm)$ の局所自由複体であるから、
命題\ref{prop:ncdp-intrinsic-degree}の仮定が満たされる。
各試験基底上で $d$ は局所定数なので、その値で分けた部分は開かつ閉である。
可換族については $E_i=\mathcal O_X(-iH)$ と
$\mathsf H=-\otimes\omega_{X/T}^{-1}$ を取り、
Kodaira消滅と滑らかな固有射のperfect・基底変更・Serre双対を用いればよい。
\end{proof}

ここで決めたのはSerre周期とdegreeの不変量である。
自己同値の最大可除指数や、各次数部分の全既約成分数を計算したものではない。

% !TeX root = ../nc-del-pezzo.tex
% Standalone section draft for papers/nc-del-pezzo.tex; 2026-11-20.
% Bibliography keys: Kontsevich, DeDekenLowen, LowenVdBDeformation.
\section{Poisson 零切断の形式近傍上の同変量子化}
\label{sec:equivariant-quantization}

この節では，重み付き標準形とその線型座標変換群を用い，
Poisson モジュライの零切断に沿う形式近傍上に
構造対象付きの Abel 圏の族を構成する．
これは量子化の一つの選択であり，全ての Abel 変形を分類する主張ではない．
Kontsevich のグラフによる formality 写像を一つ固定し，以下では同じ写像を全ての族に用いる．

次数 $d=1,2$ に応じて
\[
 S_R=\begin{cases}
 R[x_0,\ldots,x_{n-1},y,z],&\deg(x_i,y,z)=(1,2,3),\quad d=1,\\
 R[x_0,\ldots,x_n,z],&\deg(x_i,z)=(1,2),\quad d=2
 \end{cases}
\]
とし，標準式を $F=z^2-y^3-A_4(x)y-A_6(x)$ または
$F=z^2-f_4(x)$ とする．その次数を $m=6,4$ と書く．
Poisson パラメータは次数1では $D_i$，次数2では $p_{ij}$ であり，
後者には Pl\"ucker 関係を課す．相対構造定理が与える ambient の二重ベクトルを
$\widetilde\pi$ とする。具体的には、次数1では
\[
 D=\sum_iD_i\partial_{x_i},\qquad
 \widetilde\pi=F_zD\wedge\partial_y+D(F)\partial_y\wedge\partial_z
              +F_y\partial_z\wedge D,
\]
次数2では $P=\sum_{i<j}p_{ij}\partial_{x_i}\wedge\partial_{x_j}$ とおいて
\[
 V_i=\sum_jp_{ij}\partial_{x_j}f_4,\qquad
 \widetilde\pi=2zP+\Bigl(\sum_iV_i\partial_{x_i}\Bigr)\wedge\partial_z
\]
である。次数1のJacobi恒等式は三変数Jacobian括弧の恒等式の多項式的延長であり、
次数2の座標Jacobiatorは $J(x_i,x_j,z)=0$ と
$J(x_i,x_j,x_k)=2\sum_\ell q_{ijk\ell}\partial_{x_\ell}f_4$ となる。
従って任意の係数環上で
\[
 [\widetilde\pi,\widetilde\pi]=0,\qquad
 [\widetilde\pi,F]=0
\]
が成り立つ．いずれも超曲面で割る前の等式である．

\subsection{冪零係数上の有限な公式}

\begin{theorem}[任意係数の中心商量子化]
\label{thm:nilpotent-universal-quantization}
$R$ を可換 $\mathbb C$-代数，$J\subset R$ を $J^N=0$ を満たすイデアルとする．
標準形の係数と Poisson パラメータを $R$ 上に取り，
全ての Poisson パラメータが $J$ に属すると仮定する．
このとき $S_R$ と同じ次数付き $R$-加群上の結合積 $\star$ と，
次数 $m$ の中心元 $\Omega$ を次の性質を満たすように構成できる．
\begin{enumerate}
\item $a\star b\equiv ab\pmod J$，$\Omega\equiv F\pmod J$ であり，
      全ての内部次数を保つ．
\item $\Omega$ は両側非零因子である．
      $B_R=(S_R,\star)/(\Omega)$ の各次数は有限自由 $R$-加群であり，
      Hilbert 級数は可換 del Pezzo section ring と同じである．
\item この構成は $R$ の任意の係数変更と可換する．
      とくに係数変更の平坦性や被約性は要求しない．
\item 同じ式を $N$ によらず使うことにより，冪零次数を下げる還元とも両立する．
\end{enumerate}
\end{theorem}

\begin{proof}
Kontsevich の $L_\infty$ 写像を $\mathcal U$ とし，
その定数項の積を通常の積 $\mu$ とする．定義する式は
\[
 \mu_\pi=\mu+\sum_{r=1}^{N-1}\frac1{r!}
       \mathcal U_r(\widetilde\pi,\ldots,\widetilde\pi),
 \qquad
 \Omega=\sum_{r=0}^{N-1}\frac1{r!}
       \mathcal U_{r+1}(F,\widetilde\pi,\ldots,\widetilde\pi).
\]
係数環 $R$ の元を微分せず，ambient 変数だけを微分する．
$r$ 個の二重ベクトルを含む項の係数は $J^r$ に属するので，これらは有限和である．
各項は多項式係数の有限階多重微分作用素であり，任意の $\mathbb C$-代数上で定義される．
結合性は Maurer--Cartan 方程式を保つ性質から，中心性は twisted chain map と
$[\widetilde\pi,F]=0$ から従う
\cite[\S4.6.2，\S6，\S8]{Kontsevich}．
この中心性は Hochschild の零次 cocycle の等式であり，中心を coboundary で割っていない．
各グラフの添字の縮約において係数の重みと微分の重みが相殺するため，
積は内部次数零で，$\Omega$ は次数 $m$ である．
単位を保つ正規化を用いる．

平坦性には一般の冪零係数上の議論が必要である．
$F$ は $z$ に関して monic な二次式なので，各次数で
\[
 S_{R,\ell}=F S_{R,\ell-m}\oplus W_{R,\ell},
 \qquad
 W_R=\operatorname{span}_R\{\text{$z$ の指数が $0$ または $1$ の単項式}\}
\]
という有限自由加群の分解がある．
$F$ による除法の商を取る写像を
$r_\ell:S_{R,\ell}\to S_{R,\ell-m}$ と書く．
量子化後の乗算 $M_\ell(a)=\Omega\star a$ は
\[
 r_\ell M_\ell=1+K_\ell,
 \qquad K_\ell(S_{R,\ell-m})\subset J S_{R,\ell-m}
\]
を満たす．従って $K_\ell^N=0$ であり，$1+K_\ell$ は有限幾何級数で可逆である．
任意の $b\in S_{R,\ell}$ を
\[
 b=M_\ell(1+K_\ell)^{-1}r_\ell(b)
   +\bigl(b-M_\ell(1+K_\ell)^{-1}r_\ell(b)\bigr)
\]
と書けば，第二項は $\ker r_\ell=W_{R,\ell}$ に属する．
また $M_\ell(a)\in W_{R,\ell}$ なら $(1+K_\ell)a=0$ なので $a=0$ である．
従って
\[
 S_{R,\ell}=M_\ell(S_{R,\ell-m})\oplus W_{R,\ell}.
\]
これは $\Omega$ の正則性と $B_{R,\ell}\simeq W_{R,\ell}$ を同時に示す．
中心性により右側の正則性も従う．
グラフ式，monic 除法，有限幾何級数は全て係数変更と可換するので，
商の同定と各次数の自由性も任意の係数変更後に保たれる．
\end{proof}

\begin{remark}[中心的な線型座標]
複素数体上の固定した非零 Poisson 構造に適合する座標を取ると，
両次数とも $n-1$ 個の線型座標 $s_1,\ldots,s_{n-1}$ が中心的である．
残る活性な三変数の重みは，次数1で $(1,2,3)$，次数2で $(1,1,2)$ となる．
二重ベクトルに $\partial_{s_i}$ 成分がないため，グラフ式に
$s_i$ 方向の微分は現れない．従って $s_i$ は補正なしに中心的であり，
積は $R[s_1,\ldots,s_{n-1}]$-双線型である．
この座標選択は普遍構成の定義には不要である．
任意の冪零係数上の非零パラメータからこのような基底を選べるとは主張しない．
\end{remark}

\subsection{正規形の全座標変換との両立}

正規形を保つ群 $\widetilde G_d$ は ambient 変数への線型作用を持つ．
次数1では $g=(A,t)$ が $(x,y,z)\mapsto(Ax,t^2y,t^3z)$，
次数2では $g=(A,b)$ が $(x,z)\mapsto(Ax,bz)$ と作用する．
係数と Poisson パラメータを同時に変換すれば，ambient 二重ベクトルは pushforward され，
方程式の直線はそれぞれ指標 $t^6$，$b^2$ で変換される．

\begin{proposition}[厳密な同変性]
\label{prop:quantization-linear-equivariance}
定理\ref{thm:nilpotent-universal-quantization}の積と中心元の直線は
$\widetilde G_d$ に関して同変である．
従って中心商 $B_R$ には正規形の全座標変換が実際の次数付き代数同型として作用する．
\end{proposition}
\begin{proof}
Kontsevich のグラフ公式は線型座標変換に関して同変である．
これは全ての微分添字と係数添字が縮約されることから直接従い，
同論文\S7冒頭の性質 P3 としても明記される\cite[\S7，P3]{Kontsevich}．
実係数での恒等式は複素化後に成立し，各有限グラフの式は多項式恒等式なので，
群の行列成分を任意の $R$-係数に取った場合にも成立する．
同じ formality 写像の twisted 成分を中心元に用いているため，
$F$ の半不変性がそのまま $\Omega$ の半不変性になる．
従って中心イデアルは保たれる．
\end{proof}

\subsection{tails，係数変更，次数スカラーの消去}

$B=B_R$ とおく．有限個の標準変数の像が $B$ を生成する．
実際，各次数で生成写像の余核は $J$ で割ると零になり，
$J$ の冪零性から余核自身が零である．
生成元の最大重みを $w_{\max}$ とすると，右 ideal $B_{\ge q}$ は，
重みが初めて $q$ 以上となる prefix を持つ単語のうち，
その prefix 全体の有限集合で生成される．これらの重みは $q+w_{\max}-1$ 以下である．
従って各 $B_{\ge q}$ は有限生成右 ideal であり，tails torsion は局所化部分圏をなす．
これは $R$ の noetherian 性によらない．
$R$ が noetherian なら $B$ は左右 noetherian である．
これは有限の $J$-進 filtration の随伴次数環が可換 noetherian 環
$(S_{R/J}/(F\bmod J))\otimes_{R/J}\operatorname{gr}_J R$ であることから従う．
大きい圏を $\mathcal Q_R=\operatorname{QGr}B$，構造対象を
$\mathcal O_R=\pi B$ とする．

\begin{lemma}[tails の係数変更]
\label{lem:tails-arbitrary-base-change}
任意の可換 $R$-代数 $R'$ に対して，自然な同値
\[
 (\operatorname{QGr}B)_{R'}
 \simeq\operatorname{QGr}(B\otimes_RR')
\]
がある．左辺は $\operatorname{QGr}B$ 内の $R'$-作用付き対象の圏である．
この同値は構造対象と次数移動を保つ．
\end{lemma}
\begin{proof}
$\pi:\operatorname{Gr}B\to\operatorname{QGr}B$ は exact な Gabriel 商であり，
完全忠実な右随伴 $\omega$ を持つ．両関手に $R'$-作用を付けると，
再び exact な左随伴と完全忠実な右随伴が得られる．
前者の定義域は $\operatorname{Gr}(B\otimes_RR')$ である．
その核は，$B$-加群として tails torsion である $R'$-作用付き加群である．
これは $B\otimes_RR'$ に関する tails torsion と一致する．
実際，$R'$ は次数零で中心的なので，ある斉次元が $B_{\ge q}$ によって消える条件と，
$(B\otimes_RR')_{\ge q}$ によって消える条件は同じである．
商の普遍性から結論を得る．
\end{proof}

\begin{lemma}[平坦な Abel 変形]
\label{lem:tails-flat-deformation}
$R$ が noetherian なら，$\mathcal Q_R$ は $R$-flat な Grothendieck 圏であり，
$\mathcal O_R$ は $R$-flat な対象である．また
\[
 (\mathcal Q_R)_{R/J}\simeq\operatorname{QCoh}(X_{R/J}),
 \qquad \mathcal O_R\bmod J\simeq\mathcal O_{X_{R/J}}
\]
である．局所Artin基底上では，これは通常の平坦 Abel 変形を与える．
\end{lemma}
\begin{proof}
$B$ に付随する $\mathbb Z$-代数は各 Hom が $R$-flat なので，
その加群圏は $R$-flat である
\cite[arXiv v1, Proposition~3.7]{LowenVdBDeformation}．
Gabriel 商の平坦性もここでは直接確認できる．
$J_0$ が商圏の injective なら，$\pi$ の exact 性から $\omega J_0$ は加群圏の injective であり，
coflat である\cite[arXiv v1, Proposition~3.4]{LowenVdBDeformation}．
有限表示 $R$-加群 $M$ に対する内部 $\operatorname{Hom}_R(M,-)$ は
有限直和間の写像の kernel で計算される．$\pi$ はこの計算と可換するので，
\[
 \operatorname{Hom}_R(M,J_0)
 \simeq\pi\operatorname{Hom}_R(M,\omega J_0)
\]
となり，$J_0$ も coflat である．同じ判定から商圏は flat である．
さらに $B$ 自身が次数ごとに $R$-free であり，$\pi$ が exact なので，
$M\mapsto\pi(M\otimes_RB)$ は exact である．これは構造対象の平坦性を示す．
還元は補題\ref{lem:tails-arbitrary-base-change}と補題\ref{lem:weighted-free-locus}による重み付き商の自由性と相対 Serre 対応による．
Abel 変形としての tails の構成は
De Deken--Lowen \cite[Proposition~4.4]{DDL2011} とも一致する．
\end{proof}

\begin{lemma}[次数スカラーの整合的な恒等化]
\label{lem:grading-scalar-trivialization}
次数付き代数 $B$ の自己同型 $\alpha_\lambda(a)=\lambda^{\deg a}a$ に対し，
右加群の係数延長関手を
$F_\lambda(M)=M\otimes_{B,\alpha_\lambda}B$ とする．
写像
\[
 \eta_{\lambda,M}:F_\lambda(M)\longrightarrow M,
 \qquad m\otimes a\longmapsto\lambda^{\deg m}ma
\]
は次数付き加群の自然同型であり，$\eta_{\lambda\mu}$ は
$\eta_\lambda,\eta_\mu$ の合成と整合する．
この自然同型は任意の次数を保つ代数同型と両立し，tails に降下する．
$M=B$ 上では，構造対象を保つ標準同型に等しい．
\end{lemma}
\begin{proof}
斉次元 $b$ について
\[
 \lambda^{\deg m+\deg b}mba
 =\lambda^{\deg m}m\alpha_\lambda(b)a
\]
なので，写像は tensor 積の関係式を保つ．逆写像は
斉次元 $m$ を $\lambda^{-\deg m}m\otimes1$ に送る．
二つの係数延長を一つに同定する写像は
$m\otimes a\otimes b\mapsto m\otimes\alpha_\mu(a)b$ である．
両経路で得る係数は
$\lambda^{\deg m}\mu^{\deg m+\deg a}$ で一致する．
これが乗法的な coherence であり，三重以上の重なりにも同じ等式を使える．
残る両立性は次数の保存から従う．$M=B$ の場合は
$a\otimes b\mapsto\alpha_\lambda(a)b$ そのものである．
\end{proof}

\begin{theorem}[内在的 Poisson 形式近傍への降下]
\label{thm:quantization-descent}
$\widehat{\mathcal P}_{n,d,0}$ を，Poisson モジュライスタックの零切断に沿う形式完備化とする．
一つのグラフ formality 写像を固定すると，この形式スタック上に
$\mathcal O$-marked な tails の整合的な形式族を構成できる．
局所 noetherian 基底上で，その各冪零段階は平坦 Abel 族である．
この構成は任意の係数変更と両立し，零切断上で通常の
$(\operatorname{QCoh}(X),\mathcal O_X)$ を与える．
ここでは $\mathcal O$-marked 圏の射を，$\mathcal O$ を保つ完全線型同値の
自然同型類とし，必要な fppf stack 化を取る．
\end{theorem}
\begin{proof}
座標 atlas $U_d\times C_{n,d}$ 上で，Poisson パラメータのイデアルを $I$ とする．
$I^N=0$ の各近傍に定理\ref{thm:nilpotent-universal-quantization}を適用する．
得られた中心商とその tails は命題\ref{prop:quantization-linear-equivariance}により
$\widetilde G_d$-同変である．
基本線束のスカラーに当たる中央部分群 $Z_d\simeq\mathbb G_m$ は
内部次数 $\ell$ 上で $\lambda^\ell$ と作用する．
補題\ref{lem:grading-scalar-trivialization}は，この部分群の tails への作用を
構造対象を保ちながら整合的に恒等化する．
従って $G_d=\widetilde G_d/Z_d$ の torsor を局所的に持ち上げて得る圏は，
持上げの差 $Z_d$ に依存せず，三重の重なりでも整合する．
これが $G_d$ 商への降下データである．
大域的な $\widetilde G_d$-torsor の存在を仮定していない．
各 $N$ の構成とその降下データは還元と可換するので，形式族を得る．
平坦性・構造対象・係数変更は直前の補題による．
\end{proof}

\begin{remark}[一次項と結論の範囲]
平方零イデアル $J$ に Poisson パラメータが入る場合，積は
\[
 a\star b=ab+\tfrac12\widetilde\pi(\mathrm da,\mathrm db),
 \qquad a\star b-b\star a=\widetilde\pi(\mathrm da,\mathrm db)
\]
である．また $\Omega=F$ である．実際，一次補正
$\mathcal U_2(F,\widetilde\pi)$ に寄与し得るグラフは，第一種の頂点を二つ，
第二種の頂点を零個持つ．$F$ の頂点から出る辺はなく，
$\widetilde\pi$ の頂点からは二辺が出る．自己ループを認めないため，
二辺はともに $F$ の頂点へ向かわねばならず，多重辺を認めない admissible graph は存在しない．
従ってこの補正は零である．Hochschild--Kostant--Rosenberg 分解の二重ベクトル成分において，
この構成の一次写像は与えた Poisson パラメータの恒等写像である．
特に量子化操作は一次で余分な可換超曲面変形を加えない．
積と中心元の構成は選んだ formality 写像に依存する．
上の定理は全 Abel 変形との同値，量子化の一意性，有限の非零複素パラメータへの代数化，
全 nc モジュライの代数性や成分性，および内在的な逆量子化を主張しない．
\end{remark}

\par

% !TeX root = ../nc-del-pezzo.tex
\section{形式群oidと自己同型を含む局所モジュライ}
\label{sec:formal-groupoid}

本節では全Abel Artin変形を扱い，相対Serre条件を入力として用いない．
定義\ref{def:nc-dp-stack}のモジュライとの完成の一致は
後の定理\ref{thm:all-artin-are-nc-dp}で証明する．
試験環は引き続き剰余体 $\mathbb C$ の局所Artin環である．
$\mathfrak P_X(R)$ を，特殊ファイバーを $(X,0)$ と同一視した
Poisson変形の群oidとする．

\begin{lemma}[微小なAbel自己同値の消滅]
\label{lem:abel-infinitesimal-aut-rigidity}
$\mathcal C$ を $\mathbb C$-線型Grothendieck圏とし，
$\operatorname{HH}^1_{\mathrm{Ab}}(\mathcal C)=0$ と仮定する．
任意の平坦Abel Artin変形 $\mathcal D$ の自己同値で，
特殊ファイバー上で恒等に自然同型なものは，$\mathcal D$ 上でも恒等に自然同型である．
\end{lemma}
\begin{proof}
Lowen--Van den Bergh \cite[Definition~8.2, Theorem~8.17]{LowenVdBDeformation}
は，Abel変形と $\mathfrak i=\operatorname{Inj}(\mathcal C)$ の平坦線型変形の
群oid擬関手を同値にする．両者の射は自然同型類である．
同論文Appendix~BのPropositions~B.1, B.3により，線型変形は
対象集合を固定した模型で表せ，自己同値も自然同型を変更して対象上恒等にできる．
また
\[
 \operatorname{HH}^1(\mathfrak i)
 \simeq\operatorname{HH}^1_{\mathrm{Ab}}(\mathcal C)=0
\]
である\cite[Theorem~6.6]{LowenVdBHochschild}．

Artin環 $R$ の長さについて帰納する．
$0\ne I\subset R$，$\mathfrak m_RI=0$ とし，$R'=R/I$ と置く．
帰納法で自己同値 $F$ の $R'$ 上の還元は恒等に自然同型である．
各対象上のその同型を，Homの平坦性と還元の全射性で持ち上げる．
還元で可逆な射は冪零性により可逆なので，対象ごとの共役により
$F$ は $R'$ 上で厳密に恒等と仮定できる．
平坦線型変形を $\mathfrak a$ と書くと，$F-\operatorname{Id}$ は
\[
 \delta:\mathfrak i(A,B)\longrightarrow I\otimes_{\mathbb C}\mathfrak i(A,B)
\]
を与える．$\mathfrak m_RI=0$ と $I^2=0$ により，関手性は
$\delta(ba)=\delta(b)a+b\delta(a)$ と同値である．
従って $\delta$ は Hochschild $1$-cocycleである．
$I$ は有限次元なので $\operatorname{HH}^1(\mathfrak i,I\otimes\mathfrak i)=0$ であり，
ある $h_A\in I\otimes\mathfrak i(A,A)$ によって
\[
 \delta(a)=h_Ba-ah_A
\]
と書ける．$1-h_A$ は $F$ から恒等関手への自然同型を与える．
これは既存の自己同値の核を消す議論であり，任意の同値を
無条件に持ち上げられると主張しているのではない．
\end{proof}

\begin{theorem}[framed群oidにおける具体量子化の完全性]
\label{thm:quantization-framed-groupoid}
任意の $R$ に対して，$\mathfrak D_X^{\mathcal O}(R)$ の自己同型群は自明である．
$\mathfrak P_X(R)$ についても同じであり，具体量子化は
\[
 \mathfrak P_X\simeq\mathfrak D_X^{\mathcal O}
\]
というArtin基底上の群oid擬関手の同値を与える．
ここで射は全ての標識を保つ完全Abel同値の自然同型類である．
\end{theorem}
\begin{proof}
Hochschild比較とHKR分解
\cite[Theorem~6.6, Corollaries~7.7.2,7.8.2]{LowenVdBHochschild}
\cite[Corollary~0.6]{YekutieliHKR2002}から
\[
 \operatorname{HH}^1_{\mathrm{Ab}}(\operatorname{Qcoh}X)
 =H^0(X,T_X)\oplus H^1(X,\mathcal O_X)=0.
\]
第二項は補題\ref{lem:directed-line-vanishing}により零であり，
第一項は入力となる可換del Pezzo族の自己同型剛性による．
前補題と命題\ref{prop:directed-section-deformation}の標識忘却を用いると，
$\mathfrak D_X^{\mathcal O}(R)$ の自己同型は自明となる．
実際の自然自己同型についても，持ち上げた豊富列の各対象上で
$\operatorname{End}(E_{i,R})=R$ であり，隣接Homの正階数自由性から
全て同じスカラーとなる．$\mathcal O_R$ の標識を保てばそのスカラーは1である．

幾何的なPoisson自己同型が小拡大の商上恒等なら，その差は
$H^0(X,T_X)\otimes I$ に属する．これが零であることから帰納し，
$\mathfrak P_X(R)$ の自己同型も自明である．
従って両群oidからその同型類を離散群oidと見たものへの射は同値である．
定理\ref{thm:concrete-quantization-all-artin}の自然な全単射と合わせれば主張を得る．
量子化の対象の選択を変えても，同じ同型類の二対象間の射は一意だから，
必要な比較同型と基底変更の整合性は一意に定まる．
\end{proof}

\subsection{Artin近傍での通常の次数付き表示の復元}

\begin{proposition}[Artin局所での全構造対象付き同値]
\label{prop:artin-full-graded-recognition}
$B,B'$ を局所Artin環 $R$ 上の次数付き代数とし，
$(\mathcal D,\mathcal O)=(\operatorname{QGr}B,\pi B)$，
$(\mathcal D',\mathcal O')=(\operatorname{QGr}B',\pi B')$ が
$\operatorname{Qcoh}X$ の平坦Abel変形を与え，$\mathcal O,\mathcal O'$ は $R$-flatとする．
次数シフトを $s,s'$ と書き，両者は特殊ファイバーで
$-\otimes\mathcal O_X(H)$ に一致するとする．
各 $(\mathcal O,s)$，$(\mathcal O',s')$ は豊富対であり，section写像が
\[
 B_m\xrightarrow{\sim}\operatorname{Hom}_{\mathcal D}(\mathcal O,s^m\mathcal O),
 \qquad
 B'_m\xrightarrow{\sim}\operatorname{Hom}_{\mathcal D'}(\mathcal O',(s')^m\mathcal O')
 \quad(m\ge0)
\]
という次数付き環の同型を与えると仮定する．
このとき自然な全単射
\[
 \operatorname{Eq}^{\mathcal O}_R(\mathcal D,\mathcal D')
 \simeq\operatorname{Isom}_{\mathrm{gr},R}(B,B')/R^\times
\]
がある．右辺では $\lambda\in R^\times$ が次数 $m$ 上で $\lambda^m$ を掛ける．
上の仮定が保たれるArtin基底変更に関して自然である．
\end{proposition}
\begin{proof}
全 $\mathcal O$ 保存同値 $F:\mathcal D\simeq\mathcal D'$ を取る．
その特殊ファイバーは $X$ の自己同型を与え，基本線束 $H$ を保つ．
従って
\[
 G=(s')^{-1}\circ F\circ s\circ F^{-1}
\]
は特殊ファイバー上で恒等となる．
補題\ref{lem:abel-infinitesimal-aut-rigidity}は構造対象の保存を要求しないので，
$G\simeq\operatorname{Id}$ となり，$\theta:Fs\xrightarrow{\sim}s'F$ が存在する．
$\theta$ を反復して $\theta_m:Fs^m\simeq(s')^mF$ を定めると，
これは整数の加法に関して整合する．負の整数には逆同型を用いる．

標識 $\alpha:F\mathcal O\simeq\mathcal O'$ により，
$a:\mathcal O\to s^m\mathcal O$ を
\[
 \mathcal O'\xrightarrow{\alpha^{-1}}F\mathcal O
 \xrightarrow{F(a)}Fs^m\mathcal O
 \xrightarrow{\theta_m}(s')^mF\mathcal O
 \xrightarrow{(s')^m\alpha}(s')^m\mathcal O'
\]
へ送る．$\theta_m$ の整合性によりこれはsection環の次数付き同型である．
各対象のEndと隣接Homを用いた
補題\ref{lem:marking-removes-two-aut}の証明から，
$\operatorname{Nat}(\operatorname{Id}_{\mathcal D'},
\operatorname{Id}_{\mathcal D'})=R$ である．
従って二つの $\theta$ の差は一つの $\lambda\in R^\times$ であり，
$\theta_m$ の差は $\lambda^m$ となる．
標識同型 $\alpha$ のスカラー変更は始域と終域で相殺される．
よって右辺の元は $F$ の標識を保つ自然同型類だけで決まる．

逆に次数付き環同型は構造対象を保ち，次数シフトと可換するtails同値を与える．
豊富対からのsection復元により，この二つの構成は互いに逆である．
具体的には全ての $s^{-m}\mathcal O$ とそれらの間の射上で同じ作用を持ち，
豊富列からの表示とfiltered colimitにより全圏で自然同型となる．
選択の差を次数スカラーで除いているので，得られた全単射は基底変更に自然である．
\end{proof}

この命題は，相対Serre公式やVeroneseの普遍中心条件を用いない．
特に後者が破れる可換点のArtin近傍でも使える．
一方，有限の非零パラメータ上の一般nc点では
特殊ファイバーの $\operatorname{HH}^1=0$ による帰納が使えないので，
本命題をそのまま大域の全同値認識へ延長してはならない．

\subsection{有限自己同型群を含む形式商}
$\Gamma_X=\operatorname{Aut}(X)$ と置く．これは有限群である．
$\widehat{\mathfrak D}^{\mathcal O}_{[X]}(R)$ は，特殊ファイバーが
$(\operatorname{Qcoh}X,\mathcal O_X)$ に同値な平坦Abel変形を対象とし，
特殊ファイバーとの同一視自体は選ばない群oidとする．
$\mathcal O$ の標識は保持する．

\begin{theorem}[自己同型を含む形式局所商]
\label{thm:unframed-formal-quotient}
$\Gamma_X$ は枠の変更により $R_X$ に連続的に作用し，
Artin試験環上で
\[
 \widehat{\mathfrak D}^{\mathcal O}_{[X]}
 \simeq[\operatorname{Spf}R_X/\Gamma_X]
\]
となる．この作用は具体量子化による同値を通じて，
Poisson変形の特殊ファイバーの枠を変える作用と一致する．
\end{theorem}
\begin{proof}
可換点での全 $\mathcal O$ 保存完全Abel自己同値の群は
$\operatorname{Aut}(X)$ である．完全同値をcoherent部分へ制限し，反標準豊富列からの復元
（命題\ref{prop:ncdp-anticanonical-hom}）を可換点に適用して得る．
実際，反標準軌道環は通常の反標準切断環であり，そのProjは $X$ である．
$X$ の自己同型は標準的に反標準線束を保ち，
切断環の次数スカラーは構造対象付き圏の自然同型で消える．

枠を選べば対象は $\mathfrak D_X^{\mathcal O}(R)$ に入る．
枠の二つの選択の差は $\Gamma_X$ の元であり，
その作用は全Artin基底変更と両立するのでYonedaにより $R_X$ 上の作用を定める．
二対象間で特殊ファイバー上の $\gamma\in\Gamma_X$ を誘導する同値は，
枠を $\gamma$ で変更した後のframed同値に等しい．
前定理によりその射は，対応する同型類が一致するときに限り存在し，存在すれば一意である．
従って群oidは同型類集合と $\Gamma_X$ の作用群oidに等しい．
局所Artin $\mathbb C$-基底上の有限定数群torsorは自明なので，
これは表示した形式商のArtin点の群oidである．

標準形の座標変換は全ての $X$ の自己同型を表し，
定理\ref{thm:quantization-descent}の同変性と次数スカラーの整合的恒等化により，
具体量子化はこの枠の変更に両立する．
\end{proof}

\begin{proposition}[同変な形式積表示]
\label{prop:equivariant-formal-product}
$N_X=H^1(X,T_X)$ とし，$C_{n,d}$ を第\ref{sec:definitions}節のPoisson錐とする．
$R_X$ の積表示は，非標準的だが $\Gamma_X$ 同変に
\[
 R_X\simeq
 \widehat{\mathcal O}_{N_X\times C_{n,d},(0,0)}
\]
と選べる．$N_X$ 上の作用は通常の変形の接表現であり，
錐上の作用は第\ref{sec:definitions}節に書いた標準形の作用の制限である．
\end{proposition}
\begin{proof}
標準形の係数空間 $U_d$ はある線型表現 $V_d$ の開部分であり，
$X$ を表す点を $f\in U_d$ とする．$G_d$ の $f$ におけるstabilizerは
$\Gamma_X$ である．$H^0(X,T_X)=0$ により，軌道接写像
$\operatorname{Lie}(G_d)\to T_fU_d=V_d$ は単射である．
有限群の平均化により，その像に対して $\Gamma_X$ 不変な補空間 $N$ を選ぶ．
$N\simeq H^1(X,T_X)$ は接空間として同変である．
$S=(f+N)\cap U_d$ と置くと，$G_d\times S\to U_d$ の微分は
$(1,f)$ で同型であり，この写像はそこで\'{e}taleである．
形式的逆写像により，恒等に還元する座標変換で各Artin変形を
$S$ 上に一意に正規化できる．

Poissonパラメータも同時にその座標変換で移す．
$G_d\times S\times C_{n,d}\to U_d\times C_{n,d}$ は，
$G_d\times S$ 上で $c\mapsto gc$ と変数変換した後には
上の\'{e}tale写像と $C_{n,d}$ の恒等との積である．
従って錐が特異でも同じ形式的正規化が働く．
$S$ の選択は $\Gamma_X$ 不変であり，残る作用は
$N\times C_{n,d}$ 上の線型な対角作用なので主張を得る．
\end{proof}

\begin{corollary}[形式粗商の正規性と一分枝性]
\label{cor:formal-coarse-normal}
形式局所商の不変式環 $R_X^{\Gamma_X}$ は完備Noether局所環で，
正規かつCohen--Macaulayな整域である．その次元は
$d=1$ で $b_{n,1}+n$，$d=2$ で $b_{n,2}+2n-1$ である．
従って全自己同型を考慮した形式粗商も一分枝である．
\end{corollary}
\begin{proof}
前命題の表示は有限型代数の固定点での完備化から得られる．
有限群の不変式環は有限型で，元の環はその上有限である．
平均化により不変式と完備化は可換するから，
$R_X^{\Gamma_X}$ は完備Noether局所環で，$R_X$ はその上有限である．
整域性と正規性は不変式部分環へ降りる．
不変式環のパラメータ系は有限拡大 $R_X$ のパラメータ系にもなるので，
$R_X$ のCohen--Macaulay性からそこで正則列である．
平均化は不変な各パラメータによる商と可換するため，
同じ列は不変式環でも正則列となる．次元は有限整拡大で変わらない．
\end{proof}

\begin{corollary}[被覆対合だけを持つ点の形式粗商]
\label{cor:generic-nc-formal-sign-quotient}
$\operatorname{Aut}(X)=\langle\tau\rangle\simeq\mu_2$，
$\tau(z)=-z$ とする．この条件を満たす $X$ は両次数・全 $n\ge2$ で稠密開をなす
\cite[命題9.4]{PoissonDP}．このとき形式粗商は
\[
 \operatorname{Spf}\Bigl(
 \mathbb C[[t_1,\ldots,t_{b_{n,d}}]]
 \widehat\otimes\widehat{\mathbb C[C_{n,d}]^{(2)}}\Bigr)
\]
である．右の完備化は頂点で取り，$(2)$ はPoissonパラメータ次数の第二Veroneseを表す．
その接空間次元は，次数1では
\[
 b_{n,1}+\binom{n+1}{2},
\]
次数2では
\[
 b_{n,2}+\frac{n(n+1)^2(n+2)}{12}
\]
となる．特に両次数とも，全 $n\ge2$ でこの形式粗商は特異である．
\end{corollary}
\begin{proof}
被覆対合は係数空間上で恒等，Poissonパラメータ上で $-1$ に作用する
\cite[注意8.5]{PoissonDP}．従って前命題の $N_X$ 上では恒等であり，
不変式はPoissonパラメータの偶数次数部分に等しい．
元の錐の座標環は次数1で生成されるので，その第二Veroneseは
元の次数2部分で生成され，これが頂点の余接空間となる．
次数1ではその次元は $\dim\operatorname{Sym}^2\mathbb C^n=\binom{n+1}{2}$．
次数2では $v=n+1$ と置き，$\binom v2$ 個のPl\"ucker変数の二次式から
独立な $\binom v4$ 個のPl\"ucker関係を引けば
\[
 \binom{\binom v2+1}{2}-\binom v4
 =\frac{v^2(v^2-1)}{12}
 =\frac{n(n+1)^2(n+2)}{12}
\]
を得る．これらはそれぞれ錐の次元 $n$，$2n-1$ より大きい．
\end{proof}

\begin{remark}[粗商のArtin点と変形の同型類]
形式粗商のArtin点は，変形の同型類集合そのものではない．
例えば $R=\mathbb C[\epsilon]/(\epsilon^2)$ 上で
Poissonパラメータ $\epsilon Q\ne0$ と零は異なる軌道だが，
符号不変式は全て偶数次なので，ともに粗商の零点へ写る．
群oidの商表示を不変式環への置換だけで扱うと，この微小情報を失う．
本節の形式商は全Abel変形の形式モジュライである．
$\mathfrak M^{\mathcal O}_{n,d}$ の完成との同定は
定理\ref{thm:all-artin-are-nc-dp}を要し，大域の代数性・全既約成分数は残問である．
\end{remark}

% !TeX root = ../nc-del-pezzo.tex
% 2026-11-20 JST / v0.1
\section{冪零変形における重み付き多項式環の相対的性質}
\label{subsec:nilpotent-weighted-as}

ここでは通常の体上のAS正則性と，基底環に関する相対的な主張を区別する．
以下の証明は基底環の局所性，Noether性，標数零を必要としない．
内部次数の移動は $M(a)_q=M_{q+a}$ とし，
$\underline{\operatorname{Ext}}$ は内部次数を集めたgraded Extを表す．

\begin{lemma}[有限自由分解の冪零持上げ]
\label{lem:nilpotent-resolution-lift}
$R$ を可換環，$J\subset R$ を冪零イデアルとする．
$A$ を非負次数付き中心的 $R$-代数とし，各 $A_q$ は
$R$-射影加群であるとする．$M$ は各内部次数で $R$-射影な
次数付き左 $A$-加群とする．
$\overline A=A/JA$ と $\overline M=M/JM$ と置く．
$\overline M$ が，シフトを指定した有限自由分解
\[
 0\longrightarrow\overline F_N\longrightarrow\cdots
 \longrightarrow\overline F_0\longrightarrow\overline M
 \longrightarrow0,\qquad
 \overline F_p=\bigoplus_\alpha\overline A(-a_{p,\alpha})
\]
を持つなら，これを還元とする有限自由分解
\[
 0\longrightarrow F_N\longrightarrow\cdots
 \longrightarrow F_0\longrightarrow M\longrightarrow0,
 \qquad F_p=\bigoplus_\alpha A(-a_{p,\alpha})
\]
が存在する．その各短完全列は内部次数ごとに $R$-線形に分裂する．
したがって任意の環準同型 $R\to T$ による基底変換後も分解である．
右加群に対しても同じ結論が成り立つ．
\end{lemma}

\begin{proof}
$\overline F_0\to\overline M$ の基底の像を $M$ に持ち上げて
$F_0\to M$ を作る．余核 $C$ は $C=JC$ を満たすので，
$J$ の冪零性から $C=0$ である．$M$ の各次数が $R$-射影だから
$0\to Z_0\to F_0\to M\to0$ は各次数で分裂し，
$Z_0$ の各次数も $R$-射影である．特に還元後の核は
$Z_0/JZ_0$ であり，指定した分解の第零syzygyに一致する．
その生成元を持ち上げて $F_1\to Z_0$ を作ると，同じ議論で
これは全射であり，核 $Z_1$ は各次数で $R$-射影である．
この操作を反復する．最後の核 $Z_N$ は $Z_N/JZ_N=0$ を満たすため
$Z_N=0$ である．得られた全ての短完全列が $R$-分裂するので，
任意の基底変換により完全性が保たれる．
\end{proof}

\begin{lemma}[有界射影複体の完全性の検出]
\label{lem:nilpotent-complex-exactness}
$J\subset R$ が冪零で，$C^\bullet$ が有界な $R$-射影加群の複体なら，
$C^\bullet/JC^\bullet$ が完全であることから
$C^\bullet$ が $R$-線形に可縮であることが従う．
内部次数ごとに射影である次数付き複体にも，同じ主張が成り立つ．
\end{lemma}

\begin{proof}
有界で完全な射影複体 $C^\bullet/JC^\bullet$ は可縮である．
その縮約 $\overline h$ を射影性により $R$-線形写像 $h$ に持ち上げる．
$u=dh+hd$ と置くと $u\equiv1\pmod J$ である．
$u-1$ は各項を $J$ 倍の中へ写すので冪零であり，$u$ は可逆である．
$u$ は $d$ と可換だから，$h u^{-1}$ が $C^\bullet$ の縮約となる．
次数付きの場合はこの操作を各内部次数で行う．
\end{proof}

\begin{proposition}[重み付き相対正則性とaugmentation Ext]
\label{prop:nilpotent-weighted-as}
$R$ を可換環，$J\subset R$ を冪零イデアル，
$A=\bigoplus_{q\geq0}A_q$ を $A_0=R$ を満たす次数付き中心的
$R$-代数とする．各 $A_q$ は有限 $R$-自由であるとし，正整数
$w_1,\ldots,w_N$ と次数付き代数の同型
\[
 A/JA\simeq (R/J)[x_1,\ldots,x_N],\qquad \deg x_i=w_i
\]
を固定する．$W=\sum_iw_i$ と置く．次が成り立つ．
\begin{enumerate}
\item augmentation加群 $R=A/A_{>0}$ は，左右とも
\[
 K_p=\bigoplus_{\substack{I\subset\{1,\ldots,N\}\\|I|=p}}
 A\Bigl(-\sum_{i\in I}w_i\Bigr),\qquad0\leq p\leq N,
\]
を項とする有限自由分解を持つ．還元後の分解は重み付きKoszul分解である．
\item $A^{\mathrm e}=A\otimes_R A^{\mathrm{op}}$ に関して，対角加群 $A$ は
同じシフトと階数を持つ長さ $N$ の有限自由 $A^{\mathrm e}$-分解を持つ．
特に $A$ は $R$ 上homologically smoothである．
\item augmentationによる $A$-作用を右辺に入れると，左右とも
\[
 \underline{\operatorname{Ext}}_A^i(R,A)=
 \begin{cases}R(W)&i=N,\\0&i\neq N\end{cases}
\]
である．
\item これらの分解，完全性，Extの計算は任意の基底変換 $R\to T$ と両立する．
\end{enumerate}
\end{proposition}

\begin{proof}
多項式環上で $x_1,\ldots,x_N$ は基底環によらず正則列である．
そのKoszul分解を，補題\ref{lem:nilpotent-resolution-lift}で
augmentation $A\to R$ の分解に持ち上げると第一の主張を得る．
同じ補題を $A^{\mathrm e}$ とその左加群 $A$ に適用する．
還元後は $(R/J)[x_1,\ldots,x_N,y_1,\ldots,y_N]$ 上の対角であり，
$x_i-y_i$ のKoszul分解を持つ．$A$ は各次数で $R$-自由だから補題の仮定を満たし，
第二の主張を得る．

第一の分解を $\underline{\operatorname{Hom}}_A(-,A)$ で双対化する．
最上位の項は $A(W)$ である．最終微分
$K_N=A(-W)\to K_{N-1}$ の行列成分の次数はそれぞれ $w_i>0$ なので，
双対複体の最後にaugmentationを付けて複体
\[
 0\longrightarrow\underline{\operatorname{Hom}}_A(K_0,A)
 \longrightarrow\cdots\longrightarrow A(W)
 \longrightarrow R(W)\longrightarrow0
\]
が得られる．この複体の還元は，完全な双対Koszul複体に
最上位の商 $(R/J)(W)$ を付けたものである．各項は内部次数ごとに
$R$-射影だから，補題\ref{lem:nilpotent-complex-exactness}により元の複体も
$R$-線形に可縮である．これが第三の主張を与える．
反対代数に同じ議論を適用すれば右加群版も得られる．

持ち上げた分解の短完全列と，今得た双対複体はいずれも $R$-線形に分裂する．
したがって任意の基底変換で完全性を保つ．分解の項は有限自由 $A$-加群なので，
その双対化も基底変換と両立する．
\end{proof}

\begin{remark}
この命題は，$R$ 自身が特異である場合も適用できる．
したがって $A$ の通常の大域次元や通常の入射次元が有限であるとは主張していない．
ここで証明した内容は，$R$ に関する対角の完全性とaugmentation Ext条件である．
また，$R$ がNoether環なら $A$ は左右Noether環である．実際，
有限な $J$-進フィルターの随伴次数環は
$\operatorname{gr}_J(R)[x_1,\ldots,x_N]$ に同型であり，Noether性を持つ．
非Noether基底についてNoether性は主張しない．
\end{remark}

\begin{proposition}[平坦な中心超曲面商]
\label{prop:flat-central-hypersurface-ext}
命題\ref{prop:nilpotent-weighted-as}の仮定の下で，
$\Omega\in A_m$ は $m>0$ の中心的な左右正則元であるとし，
$B=A/(\Omega)$ は $R$-平坦であるとする．左右とも
\[
 \underline{\operatorname{Ext}}_B^i(R,B)=
 \begin{cases}R(W-m)&i=N-1,\\0&i\neq N-1\end{cases}
\]
が成り立つ．この計算は任意の基底変換と両立する．

さらに $\overline\Omega$ が，固定した多項式表示においてある変数に関して
monicな多項式であれば，$\Omega$ の左右正則性と $B$ の各次数の
有限 $R$-自由性は自動的に成り立つ．
\end{proposition}

\begin{proof}
正則中心元の自由分解
\[
 0\longrightarrow A(-m)\xrightarrow{\,\Omega\,}A
 \longrightarrow B\longrightarrow0
\]
から，$\underline{\operatorname{Ext}}_A^q(B,A)$ は $q=1$ のときだけ非零で，
その値は $B(m)$ である．従ってchange-of-ringsの随伴関係
\[
 \operatorname{RHom}_A(R,A)
 \simeq\operatorname{RHom}_B
 \bigl(R,\operatorname{RHom}_A(B,A)\bigr)
\]
は
\[
 \underline{\operatorname{Ext}}_B^i(R,B)
 \simeq\underline{\operatorname{Ext}}_A^{i+1}(R,A)(-m)
\]
を与える．前命題を適用すれば表示した公式が得られる．
$B$ が $R$-平坦なので，任意の基底変換後も上の短完全列は完全であり，
$\Omega$ は左右正則である．随伴関係と前命題の有限自由分解は基底変換に自然なので，
Extの計算も基底変換と両立する．

最後の主張を示す．$\overline\Omega$ がある変数に関してmonicなら，
$(R/J)[x_1,\ldots,x_N]/(\overline\Omega)$ の各内部次数は有限自由であり，
その単項式基底を選べる．次数 $q$ でこの基底を $A_q$ に持ち上げると，写像
\[
 A_{q-m}\oplus R^{h_q}\longrightarrow A_q,
 \qquad(a,c)\longmapsto\Omega a+\widetilde c
\]
は $J$ による還元後に同型となる．これは射影加群間の写像であるから，
還元後の逆写像を持ち上げ，冪零誤差を有限幾何級数で除くことにより同型と分かる．
従って $\Omega:A_{q-m}\to A_q$ は分裂単射で，その余核は $R^{h_q}$ である．
全次数を合わせれば $\Omega$ の正則性と $B$ の各次数の有限自由性が従う．
\end{proof}

\begin{remark}[中心正則性と普遍的正則性の区別]
中心正則性だけでは任意の基底変換に関する主張は得られない．
例えば $R=k[t]$, $A=R[x]$, $\Omega=tx$ は中心正則であるが，
$R\to R/(t)$ の後には $\Omega$ は零となる．上の命題では商の平坦性が
この問題を排除している．本研究の重み付き超曲面は，次数 $1$ では
重み $3$ の変数，次数 $2$ では重み $2$ の変数についてmonicな標準形を持つため，
この条件を満たす．また，それぞれ
\[
 (N,W,m)=(n+2,n+5,6),\qquad(n+2,n+3,4)
\]
であるから，中心商の唯一の非零augmentation Extは
$\underline{\operatorname{Ext}}_B^{n+1}(R,B)=R(n-1)$ となる．
\end{remark}

% !TeX root = ../nc-del-pezzo.tex
\section{Artin族の局所化，豊富性と相対perfect圏}
\label{sec:artin-ample-perfect}

本節では $R$ を剰余体 $k=\mathbb C$ を持つ局所Artin環，
$\mathfrak m^N=0$ とする．$B$ は定理\ref{thm:nilpotent-universal-quantization}の
中心超曲面商とし，$B\otimes_R k=B_0$ は固定した滑らかな $X$ の切断環である．
$T$ は任意の可換 $R$-代数，$B_T=B\otimes_RT$，
$J=\mathfrak mT$，$T_0=T/J$ と置く．
noetherian性とcoherent対象を述べる箇所では $T$ をnoetherianとする．
$\mathcal Q_T=\operatorname{QGr}B_T$，$E_{i,T}=\pi B_T(-i)$ と書く．

\subsection{有限Ore被覆と導来切断}
\begin{lemma}[斉次Ore局所化と重なり]
\label{lem:artin-ore-cover}
$B$ の正次数斉次元 $f$ で $B_0$ における像が非零なものを固定する．
その $B_T$ への像の冪は
左右の正則Ore集合をなす．局所化は任意の係数変更と可換する．
有限個のそのような元を同時に逆転する局所化は，その積を逆転する局所化と
標準的に一致し，積の順序によらない．
次数2では $f_i=x_i$，次数1ではこれらと $f_*=y\star z$ を取ると，
これらの局所化はtailsを被覆し，各局所環は強次数付きである．
\end{lemma}
\begin{proof}
$[B_T,B_T]\subset JB_T$ なので $\operatorname{ad}_f^N=0$ である．
$m\ge N-1$ に対して
\[
 b f^m=\sum_{a=0}^{N-1}(-1)^a\binom ma f^{m-a}\operatorname{ad}_f^a(b)
\]
となる．$m\ge q+N-1$ とすれば右辺は左から $f^q$ で割れ，
反対側の公式も成立するから左右のOre条件を得る．
$B_T$ は $T$-flatであり，その $J$-進随伴次数は
$B_0\otimes_k\operatorname{gr}_J T$ である．$B_0$ は整域なので，
$\bar f$ の乗法は各随伴次数で単射であり，$f$ は左右正則である．
局所化は $B_T$ の平坦な左右加群であり，内部次数ごとの $T$-加群も
局所化の直極限表示により平坦である．普遍性から任意の係数変更との両立を得る．

$g=f_1\cdots f_s$ を逆転すると，$J$ で割った可換環では各 $f_i$ が可逆である．
$J$ は局所化後も冪零なので，各 $f_i$ は局所化環自体で可逆である．
逆に各 $f_i$ を逆転すれば $g$ は可逆である．これが重なりの標準的同定を与える．
さらに乗法による $B_f\otimes_B B_g\to B_{fg}$ は同型である．
両辺は $R$-flatであり，還元は可換局所化の標準同型であるから，
有限の $\mathfrak m$-進filtrationにより元の写像も同型となる．
この同型を $T$ へ移すことで，全ての重なりにおける
局所化関手の合成の可換性も得る．

選んだ $f_i$ の可換還元の零点集合はaffine coneの頂点だけである．
次数1では $x=0$，$yz=0$，$z^2=y^3$ から $y=z=0$ が従う．
よって全ての局所化で零となる加群の $J$-進随伴次数はtails torsionであり，
有限のfiltrationにより元の加群もtails torsionである．この議論は加群の $T$-flat性を使わない．
逆向きは，tails torsionの各元が十分高い $f_i$ の冪で消えることによる．
$x_i$ の開では $x_i$ が次数1の可逆元である．$yz$ の開では $y,z$ がともに可逆であり，
$z\star y^{-1}$ が次数1の可逆元である．従って各局所化は強次数付きである．
\end{proof}

$B_{T,I}$ を添字集合 $I$ の元を同時に逆転した環とし，
次数付き加群または複体 $M$ に対して，非増大添字による有限複体
\[
 \check C_T(M)=\left[
 \bigoplus_i M\otimes_{B_T}B_{T,i}\longrightarrow
 \bigoplus_{i<j}M\otimes_{B_T}B_{T,ij}\longrightarrow\cdots\right]
\]
を次数0から置く．微分は制限写像の交代和である．複体 $M$ に対しては
有限方向の全複体を取る．重なりの同定は局所化の普遍性により整合する．

\begin{proposition}[有限\v{C}ech模型]
\label{prop:artin-cech-model}
$\pi:\operatorname{Gr}B_T\to\mathcal Q_T$ の右随伴を $\omega$ とすると，
任意の $M\in D(\operatorname{Gr}B_T)$ について
\[
 R\omega(\pi M)\simeq\check C_T(M).
\]
特に任意の整数 $i,j$ に対して
\[
 R\operatorname{Hom}_{\mathcal Q_T}(E_{j,T},E_{i,T})
 \simeq\check C_T(B_T)_{j-i}.
\]
これらの公式は任意の係数変更に対して導来的に両立する．
\end{proposition}
\begin{proof}
全ての局所化はexactである．増大複体 $M\to\check C_T(M)$ を一つの $f_i$ で局所化すると，
添字 $i$ を挿入する通常の縮約を持つ．よってそのconeのコホモロジーは
全ての局所化で零になり，前補題によりtails torsionである．

一方，局所化環上の任意の複体 $L$ を $B_T$-複体と見たものは，
tails torsionのコホモロジーを持つ複体全体に右直交する．実際，exactなOre局所化と
その制限関手の導来随伴により，そのような複体から $L$ への全ての導来Homは零である．
従って有限全複体 $\check C_T(M)$ も同じ右直交圏に入る．
これが $M$ のtails torsionに関する導来局所化を特徴付ける．
exactなGabriel商の導来右随伴は完全忠実であり，
$\ker D\pi$ はまさにtails torsionのコホモロジーを持つ複体であるから，最初の同型を得る．

$B_T(-j)$ は次数付き加群圏の射影生成元の一つである．導来随伴を適用して二番目の公式を得る．
全ての局所化環は $T$-flatであり，局所化の普遍性と有限全複体の構成は
係数変更と可換する．一般の複体にはK-flatな係数変更を用いれば同じ公式が成立する．
\end{proof}

\subsection{飽和性，ACM条件，正の周期の豊富性}
\begin{theorem}[普遍的な有向ACMとsection復元]
\label{thm:artin-saturation-acm}
全ての $i\le j$ と任意の $R\to T$ について，標準写像は準同型
\[
 (B_T)_{j-i}[0]\xrightarrow{\ \sim\ }
 R\operatorname{Hom}_{\mathcal Q_T}(E_{j,T},E_{i,T})
\]
である．従って次数付き切断環の標準写像
$B_T\to\bigoplus_{q\ge0}\operatorname{Hom}(\mathcal O_T,\mathcal O_T(q))$
は積と両立する同型であり，全てのHomと消滅は任意の係数変更で保たれる．
\end{theorem}
\begin{proof}
$q=j-i\ge0$ と置く．$B_q\to\check C_R(B)_q$ のconeは有界な $R$-flat複体である．
$k$ で割ると通常のaffine被覆による $R\Gamma(X,\mathcal O_X(qH))$ の計算になる．
可換切断環の飽和性とKodaira消滅により，還元後のconeは非輪状である．
$\mathfrak m$-進filtrationの各随伴次数は，この非輪状な $k$-複体と
$\mathfrak m^a/\mathfrak m^{a+1}$ のテンソル積なので，cone自体も非輪状である．
さらに任意の $T$ 上で同じconeは有界な $T$-flat複体であり，
$J$ で割ったものは可換模型の非輪状複体と $T_0$ のテンソル積となる．
同じ有限filtrationの議論が適用できる．
標準写像は斉次元の乗法によるものであるから，積との両立も保たれる．
\end{proof}

\begin{theorem}[任意の正周期の相対豊富性]
\label{thm:artin-positive-shift-ample}
$T$ がnoetherianなら，任意の整数 $r>0$ に対して
$(\mathcal O_T,(r))$ は相対豊富対である．
これは任意のnoetherian係数変更後にも成立する．
特に $r=n-1$ に適用できる．
\end{theorem}
\begin{proof}
$B_T$ は左右noetherianであり，$\mathcal Q_T$ は局所noetherianである．
coherent対象 $M$ の有限 $J$-進filtrationの各商は，
$X_{T_0}$ 上のcoherent層である．前命題の\v{C}ech模型は任意の加群に対してexactな局所化から作られるので，
このfiltrationを保つ．$\mathcal O_X(rH)$ はampleであるから相対Serre消滅により，
固定した $M$ に対して十分大きい $a$ で
\[
 \operatorname{Ext}^q_{\mathcal Q_T}(\mathcal O_T(-ar),M)=0
 \quad(q>0)
\]
となる．各filtration商からの帰納によりHomは有限 $T$-加群である．
全射 $M\to M'$ のcoherentな核にこの消滅を適用すれば，十分大きい $a$ でHomの全射性を得る．

さらに $M/JM$ は十分大きい $a$ で有限個の $\mathcal O_{X_{T_0}}(-arH)$ により生成される．
$JM$ への一次Extの消滅を使い，これらの生成写像を
$\mathcal O_T(-ar)$ の有限直和から $M$ へ持ち上げる．
余核 $C$ は $C/JC=0$ を満たすから $J$ の冪零性により零である．
これが豊富性の二条件を与える．議論は任意のnoetherian $T$ について同じである．
\end{proof}

\subsection{Compact対象，proper性と任意基底変更}
\begin{proposition}[Compact生成と相対proper性]
\label{prop:artin-compact-proper}
$D(\mathcal Q_T)$ は $\{E_{i,T}\}_{i\in\mathbb Z}$ によりcompact生成される．
そのcompact部分圏の全てのRHomは $T$-perfectである．
任意の $R\to T$ に対してcompact圏の導来基底変更は
$\operatorname{Perf}(\mathcal Q_T)$ とMorita同値であり，RHomと基底変更が両立する．
\end{proposition}
\begin{proof}
有限\v{C}ech局所化関手は任意の直和を保つ．従って導来右随伴 $R\omega$ の像は直和で閉じ，
$R\omega$ 自体が直和を保つ．$B_T(-i)$ は $D(\operatorname{Gr}B_T)$ のcompact生成系なので，
導来随伴により $E_{i,T}$ はcompactであり，これらが $D(\mathcal Q_T)$ を生成する．
$T=R$ におけるcompact性は，既知の持上げ定理
\cite[Propositions~5.8--5.9]{DDL2011}からも得られる．
その結果を任意基底変更へ無条件に拡張せず，ここでは\v{C}ech模型で直接確認した．

任意の整数 $q$ に対し，$\check C_R(B)_q$ は有界なflat $R$-複体であり，
還元後のコホモロジーは $H^*(X,\mathcal O_X(qH))$ なので有限次元かつ有界である．
Artin局所環上ではflat加群は自由である．還元複体をそのコホモロジーと可縮な複体に分解し，
各次数の自由基底を持ち上げると，微分は還元微分の $R$-線型延長と
$\mathfrak m$ に値を持つ摂動の和となる．縮約に対する有限幾何級数
$(1-h\delta)^{-1}=\sum_{a=0}^{N-1}(h\delta)^a$ を用いると，
$\check C_R(B)_q$ は有界有限自由 $R$-複体にホモトピー同値となる．
この有限模型は任意の $T$ へテンソルしてよく，前命題の模型に一致する．
従って全ての $R\operatorname{Hom}(E_{i,T},E_{j,T})$ は $T$-perfectであり，基底変更に両立する．
有限coneと直和因子を両変数で取ることにより，全compact対象について同じ結論を得る．

最後に，$E_i$ を対象とする小さいdg圏の任意基底変更は，
$E_{i,T}$ を対象とするdg圏に準同値である．両方がcompact圏を生成するので，
導来Morita理論\cite[\S4.3,\S8.2]{Keller1994}から主張の基底変更同値を得る．
\end{proof}

\subsection{対角dg加群のcompact性とsmooth性}
\begin{lemma}[dg 加群の compact 性の冪零持上げ]
\label{lem:dg-compact-nilpotent-lift}
$(R,\mathfrak m)$ を剰余体 $k$ の局所 Artin 環とし，
$\mathcal B$ を小さい $R$-線形 dg 圏とする．
$\mathcal B_0=\mathcal B\otimes_R^{\mathbf L}k$ と置く．
$M\in D(\mathcal B)$ に対し，
$M_0=M\otimes_R^{\mathbf L}k$ が $D(\mathcal B_0)$ で compact なら，
$M$ は $D(\mathcal B)$ で compact である．
\end{lemma}
\begin{proof}
全ての基底変更を計算するため，Hom 複体が $R$-K-flat な dg 模型を取る．
係数変更 $u:\mathcal B\to\mathcal B_0$ による restriction を
$\operatorname{res}:D(\mathcal B_0)\to D(\mathcal B)$ と書く．
この関手は直和を保ち，その左随伴は
$-\otimes_{\mathcal B}^{\mathbf L}\mathcal B_0$ である．
従って任意の対象族 $L_\alpha\in D(\mathcal B_0)$ について，随伴と
$M_0$ の compact 性から
\[
 \bigoplus_\alpha R\operatorname{Hom}_{\mathcal B}
       (M,\operatorname{res}L_\alpha)
 \xrightarrow{\sim}
 R\operatorname{Hom}_{\mathcal B}
       \left(M,\bigoplus_\alpha\operatorname{res}L_\alpha\right)
\]
を得る．

$\mathfrak m^N=0$ とし，$V_a=\mathfrak m^a/\mathfrak m^{a+1}$ と置く．
$R$ の有限フィルトレーションに対角 bimodule
$\mathcal B\otimes_R^{\mathbf L}-$ を施すと，
\[
 F_{a+1}\longrightarrow F_a
 \longrightarrow\mathcal B_0\otimes_kV_a\longrightarrow,
 \qquad F_a=\mathcal B\otimes_R^{\mathbf L}\mathfrak m^a
\]
という bimodule の三角を得る．ここで $F_0=\mathcal B$，$F_N=0$ であり，
第三項の左右作用はともに $\mathcal B_0$ を経由する．
任意の $Y\in D(\mathcal B)$ に対して
$R\operatorname{Hom}_{\mathcal B}(-,Y)$ をこの有限列に施す．
ここでbimoduleからのRHomは，残る作用を備えたdg加群を意味する．
すると $Y$ は有限回の三角によって
\[
 R\operatorname{Hom}_{\mathcal B}
    (\mathcal B_0\otimes_kV_a,Y)
\]
から構成される．これらは，第一変数に残る
$\mathcal B_0$-作用によって全て $\operatorname{res}$ の像に属する．
この構成は $Y$ に関手的で，列の長さは $Y$ に依存しない．

対象族 $Y_\alpha$ にこの列を適用し，列の各三角の直和を取る．
最初に証明した restriction の像上の比較同型と三角に関する帰納から，
\[
 \bigoplus_\alpha R\operatorname{Hom}_{\mathcal B}(M,Y_\alpha)
 \xrightarrow{\sim}
 R\operatorname{Hom}_{\mathcal B}
       \left(M,\bigoplus_\alpha Y_\alpha\right)
\]
を得る．これが compact 性である．なお coinduction
$R\operatorname{Hom}_{\mathcal B}(\mathcal B_0,-)$ 自体が直和を保つとは
仮定していない．個々の $Y_\alpha$ の有限列を直和したことだけを用いた．
\end{proof}

\begin{remark}
これは \cite[Proposition~5.8]{DDL2011} の compact 持上げの論法を
小さい dg 圏の加群へ直接適用したものである．ここではその証明を書いたため，
平坦 Abel 圏の定理を dg 圏へ無説明で引用してはいない．
\end{remark}

\begin{theorem}[全 Artin 量子化の smooth proper な perfect 圏]
\label{thm:artin-smooth-proper}
$\operatorname{Perf}(\mathcal Q_R)$ は $R$ 上 smooth かつ proper である．
この性質は任意の係数変更後にも保たれる．
\end{theorem}
\begin{proof}
全ての $E_i$ を対象とする dg 圏の $R$-K-flat な模型を $\mathcal A_R$ とする．
前命題の有限 \v{C}ech 模型により，その各 Hom は $R$-perfect であり，
\[
 \mathcal A_R\otimes_R^{\mathbf L}k\simeq\mathcal A_0
\]
は $\operatorname{Perf}(X)$ を生成する line bundle 列の dg 圏である．
従って proper 性は前命題で証明済みである．

$\mathcal B=\mathcal A_R^{\mathrm{op}}
 \otimes_R^{\mathbf L}\mathcal A_R$ とし，
$M$ を $\mathcal A_R$ の対角 bimodule とする．
その基底変更は $\mathcal A_0$ の対角 bimodule である．
$X$ は滑らかな射影多様体なので $\operatorname{Perf}(X)$ は smooth であり，
Morita 不変性からこの還元対角は $D(\mathcal B_0)$ で compact である．
補題\ref{lem:dg-compact-nilpotent-lift}を適用すると $M$ は
$D(\mathcal B)$ で compact，すなわち perfect である．
これが $\mathcal A_R$ の相対 smooth 性である．

ここから有限個の $E_i$ の直和をcompact generatorとして取れることも従う．
実際，perfectな対角の有限cell表示には有限個の外積representableしか現れない．
任意のrepresentableをこの対角とテンソルすると，proper性により係数は
$R$-perfectであるから，それらの有限個の第二因子representableの厚い閉包に入る．


perfect 対角と perfect な Hom 複体は任意の導来基底変更で保たれる．
前命題の compact 圏の基底変更同値を用いれば，任意の $R\to T$ に対して
$\operatorname{Perf}(\mathcal Q_T)$ についても同じ結論を得る．
\end{proof}

\subsection{二重tails上の平坦な対角対象}
以上のsmooth性の証明とは独立に，重み付き表示から具体的な対角対象も構成できる．
$C=B^{\mathrm{op}}\otimes_RB$ を二重次数付き代数とする．
第一・第二因子それぞれについて前節の有限Ore被覆を取り，
両因子の局所化で零となる対象を除くGabriel商を $\mathcal P_R$ とする．
より明示的には，全ての $i,j$ について
$C\to B_i^{\mathrm{op}}\otimes_RB_j$ で零となる加群を核とする．
各局所化は強二重次数付きなので，その次数 $(0,0)$ 加群圏を用いても同じ商圏を得る．

\begin{lemma}[積の平坦変形と対角]
\label{lem:artin-product-diagonal}
$\mathcal P_R$ は $\operatorname{QCoh}(X\times X)$ の平坦Abel変形である．
二重次数付き $C$-加群
\[
 \Delta_B{}_{p,q}=B_{p+q}\quad(p,q\in\mathbb Z),\qquad
 u\cdot(a^{\mathrm{op}}\otimes b)=a u b
\]
の像 $\Delta_R\in\mathcal P_R$ は $R$-flatであり，その還元は
対角の構造層 $\Delta_*\mathcal O_X$ である．$\Delta_R$ は $D(\mathcal P_R)$ のcompact対象である．
\end{lemma}
\begin{proof}
二重局所化の重なりと有限\v{C}ech模型は前節と同様に構成できる．
$C$ は各二重次数で $R$-freeであり，Gabriel商のinjective/coflat判定
（補題\ref{lem:tails-flat-deformation}の証明）により $\mathcal P_R$ は平坦である．
係数変更も同じGabriel商の議論により可換する．
可換還元は二つの自由な重み付き射影商の積なので $\operatorname{QCoh}(X\times X)$ である．

$\Delta_B{}_{p,q}$ は各次数で $R$-freeだからその商圏内の像も $R$-flatである．
還元後，$B_{0,p+q}$ は対角上の $\mathcal O_X((p+q)H)$ の切断である．
局所化して次数 $(0,0)$ を取ると，$U_i\times U_j$ 上で
$\Delta_*\mathcal O_{U_i\cap U_j}$ を得るので，還元は対角の構造層である．
$X$ は滑らかなので $\Delta_*\mathcal O_X$ は $X\times X$ 上perfectである．
$\Delta_R$ はflatであり，その導来還元は通常の還元に等しい．
従って \cite[Proposition~5.8]{DDL2011} をこの平坦Abel変形に適用すればcompact性を得る．
ここでは $\Delta_R$ 自体が相対平坦な持上げとして存在することを使っており，
任意の特殊ファイバーのperfect対象が無条件に持ち上がるとは述べていない．
\end{proof}


\begin{remark}[相対Serre公式との役割分担]
以上により，Artin段階ではsection saturation，有向ACM，任意の正周期の豊富性，
compact生成，相対perfectなRHomとsmooth proper性が得られる．
これだけではSerre関手を次数移動 $(-(n-1))[n]$ と同定していない．
その同定には相対Serre双対bimoduleの比較が必要である．
また，Artin係数での結論を有限の非零複素パラメータの代数的族へ延長したとは主張しない．
\end{remark}

% !TeX root = ../nc-del-pezzo.tex
\section{体積形式と相対 Serre 双対}
\label{sec:relative-serre-artin}

\subsection{任意の冪零係数上の無捩れな ambient 双対性}

この小節では，Kontsevich の写像と Shoikhet の鎖写像を対として固定する．
後者は前者を用いた Hochschild 鎖の $L_\infty$ 加群構造から微分形式への写像である
\cite[Theorem~2.2.7, \S2.2.8, \S3.2--3.3]{Shoikhet2003}．
その線型座標変換に関する同変性は
\cite[Theorem~3]{DolgushevChains2006}で明記されている．
以下の議論は多項式環上のグラフ公式だけを使い，
滑らかな多様体上の大域化や非線型座標変換の厳密同変性を使わない．

\begin{lemma}[ambient 括弧の相対 unimodularity]
\label{lem:ambient-unimodular-all-artin}
第\ref{sec:equivariant-quantization}節の二つの ambient 二重ベクトルは，
標準の相対体積形式
\[
 \omega=dx_0\wedge\cdots\wedge dx_{n-1}\wedge dy\wedge dz
 \quad(d=1),\qquad
 \omega=dx_0\wedge\cdots\wedge dx_n\wedge dz
 \quad(d=2)
\]
に関して $\operatorname{div}_\omega\widetilde\pi=0$ を満たす．
この等式は任意の係数環上で成立し，次数2の $P$ に対して
二つのベクトルの外積による表示を選ぶ必要はない．
\end{lemma}
\begin{proof}
$\widetilde\pi=\sum_{a<b}\pi^{ab}\partial_a\wedge\partial_b$ とし，
発散の第 $b$ 成分を $\sum_a\partial_a\pi^{ab}$ とする．
次数1では $D_i$ は相対微分で定数である．
$x_i$ 成分は $-D_iF_{zy}+D_iF_{yz}=0$，
$y$ 成分は $D(F_z)-\partial_zD(F)=0$，
$z$ 成分は $-D(F_y)+\partial_yD(F)=0$ である．
次数2では $\pi^{ij}=2zp_{ij}$ と
$\pi^{iz}=V_i=\sum_jp_{ij}\partial_jf_4$ である．
$x_i$ 成分は零であり，$z$ 成分は
$\sum_{i,j}p_{ij}\partial_i\partial_jf_4=0$ となる．
最後の等式は反対称性と混合微分の可換性による．
従って $L_{\widetilde\pi}\omega=0$ である．
\end{proof}

\begin{lemma}[鎖 formality の冪零係数版]
\label{lem:nilpotent-chain-formality}
$R$ を可換 $\mathbb C$-代数，$J\subset R$ を冪零イデアルとし，
$S_R=R[x_1,\ldots,x_N]$ に重み $w_i>0$ を与える．
$\pi\in J\bigwedge^2\operatorname{Der}_R(S_R)$ を
内部次数零の Poisson 二重ベクトルとする．
固定した Kontsevich--Shoikhet の対で量子化した代数を $A$ とすると，
相対 Hochschild 鎖複体から相対微分形式への内部次数を保つ鎖写像
\[
 \mathcal K_\pi:
 (C_\bullet(A/R),b_\star)
 \longrightarrow(\Omega^\bullet_{S_R/R},L_\pi)
\]
があり，これは擬同型である．その $J$ による還元は
$\frac1{q!}a_0\,da_1\wedge\cdots\wedge da_q$ による HKR 写像である．
構成は任意の $R\to T$ と両立する．
\end{lemma}
\begin{proof}
Shoikhet の $q$ 次の構造写像を $\mathcal K_q$ と書くと，
\[
 \mathcal K_\pi(c)
 =\sum_{q\ge0}\frac1{q!}\mathcal K_q(\pi,\ldots,\pi;c)
\]
と定める．$q$ 番目の項の係数は $J^q$ に入るため有限和である．
各グラフは ambient 変数だけを微分する有限階の式であり，
多項式を多項式微分形式へ送る．従って複素多項式上で成立する
鎖写像の恒等式は任意の $R$ 上に拡張できる．
これは一変数の $\hbar$ 進量子化から値 $\hbar=1$ を取る操作ではない．
$J$ 進フィルトレーションの各随伴次数上でこの写像は，
$J^a/J^{a+1}$ を複素ベクトル空間として係数に持つ多項式 HKR 写像である．
その擬同型性とフィルトレーションの有限性から元の写像も擬同型となる．
内部次数については各グラフの縮約が重みを保存する．
有限グラフ公式なので係数変更は平坦性なしに可換し，
係数変更後にも同じ有限フィルトレーションの議論が適用できる．
\end{proof}

\begin{theorem}[重み付き相対 Calabi--Yau 同型]
\label{thm:nilpotent-ambient-cy}
前補題の仮定に加え，$W=\sum_iw_i$ とし，
$\omega=dx_1\wedge\cdots\wedge dx_N$ に関して
$L_\pi\omega=0$ と仮定する．$A^{\mathrm e}=A\otimes_R A^{\mathrm{op}}$ と置くと，
次数付き bimodule の導来圏で
\[
 A^!:=\underline{R\operatorname{Hom}}_{A^{\mathrm e}}(A,A^{\mathrm e})
 \simeq A(W)[-N].
\]
この同型は，体積形式と鎖 formality の対を固定したとき任意の係数変更と両立する．
従って全ての基底変更後に左右の作用に Nakayama 自己同型の捻りは残らない．
\end{theorem}
\begin{proof}
$\omega$ は $L_\pi$ に関する最高次の cycle である．
補題\ref{lem:nilpotent-chain-formality}の同型により，
これに対応する一意の類
\[
 c_\pi\in HH_N(A/R)_W
\]
を得る．その還元は多項式環の標準的な体積類である．
命題\ref{prop:nilpotent-weighted-as}により対角 $A$ は
有限自由 $A^{\mathrm e}$-分解を持つ．有限双対性から
\[
 A\mathbin{\otimes^{\mathbf L}_{A^{\mathrm e}}}A
 \simeq R\operatorname{Hom}_{(A^{\mathrm e})^{\mathrm{op}}}(A^!,A)
\]
が成り立つ．右 $A^{\mathrm e}$-加群は通常の反転で bimodule と読む．
従って内部次数 $W$ の類 $c_\pi$ は，内部次数零の射
\[
 \theta_\pi:A^!\longrightarrow A(W)[-N]
\]
を与える．還元後，この射は Koszul 対角分解の標準双対同型である．
有限自由分解を使って $\theta_\pi$ を chain map として表すと，
その cone は各内部次数で有界な $R$-射影加群の複体となり，
$J$ による還元後には完全である．
補題\ref{lem:nilpotent-complex-exactness}により cone は完全であり，
$\theta_\pi$ は同型である．

$R\to T$ による係数変更に対し，$c_\pi$ の像は
$\mathcal K_{\pi_T}$ によって $\omega_T$ へ送られる．
後者も擬同型なので，この像は $c_{\pi_T}$ に等しい．
有限自由対角分解の双対性も任意の係数変更と可換する．
従って $\theta_\pi$ は基底変更と両立する．
\end{proof}

\begin{corollary}[次数1・2への適用]
\label{cor:dp-ambient-cy}
第\ref{sec:equivariant-quantization}節の量子化を
Kontsevich--Shoikhet の対から選ぶと，任意の冪零係数上で
\[
 A_R^!\simeq
 \begin{cases}
 A_R(n+5)[-n-2],&d=1,\\
 A_R(n+3)[-n-2],&d=2
 \end{cases}
\]
である．とくに定理\ref{thm:concrete-quantization-all-artin}の
全 Artin 表示には，この無捩れな相対 ambient 双対性を持つ表示を選べる．
\end{corollary}
\begin{proof}
補題\ref{lem:ambient-unimodular-all-artin}と前定理を用い，
$N=n+2$ とそれぞれの重みの和を代入する．
\end{proof}

\begin{remark}[既知の unimodularity 定理と射影的双対性]
Dolgushev は滑らかな affine 多様体の一変数形式量子化について，
Van den Bergh 双対化 bimodule が無捩れになることと
形式 Poisson 構造の unimodularity の同値を証明した
\cite[Theorem~3]{DolgushevModular2009}．
本稿ではその一変数定理を任意の非被約基底へ直接適用せず，
鎖 formality と有限自由対角分解による上の証明を用いた．
これは既知の原理の相対・重み付きの適用であり，原理自体を新規結果とはしない．

$\Omega$ が次数 $m$ の正則中心元なら，二項分解によって
\[
 \underline{R\operatorname{Hom}}_{A_R}
 (B_R,A_R(-W)[N])\simeq B_R(m-W)[N-1]
 =B_R(-(n-1))[n+1]
\]
を得る．ここでは両辺を中心商上の bimodule として読む．
この随伴複体の計算だけで，相対 balanced dualizing complex の全公理，
射影 tails の相対 Serre 双対，または compact 生成を証明したことにはならない．
以下では smooth proper な相対 dg 圏が得られた場合に，
Serre 関手の無捩れな同定を別の剛性論法で行う．
\end{remark}

\subsection{対角 bimodule の剛性による Serre 関手の同定}

この小節で dg 代数の包絡代数を書くときは
$A^{\mathrm e}=A\otimes_R^{\mathbf L}A^{\mathrm{op}}$ を意味する．
同値に，$R$-cofibrant な模型を選んで通常のテンソル積で計算してよい．

\begin{lemma}[対角の二つの持上げの同型]
\label{lem:diagonal-bimodule-rigidity}
$(R,\mathfrak m)$ を剰余体 $\mathbb C$ の局所 Artin 環とし，
$A$ を $R$ 上 smooth proper な dg 代数とする．
$A_0=A\otimes_R^{\mathbf L}\mathbb C$ と置き，$HH^1(A_0)=0$ と仮定する．
$P,Q\in\operatorname{Perf}(A^{\mathrm e})$ の還元に
$P_0\simeq A_0\simeq Q_0$ を固定する．
この還元上の任意の同型 $P_0\to Q_0$ は，
$P\simeq Q$ へ持ち上がる．
\end{lemma}
\begin{proof}
$P$ の完全性と $A$ の proper 性から
$K=R\operatorname{Hom}_{A^{\mathrm e}}(P,Q)$ は $R$-perfect である．
これは任意の導来基底変更と可換し，
$K\otimes_R^{\mathbf L}\mathbb C\simeq
 R\operatorname{Hom}_{A_0^{\mathrm e}}(A_0,A_0)$ である．
$R/\mathfrak m^{a+1}\to R/\mathfrak m^a$ の核を
$I_a=\mathfrak m^a/\mathfrak m^{a+1}$ と置く．
$\mathfrak m I_a=0$ なので，導来テンソル積の結合律から
\[
 I_a\otimes_R^{\mathbf L}K
 \simeq I_a\otimes_{\mathbb C}
       (\mathbb C\otimes_R^{\mathbf L}K)
\]
である．従って写像持上げの障害は，
この小拡大を $K$ とテンソルした三角から
\[
 H^1(I_a\otimes_R^{\mathbf L}K)
 =I_a\otimes_{\mathbb C}HH^1(A_0)=0
\]
に入る．従って $H^0(K)$ の元として所望の射を帰納的に持ち上げられる．
ここでは $P,Q$ の存在は仮定済みであり，その間の射だけを持ち上げている．
未構成の対象の持上げに現れる $HH^2$ の障害の消滅を主張していない．
その cone は $R$-perfect で還元が零なので，導来 Nakayama により零である．
よって得られた射は同型である．
\end{proof}

\begin{proposition}[smooth proper な Artin 持上げの相対 Serre 関手]
\label{prop:serre-from-smooth-proper-lift}
$R$ を前補題の局所 Artin 環とし，$\mathcal C_R$ を
$R$ 上 smooth proper な小さい dg 圏で，その完全包が compact generator を持つものとする．
以下ではその冪等完備な完全包を同じ記号で書く．
その還元が $\operatorname{Perf}(X)$ であるとする．
$\mathsf T_R$ を $\mathcal C_R$ の dg 自己同値で，
還元が $-\otimes\mathcal O_X(H)$ になるものとする．
$X$ が次数1または2の滑らかな del Pezzo $n$-fold なら，
相対 Serre 関手 $\mathsf S_R$ が存在し，
\[
 \mathsf S_R\simeq\mathsf T_R^{-(n-1)}[n].
\]
この同型を一つ選ぶと，任意の $R\to T$ に対して導来基底変更された
smooth proper 圏の相対 Serre 同型を与える．
$\mathsf T_R$ が Abel 圏の次数シフト $(1)$ から来る場合，
$\mathsf S_R^{-1}[n]$ は完全自己同値 $(n-1)$ の制限であり，
$E_i=\pi B_R(-i)$ に対して $E_i\mapsto E_{i-(n-1)}$ となる．
\end{proposition}
\begin{proof}
compact generator による Morita 模型を $A$ とし，
$A^*=R\operatorname{Hom}_R(A,R)$ と置く．
$A$ が proper なので $A^*$ は $R$-perfect である．
$A$ が smooth なので，$R$-perfect な $A$-加群は $A$-perfect である．
実際，有限な perfect $A^{\mathrm e}$-分解で対角を表し，その加群をテンソルすると，
$R$-perfect な係数を持つ $A$ の有限な cone 表示を得る．
従って $\mathsf S_R=-\otimes_A^{\mathbf L}A^*$ は
$\operatorname{Perf}(A)$ の自己関手である．
テンソルと Hom の随伴関係から
\[
 R\operatorname{Hom}_R(R\operatorname{Hom}_A(M,N),R)
 \simeq R\operatorname{Hom}_A
 (N,M\otimes_A^{\mathbf L}A^*)
\]
を得る．これは $M=A$ の場合には通常の随伴同型であり，
有限な cone と直和因子により全ての $M\in\operatorname{Perf}(A)$ に拡張する．
$\mathsf S_R$ の右随伴は
$R\operatorname{Hom}_A(A^*,-)$ である．
$A^*$ と入力が $A$-perfect なのでその値は $R$-perfect であり，
smooth 性から $A$-perfect でもある．従って右随伴も完全圏を保つ．
この右随伴とその随伴写像は基底変更と両立する．
その単位・余単位は特殊ファイバーで通常の Serre 同値のものになるため，
$R$-perfect な cone に導来 Nakayama を適用すれば同型となる．
従って $\mathsf S_R$ は相対 Serre 自己同値である．

$\mathsf S_R\mathsf T_R^{n-1}[-n]$ を表す bimodule を $P$ とする．
その右 $A$-加群としての値は生成元の像なので $A$-perfect，従って
$R$-perfect である．$A^{\mathrm e}$ 自身も $R$ 上 smooth なので，
（その対角は $A$ と $A^{\mathrm{op}}$ のperfect対角の外積である），
同じ対角分解の議論により $P$ は $A^{\mathrm e}$-perfect である．
$\omega_X\simeq\mathcal O_X(-(n-1)H)$ により $P_0\simeq A_0$ である．
HKR と $H^0(X,T_X)=H^1(X,\mathcal O_X)=0$ により
\[
 HH^1(A_0)=HH^1(X)
 =H^0(X,T_X)\oplus H^1(X,\mathcal O_X)=0.
\]
$H^0(T_X)=0$ は \cite[補題8.7]{PoissonDP} の入力であり，
$H^1(\mathcal O_X)=0$ は Kodaira 消滅による．
補題\ref{lem:diagonal-bimodule-rigidity}を $Q=A$ に適用すると
$P\simeq A$ であり，所望の Serre 同型を得る．

$A$ の $R$-完全性により $A^*$ の構成は任意の基底変更と可換する．
perfect bimodule，その合成，同型も導来基底変更で保たれるので，
選んだ Serre 同型は任意の $T$ 上へ移る．
最後の主張は次数シフトの合成と $E_i$ の定義による．
\end{proof}

\begin{remark}[perfect圏と有界coherent導来圏]
前命題のsmooth proper性とcompact生成の入力は、
本稿のArtin中心商については定理\ref{thm:artin-smooth-proper}で満たされる．
相対 Serre 同型の適用先は $\operatorname{Perf}$ であり，
特異な Artin 基底上の全 $D^b(\operatorname{coh})$ ではない．
例えば $R=\mathbb C[\epsilon]/(\epsilon^2)$ 上の自明族でも，
$\mathcal O_{X\times\operatorname{Spec}R}\otimes_R\mathbb C$ は coherent だが，
$R$ 上の射影次元が無限なので相対 compact ではない．
反標準列の豊富性とsection saturationは、それぞれ
定理\ref{thm:artin-positive-shift-ample}と定理\ref{thm:artin-saturation-acm}による．
\end{remark}

\subsection{全Abel変形と非可換多様体のArtin完成の一致}

\begin{lemma}[局所Artin基底上のstack化と枠付き変形]
\label{lem:artin-stackification-comparison}
$\mathcal F$ を局所noetherianな複素スキームのfppf site上の
群oid擬関手，$\mathcal F^a$ をそのstack化，$x\in\mathcal F(\mathbb C)$ とする．
$\mathcal F_x(R)$ は $y\in\mathcal F(R)$ と
$y\otimes_R\mathbb C\simeq x$ という枠の組の群oidとする．
剰余体 $\mathbb C$ の局所Artin環上で，ある完備局所
$\mathbb C$-代数 $A$ に対して自然な同値
\[
 \mathcal F_x(R)\simeq
 h_A(R):=\operatorname{Hom}_{\mathrm{cont,loc},\mathbb C}(A,R)
\]
が存在すると仮定する．右辺は離散群oidとみなす．
このとき自然な関手
$\mathcal F_x(R)\to(\mathcal F^a)_x(R)$ は同値である．
\end{lemma}
\begin{proof}
まず $R$ 上の任意のfppf被覆は，有限忠実平坦な局所Artin
$\mathbb C$-代数 $S$ による一つの被覆で細分できる．
実際，fppf被覆は局所準有限なfppf被覆で細分できる
\cite[Tag 0572]{StacksProject}．
$\operatorname{Spec}R$ の一点を覆う細分の非空なaffine開を取る．
その座標環は $\mathbb C$ 上有限型かつ零次元なので有限次元である．
その非零な局所因子を $S$ とすれば，$R$ 上有限平坦かつ非零であるから
忠実平坦であり，剰余体は $\mathbb C$ である．
有限個の局所Artin $\mathbb C$-代数の $R$ 上のテンソル積も
局所Artin $\mathbb C$-代数である：各極大イデアルから生成される
イデアルは冪零で，その商は $\mathbb C$ になる．
従って被覆の全ての\v{C}ech重なりも同じ試験環の範囲に入る．
また $R=\mathbb C$ の被覆は $\mathbb C$-点により分裂するので，
$\mathcal F(\mathbb C)\to\mathcal F^a(\mathbb C)$ は同値である．

忠実平坦性から $R\to S\rightrightarrows S\otimes_RS$ は等化子である．
従って $h_A$ の元とその等式はこの被覆について降下する．
局所性も降下するので，得られる準同型 $A\to R$ は局所的かつ連続である．
$\mathcal F_x$ は離散群oid $h_A$ に同値だから，その降下は有効であり，
同じ同型類の二対象間にはただ一つの射がある．

これがstack化の射にも適用できることを先に確認する．
$\mathcal F_x(R)$ の二対象間の $\mathcal F^a$ における枠付き射は，
被覆上で $\mathcal F$ の射として表せる．上の $S$ に細分すると
両者の $h_A(S)$ での値は一致するので，忠実平坦性から
$h_A(R)$ でも一致し，$\mathcal F_x(R)$ に一意の射がある．
stack化された射は局所的にこの射と一致するため全体でも一致する．
よって自然な関手は完全忠実である．

次に $(\mathcal F^a)_x(R)$ の対象を $S$ 上で
$\mathcal F_x(S)$ の対象として表す．特殊ファイバーの枠は，
$\mathcal F(\mathbb C)\simeq\mathcal F^a(\mathbb C)$ により
元の群oid内で選べる．重なり上の降下同型は，先に示した完全忠実性によって
$\mathcal F_x(S\otimes_RS)$ 内の射へ一意に持ち上がる．
従って $h_A(S)$ の値は降下して $h_A(R)$ の元を与える．
対応する $\mathcal F_x(R)$ の対象は，stack化後に元の対象と同型となる．
これで本質的全射も得た．
\end{proof}


\begin{theorem}[全Artin変形の非可換del Pezzo性]
\label{thm:all-artin-are-nc-dp}
$X$ を次数 $d\in\{1,2\}$ の滑らかな複素del Pezzo $n$-fold，$n\ge2$ とする．
その構造対象付き平坦Abel Artin変形は全て
$\mathfrak M^{\mathcal O,\mathrm{hyp}}_{n,d}$ に属する．
そのperfect圏は相対的にsmooth properであり、ある通常の次数シフトに対して
\[
 \mathsf S_R\simeq (-(n-1))[n],\qquad
 \mathsf S_R^{-1}[n]\simeq(n-1)
\]
となる．これらの性質は任意の許される基底変更後にも成立する．

従って，可換点 $[X]$ におけるArtin完成を変形群oidとして読むと，
自然な包含は同値
\[
 \widehat{\mathfrak M}^{\mathcal O,\mathrm{hyp}}_{n,d,[X]}
 \simeq\widehat{\mathfrak M}^{\mathcal O}_{n,d,[X]}
 \simeq\widehat{\mathfrak A}^{\mathcal O}_{n,d,[X]}
 \simeq\widehat{\mathfrak D}^{\mathcal O}_{[X]}
 \simeq[\operatorname{Spf}R_X/\operatorname{Aut}(X)]
\]
を与える．特殊ファイバーの枠を固定した版は $\operatorname{Spf}R_X$ である．
また，具体量子化表示 $B,B'$ 間の全構造対象付き同値は
\[
 \operatorname{Eq}^{\mathcal O}_R(\operatorname{QGr}B,\operatorname{QGr}B')
 \simeq\operatorname{Isom}_{\mathrm{gr},R}(B,B')/R^\times
\]
で認識される．
\end{theorem}
\begin{proof}
定理\ref{thm:concrete-quantization-all-artin}によって任意のAbel変形に
通常の重み付き中心超曲面表示を選ぶ．
定理\ref{thm:artin-saturation-acm}と定理\ref{thm:artin-positive-shift-ample}により，
全シフト列は許容表示であり，有向ACMと正周期の豊富性を満たす．
定理\ref{thm:artin-smooth-proper}によりcompact生成と相対smooth proper性，
perfectなRHomと任意基底変更との両立を得る．有限compact generatorも同定理の証明にある．
命題\ref{prop:serre-from-smooth-proper-lift}を次数シフト $(1)$ に適用すれば
表示した相対Serre公式が成り立つ．
$E_i=\pi B(-i)$ と $\mathsf H=(n-1)$ は必要なSerre周期を整合的に満たす．

任意の基底変更の幾何学的ファイバーでは冪零パラメータが消えるため，
元の滑らかな $X$ の体拡大となり，ホモロジー次元は $n$ である．
ambient代数も重み付き多項式環の体拡大となり，所定の重み・大域次元・
Gorensteinパラメータを持つAS正則代数である．
中心元の普遍的正則性と商の次数別平坦性は
定理\ref{thm:nilpotent-universal-quantization}に含まれる．
よって定義\ref{def:nc-dp-stack}と超曲面部分問題の全条件が満たされる．

これらの部分問題は全 $\mathcal O$ 保存同値についてfullであり，
全Artin対象が既に最小の部分問題に入るから，stack化前の
枠付き変形群oidは全て $\mathfrak D_X^{\mathcal O}$ と同値である．
定理\ref{thm:quantization-framed-groupoid}と
定理\ref{thm:abel-formal-ring}によれば，これは離散群oid
$h_{R_X}$ に同値である．
補題\ref{lem:artin-stackification-comparison}をそれぞれの
stack化前の部分問題に適用すると，fppf stack化後も
枠付きArtin完成は変わらない．
枠を忘れる操作は $\operatorname{Aut}(X)$ による商であり，
定理\ref{thm:unframed-formal-quotient}の形式商表示を得る．
この議論は正次元の基底上の変形分類や降下の比較を仮定しない．

次数付き同型による認識は、飽和性と豊富性を
命題\ref{prop:artin-full-graded-recognition}へ適用して得る．
\end{proof}

\begin{remark}[局所定理と大域的な成分数]
この定理により，定義した非可換多様体の可換点での形式分枝は一つであり，
形式粗商は正規Cohen--Macaulay整域となる．
被覆対合だけを持つ可換点の粗商の特異性は
系\ref{cor:generic-nc-formal-sign-quotient}で計算される．
一方，有限型atlas，形式量子化の代数化，一般nc点の全同値認識，
可換locusに接しない成分の有無は未解決である．
本定理はそれらを仮定せずにArtin完成を同定するものである．
\end{remark}

% !TeX root = ../nc-del-pezzo.tex
% 2026-11-20 JST / v0.1
\section{固定した非零方向の形式一般ファイバーと中心}
\label{subsec:formal-generic-center}

ここでは滑らかな複素多様体 $X$ と非零 Poisson 構造 $\pi$ を固定する．
冪零係数環上の任意の非零パラメータに対して適合座標が存在すると主張していない．
適合する複素線型座標を選び，中心的な一次座標を
$s_1,\ldots,s_k$，$k=n-1$ と書く．
$P=\mathbb C[s_1,\ldots,s_k]$，$K=\operatorname{Frac}(P)$ とする．
残る変数を次数1では $(u,y,z)$，次数2では $(u,v,z)$ とすると，
可換切断環 $B_0$ は
\[
 B_0=\begin{cases}
 P[u,y,z]/(z^2-y^3-A_4(u,s)y-A_6(u,s)),&d=1,\\
 P[u,v,z]/(z^2-f_4(u,v,s)),&d=2
 \end{cases}
\]
であり，Poisson 括弧はこれらの活性な三変数の Jacobian 括弧である．

\begin{lemma}[古典切断環の中心]
\label{lem:classical-casimir-center}
$B_0\otimes_PK$ は幾何学的整な $K$-代数であり，
\[
 Z_{\mathrm{Pois}}(\operatorname{Frac}B_0)=K,
 \qquad Z_{\mathrm{Pois}}(B_0)=P
\]
である．
\end{lemma}
\begin{proof}
次数1では $y^3+A_4(u,s)y+A_6(u,s)$ の $y$-次数が3であるため，
$\overline K[u,y]$ において平方ではない．UFD であることから分数体でも平方でなく，
$z$ の二次式は $\overline K[u,y,z]$ で既約である．

次数2で，反対に $f_4$ が $\overline K[u,v]$ の平方であると仮定する．
平方根のある非零係数を1に正規化すると，平方根の一意性を符号を除いて用いることにより，
正規化した多項式は $K$ 上に降下する．従って
$f_4=a g^2$，$a\in K^\times$，$g\in K[u,v]$ と書ける．
$g$ が $u,v$ に関して非定数なら，$P$ の UFD 性と Gauss の補題により，
$f_4\in P[u,v]$ は $u,v$ に関して非定数な既約因子を重複して持つ．
これは分岐四次超曲面の滑らかさに反する．
$g$ が定数なら $f_4\in P$ であり，その斉次性から
$s_1=\cdots=s_k=0$ で定まる活性変数の射影直線上で $f_4$ と全偏微分が消える．
これも滑らかさに反する．従って $f_4$ は平方でなく，
ここでも $z$ の二次式は $\overline K[u,v,z]$ で既約となる．

$L=\operatorname{Frac}B_0$ とすると $\operatorname{trdeg}_K L=2$ である．
いずれの場合も活性な最初の二変数の括弧は $2z\ne0$ なので，
$\Omega^1_{L/K}\to\operatorname{Der}_K(L)$ は同型となる．
従って Poisson 中心の元 $a\in L$ は $\mathrm da=0$ を満たし，標数零より $K$ 上代数的である．
幾何学的整性により $K$ は $L$ の中で代数的に閉じているため，中心は $K$ である．
最後に $F$ は $z$ について monic なので，$B_0$ は
活性な最初の二変数の単項式と $z$ の指数0または1で与えられる $P$-自由基底を持つ．
この基底は $1$ を含むため，$B_0\cap K=P$ である．
\end{proof}

\begin{theorem}[形式一般ファイバーの普遍中心条件]
\label{thm:formal-generic-center}
$\hbar\pi$ に対して前述の量子化を各冪零段階で行い，
内部次数ごとに逆極限を取った中心商を $B_\hbar$ とする．
$R_\hbar=\mathbb C[[\hbar]]$，$K_\hbar=\mathbb C((\hbar))$，
$P_\hbar=R_\hbar[s_1,\ldots,s_k]$ と書く．このとき
\begin{enumerate}
\item $B_\hbar$ は中心的な $P_\hbar$ 上の次数付き自由加群である．
\item 各内部次数 $\ell$ に対して
\[
 Z(B_\hbar)_\ell=P_\ell[[\hbar]]
\]
が成り立つ．従って $B_\eta=B_\hbar\otimes_{R_\hbar}K_\hbar$ とすると
\[
 Z(B_\eta)=K_\hbar[s_1,\ldots,s_k].
\]
\item 任意の可換 $K_\hbar$-代数 $T$ と任意の $r\ge0$ に対して
\[
 Z(B_\eta\otimes_{K_\hbar}T)_r
   =T[s_1,\ldots,s_k]_r
\]
である．また $B_\eta\otimes_{K_\hbar}T$ は $T[s_1,\ldots,s_k]$ 上自由である．
\end{enumerate}
\end{theorem}
\begin{proof}
各内部次数の逆極限は有限自由 $R_\hbar$-加群である．
二重ベクトルに $\partial_{s_i}$ 成分がないため，グラフ積は
$P_\hbar$-双線型であり，$s_i$ は補正なしに中心に残る．
monic 除法の写像と中心元の乗算も $P_\hbar$-線型である．
定理\ref{thm:nilpotent-universal-quantization}の分裂に現れる逆行列
$(1+K_\ell)^{-1}$ を $\hbar$-進極限で取ると，これも $P_\hbar$-線型となる．
従って活性変数の単項式と $z$ の指数0または1による同じ基底が，
$B_\hbar$ の $P_\hbar$-自由基底を与える．
ここで各次数に現れる単項式は有限個であり，逆極限を取るのは内部次数ごとである．

$b\in Z(B_\hbar)_\ell$ を還元すると，交換子の一次項から
$b\bmod\hbar$ は $B_0$ の Poisson 中心に入る．
補題\ref{lem:classical-casimir-center}により，これはある $p_0\in P_\ell$ に等しい．
$p_0$ 自身が量子中心に属するので，$b-p_0=\hbar b_1$ と書ける．
$\hbar$ は非零因子だから $b_1$ も中心的である．
反復して $b=\sum_{j\ge0}\hbar^jp_j$，$p_j\in P_\ell$ を得る．
逆の包含は既に示したため，次数別の中心の式が従う．
$\hbar$ の非零因子性により，局所化後に中心となる元は
$\hbar$ の冪を掛ければ局所化前から中心に入る．これで一般ファイバーの中心を得る．

$B_\hbar$ は標準の有限個の重み付き変数で代数として生成される．
実際，各次数で対応する有限個の単語からの生成写像は $\hbar$ による還元後全射であり，
有限 $R_\hbar$-加群に対する Nakayama の補題を使える．
従って $B_\eta$ の中心の次数 $r$ は，これらの有限個の生成元との交換子写像
\[
 (B_\eta)_r\longrightarrow\bigoplus_j(B_\eta)_{r+w_j}
\]
の核である．これは有限次元 $K_\hbar$-ベクトル空間の間の写像である．
任意の $K_\hbar$-代数は $K_\hbar$ 上平坦なので，その核は任意の係数変更と可換する．
自由性の係数変更と合わせて最後の主張を得る．
\end{proof}

\begin{corollary}[Veronese 持上げの代数的仮定を満たす形式一般例]
\label{cor:formal-generic-veronese}
$n\ge3$ なら，任意の固定した滑らかな $(X,\pi)$，$\pi\ne0$ に対する $B_\eta$ は，
$k=n-1\ge2$ と任意の $r\ge2$ について
定理\ref{thm:ncdp-veronese-lift}の中心・平坦性仮定を満たす．
特に $r=n-1$ の Veronese 持上げに適用できる．
$n=2$ では反標準指数 $r=1$ なので持上げ自体が不要である．
\end{corollary}

これは $\mathbb C((\hbar))$ 上の例であり，有限の非零複素数への特殊化を構成したものではない．
また本節は相対 Serre 同型，saturation，Nakayama の捻りの消去を証明していないので，
全ての $\mathcal O$ 保存 Abel 同値の認識にはそれらの条件が依然として必要である．
$\hbar$ が冪零となる Artin 商では中心が増大し得るため，
上の普遍中心条件を零切断の全 Artin 近傍へ延長しない．

% !TeX root = ../nc-del-pezzo.tex
\section{有限パラメータ量子化と普遍中心}
\label{sec:finite-auxiliary-quantization}

本節では，特殊な非空な部分族について形式パラメータを有限の値へ特殊化できることを証明する．
方法は，通常の射影的な次数とは別に正の補助重みを与え，
Poisson括弧が補助次数を厳密に下げることを用いる．
これにより各グラフ式は多項式となる．
一般の楕円的な最高次項はこの条件を満たさず，
全Poisson方向の代数化や有限パラメータでの全モジュライ比較は結論しない．

\subsection{有限性と普遍的なPBW基底}

$R$ を任意の可換 $\C$-代数とし，$x,y,z$ に正整数の補助重み $(a,b,c)$ を与える．
$R$ の補助重みは零とする．以下の仮定を置く：
\begin{equation}
 F=z^2/2+V(x,y)\in R[x,y,z],\qquad
 \deg_{\mathrm{aux}}V\le L:=2c<a+b+c,
 \qquad \delta:=a+b+c-L>0.
 \label{eq:negative-auxiliary-hypothesis}
\end{equation}
$R$ には係数パラメータと，後でCasimirとする中央変数を含めてよい．
微分は活性変数 $x,y,z$ だけについて行う．
Jacobian二重ベクトル
\begin{equation}
 \pi_F=F_z\,\partial_x\wedge\partial_y
       +F_x\,\partial_y\wedge\partial_z
       +F_y\,\partial_z\wedge\partial_x
 \label{eq:finite-jacobian-bivector}
\end{equation}
はPoissonであり，$F$ はCasimirである．
グラフformality写像を一つ固定し，その正規化を
$[u,v]_{\star}=\hbar\{u,v\}_{\pi_F}+O(\hbar^2)$ とする．

\begin{theorem}[負の補助次数による代数化]
\label{thm:negative-auxiliary-algebraization}
仮定\eqref{eq:negative-auxiliary-hypothesis}の下で，次が成立する．
\begin{enumerate}
\item 固定したグラフ式による積
\[
 u\star v=uv+\sum_{r\ge1}\hbar^r B_r(u,v)
\]
は $R[\hbar,x,y,z]$ 上の結合的な $R[\hbar]$-代数 $A$ を定める．
補助次数がそれぞれ $p,q$ 以下の入力に対して
\[
 \deg_{\mathrm{aux}}B_r(u,v)\le p+q-r\delta
\]
であり，特に $r>(p+q)/\delta$ の項は零である．
\item 同じformalityのtwistにより中心元
\[
 \Omega=F+\sum_{r\ge1}\hbar^r\Omega_r\in A,
 \qquad
 \deg_{\mathrm{aux}}\Omega_r\le L-r\delta
\]
が定まり，$r>L/\delta$ の項は零である．
その最高補助次数部分は
$F_L=z^2/2+V_L(x,y)$ である．係数 $1/2$ は固定単元である．
\item $A$ は三生成元と三つの交換子関係による有限表示を持つ．
順序付きstar単項式
\[
 x^{\star i}\star y^{\star j}\star z^{\star k}
 \qquad (i,j,k\ge0)
\]
が $R[\hbar]$-基底をなし，
$\operatorname{gr}_{\mathrm{aux}}A\cong R[\hbar,x,y,z]$ である．
\item $\Omega$ は左右の非零因子であり，中心商
$B=A/(\Omega)$ は有限表示を持つ．その $R[\hbar]$-基底は
\[
 x^{\star i}\star y^{\star j}\star z^{\star\epsilon},
 \qquad i,j\ge0,\quad\epsilon\in\{0,1\},
\]
である．さらに
\[
 \operatorname{gr}_{\mathrm{aux}}B
 \cong R[\hbar,x,y,z]/(F_L).
\]
\item これらの構成，基底，有限表示は任意の係数変更と両立する．
特に $\hbar$ は任意の値へ特殊化できる．
係数環がnoetherianなら，$A$ と $B$ は左右noetherianである．
\end{enumerate}
\end{theorem}

\begin{proof}
\eqref{eq:finite-jacobian-bivector}の各項の補助次数は $-\delta$ 以下である．
定数係数のグラフ多重微分作用素は重み付き変数の拡大と可換するため，
$r$ 個の $\pi_F$ を含む積の項は入力の補助次数を少なくとも $r\delta$ 下げる．
多項式の補助次数は非負なので，(1)の有限性が従う．
結合律は形式グラフ恒等式であるが，三つの固定した多項式に適用した両辺は
$\hbar$ の多項式である．したがって形式的な恒等式は
$R[\hbar,x,y,z]$ における恒等式となり，(1)を得る．

$[\pi_F,F]=0$ なので，twisted formality写像はPoisson中心の元 $F$ を
Hochschild零次cocycle，すなわち中心元へ送る\cite{Kontsevich}．
その第 $r$ 項は $F$ 一つと $\pi_F$ を $r$ 個含むため，
同じ重み評価により補助次数が $L-r\delta$ 以下となる．
これは有限和であり，形式中心性は実際の多項式中心性を与える．
$r\ge1$ の項は補助次数が $L$ 未満なので，最高次記号は $F_L$ のままである．

補助次数による増大filtrationを $A$ に入れると，
$u\star v-uv$ は $\deg_{\mathrm{aux}}u+\deg_{\mathrm{aux}}v$ より低い．
したがって順序付きstar単項式は通常の順序付き単項式に対して三角的であり，
対角成分は1である．各補助次数以下の多項式は
$R[\hbar]$ 上有限自由なので，この三角変換は有限な帰納法で可逆である．
これで(3)の基底とassociated gradedの同定を得る．

有限表示もこの基底から直接得られる．
$[y,x]_\star,[z,x]_\star,[z,y]_\star$ を順序付きstar基底で表すと，
それぞれ補助次数 $a+b,a+c,b+c$ より低い有限な多項式になる．
自由代数の対応する三つの交換子関係による商を $A'$ とする．
逆順の隣接生成元を順序付きの積と低い補助次数の項へ置き換えると，
補助次数と逆転数の辞書式順序が下がるため，置換は停止する．
従って順序付き単項式は $A'$ を張る．
自然な全射 $A'\to A$ によるそれらの像は既に基底なので，
関係の重なりから追加の関係が生じることはなく，$A'\cong A$ である．

次に，$2F_L$ は $z$ についてmonicなので，
$R[\hbar,x,y,z]$ の非零因子である．
最高補助次数記号を取ると，$\Omega$ による左右の乗法も単射となる．
$W\subset A$ を $z$ の指数が0または1の順序付きstar単項式の張る自由部分加群とする．
任意の元の最高次記号を $F_L$ でmonic除法し，商と余りを $A$ へ持ち上げる．
商に $\Omega$ を掛けたものと $W$ の余りを引くと補助次数が下がるため，
帰納法から $A=A\star\Omega+W$ を得る．
もし $0\ne w=q\star\Omega\in W$ なら，最高次記号について
$w_{\mathrm{top}}=q_{\mathrm{top}}F_L$ となる．
左辺は $z$ の次数が2未満であり，monic除法の一意性に反する．
従ってこの和は直和であり，(4)の基底とassociated gradedの記述を得る．
中心性により二側idealも $A\star\Omega$ と一致する．

グラフ式の微分は活性変数だけに作用し，係数は固定した複素グラフ重みである．
三角変換とmonic除法はいずれも有限で，除算を要しない．
従って全構成は任意の係数変更と可換する．
noetherian性は，非負整数の補助次数filtrationのassociated gradedが
noetherian可換環であることと，filtered--gradedの標準的な帰納法から従う．
\end{proof}

\begin{remark}[通常の次数と有限パラメータ]
\label{rem:finite-ordinary-grading}
別に正整数の通常の次数 $(w_x,w_y,w_z)$ があり，
$F$ が通常の次数 $D=w_x+w_y+w_z=2w_z$ の斉次式であるとする．
係数環 $R$ も通常の次数を持ってよい．
このとき $\pi_F$ の通常の次数は零であり，
$\deg\hbar=0$ とすればstar積は通常の次数を保ち，$\Omega$ も次数 $D$ である．
補助次数は有限性とPBWのために用い，射影的な次数を置き換えない．
また，ここでの $\hbar$ の特殊化は，各式が有限多項式であることを証明した後に行うものである．
一般の形式級数へ $\hbar=1$ を代入する操作は用いていない．
\end{remark}

\subsection{分離型potentialの明示的な中心元}
\label{subsec:explicit-finite-potential}

グラフ式とは別に，関係式と中心元の全係数を有限に書ける模型を構成する．
両者の量子化が同じ形式座標で一致するとは仮定しない．
$R$ を任意の可換 $\C$-代数，$h\in R$ とし，
\[
 P(x)=\sum_{i=0}^{5}p_i x^i,\qquad
 Q(y)=\sum_{j=0}^{3}q_j y^j
\]
と置く．係数は全て中央とする．巡回potential
\[
 \Phi=xyz-yxz-h\bigl(z^2/2+P(x)+Q(y)\bigr)
\]
のJacobian代数を
\begin{equation}
 A_R=R\langle x,y,z\rangle/
 ([x,y]-hz,\ [y,z]-hP'(x),\ [z,x]-hQ'(y))
 \label{eq:finite-potential-relations}
\end{equation}
と書く．$\operatorname{Sym}$ は可換単項式を，その文字の全ての異なる順序の単語の平均へ送る
Weyl対称化であり，係数上線型に延長する．

\begin{theorem}[両次数の普遍PBW族と有限Casimir公式]
\label{thm:explicit-finite-pbw-casimir}
代数\eqref{eq:finite-potential-relations}は
$x^a y^b z^c$ ($a,b,c\ge0$) を $R$-基底に持つ．さらに
\begin{align}
 \Omega={}&\frac{z^2}{2}+P(x)+Q(y)
       -\frac{h^2}{12}\operatorname{Sym}\bigl(P''(x)Q''(y)\bigr)\notag\\
 &+\frac{2h^4}{3}p_5(3q_1q_3-q_2^2)x
       -\frac{3h^6}{2}p_5q_3^2P'(x)
 \label{eq:finite-casimir-explicit}
\end{align}
は中心元かつ左右非零因子である．商 $B_R=A_R/(\Omega)$ は
\[
 \{x^a y^b z^\epsilon\mid a,b\ge0,\ \epsilon=0,1\}
\]
を $R$-基底に持つ．これらの主張は任意の係数変更と両立し，
$h$ の冪零性や可逆性を仮定しない．$R$ がnoetherianなら両代数は左右noetherianである．

$p_5=0$ の場合，\eqref{eq:finite-casimir-explicit}の最後の二項は消える．
これが次数2の模型であり，$p_5$ を許す場合を次数1に用いる．
\end{theorem}

\begin{proof}
まず補助重みを $(x,y,z)=(6,10,15)$，$R$ の重みを0とする．
次の規則を，重みの後に $z>y>x$ の辞書式で比較する整礎的な単語順序に対して用いる：
\[
 yx\longmapsto xy-hz,\qquad
 zx\longmapsto xz+hQ'(y),\qquad
 zy\longmapsto yz-hP'(x).
\]
交換項は同じ重みで辞書式に下がり，他の項は重みが真に下がる．
唯一の重なり $zyx$ の二つの簡約は
\[
 xyz-hz^2+h\,yQ'(y)-hP'(x)x,
 \qquad
 xyz-hxP'(x)+hQ'(y)y-hz^2
\]
で一致する．包含重なりはない．従って各単語には一意な簡約形があり，
係数1のDiamond補題によって順序付き単項式が基底となる．
$p_5=0$ の場合は補助重み $(3,4,6)$ でも同じ証明が成立する．

中心性を有限な恒等式で確認する．
\[
 H=\frac{z^2}{2}+P+Q-\frac{h^2}{12}\operatorname{Sym}(P''Q''),\quad
 G=\frac{2h^4}{3}p_5(3q_1q_3-q_2^2)x
       -\frac{3h^6}{2}p_5q_3^2P'
\]
と置く．$Q'''=6q_3$ が中央であることと三つの関係式から
\[
 \left[x,\frac{z^2}{2}+P+Q\right]
       =-\frac{h^2Q'''}{12}[y,P'],\qquad
 [x,\operatorname{Sym}(P''Q'')]=-Q'''[y,P']
\]
を得るので $[x,H]=[x,G]=0$ である．
もう一つの交換子については
\begin{equation}
 [y,H]=\frac{2h^5}{3}p_5(3q_1q_3-q_2^2)z
          +\frac{3h^6}{2}p_5q_3^2[y,P']=-[y,G].
 \label{eq:finite-casimir-residual}
\end{equation}
これらは $P$ の次数5以下の各単項式と $Q$ の次数3以下の各単項式について，
上の三規則で展開すれば得られる多項式恒等式である．
対称化に現れる混合項は $\operatorname{Sym}(x^i y)$，$0\le i\le3$ だけであり，
\[
 \operatorname{Sym}(x^i y)=\frac{1}{i+1}\sum_{j=0}^{i}x^j yx^{i-j}
\]
を用いれば有限個の展開になる．
$p_5=0$ の場合の $y$-交換子は，$E=[x,Q']$，$\{u,v\}_+=uv+vu$ と書けば
\[
 [y,z^2/2+P+Q]=-h^2\bigl(p_3E/2+p_4\{x,E\}_+\bigr),
 \quad
 [y,\operatorname{Sym}(P''Q'')]=-6p_3E-12p_4\{x,E\}_+
\]
で相殺する．$p_5x^5$ の追加分を同じ規則で簡約した残差が
\eqref{eq:finite-casimir-residual}である．
全係数を独立変数にした有理数係数の展開は
\texttt{computations/verify\_finite\_families.py}にも収録した．

以上より $[x,\Omega]=[y,\Omega]=0$．普遍係数環
$\mathbb Q[h,p_0,\ldots,p_5,q_0,\ldots,q_3]$ 上ではPBW基底により $h$ は非零因子なので，
\[
 h[z,\Omega]=[[x,y],\Omega]=0
\]
から $[z,\Omega]=0$ を得る．この普遍多項式恒等式を特殊化するので，
$h$ にtorsionがある係数環でも中心性が成立する．

補助次数に関して全ての補正項は30未満であり，最高記号は
\[
 g=z^2/2+p_5x^5+q_3y^3
\]
である．$p_5=0$ の模型を重み $(3,4,6)$ で見ると，最高記号は
$g=z^2/2+p_4x^4+q_3y^3$ となる．
どちらも $2g$ は $z$ にmonicなので，任意の $R$ 上で正則である．
従って $\Omega$ の乗法はstrictかつ単射であり，中央性から
\[
 \operatorname{gr}B_R=R[x,y,z]/(g)
\]
となる．monic除法の基底を有限なfiltration帰納法で持ち上げれば，
主張の商基底を得る．同じ証明が係数変更後にも成立する．
noetherian性は可換なassociated gradedから従う．
\end{proof}

\begin{remark}[通常の射影的次数]
\label{rem:explicit-finite-internal-grading}
$R=\Lambda[s_1,\ldots,s_{n-1}]$ とし，$\deg s_i=1$，$\deg h=0$ とする．
次数2では $p_5=0$，$(\deg x,\deg y,\deg z)=(1,1,2)$，
$\deg p_i=4-i$，$\deg q_j=4-j$ とする．
次数1では $(\deg x,\deg y,\deg z)=(1,2,3)$，
$\deg p_i=6-i$，$\deg q_j=6-2j$ とする．
各係数は指定次数の $s$ の斉次式であり，その係数と $h$ を
$\Lambda$ の有限個の独立変数に取れる．
\eqref{eq:finite-casimir-explicit}はそれぞれ通常の次数4，6で斉次となる．
従ってこれらは全 $n\ge2$ に対する有限表示次数付き代数の有限型係数族である．
$h=0$ の商は $z^2/2+P+Q=0$，一次commutatorはこの式のJacobian括弧である．
\end{remark}

\subsection{全次元の次数1・2への適用}

$k=n-1\ge1$，$P=\C[s_1,\ldots,s_k]$ とし，$s_i$ の通常の次数を1，補助次数を0とする．
係数をさらに有限個のパラメータで動かす場合は，その多項式環 $P_0$ を付けて
$R=P_0[s_1,\ldots,s_k]$ とする．

\begin{corollary}[両次数の有限型係数族]
\label{cor:finite-dp-coefficient-families}
次の二種類の斉次式に定理\ref{thm:negative-auxiliary-algebraization}を適用できる．
\begin{enumerate}
\item 次数2では通常の重みを
$(\deg x,\deg y,\deg z)=(1,1,2)$，補助重みを $(3,4,6)$ とし，
\[
 F=z^2/2+a_0x^4+\sum_{i+j\le3}f_{ij}(s)x^iy^j,
 \qquad \deg f_{ij}=4-i-j.
\]
$a_0$ は通常の次数零の係数である．
この場合 $L=12$，$\delta=1$ である．
\item 次数1では通常の重みを $(1,2,3)$，補助重みを $(6,10,15)$ とし，
\[
 F=z^2/2+y^3+A_4(x,s)y+A_6(x,s),
\]
ただし $A_i$ は通常の次数 $i$ の斉次式で，
$\deg_xA_4\le3$，$\deg_xA_6\le5$ とする．
この場合 $L=30$，$\delta=1$ である．
\end{enumerate}
全ての係数と全ての $\hbar$ に対して中心商は次数ごとに係数環上自由であり，
通常のHilbert級数はそれぞれ
\[
 H_{n,2}(t)=\frac{1-t^4}{(1-t)^{n+1}(1-t^2)},
 \qquad
 H_{n,1}(t)=\frac{1-t^6}{(1-t)^n(1-t^2)(1-t^3)}
\]
である．全係数を独立な座標とすると，これは有限個のパラメータによる
有限表示次数付き代数とその中心商の族となる．
\end{corollary}

\begin{proof}
次数2では $3i+4j\le12$ が $i+j\le3$ の全ての組と $(i,j)=(4,0)$ について成立する．
次数1では $y^3$ と $x^5$ の補助次数は30，
$x^iy$，$i\le3$，の補助次数は28以下である．
従って定理の仮定を満たす．
通常の次数に関しても注意\ref{rem:finite-ordinary-grading}の条件が成立する．
基底を $s^\alpha x^iy^jz^\epsilon$，$\epsilon=0,1$，で数えればHilbert級数を得る．
\end{proof}

\begin{proposition}[滑らかな可換点の非空性]
\label{prop:finite-family-smooth-points}
両次数・全 $n\ge2$ について，系\ref{cor:finite-dp-coefficient-families}の
$\hbar=0$ のファイバーには滑らかなdel Pezzo $n$-foldが存在する．
具体例は
\begin{align*}
 F_{n,2}&=z^2/2+x^4+y^3s_1+\sum_{i=1}^{n-1}s_i^4,\\
 F_{n,1}&=z^2/2+y^3+x^5s_1+\sum_{i=1}^{n-1}s_i^6
\end{align*}
で与えられる．従ってそれぞれの係数空間には，滑らかな可換点からなる非空なZariski開集合がある．
\end{proposition}

\begin{proof}
次数2では $F_z=z$，$F_x=4x^3$，
$F_y=3y^2s_1$，$F_{s_1}=y^3+4s_1^3$，
$F_{s_i}=4s_i^3$ ($i\ge2$) である．
これらが全て零なら全座標が零である．
次数1でも $y,z,s_i$ ($i\ge2$) に関する偏微分からそれらの座標が零となり，
$5x^4s_1=x^5+6s_1^5=0$ から $x=s_1=0$ を得る．
従って両方のaffine coneは原点外で滑らかである．
次数2の重み付き射影空間の非自明な安定化群を持つ点は $z$ 軸上にあり，
次数1では $y$ 軸または $z$ 軸上にある．
どちらの超曲面もこれらの点を含まないので，通常の射影多様体として滑らかである．
重み付き超曲面の随伴公式と次数公式により
$-K_X=(n-1)H$，$H^n=2$ または1となる．
最後の主張は，固定された重み付き射影空間内の超曲面の滑らかさが係数の開条件であることによる．
\end{proof}

\subsection{分離型有限パラメータ模型の相対bimodule CY性}
potentialから対角複体を作る構成と巡回Hessianの対称性は
Ginzburgの構成である\cite[Proposition 5.1.9, Theorem 5.3.1,
Proposition 8.1.3(ii)]{Ginzburg2007}．
以下では任意の可換係数環を扱うため，複体・完全性・自己双対性を直接記述する．
可換 $\mathbb C$-代数 $R$ と $h\in R$ に対し，
\[
 \Phi=xyz-yxz-h\bigl(z^2/2+P(x)+Q(y)\bigr),\qquad
 A=R\langle x,y,z\rangle/(\partial_x\Phi,\partial_y\Phi,\partial_z\Phi)
\]
と置く．以下では $\deg P\le4$, $\deg Q\le3$ と補助重み $(3,4,6)$，
または $\deg P\le5$, $\deg Q\le3$ と補助重み $(6,10,15)$ を用いる．
$R$ の補助重みは0である．PBW定理により $A$ は $R$ 上自由で，
補助随伴次数環は $R[x,y,z]$ である．

内部次数も考える場合，$(x,y,z)$ の内部重みを $(a,b,c)$ とし，
$m=a+b+c$，$\Phi$ は $R$ の次数を含めて内部斉次次数 $m$ と仮定する．
対象となる二つの場合は
\[
 (a,b,c;m)=(1,1,2;4),\qquad (1,2,3;6)
\]
である．内部次数を忘れた主張にはこの斉次性の仮定は不要である．

\begin{proposition}\label{prop:finite-separated-relative-cy}
$A$ は $R$ 上homologically smoothであり，
\[
 A^!_R:=R\!\operatorname{Hom}_{A\otimes_R A^{\mathrm{op}}}
          (A,A\otimes_R A^{\mathrm{op}})
       \simeq A(m)[-3]
\]
という相対bimodule CY同型を持つ．
この同型は選んだ生成元の順序とpotentialから構成され，
任意の係数変更と両立する．$h$ は任意であり，冪零性を仮定しない．
\end{proposition}
\begin{proof}
$x_1=x,x_2=y,x_3=z$，$w_1=a,w_2=b,w_3=c$，
$r_i=\partial_{x_i}\Phi$ とする．
自由代数の非可換微分を
\[
 \Delta_j(u)=\sum_{u=vx_jw}v\otimes w^{\mathrm{op}}
\]
で定め，$R$-線型に延長する．和は単語中の $x_j$ の各出現を走る．
以下ではその $A^e_R=A\otimes_R A^{\mathrm{op}}$ への像を用いる．
$\delta_i=x_i\otimes1-1\otimes x_i^{\mathrm{op}}$ と置く．

有限自由 $A^e_R$-加群からなる列
\[
 0\longrightarrow A^e_R(-m)
 \xrightarrow{d_3}
 \bigoplus_i A^e_R(-(m-w_i))
 \xrightarrow{d_2}
 \bigoplus_i A^e_R(-w_i)
 \xrightarrow{d_1} A^e_R
 \xrightarrow{\mu} A\longrightarrow0
\]
を考える．$e_i,f_i,\omega$ をそれぞれ第1，第2，第3項の基底とし，
\[
 d_1(e_i)=\delta_i,\qquad
 d_2(f_i)=\sum_j\Delta_j(r_i)e_j,\qquad
 d_3(\omega)=\sum_i\delta_i f_i
\]
と定める．ここで $u\otimes v^{\mathrm{op}}$ はbimoduleの元に
$u(-)v$ として作用する．$d_1d_2=0$ は $r_i=0$ と
\[
 \sum_j\Delta_j(r_i)\delta_j
       =r_i\otimes1-1\otimes r_i^{\mathrm{op}}
\]
から従う．また巡回微分の恒等式
$\sum_i[x_i,r_i]=0$ を非可換微分すると，関係式 $r_i=0$ の下で
$\sum_i\delta_i\Delta_j(r_i)=0$ を得るので $d_2d_3=0$ である．
本件の $d_2$ の係数行列 $H_{ji}=\Delta_j(r_i)$ は特に
\[
 H=\begin{pmatrix}
 -h\Delta_xP'(x)&\delta_z&-\delta_y\\
 -\delta_z&-h\Delta_yQ'(y)&\delta_x\\
 \delta_y&-\delta_x&-h
 \end{pmatrix}
\]
である．$\Delta_xP'$，$\Delta_yQ'$ はそれぞれ左右の同じ生成元に関する
差商であり，両tensor因子の交換に不変である．

この複体に補助filtrationを入れる．$x_i$ の補助重みを $v_i$，
$W=v_1+v_2+v_3$ とし，$e_i,f_i,\omega$ の補助次数を
それぞれ $v_i,W-v_i,W$ とする．
$xyz-yxz$ の補助次数は $W$，残るpotentialの項は高々 $W-1$ である．
従って随伴次数で $r_i$ は順に $[y,z],[z,x],[x,y]$ となり，
上の複体は $R[x,y,z]$ の対角を解消するKoszul複体になる．
具体的には $R[x_L,y_L,z_L,x_R,y_R,z_R]$ の正則列
$x_L-x_R,y_L-y_R,z_L-z_R$ のKoszul複体である．
この列は任意の $R$ 上で正則である．
全てのfiltrationは下に有界で，元は有限filtration次数を持つため，
最高次数のcycleを随伴次数で境界として持ち上げ，次数を有限回下げることにより，
元の複体も完全である．従ってこれは長さ3の有限自由対角解消である．

最後に $\tau:A^e_R\to A^e_R$ を
$\tau(u\otimes v^{\mathrm{op}})=v\otimes u^{\mathrm{op}}$ とする．
これはanti-involutionであり，$\tau(\delta_i)=-\delta_i$．
巡回単語の二つの切断箇所を交換することから，巡回Hessianの恒等式
\[
 \tau\bigl(\Delta_j(\partial_i\Phi)\bigr)
       =\Delta_i(\partial_j\Phi)
\]
を得る．従って上の解消を $\operatorname{Hom}_{A^e_R}(-,A^e_R)$ で双対化し，
$\tau$ によって右加群を左加群に読み替えると，逆順の同じ複体を内部次数 $m$ だけ
ずらしたものになる．双対複体の次数 $0,1,2,3$ において
それぞれ符号 $(1,-1,-1,1)$ を入れれば上の微分規約と一致する．
よって双対複体のコホモロジーは次数3の $A(m)$ だけであり，主張の同型を得る．
非可換微分，巡回Hessian，補助filtrationとPBW基底は全て係数変更で保たれる．
変更後の任意の可換係数環に同じ完全性の証明を適用できるので，任意基底変更との両立も従う．
\end{proof}

\begin{proposition}\label{prop:finite-separated-absolute-cy}
$\Lambda$ を任意の可換 $\mathbb C$-代数，
$P_0=\Lambda[s_1,\ldots,s_k]$ とし，$\deg s_i=1$ とする．
前命題の係数環を $R=P_0$ とし，$P,Q$ の係数はこの中央変数の斉次式であるとする．
このとき $A$ は $\Lambda$ 上homologically smoothであり，
\[
 A^!_\Lambda\simeq A(m+k)[-3-k]
\]
である．従って $k=n-1$ のとき
\[
 A^!_\Lambda\simeq
 \begin{cases}
   A(n+3)[-n-2],&d=2,\\
   A(n+5)[-n-2],&d=1.
 \end{cases}
\]
これらは任意の $\Lambda$-代数への基底変更と両立する．
\end{proposition}
\begin{proof}
$E=A\otimes_\Lambda A^{\mathrm{op}}$，
$\overline E=A\otimes_{P_0}A^{\mathrm{op}}$ と置く．
$s_i$ は $A$ の中央元なので
\[
 \overline E=E/(t_1,\ldots,t_k),\qquad
 t_i=s_i\otimes1-1\otimes s_i^{\mathrm{op}}
\]
であり，各 $t_i$ は $E$ の中央元である．
PBW基底により $A$ は $P_0$ 上自由であるから，$E$ は
$P_0\otimes_\Lambda P_0$ 上自由である．
多項式環 $P_0\otimes_\Lambda P_0$ における差の列 $(t_i)$ は正則なので，
$E$ においても正則である．従って $\overline E$ は $E$ 上長さ $k$ の有限自由
Koszul解消を持ち，
\[
 R\!\operatorname{Hom}_{E}(\overline E,E)
      \simeq\overline E(k)[-k]
\]
を得る．この同型は $\overline E$--$E$ bimoduleの同型である．
前命題から $A$ は $\overline E$ 上perfectであり，
$\overline E$ 自体も $E$ 上perfectなので，$A$ は $E$ 上perfectである．
導来随伴と前命題によって
\begin{align*}
 R\!\operatorname{Hom}_E(A,E)
 &\simeq R\!\operatorname{Hom}_{\overline E}
       \bigl(A,R\!\operatorname{Hom}_E(\overline E,E)\bigr)\\
 &\simeq R\!\operatorname{Hom}_{\overline E}(A,\overline E)(k)[-k]\\
 &\simeq A(m+k)[-3-k]
\end{align*}
を得る．差のKoszul列の順序 $s_1,\ldots,s_k$ を固定すれば，
全ての同型は係数変更と両立する．
\end{proof}

\begin{remark}
上の結果は有限非零パラメータを含むambient代数のbimodule CY性である．
中央超曲面商のtailsのsmooth proper性やSerre公式を，この結果だけから結論しない．
\end{remark}

\subsection{有限非零パラメータの普遍中心とAbel族}
\label{subsec:finite-universal-center}

系\ref{cor:finite-dp-coefficient-families}の最高補助次数部分を次の形に固定する：
\[
 F_0=\begin{cases}
 z^2/2+x^4+s_1y^3,&d=2,\\
 z^2/2+y^3+s_1x^5,&d=1.
 \end{cases}
\]
残る許された低い補助次数の係数は自由に動かしてよい．
この規格化を施した係数空間を $T_d=\operatorname{Spec}\C[c_\alpha]$ と書く．
定理\ref{thm:explicit-finite-pbw-casimir}の分離型族でも，同じ規格化
$(p_4,q_3)=(1,s_1)$ または $(p_5,q_3)=(s_1,1)$ を用いる．
以下の中心の定理は，この二つの量子化のいずれにも成立する．

\begin{lemma}[最高記号のPoisson中心]
\label{lem:finite-symbol-poisson-center}
$P_0=\C[s_1,\ldots,s_{n-1}]$ とし，
$T=P_0[x,y,z]/(F_0)$ にJacobian括弧を入れる．
このとき $T$ は整域であり，Poisson中心は正確に $P_0$ である．
\end{lemma}
\begin{proof}
$K=\operatorname{Frac}P_0$ とする．$K$ の代数閉包上でも
$F_0$ は既約である．実際，$z$ に関する二次式が可約なら，残る $x,y$ の多項式が
平方でなければならないが，その $y$-次数は3である．
monic基底により $T\to T\otimes_{P_0}K$ は単射なので，$T$ も整域である．
$L=\operatorname{Frac}(T\otimes_{P_0}K)$ とすると，幾何学的整性から
$K$ は $L$ 内で相対代数閉である．Jacobian括弧は
$\{x,y\}=z\ne0$ なので，二次元の $L/K$ 上でrank2である．
Poisson中心の元 $f\in L$ は $df=0\in\Omega^1_{L/K}$ を満たし，
標数0なので $K$ 上代数的，従って $K$ に属する．
最後に，$x^ay^bz^\epsilon$，$\epsilon=0,1$，は $P_0$ 上でも $K$ 上でも基底なので，
$T\cap K=P_0$ である．
\end{proof}

\begin{theorem}[全ての有限非零パラメータにおける普遍中心]
\label{thm:finite-parameter-universal-center}
上の規格化を施した族について
$S=\C[c_\alpha,h,h^{-1}]$ と置く．任意の可換 $S$-代数 $U$ に対して
\[
 Z(B_S\otimes_S U)=U[s_1,\ldots,s_{n-1}]
\]
である．特に，各複素係数と各 $h\ne0$ で同じ中心公式が成立する．
この結論には追加の例外パラメータの除外はない．
\end{theorem}
\begin{proof}
まず $S$ の複素閉点を固定する．補助filtrationについて
\[
 \operatorname{gr}B_h=P_0[x,y,z]/(F_0),\qquad
 \operatorname{gr}_{-1}[u,v]=h\{\operatorname{gr}u,\operatorname{gr}v\}_{F_0}
\]
が成り立つ．グラフ族では二次以上のグラフ項は補助次数を2以上下げるためであり，
分離型族では三つの交換関係から直接従う．
中心元 $b$ の最高記号は，$h\ne0$ よりPoisson中心に属する．
補題\ref{lem:finite-symbol-poisson-center}から，この記号は補助次数0でなければならない．
基底の $x,y,z$ は正の補助重みを持つので，$b$ 自身が $P_0$ に属する．
逆包含は中央係数としての構成から従う．

任意の基底変更には，内部次数ごとの有限行列を用いる．
$C_m=(S[s])_m$ は基底の一部として $(B_S)_m$ の直和因子であり，
両者は有限自由 $S$-加群である．三生成元との交換子写像が誘導する
\[
 \overline c_m:(B_S)_m/C_m
 \longrightarrow\bigoplus_{u\in\{x,y,z\}}(B_S)_{m+\deg u}
\]
は，全ての複素閉点で単射である．従って，その列数と同じ大きさの小行列式が生成する
idealはどの極大idealにも含まれず，単位idealである．
各小行列式が可逆な開集合では写像はsplit injectionとなるので，
任意の $S\to U$ 後も単射である．
従って中心の内部次数 $m$ 部分は正確に $U[s]_m$ となる．
斉次な生成元と可換する条件は次数別に分かれるので，全中心について主張を得る．
次数ごとに開被覆を用いてよく，全次数に共通の新たな分母は不要である．
\end{proof}

\begin{corollary}[有限パラメータでのVeronese持上げの代数的仮定]
\label{cor:finite-parameter-veronese-input}
$n\ge3$ とする．定理\ref{thm:finite-parameter-universal-center}の族は，
二つ以上の中心一次変数を持ち，その中心多項式環上自由である．
従って定理\ref{thm:ncdp-veronese-lift}の普遍中心と平坦性に関する仮定を
有限非零パラメータ上で満たす．従ってこのような二つの族の任意の $r\ge2$ の
Veronese次数付き同型はfppf局所的に元の次数付き環へ持ち上がり，
その持上げ層は $\mu_r$-torsorである．中心の追加標識は不要である．
\end{corollary}
\begin{proof}
中心一次変数は $n-1\ge2$ 個であり，中心上の自由基底は
$x^ay^bz^\epsilon$，$\epsilon=0,1$，である．普遍中心は前定理による．
\end{proof}

\begin{proposition}[有限型の平坦Abel族]
\label{prop:finite-pbw-abel-family}
上のグラフ族または分離型族の係数環をnoetherianな $\C$-代数 $\Lambda$ とする．
$\QGr B_\Lambda$ と $\pi B_\Lambda$ は，それぞれ
$\Lambda$-flatなGrothendieck圏とflatな構造対象を与え，
許された任意の基底変更と両立する．
$h=0$ の滑らかな可換点では $(\Qcoh X,\OO_X)$ を与え，
その点におけるArtin完成は $\mathfrak M^{\OO,\mathrm{hyp}}_{n,d}$ に入る．
\end{proposition}
\begin{proof}
各次数が $\Lambda$ 上自由であり，$B_\Lambda$ は有限生成noetherian次数付き代数である．
補題\ref{lem:tails-arbitrary-base-change}と\ref{lem:tails-flat-deformation}の証明における
作用付き加群のGabriel商およびinjectiveのcoflat性の議論は，
ここでは冪零性を一切使わず適用できる．
可換点の同定には滑らかな重み付き模型の自由locusを用いる．
最後の主張は定理\ref{thm:all-artin-are-nc-dp}による．
\end{proof}

\subsection{先行研究との関係と適用範囲}

\begin{remark}[既存のfiltered量子化との関係]
Etingof--Ginzburg\cite[Proposition~7.3.1(ii), Corollary~7.4.1(ii)]{EG2010}は，
homogeneous CY potentialへの低次数potentialの付加がPBW変形を与えることを証明する．
同\cite[Proposition~8.1.1(i)]{EG2010}は，形式potentialのCasimir中心元の持上げを与える．
本節ではそれらと共通する負次数の機構をグラフ式に直接適用し，
中央多項式係数，任意の係数変更，中心商の自由基底を一つの証明で扱った．
負次数の量子化という一般原理の新規性は主張しない．
また，同\cite[Definition~3.4.3, Lemma~7.1.5]{EG2010}のgenericおよびCY条件は
可算個の開条件に基づくので，それだけから有限個の分母で定義される開集合を推論しない．
\end{remark}

\begin{proposition}[一般の楕円的最高次項に対する補助重みの障害]
\label{prop:elliptic-top-auxiliary-obstruction}
次数2の活性式が $x^4,y^4,z^2$ を全て非零係数で含む場合，
または次数1の活性式が $x^6,y^3,z^2$ を全て非零係数で含む場合，
正の補助重みに関してその式の次数を重みの和より厳密に小さくすることはできない．
\end{proposition}

\begin{proof}
次数2で $D=\max(4a,4b,2c)$ とすると，
$a+b+c\le D/4+D/4+D/2=D$ である．
次数1では $D=\max(6a,3b,2c)$ とし，
$a+b+c\le D/6+D/3+D/2=D$ を得る．
他の単項式を加えても式の補助次数は下がらない．
\end{proof}

次数2の上の族では $s=0$ へのrestrictionが $z^2/2+a_0x^4$ となる．
次数1では $s=0$ へのrestrictionが $z^2/2+y^3$ となる．
従ってこれらは特別な活性平面を持つ族であり，一般のPoisson方向を被覆するとは主張しない．
有限パラメータの商 $B$ が存在し平坦であることと，
そのtailsが本稿の $\mathfrak M^{\OO}_{n,d}$ の全条件を満たすことも区別する．
滑らかな可換点におけるArtin厚化には定理\ref{thm:all-artin-are-nc-dp}を適用できるが，
有限の非零 $\hbar$ における相対Serre，飽和性，smooth proper性は別途検証を要する．

% !TeX root = ../nc-del-pezzo.tex
\section{有限パラメータ族の局所コホモロジーと相対Serre双対}
\label{sec:finite-tails}

本節では $n\ge2$ とし，$\Lambda$ をnoetherianな可換 $\mathbb C$-代数とする．
定理\ref{thm:explicit-finite-pbw-casimir}の分離型有限パラメータ族のambientを $A$，中央元を $\Omega$，
$B=A/(\Omega)$ と書く．記号を
\[
 N=n+2,\qquad d=N-1=n+1,\qquad
 (l,D)=\begin{cases}(n+3,4)&\text{次数2},\\(n+5,6)&\text{次数1},\end{cases}
 \qquad r=l-D=n-1>0
\]
とする．定理\ref{thm:explicit-finite-pbw-casimir}，
命題\ref{prop:finite-separated-absolute-cy}によって，次の性質が得られている：
$A,B$ は左右noetherianで各内部次数が有限自由，$A_0=B_0=\Lambda$，
$A$ は $\Lambda$ 上homologically smoothで
\[
 A^!_\Lambda\simeq A(l)[-N],
\]
$\Omega$ は次数 $D$ の中央正則元であり，これらは任意の係数変更で保たれる．
以下の加群計算は可換係数環一般でも成り立つが，
coherent tailsの主張は本稿のsiteに合わせて局所noetherian基底に限定する．

次数付き $\Lambda$-複体の相対双対を
\[
 (D_\Lambda M)_j=R\!\operatorname{Hom}_\Lambda(M_{-j},\Lambda)
\]
と書く．各次数が有限自由の代数の場合には通常の次数別双対 $M^\vee$ と一致する．
左右の式は反対環にも適用する．

\subsection{有限な係数加群に対するambient双対}
\begin{lemma}\label{lem:relative-proper-module-duality}
$M$ を次数付き左 $A$-加群で，$\Lambda$ 上perfectなものとする．
このとき $M$ は $A$ 上perfectで，右 $A$-加群として自然な同型
\[
 R\!\operatorname{Hom}_A(M,A)
       \simeq D_\Lambda M(l)[-N]
\]
を持つ．$M$ がbimoduleなら，残る作用についても自然である．
同様に，$M$ が $B$-加群で $\Lambda$ 上perfectなら
\[
 R\!\operatorname{Hom}_B(M,B)
       \simeq D_\Lambda M(r)[-d].
\]
後者では $M$ が $B$ 上perfectであるとは主張しない．
\end{lemma}
\begin{proof}
$A$ の有限projective対角解消を右 $A$ 上で $M$ とtensorすると，
$M$ は有限個の $A\otimes_\Lambda M$ から作られるので $A$ 上perfectである．
対角の双対可能性，tensor--Hom随伴，$M$ の $\Lambda$-perfect性から
\begin{align*}
 R\!\operatorname{Hom}_A(M,A)
 &\simeq R\!\operatorname{Hom}_{A^e_\Lambda}
       \bigl(A,R\!\operatorname{Hom}_\Lambda(M,A)\bigr)\\
 &\simeq A^!_\Lambda\otimes^L_{A^e_\Lambda}
             (A\otimes_\Lambda D_\Lambda M)\\
 &\simeq D_\Lambda M(l)[-N]
\end{align*}
を得る．最後のtensorでは $A\otimes_\Lambda D_\Lambda M$ を
$A^{\mathrm{op}}$ からの誘導加群として見ればよい．
この議論はReyes--Rogalski--Zhangの体上の議論
\cite[Lemma 1.2]{RRZ2014}と同じ随伴であり，ここでは係数環についてのperfect性を明示した．

中央正則元の二項解消から
$R\!\operatorname{Hom}_A(B,A)\simeq B(D)[-1]$．
$B$-加群 $M$ についてのchange of ringsにより
\[
 R\!\operatorname{Hom}_A(M,A)
 \simeq R\!\operatorname{Hom}_B(M,B)(D)[-1]
\]
なので第二の式を得る．全ての同型は自然で，次数付きbimoduleの作用と両立する．
\end{proof}

\begin{corollary}\label{cor:finite-relative-augmentation-ext}
任意の許される基底変更後にも
\[
 R\!\operatorname{Hom}_A(\Lambda,A)=\Lambda(l)[-N],\qquad
 R\!\operatorname{Hom}_B(\Lambda,B)=\Lambda(r)[-d]
\]
が左右で成り立つ．幾何学的体ファイバーでは $A$ は大域次元 $N$，
AS index $l$ のAS正則代数であり，$B$ はinjective dimension $d$，
AS index $r$ のAS--Gorenstein代数である．
\end{corollary}
\begin{proof}
相対Ext式は前補題で $M=\Lambda$ としたものである．
体ファイバーのambientは連結次数付きCYなので
\cite[Lemma 1.2]{RRZ2014}によりAS正則である．
Hilbert級数からGK次元も有限である．中央元のchange of ringsを任意の $B$-加群に適用すると，
$\operatorname{injdim}_B B\le N-1$ を得る．
augmentation Extの非消滅から等号であり，左右のAS--Gorenstein条件が従う．
非正則な $\Lambda$ 上で絶対的な有限大域次元を主張するものではない．
\end{proof}

\subsection{相対局所コホモロジー}
$C$ を $A$ または $B$ とし，$\Gamma_C$ を内部次数のirrelevant idealによるtorsion関手とする．
\[
 R\Gamma_C(M)=\mathop{\mathrm{colim}}_q
       R\!\operatorname{Hom}_C(C/C_{\ge q},M)
\]
を用いる．ここでHomは全ての内部次数を含む．

\begin{theorem}\label{thm:finite-relative-local-cohomology}
次数付きbimoduleとして
\[
 R\Gamma_A(A)\simeq A^\vee(l)[-N],\qquad
 R\Gamma_B(B)\simeq B^\vee(r)[-d].
\]
さらに $\Gamma_B$ のcohomological dimensionは高々 $d$ で，
\[
 R\Gamma_B(M)\simeq R\Gamma_B(B)\otimes_B^L M
\]
が左加群に対して自然に成り立つ．右加群には左右を交換した式を用いる．
これらの関手は直和を保つ．構造加群 $A,B$ の式は任意の基底変更と両立する．
\end{theorem}
\begin{proof}
$A/A_{\ge q}$ と $B/B_{\ge q}$ は $\Lambda$ 上有限自由である．
前補題を適用して $q$ のcolimitを取ると，最初の二つの式を得る．
次数ごとに双対が安定するので，この計算は任意の基底変更と両立する．

$A/A_{\ge q}$ は $A$ 上perfectであるため，
\[
 R\Gamma_A(M)
 \simeq\left(\mathop{\mathrm{colim}}_q
 R\!\operatorname{Hom}_A(A/A_{\ge q},A)\right)\otimes_A^L M.
\]
従って $R\Gamma_A$ はsmashingであり，直和を保つ．
有限projective対角解消をtensorすることによって各 $A/A_{\ge q}$ のprojective dimensionには
$q$ によらない上界があるので，$\Gamma_A$ は有限cohomological dimensionを持つ．
上の具体式 $A^\vee(l)[-N]\otimes_A^LM$ によって，その上界を $N$ にできる．

次に $B$-加群上で $R\Gamma_A=R\Gamma_B$ であることを確認する．
$I$ を次数付きinjective $B$-加群とすると，change of ringsにより
\[
 \operatorname{Ext}^i_A(A/A_{\ge q},I)
 \cong\operatorname{Hom}_B
   \bigl(\operatorname{Tor}_i^A(B,A/A_{\ge q}),I\bigr).
\]
$\operatorname{Tor}_i$ は $i>1$ で零であり，$i=1$ の項は
$\Omega$ の正則性から内部次数 $q,\ldots,q+D-1$ にのみ支持を持つ．
$q'\ge q+D$ なら切断写像が誘導する $\operatorname{Tor}_1$ の写像は零である．
従って $q$ のcolimitを取ると $I$ は $\Gamma_A$-acyclicとなる．
両torsion関手は $B$-加群上で一致するので，injective $B$-解消により
導来関手も一致する．まず下に有界な複体で比較し，有限cohomological dimensionにより
非有界導来圏へ延長する．

従って導来tensorの結合性から
\begin{align*}
 R\Gamma_B(M)
 &\simeq R\Gamma_A(A)\otimes_A^L M\\
 &\simeq\bigl(R\Gamma_A(A)\otimes_A^L B\bigr)\otimes_B^L M\\
 &\simeq R\Gamma_B(B)\otimes_B^L M.
\end{align*}
これは直和を保つ．$B^\vee(r)[-d]$ の式から，離散加群 $M$ の局所コホモロジーは
次数 $d$ より上で零である．

後で用いる局所化の冪等性も確認できる．内部次数が上に有界なtorsion加群 $T$ に対し，
$A/A_{\ge q}$ の自由解消を，零次項は $A$，それ以外の項の生成次数は全て $q$ 以上となるように取る．
固定した内部次数では，$q\to\infty$ により正の解消次数から $T$ へのHomは零になるので，
$R\Gamma_A(T)=T$ である．一般のtorsion加群は，内部次数が上に有界な有限生成torsion部分加群の
filtered colimitである．$A/A_{\ge q}$ のperfect性から $R\Gamma_A$ はfiltered colimitを保つため，
一般のtorsion加群にも同じ結論が成立し，有限cohomological dimensionにより
torsionコホモロジーを持つ複体にも延長される．$B$ に対しては比較同型を使う．
また $R\Gamma_B(M)$ の各コホモロジーはtorsionである：
切断からのExtにおける残りの $B$-作用は $B/B_{\ge q}$ を経由するためである．
従って $R\Gamma_B$ はtorsion導来部分圏への冪等なcolocalizationであり，
その局所化三角は導来商 $D(\operatorname{QGr}B)$ の局所化三角と一致する．
\end{proof}

\begin{corollary}[相対局所双対]\label{cor:finite-relative-local-duality}
自然な同型
\[
 D_\Lambda R\Gamma_B(M)
       \simeq R\!\operatorname{Hom}_B(M,B(-r)[d])
\]
が成り立つ．ここでも右加群には左右を交換した式を用いる．
これは相対局所双対であり，非正則な $\Lambda$ 上での絶対balanced dualizing complexの
全公理を意味するとは限らない．
\end{corollary}
\begin{proof}
前定理のsmashing表示に導来tensor--Hom随伴を適用する．
$B$ の各次数は有限自由なので
$D_\Lambda(B^\vee(r)[-d])=B(-r)[d]$ である．
\end{proof}

\subsection{飽和性，ACM，豊富列，compact生成}
$\mathcal Q=\operatorname{QGr}B$，$\mathcal O(t)=\pi B(t)$ とする．
$\pi$ の右随伴を $\omega$ と書く．局所コホモロジーと導来sectionの三角
\[
 R\Gamma_B(M)\longrightarrow M\longrightarrow R\omega\pi M
\]
を用いる．これは次数付きnoetherian torsion商の通常の局所化三角である．

\begin{theorem}\label{thm:finite-tails-cohomology}
任意の整数 $t$ について
\[
 R\!\operatorname{Hom}_{\mathcal Q}(\mathcal O,\mathcal O(t))
 \simeq
 \begin{cases}
  B_t,&t\ge0,\\
  \operatorname{Hom}_\Lambda(B_{-t-r},\Lambda)[-n],&t\le-r,\\
  0,&-r<t<0.
 \end{cases}
\]
従って切断環は全次数で飽和し，有向ACM条件が成立する．
これらのRHomは $\Lambda$ 上perfectで，任意の許される基底変更と両立する．
\end{theorem}
\begin{proof}
局所化三角の内部次数 $t$ を取る．
$R\Gamma_B(B)$ の唯一のコホモロジーは次数 $d=n+1$ にあり，
その内部次数 $t$ 部分は $\operatorname{Hom}_\Lambda(B_{-t-r},\Lambda)$．
$r>0$ により $B_t$ とこの双対は同時に非零にならない．
従って各 $t$ の複体は上記の一つのコホモロジー群だけを持つ．
基底変更は有限自由加群の双対と局所コホモロジーの構成に両立する．
\end{proof}

\begin{proposition}\label{prop:finite-full-shift-ample}
任意のnoetherian次数付き $B$-加群 $M$ について，
$H^i(\mathcal Q,\pi M(t))$ は有限 $\Lambda$-加群であり，
$i>0$ では $t\gg0$ に対して零である．
$B$ は左右で相対 $\chi$ 条件を満たす：全ての $q,i$ について
$\underline{\operatorname{Ext}}^i_B(B/B_{\ge q},M)$ は有限 $\Lambda$-加群である．
従って全シフト列 $(\mathcal O(-j))_{j\in\mathbb Z}$ は相対豊富列であり，
この性質は許される基底変更後も成立する．
\end{proposition}
\begin{proof}
$M$ の各項が有限生成次数付き自由加群である解消 $F^\bullet$ を選ぶ．
局所コホモロジーは $B^\vee(r)[-d]\otimes_BF^\bullet$ で計算できる．
固定したコホモロジー次数の項は有限個の $B^\vee$ の次数シフトの和なので，
各内部次数で有限 $\Lambda$-加群であり，内部次数について右有界である．
その部分商にも同じ性質がある．$i>0$ のsection cohomologyは
$H^{i+1}_{\ge1}(M)$ から計算されるので有限性と消滅を得る．
$H^0$ の有限性は局所化の完全列から従う．

$T_q=B/B_{\ge q}$ はtorsionなので，導来商の随伴から
\[
 R\!\underline{\operatorname{Hom}}_B(T_q,M)
 \simeq R\!\underline{\operatorname{Hom}}_B(T_q,R\Gamma_B(M)).
\]
$T_q$ の有限型自由解消と，有限範囲 $0\le j\le d$ にのみ支持を持つ
$H^jR\Gamma_B(M)$ に対するスペクトル系列を用いると，このExtは
内部次数について右有界で，各内部次数が有限 $\Lambda$-加群となる．
一方，左辺を $M$ から直接計算すれば，$M$ の内部次数が左有界なので
同じExtは左有界でもある．従って全内部次数の和が有限 $\Lambda$-加群であり，
相対 $\chi$ 条件を得る．反対環でも同じである．

任意の大きい $q$ について $M_{\ge q}$ は有限生成であり，
$\pi M_{\ge q}=\pi M$ だから，十分大きい添字からの有限和による全射を得る．
全射に対するHomの全射性は，その核への $H^1$ の消滅から従う．
任意の局所noetherian係数変更後も同じ証明を適用できる．
これは全シフトの豊富性であり，任意の正周期部分列の豊富性は別途確認する．
\end{proof}

\begin{proposition}\label{prop:finite-tails-compact-proper}
$\mathcal O(j)$ は $D(\mathcal Q)$ のcompact生成系である．
従ってそのthick閉包は $D(\mathcal Q)^c$ に一致し，このperfect圏は
$\Lambda$ 上properである．許される係数変更はcompact対象を保ち，
compact対象間のRHomは基底変更と可換する．
\end{proposition}
\begin{proof}
$R\Gamma_B$ が直和を保つので，局所化三角から $R\omega\pi$ も直和を保つ．
$\pi R\omega=\mathrm{Id}$ を用いると $R\omega$ 自体も直和を保つ．
$B(j)$ は $D(\operatorname{Gr}B)$ でcompactなので，随伴によって $\mathcal O(j)$ はcompact．
これら全てからの導来Homが零なら $R\omega$ を施した対象は $D(\operatorname{Gr}B)$ で零，
従って元の対象も零である．これがcompact生成を与える．
生成対象間のRHomは前定理で有限自由の一つのシフトとして計算した．
有限coneと直和因子に閉じることにより，任意のcompact対のRHomは $\Lambda$ 上perfectとなる．
係数変更は $\mathcal O(j)$ を対応する $\mathcal O_T(j)$ へ送り，
生成対象間で既にRHom比較が同型なので，全compact対へ拡張する．
\end{proof}

\subsection{有限個のtwistからなる生成元}
\begin{proposition}[有限windowとproperなdg模型]
\label{prop:finite-window-proper}
有限直和
\[
 E=\bigoplus_{j=0}^{l-1}\mathcal O(j)
\]
は $D(\mathcal Q)$ のcompact生成元であり，
\[
 D(\mathcal Q)^c=\operatorname{thick}(E).
\]
従って $C=R\!\operatorname{End}_{\mathcal Q}(E)$ は $\Lambda$-perfectなdg代数で，
$\operatorname{Perf}(C)\simeq D(\mathcal Q)^c$ である．
任意の許されたnoetherian係数変更 $\Lambda\to T$ に対し
\[
 C\otimes_\Lambda^L T
   \simeq R\!\operatorname{End}_{\operatorname{QGr}B_T}(E_T)
\]
が成り立つ．
\end{proposition}
\begin{proof}
命題\ref{prop:finite-tails-compact-proper}により，全てのtwistはcompact生成系である．
従って有限windowから全てのtwistが生成されることを示せばよい．
$\Lambda$ の有限自由右 $A$-分解を，その内部次数も指定して構成する．
中央変数 $s_i$ の補助重みを1，活性変数の補助重みを
次数2では $(6,8,12)$，次数1では $(12,20,30)$ とする．
これは中央変数も正の重みを持つ，この段落だけの補助filtrationである．
係数 $p_i(s),q_j(s)$ の内部次数を含めて数えると，
$z^2/2+P+Q$ の最高補助次数はそれぞれ25，61以下であり，
活性三変数の重みの和26，62より小さい．
従って全ての交換子関係の補正項は補助次数を真に下げる．
既知のPBW基底により
\[
 \operatorname{gr}_{\mathrm{aux}} A
   =\Lambda[s_1,\ldots,s_{n-1},x,y,z]
\]
であり，この等式は内部次数も保つ．

右辺のaugmentationのKoszul分解をfilteredに持ち上げる．
その自由加群の基底と内部shiftを固定し，微分を順に持ち上げる．
微分の合成に残る誤差の最高記号はKoszul複体のcycleであり，
その完全性を用いて同じ内部次数のboundaryを引けば，誤差の補助次数が下がる．
全ての変数の補助重みは正なので，この補正は有限回で終わる．
こうして得られた複体のcycleにも同じ最高記号の除去を行えば完全性が従い，
最後の項の単射性もassociated gradedの単射性から従う．
Koszul分解は任意の $\Lambda$ 上で完全であり，
この持上げに $\Lambda$ の体性や正則性は不要である．
まず係数を独立変数とする普遍係数環上でこの有限回の持上げを行う．
任意の係数変更後にもそのassociated gradedは同じKoszul分解であるため，
同じ最高記号の帰納法が完全性を示す．
従って一度構成した有限自由分解は任意の係数変更に対して解消のままである．
得られた有限自由分解 $K_\bullet$ のshiftは生成元の内部重みの部分和のままである．
従って両端は $A(-l)$ と $A$ がそれぞれ一つであり，
他の全ての項は $1\le e\le l-1$ を満たす $A(-e)$ の有限直和である．

中心正則元の二項分解
$0\to A(-D)\xrightarrow{\Omega}A\to B\to0$ を用いると，
\[
 \operatorname{Tor}^A_i(\Lambda,B)=
 \begin{cases}
 \Lambda&i=0,\\
 \Lambda(-D)&i=1,\\
 0&i\ne0,1.
 \end{cases}
\]
ここでは $\Omega$ の正の内部次数によりaugmentation上の作用が零である．
両方のTorは次数付きtorsion加群なので，
$\pi(K_\bullet\otimes_A B)$ は有限の完全複体になる．
従って $\mathcal O(-l)$ は
$\mathcal O(-l+1),\ldots,\mathcal O$ のthick閉包に属し，
同様に $\mathcal O$ は他の端点と中間項のthick閉包に属する．
この完全複体をtwistし，両方向に帰納すれば，
長さ $l$ の任意の連続したtwistの組から全てのtwistが生成される．
特に $E$ はcompact生成元である．
compact対象のthick生成とdg Morita同値は
Kellerの定理\cite[\S4.3,\S8.2]{Keller1994}を適用して得られる．
$C$ の相対perfect性と基底変更は命題\ref{prop:finite-tails-compact-proper}を
有限和 $E$ に適用したものである．

\end{proof}

\subsection{perfect圏上の相対Serre双対}
\begin{theorem}\label{thm:finite-tails-relative-serre}
$\operatorname{Perf}(\mathcal Q/\Lambda)=D(\mathcal Q)^c$ 上の相対Serre関手は
\[
 S_\Lambda\simeq(-r)[n]
\]
である．この同型は許される任意の基底変更と両立する．
この主張に $\mathcal Q$ のsmoothnessは不要である．
\end{theorem}
\begin{proof}
まず左と右のtorsion複体を $E_\ell,E_r$ と書く．
両者のコホモロジーは $B^\vee(r)[-d]$ なので，左右どちらについてもtorsionである．
左右のsmashing表示を使うと，随伴の自然な写像は
\[
 E_\ell\otimes_B^L E_r\xrightarrow{\sim}E_\ell,
 \qquad E_\ell\otimes_B^L E_r\xrightarrow{\sim}E_r
\]
を与える．これらは $B$ への写像と両立するので，左右共通の
torsion bimoduleを $E$ と書ける．局所化bimoduleを
\[
 E\longrightarrow B\longrightarrow L
\]
で定める．すると左右とも $R\omega\pi$ は $L$ との導来tensorであり，
$L\otimes_B^LL\simeq L$ である．torsion複体の導来部分圏は
このsmashing局所化の核に一致する．
$L$ は対応する局所化dg代数として取ることができ，
$\mathcal O(i)$ 全体のdg部分圏のHom複体は
\[
 R\!\operatorname{Hom}_{\mathcal Q}(\mathcal O(i),\mathcal O(j))
       =L_{j-i}
\]
であり，合成は $L$ の積で与えられる．

局所化三角を次数別に $\Lambda$-双対化すると，bimoduleの三角
\[
 D_\Lambda L\longrightarrow B^\vee
       \longrightarrow B(-r)[d]
\]
を得る．$D_\Lambda L$ は左右ともlocalである．実際，任意のtorsion複体 $T$ に対し
\[
 R\!\operatorname{Hom}_B(T,D_\Lambda L)
       \simeq D_\Lambda(L\otimes_B^LT)=0
\]
であり，反対側も同じである．従って上の三角に左右から $L$ をtensorすると，
第一項はそのまま残り，第二項は
$B^\vee=E(-r)[d]$ なので消える．これにより，局所化bimoduleとしての同型
\[
 D_\Lambda L\simeq L(-r)[d-1]=L(-r)[n]
\]
を得る．このbimoduleの同型は上のdg部分圏における合成と両立するので，
単にHom群の次元が一致するだけではなく，coherentなSerre対を与える．
compact生成とproper性から，そのperfect加群へ拡張すると，任意のcompact対象 $E',F$ について
\[
 R\!\operatorname{Hom}_{\mathcal Q}(E',F)^\vee_\Lambda
       \simeq R\!\operatorname{Hom}_{\mathcal Q}(F,E'(-r)[n])
\]
を得る．$(-r)[n]$ は自己同値なので相対Serre関手である．
生成系のRHomは次数ごとに有限自由の一つのシフトであり，
局所化三角・双対・traceはその基底変更と両立するので，この双対も基底変更と両立する．

体上のAS--Gorenstein代数に対する局所双対とprojective Serre双対は
\cite[Theorem 4.2, Corollary 4.3]{YZ1997}の範囲にある．
ここではその絶対dualizing定理を非正則基底に適用する代わりに，
相対CYからの局所コホモロジーと局所化bimoduleを直接計算した．
\end{proof}

\subsection{正周期の生成を有限個の次数で検査する}
\label{subsec:finite-periodic-generation}

ここでは基本シフトから反標準シフトへ移る際の生成条件を扱う．
重み付き代数では，すべてのシフトを使う豊富性から，
固定周期の部分列の豊富性を無条件に推論することはできない．
次の命題は，その生成条件を有限個の行列の全射性に帰着する．

\begin{proposition}[正周期の生成条件の開性]
\label{prop:finite-periodic-generation}
$S$ を有限型可換 $\C$-代数，$B=\bigoplus_{m\ge0}B_m$ を
$B_0=S$ を満たす次数付き $S$-代数とする．各 $B_m$ は有限表示
$S$-加群であり，$B$ は重み $1\le w_\ell\le w$ の有限個の
斉次生成元 $b_\ell$ により生成されるとする．$r\ge1$ を固定する．
$s\in\Spec S$ を複素閉点とし，$B_s$ は右noetherianと仮定する．
$\operatorname{qgr}B_s$ において，
各 $j=0,\ldots,r-1$ に対し有限個の整数 $a_{j\ell}$ と全射
\[
 \bigoplus_\ell\mathcal O_s(-r a_{j\ell})
 \longrightarrow\mathcal O_s(-j)
\]
が存在するとする．これらの射は次数 $r a_{j\ell}-j$ の
$B_s$ の元から誘導され，$a_{0\ell}>0$ と仮定する．
このとき $s$ のZariski開近傍 $V$ 上へこれらの射を持ち上げられ，
任意の係数変更後にも $\QGr$ における全射となる．
任意の許されたnoetherian基底変更後でも，任意の $\mathcal O(-u)$ は，
任意に大きい $a$ に対する $\mathcal O(-ra)$ の有限和の商となる．
\end{proposition}

\begin{proof}
射を定める有限個の元を対応する $B_{r a_{j\ell}-j}$ へ持ち上げ，
次数付き右加群の写像
\[
 f_j:\bigoplus_\ell B(-r a_{j\ell})\longrightarrow B(-j)
\]
を作る．$Q_j=\operatorname{coker} f_j$ と置く．
各次数 $(Q_j)_m$ は有限表示 $S$-加群である．
特殊ファイバーの $Q_j\otimes_S\C(s)$ は有限生成torsion加群だから，
ある $D>r$ に対して全ての $j$ と $m\ge D$ でその次数 $m$ 部分が零である．
有限個の加群 $(Q_j)_m$，$D\le m<D+w$，の支持を除くと，
$s$ を含む開集合 $V$ 上でこれらの加群自身が零となる．
$Q_j$ は次数 $j$ の一つの元で生成される右 $B$-加群であるから，
$m>j$ に対して
\[
 (Q_j)_m=\sum_\ell (Q_j)_{m-w_\ell}b_\ell
\]
である．従って $m$ について帰納すると $m\ge D$ で全て零となる．
この消滅は任意の係数変更で保たれる．したがって各 $f_j$ は
$\QGr$ において全射を誘導する．

$j=0$ の全射を次数 $-rt$ だけ移動すると，
$\mathcal O(-rt)$ を $\mathcal O(-r(t+a_{0\ell}))$ の有限和で覆える．
$a_{0\ell}>0$ なのでこれを反復すれば，始点の添字を任意に大きくできる．
一般の $u$ を $u=qr+j$，$0\le j<r$ と書き，$j$ の全射を $-qr$ だけ
移動してからこの反復を行えば最後の主張を得る．
\end{proof}

\begin{corollary}[有限開近傍上の相対反標準豊富性]
\label{cor:finite-periodic-ampleness}
前命題の $V$ 上で，許された全てのnoetherian係数変更 $T$ に対し
$B_T$ が右noetherianであり，全シフトについてcoherent対象のSerre消滅を持つとする．
このとき，全てのそのような $T$ において
$(\mathcal O_T,(r))$ は豊富な対である．

特に，分離型有限パラメータ族の有限型係数空間で，$s$ が $h=0$ の
滑らかな可換del Pezzo模型を与えるとする．このとき $s$ のZariski開近傍 $V$ 上で
反標準列 $(\mathcal O,(n-1))$ は相対的に豊富であり，
許された任意のnoetherian基底変更後も豊富である．
\end{corollary}
\begin{proof}
任意のそのような係数変更 $T$ を固定する．各coherent対象は
有限個の $\mathcal O_T(-u)$ の和の商であり，
命題\ref{prop:finite-periodic-generation}の普遍的な全射を反復すれば，
任意に大きい $a$ に対する $\mathcal O_T(-ra)$ の有限和で覆える．
これが豊富性の第一条件である．
coherent対象の全射の核を $K$ とすると，
\[
 \operatorname{Ext}^1(\mathcal O_T(-ra),K)=H^1(K(ra))=0
 \qquad(a\gg0)
\]
であるから，十分大きい $a$ で
$\operatorname{Hom}(\mathcal O_T(-ra),-)$ はその全射を保つ．
これが第二条件であり，各基底変更後のSerre消滅を適用したので，
主張はファイバーだけに限られない．

分離型族では命題\ref{prop:finite-full-shift-ample}が
任意のnoetherian基底変更後のSerre消滅を与える．
滑らかな可換点では $\mathcal O_X(r)$ が豊富なので，
各 $\mathcal O_X(ra-j)$ は $a\gg0$ で大域生成される．
定理\ref{thm:finite-tails-cohomology}の飽和性により，
これらの切断は $B_{ra-j}$ であり，前命題の有限個の全射を選べる．
有限開近傍上の生成と全基底変更後のSerre消滅を合わせて結論を得る．
\end{proof}

\begin{remark}[有限パラメータ族で得られた条件と残る問題]
\label{rem:finite-tails-scope}
分離型族では全ての有限 $h$ において，飽和性，有向ACM，相対 $\chi$，
全シフトの豊富性，compact生成，proper性，無捩れな相対Serre公式が得られた．
反標準周期 $r=n-1$ 自体の相対豊富性は，滑らかな可換点の有限な開近傍で得られる．
いずれも本節で指定した任意のnoetherian基底変更に対する主張である．

一方，$C=R\!\operatorname{End}(E)$ の対角が $C^{\mathrm e}$ 上perfectであること，
すなわちcompact圏のhomological smoothness，およびcoherent圏全体の有限ホモロジー次元は
本節の結論に含まれない．滑らかな可換特殊ファイバーや全Artin近傍のsmoothnessだけで，
有限型基底のZariski開近傍のsmoothnessを結論しない．
これらを完成するには，この族の対角について有限完全分解を構成するか，
適用条件を満たす追加の開性定理を検証する必要がある．
AS--Gorenstein商のperfect圏のSerre公式から，全 $D^b(\mathrm{coh})$ のsmoothnessを
推論してはならない．
\end{remark}

% !TeX root = ../nc-del-pezzo.tex
\section{次数2の一般楕円型最高次項を持つ有限型量子化}
\label{sec:elliptic-degree-two-finite}

本節では活性重みを $(1,1,2)$ とする。
負の補助次数を仮定せず，三次Sklyanin代数の具体的な正則領域を用いる。
結論は，有限型基底上の有限表示，普遍的平坦性，相対ambient CY性，
および普遍的に正則な中心商である。
有限非零パラメータでのtailsのsmooth proper性や，
全 $\mathcal O$ 保存同値によるモジュライのatlasを結論しない。
また，全複素ファイバーのNoether性と，係数環上の族の強Noether性は区別する。

\subsection{三次Sklyanin代数を与える最高次potential}

巡回語の空間で
\[
 U_4=\frac{\alpha}{4}(x^4+y^4)
       +\beta\left(\frac23x^2y^2+\frac13xyxy\right),\qquad
 u_4=\frac{\alpha}{4}(x^4+y^4)+\beta x^2y^2
\]
と置く。前者の可換化は後者である。
各単項式 $x^iy^j$ の完全対称化の巡回類を $\operatorname{Sym}(x^iy^j)$ と書き，
\[
 U=U_4+\sum_{i+j\le3}c_{ij}\operatorname{Sym}(x^iy^j),\qquad
 F=\frac12z^2+u_4+\sum_{i+j\le3}c_{ij}x^iy^j
\]
を考える。$c_{00}$ は関係式に現れないが，中心元の定数項の規格化に用いる。
potentialと代数を
\begin{equation}\label{eq:elliptic-potential}
 \Phi=xyz-yxz-h\left(\frac12z^2+U\right),\qquad
 A=R\langle x,y,z\rangle/
  \bigl([x,y]-hz,\ [y,z]-h\partial_xU,\ [z,x]-h\partial_yU\bigr)
\end{equation}
で定める。ここで $R$ の元は中央である。
$c_{ij}=0$ とした代数を $A^{\mathrm{top}}$ と書く。

\begin{lemma}\label{lem:elliptic-top-as-open}
次の多項式を置く。
\[
 a=1+\frac23\beta h^2,\quad b=-2+\frac23\beta h^2,\quad
 c=\alpha h^2,\qquad
 \Delta_{\mathrm{AS}}=
   \left(1+\frac23\beta h^2\right)(-3+4\beta h^2).
\]
$T=\mathbb C[\alpha,\beta,h,\Delta_{\mathrm{AS}}^{-1}]$ 上の
$A^{\mathrm{top}}$ は次数ごとに有限局所自由で，任意の係数変更に対し
\[
 H_{A^{\mathrm{top}}}(t)=\frac1{(1-t)^2(1-t^2)}
\]
を持つ。全複素ファイバーは3次元のAS正則整域である。
特に $h=0$ の全点がこの開集合に含まれる。
\end{lemma}
\begin{proof}
$h\ne0$ では $z=h^{-1}(xy-yx)$ を消去でき，二つの関係式は
\[
 a(y^2x+xy^2)+byxy+cx^3=0,\qquad
 a(x^2y+yx^2)+bxyx+cy^3=0
\]
となる。三次type-A代数の非正則パラメータは
\[
 \{[a:b:c]\mid a^2=b^2=c^2\}
   \cup\{[0:0:1],[0:1:0]\}
\]
に限られる
\cite[Example 3.11, (75)--(76)]{BelmansRaedschelders2020}。
$a(a^2-b^2)=\Delta_{\mathrm{AS}}$ であるため，指定した開集合はこれらを避ける。
AS正則性と標準Hilbert級数を得る。
Noether性は次数一生成の3次元正則代数の既知の性質であり，
Artin--Tate--Van den Bergh\cite[Section 2, p.~338]{ATV1991}に明記されている。
さらに大域次元とGK次元がともに3なので，
同論文のTheorem 3.9から整域性が従う。
$h=0$ では三生成元表示そのものから $A^{\mathrm{top}}=\mathbb C[x,y,z]$ である。
この点で $z$ を消去してはならない。

各次数の普遍代数は有限個の単語と関係式による有限表示 $T$-加群である。
全閉点でファイバー次元が一定なので，その関係行列の階数は一定である。
$T$ は被約であるから，必要な階数を超える小行列式は恒等的に零であり，
所要階数の非零小行列式による開被覆上で有限自由性を得る。
従って各次数は有限局所自由であり，その表示は任意の基底変更と両立する。
\end{proof}

\begin{lemma}\label{lem:elliptic-top-relative-cy}
上の $A^{\mathrm{top}}$ は $T$ に相対的な3-Calabi--Yau代数であり，
potential対角複体は任意の係数変更後も完全である。
内部次数を含む相対双対は
\[
 R\!\operatorname{Hom}_{(A^{\mathrm{top}})^e_T}
       (A^{\mathrm{top}},(A^{\mathrm{top}})^e_T)
       \simeq A^{\mathrm{top}}(4)[-3]
\]
である。
\end{lemma}
\begin{proof}
各複素ファイバーでpotential対角複体を考える。
$h=0$ では三変数の対角Koszul複体である。
$h\ne0$ では $z$ の消去と，対角複体内の係数 $-h$ の可縮対の消去により，
二生成元potential
\[
 W=a x^2y^2+\frac b2xyxy+\frac c4(x^4+y^4)
\]
の対角複体になる。
このpotentialの巡回微分は上の二つの三次関係式である。
この二生成元代数を $S$ と書く。
AS正則三次代数の標準極小解消
\[
 0\longrightarrow S(-4)\longrightarrow S(-3)^2
  \longrightarrow S(-1)^2\longrightarrow S\longrightarrow\mathbb C
  \longrightarrow0
\]
において，次数4のsyzygyは1次元である。
巡回恒等式 $[x,\partial_xW]+[y,\partial_yW]=0$ がその非零生成元を与えるため，
potential対角複体を片側のaugmentationでテンソルした複体はこの極小解消である。
下に有界な次数付き自由複体に対するgraded Nakayamaにより，対角複体自体も完全である。
巡回Hessianのtensor flip対称性は無捩れな自己双対性を与える。
これはGinzburgのpotential複体の標準的な証明である
\cite[Proposition 5.1.9, Theorem 5.3.1, Proposition 8.1.3(ii)]{Ginzburg2007}。

普遍複体の各内部次数は有限射影 $T$-加群の有界複体であり，
augmentation項も補題\ref{lem:elliptic-top-as-open}により有限射影である。
全閉点で完全であるため，これらは局所的に分裂完全である。
従って任意の $T$-代数への係数変更後も完全である。
Hessianの自己双対性は係数環上の恒等式なので，任意の係数変更に両立する。
\end{proof}

\subsection{任意の低次係数と有限行列による中心持上げ}

以下の補題は，複素数上のfiltered PBW定理を
任意の中央係数へ移すために必要な基底変更の議論を明記するものである。

\begin{lemma}\label{lem:elliptic-universal-lower-pbw}
$S=T[c_{ij}:i+j\le3]$ とし，$A_S$ を\eqref{eq:elliptic-potential}で定める。
$S$ の活性filtration次数を0，$(x,y,z)$ の活性次数を $(1,1,2)$ とすれば，
\[
 \operatorname{gr}_{\mathrm{act}}A_S
       \simeq A^{\mathrm{top}}\otimes_T S.
\]
各filtration部分および各随伴次数部分は有限射影 $S$-加群である。
この同型と射影性は，任意の可換 $S$-代数への係数変更に両立する。
また $A_S$ は $S$ 上相対3-Calabi--Yauである。
\end{lemma}
\begin{proof}
複素数値の最高次パラメータを固定すると，
補題\ref{lem:elliptic-top-relative-cy}と
Etingof--GinzburgのProposition 7.3.1(ii)から，
全ての低次potentialについてPBW性とCY性が従う
\cite{EG2010}。
ここでは同論文の「generic」を有限開条件と読み替えていない。
最高次potentialがCYであることを既に具体的に証明した上で，
固定したCY topに対する全lower termsの定理を使っている。

活性次数の上限 $N$ を固定する。
最高次代数の次数 $0,\ldots,N$ の各部分を局所的に自由化し，
その基底を同じ次数の単語の線型結合で持ち上げる。
最高次関係式による次数降下から，持上げた基底は
$F_NA_S$ を $S$ 上生成する。
これらの間の $S$-線型関係があれば，任意の複素最高次パラメータと
任意の複素低次係数へ特殊化する。
固定topのPBW定理により全ての係数が零となる。
普遍係数環は被約かつ有限型なので，元の係数多項式も零である。
従ってこの生成系は基底である。
各 $N$ で得られる議論を貼り合わせると，主張の随伴次数同型と射影性を得る。
filtrationの短完全列は射影加群の列として分裂するため，
任意の非被約係数への基底変更も許される。

相対CY性には同じpotential対角複体を使う。
活性filtrationを入れた随伴次数複体は，
$A^{\mathrm{top}}$ の完全な対角複体の係数変更である。
filtrationは下に有界で各元は有限次数を持つため，
最高次数のcycleを境界として持ち上げる次数降下により元の複体も完全である。
この議論とHessianの自己双対性は任意の基底変更に両立する。
\end{proof}

\begin{theorem}\label{thm:elliptic-finite-central-family}
$\alpha_0(4\beta_0^2-\alpha_0^2)\ne0$ とし，
$p=(\alpha_0,\beta_0,0)$ を可換最高次パラメータとする。
$p$ のある主開近傍 $V\subset\operatorname{Spec}T$ が存在し，
その座標環を $T_V$ と置くと，次が成立する。
\begin{enumerate}
 \item 全低次係数環 $S_V=T_V[c_{ij}:i+j\le3]$ 上の $A$ は有限表示で，
       活性filtrationに関して普遍的PBW性と相対CY性を持つ。
 \item $A$ に活性filtration次数4の中心元 $\Omega$ が存在し，
       $\Omega\bmod h=F$ である。
       必要な分母は最高次パラメータ $\alpha,\beta,h$ だけに依存し，
       低次係数 $c_{ij}$ に依存しない。
 \item $\Omega$ は任意の基底変更後も左右の非零因子であり，
       $B=A/(\Omega)$ は $S_V$ 上平坦である。
       より正確には，各活性随伴次数部分は有限射影で，
       \[
        \operatorname{gr}_{\mathrm{act}} B
         \simeq (A^{\mathrm{top}}/(\Omega^{\mathrm{top}}))\otimes_{T_V}S_V,
        \qquad
        H_{A^{\mathrm{top}}/(\Omega^{\mathrm{top}})}(t)
          =\frac{1-t^4}{(1-t)^2(1-t^2)}.
       \]
       ここで $\Omega^{\mathrm{top}}$ は $\Omega$ の最高活性部分である。
 \item $\deg c_{ij}=4-i-j$，$\deg(\alpha,\beta,h)=0$ とすると，
       $\Omega$ は全内部次数4の斉次元に選べる。
\end{enumerate}
\end{theorem}
\begin{proof}
PBW性とCY性は前の補題による。
近傍を必要に応じて $\alpha(4\beta^2-\alpha^2)\ne0$ の領域へ制限する。
中心元には活性次数6以下の有限な線型代数だけを使う。
各次数 $d\le6$ において，整列語 $x^iy^jz^k$ と，
非整列語 $w$ から整列語への隣接交換を集めた関係列を用いる。
$h=0$ で後者は $w-\operatorname{ord}(w)$ であり，
非整列語成分に関する行列は単位行列である。
一般の $h$ では同じ行列は $I+hK_d$ になる。
有限個の行列式 $\det(I+hK_d)$ を可逆化すれば，
次数6以下で整列語を基底にできる。
これらの行列式は最高次パラメータだけの関数で，$h=0$ で1である。
低次係数を戻したときの関係行列は活性次数に関して三角行列となり，
対角ブロックは最高次の場合と同じである。
従って同じ局所化で $F_6A$ も整列語基底を持つ。

$A/hA$ が可換であり，$A$ は $h$ に関して捩れを持たないので，
\[
 D:F_4A\longrightarrow F_5A\oplus F_5A\oplus F_6A,
 \qquad
 D(v)=h^{-1}([v,x],[v,y],[v,z])
\]
は正則な係数を持つ線型写像である。
ここで $h^{-1}$ は商環への移行ではなく，$hA$ 内の元の一意的な割算を表す。
基底表示では分子の全係数が $h$ で割り切れる。
$F_4A$ の階数は $1+2+4+6+9=22$ である。

$h=0$ では $D$ はJacobian Poisson括弧による写像であり，
\[
 \ker D=\mathbb C\cdot1\oplus\mathbb C\cdot F
\]
となる。
この低次の主張には全Poisson中心の定理は不要である。
実際 $v=g_0(x,y)+zg_1(x,y)+kz^2$ と書き，
$\{v,x\}=\{v,y\}=0$ の $z$ に関する係数を比較すると，
$U_y^{\mathrm{ab}}g_1=U_x^{\mathrm{ab}}g_1=0$ から $g_1=0$，
続いて $\partial_xg_0=2k\partial_xU^{\mathrm{ab}}$，
$\partial_yg_0=2k\partial_yU^{\mathrm{ab}}$ を得る。
従って $v=2kF+\text{constant}$ である。
最高次係数 $\alpha\ne0$ により，この比較に用いた偏微分は非零である。

まず低次係数を全て零にする。
$D$ は活性次数を，各出力で生成元の重みだけ上げる。
従って，出力座標に「出力次数から生成元の重みを引いた次数」を付けると，
$D$ は次数ごとのブロックに分かれる。
次数 $1,2,3,4$ でそれぞれ $2,4,6,8$ 本の独立な行を選び，
定数項と $z^2$ の係数を読む二つの行を加える。
こうして $p$ で可逆な22次正方行列 $M$ を得る。
有限個の分母を払った $\det M$ の非零開集合にさらに制限する。
低次係数を戻しても，行列は活性次数に関して三角で，
対角ブロックは低次係数を含まない。
従ってその行列式は最高次パラメータだけに依存する。

選んだ20本の $D$ の行を零とし，定数項を $c_{00}$，
$z^2$ の係数を $1/2$ と指定した22個の線型方程式を解く。
その唯一の解を $\Omega$ とする。
ここでまだ選んでいない交換子の行も零になることを示す。
$h\ne0$ の任意の複素特殊化では，
Etingof--GinzburgのTheorem 3.3.2(ii)により
$F_4A$ に非定数の中心元がある。
従って $\ker D$ は少なくとも2次元である。
一方，$M$ の可逆性から，定数項と $z^2$ 係数による
$\ker D\to\mathbb C^2$ は単射である。
よってこれは同型であり，規格化された解 $\Omega$ は全交換子方程式を満たす。
$h=0$ では上の直接計算により解は $F$ である。
全複素特殊化で残りの行が零であり，係数環が被約なので，
全交換子方程式は普遍係数環上でも恒等的に零である。
中心性はその後の任意の非被約基底変更にも保たれる。

最高活性部分 $\Omega^{\mathrm{top}}$ は低次係数に依存せず，
整列基底での $z^2$ の係数が $1/2$ であるため全複素点で非零である。
最高次代数の全複素ファイバーは整域だから，
各次数の乗法写像
\[
 A^{\mathrm{top}}_{d-4}\xrightarrow{\ \Omega^{\mathrm{top}}\ }
 A^{\mathrm{top}}_d
\]
は全ファイバーで単射である。
有限局所自由加群の写像に対する小行列式の判定から，
これは局所的に分裂単射で，余核も有限局所自由である。
従って任意の係数変更後も単射である。
活性filtrationで最高項を比較すると，$\Omega$ による乗法は普遍的に単射かつstrictで，
その余核の随伴次数は表示した最高次商になる。
随伴次数部分が射影であることから商のfiltrationも分裂し，普遍平坦性を得る。

最後に，低次係数へ重み $4-i-j$ を与えた同時スカラー作用を考える。
関係式と全中心方程式はこの作用に同変であり，
定数項 $c_{00}$ と $z^2$ 係数 $1/2$ の規格化も次数4に適合する。
解の一意性により $\Omega$ は内部斉次次数4を持つ。
\end{proof}

\begin{remark}\label{rem:elliptic-center-scope}
上の議論は，中心元の構成に必要な次数6以下の有限行列を使っている。
全次数にわたる中心の階数一定性を仮定していない。
最高次代数では楕円自己同型の有限位数パラメータで高次の中心元が増え得るので
\cite[Theorem II]{ATV1991}，
$Z(A)=S_V[\Omega]$ を有限開集合全体で結論しない。
また，各次数の局所自由性は，全次数を同時に整列単項式で自由化する
一つの主開集合の存在を意味しない。
本証明で整列基底を固定するのは中心方程式に必要な次数6以下だけである。
\end{remark}

\subsection{最高次中心元の明示式}

有限行列による持上げとは独立に，最高次中心元の具体式も得られる。
\[
 \begin{aligned}
 C={}&b(a^2-c^2)(xyxy+yxyx)
      +a(b^2-a^2)(yx^2y+xy^2x)\\
     &+a(c^2-a^2)(x^2y^2+y^2x^2)
      +c(a^2-b^2)(x^4+y^4)
 \end{aligned}
\]
はDe Laetの四次中心元である
\cite[Section 5.2, Theorem 5.3]{DeLaet2017}。
$C_0=C|_{a=1,b=-2,c=0}$ と置く。
\[
 \widetilde\Omega^{\mathrm{top}}
  =-\frac18\left(
       \frac{C-C_0}{h^2}-4z^2+x\partial_xU_4+y\partial_yU_4
      \right)
\]
は最高次代数の中心元で，
$\widetilde\Omega^{\mathrm{top}}\bmod h=\frac12z^2+u_4$ を満たす。
ここで $(C-C_0)/h^2$ は自由代数の係数を多項式として割った式であり，
$h$ を可逆化した式ではない。
実際，最高次代数で
\[
 C_0=h^2\bigl(-4z^2+x\partial_xU_4+y\partial_yU_4\bigr)
\]
が成り立つため，$h\ne0$ では
$\widetilde\Omega^{\mathrm{top}}=-C/(8h^2)$ である。
$h$-捩れがないことから中心性は $h=0$ を含む族で成立する。
中心性は次の自由代数内の多項式恒等式からも直接確認できる。
$r_1=a(y^2x+xy^2)+byxy+cx^3$，$r_2=a(x^2y+yx^2)+bxyx+cy^3$ と置くと，
\begin{align*}
 [C,x]={}&-(a^2+c^2)r_1x^2-2ac\,r_1y^2
          +(a^2+b^2)xr_1x-(b^2-c^2)x^2r_1\\
         &+2ac\,y^2r_1-2ab\,r_2xy-(a^2+b^2)r_2yx
          +(a^2+b^2)xr_2y+2ab\,yxr_2.
\end{align*}
$[C,y]$ の式は $x,y$ および $r_1,r_2$ を同時に交換して得られる。
従って最高次代数の関係式の下で両交換子は零となる。

表示した自由語での $z^2$ の係数は $1/2$ であるが，
純粋な $x,y$ の語を整列基底へ直すと $z^2$ 項が生じ得る。
従って定理\ref{thm:elliptic-finite-central-family}の規格化中心元と比較する際は，
整列基底での $z^2$ 係数を用い，必要ならその非零開集合へ制限してスカラー倍する。
この係数は $h=0$ で $1/2$ なので，追加の制限も有限な開条件である。

\subsection{一般次元への中央係数の付加と半古典極限の範囲}

\begin{corollary}\label{cor:elliptic-general-dimension}
$n\ge2$ とし，中心一次変数 $s_1,\ldots,s_{n-1}$ を加える。
各 $c_{ij}$ を $s$ の任意の次数 $4-i-j$ の斉次多項式で置き換える。
その係数を独立な有限個のパラメータとし，
得られる有限型係数環を $K$ と書く。このとき $K$ 上で
普遍的に次数ごとに平坦な有限表示代数 $B=A/(\Omega)$ を得る。
\[
 H_{B_\xi}(t)=\frac{1-t^4}{(1-t)^{n+1}(1-t^2)},\qquad
 B/(h)\simeq (K/(h))[s_1,\ldots,s_{n-1},x,y,z]/(F)
\]
である。ここで $\xi$ は $\operatorname{Spec}K$ の任意の幾何点である。
最初の交換子は
\[
 \{x,y\}=z,\qquad \{y,z\}=\partial_xF,
 \qquad\{z,x\}=\partial_yF,\qquad\{s_i,-\}=0
\]
を与える。最高次係数にだけ依存する局所化なので，この置換は全て許される。
\end{corollary}
\begin{proof}
これは定理の任意基底変更と内部斉次性の適用である。
活性随伴次数で中心変数の多項式環と最高次商のテンソル積になり，
各内部次数では有限個の射影加群の直和となる。
Hilbert級数と普遍的次数別平坦性が従う。
複素係数を固定すると，ambient代数の活性随伴次数環は
Noether代数 $A^{\mathrm{top}}$ の中央多項式拡大なので，
そのambient代数と中心商 $B$ もNoetherである。
交換子の第一項は関係式の $h$ に関する一次部分から直ちに得られる。
\end{proof}

\begin{proposition}\label{prop:elliptic-semiclassical-coverage}
平方完成と方程式の非零定数倍により
$F=z^2/2+f_4(s,x,y)$ と正規化した次数2の可換Jacobian模型で，
選んだ活性平面上の二変数四次最高次項が
異なる四つの零点を持つと仮定する。
活性変数の $\operatorname{GL}_2$ 変換とPoissonパラメータの非零スカラー倍の後，
その模型は上の族の $h=0$ の特殊化と最初の交換子として実現される。
この条件は可換Poisson対のパラメータ空間の稠密開条件である。
\end{proposition}
\begin{proof}
異なる四点からなる $\mathbb P^1$ の集合で，二組への分割を選ぶ。
その二組をそれぞれ交換する射影変換は位数2である。
二つの不動点を $0,\infty$ に移すと，この変換は $t\mapsto-t$ になり，
四点は $\{r,-r,q,-q\}$ となる。
従って四次式は $Ax^4+Bx^2y^2+Cy^4$ と書け，$AC\ne0$ である。
対角変換によって $x^4$ と $y^4$ の係数を等しくできるため，
$u_4=\alpha(x^4+y^4)/4+\beta x^2y^2$ の形を得る。
四つの零点が異なる条件は $\alpha(4\beta^2-\alpha^2)\ne0$ である。
座標の選択は有限な補助選択とフレームの選択を伴うものであり，
一意的な大域正規形を主張していない。

受動的座標変換 $(x,y)^t=M(x',y')^t$ で $z$ を固定すると，
Jacobian括弧は
\[
 \pi_F\longmapsto (\det M)^{-1}\pi_{F\circ M}
\]
と変換される。従って $h'=h/(\det M)$ と置けば同じ第一交換子を表す。
最高次項以外の全ての項は $c_{ij}(s)$ に自由に含められる。
必要な開集合は各可換最高次パラメータ $(\alpha,\beta,0)$ を含むので，
任意のこのような古典模型の近傍で上の構成が使える。

四次式の判別式非零は開条件であり，例えば
\[
 F=\frac12z^2+x^4+y^4+\sum_{i=1}^{n-1}s_i^4
\]
が全 $n\ge2$ でその条件を満たす。
この例のaffine coneは原点外で滑らかであり，
重み付き射影空間の $z$ 軸上の特異点も超曲面上にない。
従って滑らかな次数2 del Pezzo $n$-foldを与える。
判別式非零条件は，係数が自由な各活性平面チャートで非空な稠密開集合である。
\end{proof}

\begin{remark}
この最後の主張は，古典的Poisson対とその第一交換子の範囲に関するものである。
有限非零 $h$ の族が $\mathfrak M^{\mathcal O}_{n,2}$ のatlasとなること，
その形式完成が全ての非可換変形を被覆すること，
あるいは全 $\mathcal O$ 保存同値がこの表示の座標変換で与えられることは，
本節からは結論しない。
同様に，全複素ファイバーがNoetherである事実だけから，
この有限型係数環上の族が強Noetherであるとは結論しない。
強Noether性と相対tailsは，第\ref{sec:elliptic-relative}節で
最高次中心商の相対曲線表示を構成してから証明する。
\end{remark}

% !TeX root = ../nc-del-pezzo.tex
\section{楕円型族の相対曲線表示と強Noether性}
\label{sec:elliptic-relative}

第\ref{sec:elliptic-degree-two-finite}節の次数2の族について，
可換楕円型点を含む有限開近傍上で強Noether性を証明する．
ここで強Noether性とは，係数環から任意のnoetherianな可換環への変更後にも
左右noetherianであることをいう．
体上のSklyanin代数のNoether性を特殊ファイバーごとに使うのではなく，
最高次中心商を相対twisted homogeneous coordinate ringとして構成する．
この構成は $h=0$ で二重対角線へ退化するpoint curveを一度blow upすることにより，
重み2の生成元 $z$ を保持したまま可換点まで延長される．

\subsection{二重対角線の解消とpoint shift}

以下では $a,b,c,T$ は補題\ref{lem:elliptic-top-as-open}の記号とする．
特に
\[
 a=1+\frac23\beta h^2,\qquad
 b=-2+\frac23\beta h^2,\qquad c=\alpha h^2,
 \qquad 2a+b=2\beta h^2.
\]
可換点 $p=(\alpha_0,\beta_0,0)$ は
$\alpha_0(4\beta_0^2-\alpha_0^2)\ne0$ を満たすとし，
定理\ref{thm:elliptic-finite-central-family}の有限開近傍と交わって作業する．
以下の近傍の縮小は最高次パラメータだけに依存する．

\begin{lemma}[可換極限を持つ相対point curve]
\label{lem:elliptic-resolved-point-curve}
$p$ のあるアフィン開近傍 $V$ 上に，滑らかで射影的な種数1の曲線族
$f:\mathcal E\to V$，自己同型 $\sigma$，および相対次数2の線束
$L$ が存在し，次を満たす．
\begin{enumerate}
 \item $h\ne0$ では $\mathcal E$ は三次関係式のpoint curveであり，
       $\sigma$ は隣接する二点から次の二点へ移るshiftである．
 \item $h=0$ では
       \[
        \mathcal E_0:
        \frac12z^2+\frac\alpha4(x^4+y^4)+\beta x^2y^2=0
        \quad\subset\mathbb P(1,1,2),
        \qquad L_0=\mathcal O_{\mathcal E_0}(1),\quad \sigma_0=1.
       \]
 \item $x,y$ は $f_*L$ の基底であり，
       $z$ は $L\otimes\sigma^*L$ の切断として定義され，
       twisted積について $xy-yx=hz$ が成り立つ．
\end{enumerate}
\end{lemma}
\begin{proof}
$\mathbb P^1_T\times_T\mathbb P^1_T$ の座標を
$(u:v),(U:V)$ とする．三次関係式を三点について多重線型化すると，
第三点は次の行列の核で与えられる：
\[
 M=\begin{pmatrix}
 avV+cuU&auV+bvU\\
 avU+buV&auU+cvV
 \end{pmatrix}.
\]
その行列式は
\begin{equation}\label{eq:elliptic-point-determinant}
 H=ac(u^2U^2+v^2V^2)-ab(u^2V^2+v^2U^2)+(c^2-b^2)uvUV.
\end{equation}
$\delta=uV-vU$ と置くと，厳密な恒等式
\begin{align}
 H&=-ab\delta^2+h^2Q,\label{eq:elliptic-point-blowup}\\
 Q&=\alpha a(u^2U^2+v^2V^2)+(\alpha^2h^2-2\beta b)uvUV\notag
\end{align}
を得る．従って $h=0$ のpoint curveは二重対角線 $2\delta^2=0$ である．

周囲を理想層 $(h,\delta)$ でblow upし，$H=0$ のstrict transformを
$\mathcal E$ とする．これは基底上射影的である．
$h$-chartで $w=\delta/h$ と書く．$u=U=1$ のアフィン表示では
$t=v/u$，$q=V/U=t+hw$ であり，strict transformの式は
\begin{equation}\label{eq:elliptic-resolved-affine}
 G=-ab w^2+\alpha a(1+t^2q^2)+(\alpha^2h^2-2\beta b)tq=0.
\end{equation}
このchartの周囲は座標 $(t,w)$ を持つ基底上の滑らかな曲面である．
$h=0$ の式は
\[
 2w^2+\alpha(1+t^4)+4\beta t^2=0.
\]
逆座標のchartでは $t'=1/t$，$q'=1/q$，$w'=-w/(tq)$ で貼り合わさる．
ここで $w'=(q'-t')/h$ としたので，大域切断 $\delta/h$ の局所成分は
元のchartで $w$，逆chartで $-w'$ である．
従って特殊ファイバーは表示した重み付き四次曲線であり，判別式の仮定から滑らかである．
blow upのもう一方のchartでは $h=k\delta$ と書き，strict transformは
$-ab+k^2Q=0$ となる．$h=0$ で $k=0$ なら $ab=-2$ と矛盾するので，
特殊ファイバーの全点は先の $h$-chartに入る．
特に，特殊ファイバーを見落としたり，二重対角線そのものを楕円曲線と扱ったりしていない．

特殊ファイバーの各点で相対Jacobian判定を適用すると，$f$ はその近くで滑らかである．
非滑らかな点の閉集合の像はproper性により閉で，$p$ を含まない．
この像を除き，さらに必要なら幾何学的連結性の開近傍へ制限すると，
全幾何ファイバーが滑らかで幾何学的整な種数1曲線となる．
以下でも同じように有限個の閉集合のproperな像を除き，最後に $p$ を含む主開近傍を選ぶ．

行列 $M$ は特殊ファイバー上で階数1である．例えば有限座標では
$M|_{h=0}=\bigl(\begin{smallmatrix}t^2&-t\\-t&1\end{smallmatrix}\bigr)$ であり，
無限遠のchartでも同様である．階数0の閉集合の像を除くと，
その核は第三点への射を定める．有限座標ではその座標は
\[
 r=-\frac{at+bq}{a+ctq}.
\]
$\mathcal H(t,q)=H(1,t,1,q)$ と置くと，多項式恒等式
\begin{equation}\label{eq:elliptic-shift-identity}
 (a+ctq)^2\mathcal H\left(q,-\frac{at+bq}{a+ctq}\right)
       =(a^2-bcq^2)\mathcal H(t,q)
\end{equation}
が成り立つ．左辺は分母を払った多項式として読む．
従って $(t,q)\mapsto(q,r)$ は $h\ne0$ の曲線を保つ．
さらに $q=t+hw$ を代入すると
\begin{equation}\label{eq:elliptic-resolved-shift}
 t_{\mathrm{new}}=q,\qquad
 w_{\mathrm{new}}=\frac{r-q}{h}
       =\frac{aw-h(2\beta q+\alpha tq^2)}{a+ctq}.
\end{equation}
分母は $h=0$ で1なので，この式は有限座標の特殊ファイバー全体で正則であり，
そこでは $(t_{\mathrm{new}},w_{\mathrm{new}})=(t,w)$ となる．
$x,y$ を同時に交換した逆座標の式が無限遠を覆う．
また第三点への射は $h=0$ で第二点と一致するから，
新しい行列式 $\delta_{\mathrm{new}}$ は $h$ で割り切れる．
これがblow upの $h$-chartへの持上げを与え，上の局所式は貼り合わさる．
式\eqref{eq:elliptic-shift-identity}と $h$-平坦性により像はstrict transformに入る．
三次関係式は語順の反転で保たれる．従って二点の交換を $\iota$ と書くと，
$\iota\sigma\iota$ が逆shiftである．核の一意性による逆写像の恒等式は
$h\ne0$ で成立し，平坦な族上の等式として全体へ延長される．
以上で $\sigma\in\Aut_V(\mathcal E)$ と $\sigma_0=1$ を得る．

二つの射影を $p_1,p_2$ と書き，$L=p_1^*\mathcal O_{\mathbb P^1}(1)$ と置く．
$p_1\sigma=p_2$ なので，$\sigma^*L=p_2^*\mathcal O(1)$ である．
切断 $\delta/h$ は $L\otimes\sigma^*L$ の大域切断となる．
これを $z$ と書き，$u,v$ の引戻しを $x,y$ と書くと，twisted積では
$xy-yx=uV-vU=hz$ である．
特殊ファイバーで $p_1$ は次数2の写像なので $L$ の相対次数は2である．
種数1曲線の正次数線束の $H^1$ は零であり，$f_*L$ は階数2で基底変更に両立する．
$x,y$ は特殊ファイバーで基底なので，必要なら二次元の行列の非零行列式による
有限開近傍へ制限し，族全体で基底にできる．
\end{proof}

\subsection{最高次中心商と相対twisted homogeneous coordinate ring}

\begin{theorem}\label{thm:elliptic-relative-thcr}
前補題の $V$ をさらに有限回縮小し，$T_V=\Gamma(V,\mathcal O_V)$ と置く．
$L_m=L\otimes\sigma^*L\otimes\cdots\otimes(\sigma^{m-1})^*L$，$L_0=\mathcal O$ とする．
積 $s\cdot t=s\otimes(\sigma^m)^*t$ によって定まる次数付き代数を
\[
 D=\bigoplus_{m\ge0}H^0(\mathcal E,L_m)
\]
と書く．このとき同型
\[
 A^{\mathrm{top}}/(\Omega^{\mathrm{top}})\simeq D
\]
が存在し，任意の可換 $T_V$-代数への係数変更に両立する．
特に，この最高次中心商と $A^{\mathrm{top}}$ は $T_V$ 上強Noetherである．
\end{theorem}
\begin{proof}
各 $m>0$ について $L_m$ のファイバー次数は $2m$ なので，
$H^1(\mathcal E_s,(L_m)_s)=0$ である．
従って $D_m$ は階数 $2m$ の有限局所自由加群で，任意基底変更に両立する．
$D_0=T_V$ についても，properで幾何学的連結な滑らかな曲線族の
$f_*\mathcal O=\mathcal O_V$ とその基底変更を用いる．

前補題の $x,y,z$ は $A^{\mathrm{top}}$ の関係式を満たす．
第一の関係式は既に確認した．残りの二つは $h\ne0$ で $z=(xy-yx)/h$ を
消去すると，行列 $M$ を定めた二つの三次関係式そのものである．
$D$ の各次数が $T_V$-平坦なので，これらの等式は $h=0$ まで成立する．
従って準同型 $A^{\mathrm{top}}\to D$ を得る．

この写像 $A^{\mathrm{top}}\to D$ は，
$h=0$ の次数1と2で同型である．次数2では
$x^2,xy,y^2,z$ が重み付き楕円曲線の四つの基底切断となる．
有限個の行列式を可逆化すれば，この写像は $V$ 上で次数1と2において全射となる．
さらに $m\ge2$ に対して
\[
 D_m\otimes_{T_V}D_1\longrightarrow D_{m+1}
\]
は全射である．これは次数2の線束 $(\sigma^m)^*L$ の評価完全列を $L_m$ で
tensorした後に $f_*$ を取って分かる．核の線束は基底からの線束のtensor因子を除けば
$L_m\otimes((\sigma^m)^*L)^{-1}$ で，そのファイバー次数は $2m-2>0$ である．
従ってその $R^1f_*$ は零で，この全射は任意基底変更後にも成立する．
ゆえに $D$ は次数1と2で生成され，写像は全次数で全射である．
必要な開条件は次数1と2だけから選んでおり，無限個の行列式の同時可逆化ではない．

この写像が $\Omega^{\mathrm{top}}$ を零に送ることを確かめる．
まず $T_V$ の分数体上で考える．整な曲線のtwisted homogeneous coordinate ringに
次数 $e>0$ の非零中心切断 $g$ があれば，$\sigma^e=1$ である．
実際，同じ次数 $m$ の二つの非零切断 $s,t$ と $g$ の可換性を比較すると
$\sigma^e(s/t)=s/t$ となる．$m\ge2$ なら $L_m$ の次数は4以上であり非常に豊富なので，
これらの比は曲線の函数体を生成する．従って $\sigma^e$ は恒等写像である．
一方，式\eqref{eq:elliptic-resolved-shift}から
\[
 t\circ\sigma^4=t+4hw\pmod {h^2}
\]
である．$w$ は特殊楕円曲線上の零でない有理函数であり，標数は0であるから，
$\sigma^4$ は族の一般点で恒等写像ではない．
よって一般点の $D$ には非零の次数4中心切断がなく，
$\Omega^{\mathrm{top}}$ の像は零である．$D_4$ は整域 $T_V$ 上局所自由なので，
同じ消滅が $V$ 上で成立する．ここでは中心全体が全ての特殊化で一定であるとは主張していない．

以上により $\Omega^{\mathrm{top}}$ を零に送る全射
$A^{\mathrm{top}}\to D$ を得た．従って中心商からも全射が誘導される．
源と標的は各次数で同じ階数の有限局所自由加群だから，これは同型である．
同型と切断環の任意基底変更の性質を合わせれば，主張した普遍的な比較を得る．

任意のnoetherianな可換 $T_V$-代数 $R$ を固定する．上の基底変更により
$D_R$ は $(\mathcal E_R,L_R,\sigma_R)$ のtwisted homogeneous coordinate ringである．
$L_R$ は全ファイバーで正次数なので相対的に豊富である．
$\sigma_R$ は基底上の自己同型で，各幾何ファイバーは整な曲線だから，
全ての基底へ収縮される積分曲線上で線束の次数を保存する．
従って相対的な数値同値類の群への作用は恒等である．
Keelerの相対 $\sigma$-ampleness判定
\cite[Theorems 7.13--7.14]{Keeler2003}から，$L_R$ は左右 $\sigma_R$-ampleであり，
$D_R$ は左右noetherianである．この定理は基底を体に限定しない．
非被約な $R$ でも，数値同値と豊富性の判定は幾何ファイバー上で同じように適用できる．

最後に $A_R^{\mathrm{top}}$ のNoether性を持ち上げる．
$\Omega^{\mathrm{top}}$ は任意基底変更後も正則で，正の次数4を持つ中央元なので，
その冪による随伴次数環は $D_R[\zeta]$ である．
これはnoetherianである．斉次イデアルの随伴次数イデアルの有限個の斉次生成元を
$A_R^{\mathrm{top}}$ へ持ち上げると，最初の非零 $\Omega$-進記号を順に消去できる．
固定した内部次数では $\Omega$-進次数に上限があるため，この消去は有限回で終わる．
従って全斉次イデアルが有限生成となる．一般のイデアルも内部次数の最高項を
用いる通常の帰納法で有限生成となる．反対環にも同じ証明を適用でき，
$A_R^{\mathrm{top}}$ の左右Noether性を得る．
\end{proof}

\begin{corollary}[任意lower項と全次元での強Noether性]
\label{cor:elliptic-strong-noetherian}
前定理の有限開集合上で，任意のlower係数を持つ
定理\ref{thm:elliptic-finite-central-family}の $A,B$ は強Noetherである．
さらに系\ref{cor:elliptic-general-dimension}で
$c_{ij}=c_{ij}(s_1,\ldots,s_{n-1})$ と置いた全 $n\ge2$ の斉次族も，
その有限型係数環上で強Noetherである．
\end{corollary}
\begin{proof}
任意のnoetherianな係数環 $R$ 上で，補題\ref{lem:elliptic-universal-lower-pbw}は
$\operatorname{gr}_{\mathrm{act}}A=A^{\mathrm{top}}_R$ を与える．
右辺がnoetherianなので，増大filtrationの最高記号を消去する帰納法から
$A$ は左右noetherianであり，その商 $B$ も同様である．
中央変数を加える場合は，まず係数環をnoetherian環 $R[s_1,\ldots,s_{n-1}]$ として
同じ証明を適用する．これは任意のnoetherian係数変更後にも使えるので，主張は強Noether性である．
局所化したのはtop係数だけであり，lower多項式の選択に新たな制限は課していない．
\end{proof}

\subsection{相対tailsと有限生成元}

\begin{theorem}\label{thm:elliptic-relative-tails}
前系の全次元の斉次族をnoetherianな可換係数環 $\Lambda$ へ変更したものを考える．
$A$ をambient，$B=A/(\Omega)$，$\mathcal Q=\QGr B$，
$r=n-1$，$l=n+3$ とする．このとき
\begin{enumerate}
 \item $A$ は $\Lambda$ 上homologically smoothで，
       $A^!_\Lambda\simeq A(n+3)[-n-2]$ である．
 \item $B$ の左右の相対局所コホモロジーは
       $R\Gamma_{\ge1}(B)\simeq B^\vee(r)[-n-1]$ である．
       第\ref{sec:finite-tails}節の飽和性，有向ACM，相対 $\chi$，
       全シフトの豊富性，compact生成，proper性，相対Serre公式が成立する．
       特にcompact圏のSerre関手は $S_\Lambda\simeq(-r)[n]$ である．
 \item $E=\bigoplus_{j=0}^{n+2}\mathcal O(j)$ はcompact生成元であり，
       $R\!\operatorname{End}(E)$ は $\Lambda$ 上perfectである．
\end{enumerate}
これらの構成・公式は許される任意のnoetherian係数変更と両立する．
さらに，滑らかな可換del Pezzo模型を与える係数点の有限開近傍では，
正周期 $r=n-1$ の列自体も相対的に豊富である．
\end{theorem}
\begin{proof}
まず $P_0=\Lambda[s_1,\ldots,s_{n-1}]$ と置く．
補題\ref{lem:elliptic-universal-lower-pbw}により，$A$ は $P_0$ 上平坦で
相対3-CYであり，相対双対の内部shiftは4である．
命題\ref{prop:finite-separated-absolute-cy}の証明は，係数環上の自由性を平坦性に
弱めても成り立つ．実際，$A\otimes_\Lambda A^{\mathrm{op}}$ は
$P_0\otimes_\Lambda P_0$ 上平坦であり，中央差
$s_i\otimes1-1\otimes s_i$ の正則性とそのKoszul解消を保つ．
この長さ $n-1$ の解消と相対potential対角解消を組み合わせ，導来随伴を適用すると
\[
 A^!_\Lambda\simeq A(4+n-1)[-3-(n-1)]
\]
を得る．中央差の順序を固定しているので追加の基底線束やNakayama捻りは現れず，
構成は任意基底変更に両立する．

系\ref{cor:elliptic-strong-noetherian}により，$A,B$ は全ての許される基底上で
左右noetherianである．各内部次数は有限射影で，$\Omega$ は普遍的に正則な次数4の中央元である．
従って第\ref{sec:finite-tails}節の相対双対と局所コホモロジーの証明を適用できる．
その証明で次数部分の有限自由性を用いた箇所は，有限射影性だけで十分である：
切断加群 $A/A_{\ge q}$，$B/B_{\ge q}$ は係数環上perfectであり，
次数別双対は基底変更に可換する．
ambientの有限projective対角解消から有限係数加群の双対性を得て，
中央元のchange of ringsと切断次数についての帰納極限を取ると，
\[
 R\Gamma_{\ge1}(B)=B^\vee(n-1)[-n-1]
\]
となる．torsionの比較で現れる $\operatorname{Tor}_1$ は長さ4の次数帯に支持を持ち，
その遷移は十分離れた切断間で零となる．従って同節のsmashing表示も成立する．
これに続く局所化三角，有限型自由解消，左右の局所化bimoduleの双対の議論は
分離型の関係式に依存しない．飽和性，ACM，$\chi$，Serre消滅，compact生成，
proper性，および $S=(-r)[n]$ が従う．

有限windowについては，有限なaugmentation解消を別に確認する．
整数 $M>1$ を固定し，中央変数の補助重みを1，
$(x,y,z)$ の補助重みを $(M,M,2M)$ とする．
$c_{ij}(s)$ の内部次数は $4-i-j$ なので，lower potentialの補助次数は
\[
 M(i+j)+(4-i-j)<4M\qquad (i+j\le3).
\]
従ってlower関係式はこのfiltration次数を真に下げ，
\[
 \operatorname{gr}_{\mathrm{aux}}A
       =\Lambda[s_1,\ldots,s_{n-1}]\otimes_\Lambda A^{\mathrm{top}}_\Lambda.
\]
実際，内部次数 $d$ を固定すると補助次数は $d+(M-1)\deg_{\mathrm{act}}$ であり，
この随伴次数の同型は既に証明した活性PBW同型からも従う．
最高次代数のaugmentation解消は，potential対角複体の片側をaugmentationした
\[
 0\longrightarrow A^{\mathrm{top}}(-4)
 \longrightarrow A^{\mathrm{top}}(-3)^2\oplus A^{\mathrm{top}}(-2)
 \longrightarrow A^{\mathrm{top}}(-1)^2\oplus A^{\mathrm{top}}(-2)
 \longrightarrow A^{\mathrm{top}}\longrightarrow\Lambda\longrightarrow0
\]
であり，補題\ref{lem:elliptic-top-relative-cy}により普遍的に完全である．
これと $s_1,\ldots,s_{n-1}$ のKoszul解消をtensorする．
得られる自由解消を，内部shiftと自由基底を固定したままfilteredに持ち上げる．
微分の合成の誤差を最高記号の完全性によって消去し，次にcycleの最高記号を消去する
命題\ref{prop:finite-window-proper}の手順がそのまま適用できる．
全補助重みが正なので，各基底ベクトルに対する補正は有限回で終わる．
普遍係数環上で一度構成すれば，任意係数変更後にも随伴次数の完全性が保たれる．
解消の両端は $A(-n-3)$ と $A$ が一つずつであり，
他の項の内部shiftは $1,\ldots,n+2$ に入る．

この解消を $B$ とtensorした際のTorは
$\Lambda$ と $\Lambda(-4)$ だけであり，いずれもtorsionである．
従ってtailsでは有限完全複体となり，長さ $n+3$ の連続したtwistから
両方向に全てのtwistを生成できる．これで表示した有限compact生成元を得る．
最後に，滑らかな可換点では $\mathcal O(r)$ が豊富である．
命題\ref{prop:finite-periodic-generation}の有限個の生成条件と，上で得た
全noetherian係数変更後のSerre消滅を用いれば，有限開近傍上で周期 $r$ の相対豊富性を得る．
\end{proof}

\begin{remark}[適用範囲]
\label{rem:elliptic-relative-scope}
相対曲線表示を用いたのは活性最高次中心商である．
一般lower項を含む $n$ 次元のtailsそのものを曲線のcoherent圏と同一視していない．
また $\sigma$ が有限位数となる特殊点を除外しておらず，中心全体の次元が跳ぶ可能性も残る．
上の結論は各古典楕円型最高次点の有限開近傍で成立し，
初めに取ったAS正則領域全体に一つの曲線表示が存在するという主張ではない．
compact圏の対角のperfect性としてのsmoothness，
全 $\mathcal O$ 保存同値の認識，および全モジュライの有限型atlasは，
相対proper性とSerre公式から自動的に従うものではない．
\end{remark}

% !TeX root = ../nc-del-pezzo.tex
\section{次数1の一般楕円型最高次項を持つ有限型量子化}
\label{sec:elliptic-degree-one-finite}

本節では活性重みを $(1,2,3)$ とし，Weierstrass型最高次項から始める。
Stephensonの重み付き楕円代数の分類を用いると，
最高次potentialがCYとなる有限開集合を明示できる。
その上で固定CY topに対するfiltered PBW定理を適用し，
全低次係数と任意の係数変更を許す族を得る。

\subsection{Weierstrass型potentialと明示的なAS正則開集合}

巡回語の空間で
\[
 U_6=\frac13y^3+\alpha x^4y+\beta x^6,\qquad
 U=U_6+\sum_{i+2j<6}c_{ij}\operatorname{Sym}(x^iy^j),
 \qquad \deg(x,y,z)=(1,2,3)
\]
と置く。ここで $\operatorname{Sym}$ は完全対称化の巡回類であり，
和には定数項を含めて12個の係数がある。
可換化を $U^{\mathrm{ab}}$ と書き，
\begin{equation}\label{eq:degree-one-elliptic-potential}
 \Phi=xyz-yxz-h\left(\frac12z^2+U\right),\qquad
 A=R\langle x,y,z\rangle/
 \bigl([x,y]-hz,\ [y,z]-h\partial_xU,\ [z,x]-h\partial_yU\bigr)
\end{equation}
と定める。全係数は中央である。全 $c_{ij}=0$ とした代数を
$A^{\mathrm{top}}$ と書く。
古典最高次方程式
\[
 F_6=\frac12z^2+\frac13y^3+\alpha x^4y+\beta x^6=0
 \quad\text{in }\mathbb P(1,2,3)
\]
が滑らかな楕円曲線を定める条件は
$\Delta=4\alpha^3+9\beta^2\ne0$ である。

\begin{lemma}\label{lem:degree-one-stephenson-normal-form}
次の係数環と有限被覆を考える。
\[
 T_0=\mathbb C[\alpha,\beta,h,\Delta^{-1}],\qquad
 T'_0=T_0[k,\rho]/(k^2+\alpha,\ \rho^3-6\beta+4k^3).
\]
$T'_0$ は $T_0$ 上階数6の有限自由かつ忠実平坦な代数であり，
$\rho$ は可逆である。
$h\ne0$ において $z$ を消去し，$y=\rho Y+kx^2$ と変数変換すると，
$A^{\mathrm{top}}$ は二生成元代数
\begin{equation}\label{eq:degree-one-stephenson-relations}
 \begin{gathered}
 S(a,b,c)=T'_0[h^{-1}]\langle x,Y\rangle/(f_5,f_4),\\
 \begin{aligned}
 f_5&=a(Y^2x+xY^2)+cYxY+bx^5,\\
 f_4&=a(x^2Y+Yx^2)+cxYx+bY^2.
 \end{aligned}
 \end{gathered}
\end{equation}
と同型になる。ここで
\[
 a=1+kh^2,\qquad b=\rho h^2,\qquad c=-2.
\]
これらの関係式は巡回potential
\[
 W_{a,b,c}=a x^2Y^2+\frac c2xYxY+
             b\left(\frac13Y^3+\frac16x^6\right)
\]
の巡回微分である。
\end{lemma}
\begin{proof}
二つの定義方程式は順に $k$ と $\rho$ についてmonicである。
従って $k^\epsilon\rho^j$，$0\le\epsilon\le1$，$0\le j\le2$ が
$T_0$-基底を与え，有限自由性と忠実平坦性が従う。
$k$ を可逆化する必要はなく，$\alpha=0$ 上の $k=0$ も含まれる。
また
\begin{equation}\label{eq:degree-one-cover-unit}
 (6\beta-4k^3)(6\beta+4k^3)
   =36\beta^2+16\alpha^3=4\Delta
\end{equation}
なので $\rho^3=6\beta-4k^3$，従って $\rho$ は可逆である。

$h\ne0$ では $z=h^{-1}[x,y]$ を消去できる。
得られる二関係式は
\[
 W=-\frac12[x,y]^2+h^2U_6
   =x^2y^2-xyxy+h^2U_6
\]
の巡回微分である。最後の等号は巡回語の空間での等号である。
変数変換 $y=\rho Y+kx^2$ は重みを保つ自由代数の自己同型である。
$[x,y]=\rho[x,Y]$ および $k^2+\alpha=0$ を用いると，巡回語の空間で
\[
 U_6(x,\rho Y+kx^2)
  =\frac{\rho^3}{3}Y^3+k\rho^2x^2Y^2+
       \left(\beta-\frac23k^3\right)x^6.
\]
$\rho^3=6\beta-4k^3$ より
\begin{equation}\label{eq:degree-one-potential-normal-form}
 W(x,\rho Y+kx^2)=\rho^2W_{a,b,c}(x,Y).
\end{equation}
関係式のイデアルについても，$R_4=\partial_yW$，
$R_5=\partial_xW$ と置けば，自由代数内で
\[
 \begin{split}
 R_4(x,\rho Y+kx^2)&=\rho f_4,\\
 R_5(x,\rho Y+kx^2)+k\rho(xf_4+f_4x)&=\rho^2f_5
 \end{split}
\]
である。$\rho$ の可逆性から主張の同型を得る。
\end{proof}

\begin{lemma}\label{lem:degree-one-elliptic-top-as-open}
\[
 D_1=(1+\alpha h^4)(5+\alpha h^4),\qquad
 T=\mathbb C[\alpha,\beta,h,(\Delta D_1)^{-1}]
\]
と置く。$T$ 上の $A^{\mathrm{top}}$ は次数ごとに有限局所自由で，
任意の係数変更後もHilbert級数
\[
 H_{A^{\mathrm{top}}}(t)=\frac1{(1-t)(1-t^2)(1-t^3)}
\]
を持つ。全幾何学的ファイバーは大域次元3のAS正則Noether整域である。
この開集合は $\Delta\ne0$ の全ての $h=0$ 点を含む。
\end{lemma}
\begin{proof}
StephensonのTheorem 1.1\cite[p.~2318]{Stephenson1997}は，
代数閉体・標数0の下で，\eqref{eq:degree-one-stephenson-relations}の代数が
次の11点を除いた全パラメータでAS正則Noether整域となり，
上記のHilbert級数を持つことを示す。
\[
 \begin{split}
 \mathcal S={}&\{(0:0:1),(0:1:0)\}
    \cup\{(1:b:-1)\mid b^3=-1\}\\
   &\cup\left\{(1:b:c)\ \middle|\
          c^2-c-1=0,\ b^3=-c^{-2}\right\}.
 \end{split}
\]
従って $a(a+c)(c^2-ac-a^2)\ne0$ は正則性の十分条件である。
補題\ref{lem:degree-one-stephenson-normal-form}のパラメータでは
\begin{equation}\label{eq:degree-one-as-discriminant}
 a(a+c)(c^2-ac-a^2)
  =-(1+\alpha h^4)(5+\alpha h^4)=-D_1.
\end{equation}
実際 $a(a+c)=k^2h^4-1=-1-\alpha h^4$，
$c^2-ac-a^2=5-k^2h^4=5+\alpha h^4$ である。
有限忠実平坦被覆上の同型から，$h\ne0$ の全幾何学的ファイバーに
Stephensonの定理を適用できる。
$h=0$ では三生成元表示\eqref{eq:degree-one-elliptic-potential}を用い，
直接 $A^{\mathrm{top}}=\mathbb C[x,y,z]$ を得る。
ここで $D_1|_{h=0}=5$ である。

各次数の普遍代数は有限個の単語と斉次関係式による有限表示 $T$-加群である。
上のファイバー計算から，その関係行列は全閉点で一定の階数を持つ。
$T$ は被約な有限型 $\mathbb C$-代数なので，
その階数を超える小行列式は恒等的に零である。
所要階数の小行列式による開被覆上では余核が自由となり，
各次数の有限局所自由性が従う。
表示とその余核は任意の係数変更に両立するので，
非被約係数への基底変更後も同じ階数を持つ。
\end{proof}

\subsection{相対CY性と全低次係数の普遍PBW性}

\begin{lemma}\label{lem:degree-one-elliptic-top-relative-cy}
$A^{\mathrm{top}}$ のpotential対角複体は $T$ 上完全であり，
任意の係数変更後も完全である。内部次数を含む相対双対は
\[
 R\!\operatorname{Hom}_{(A^{\mathrm{top}})^e_T}
       (A^{\mathrm{top}},(A^{\mathrm{top}})^e_T)
       \simeq A^{\mathrm{top}}(6)[-3]
\]
となる。従って $A^{\mathrm{top}}$ は相対3-Calabi--Yau代数である。
\end{lemma}
\begin{proof}
まず複素ファイバーで完全性を示す。
$h=0$ ではpotential対角複体は重み $(1,2,3)$ の対角Koszul複体である。
$h\ne0$ では，$z$ の消去に対応する係数 $-h$ の可縮対を除くことで，
二生成元potential $W=-[x,y]^2/2+h^2U_6$ の対角複体を得る。
この二生成元代数を $S$ と書く。
補題\ref{lem:degree-one-stephenson-normal-form}と
StephensonのTheorem 5.7，(5.5)\cite[pp.~2336--2338]{Stephenson1997}より，
$S$ のaugmentationは極小解消
\[
 0\longrightarrow S(-6)\longrightarrow S(-5)\oplus S(-4)
  \longrightarrow S(-1)\oplus S(-2)\longrightarrow S
  \longrightarrow\mathbb C\longrightarrow0
\]
を持つ。重みを保つ変数変換の下でもこの次数列は変わらない。
次数6の極小syzygyは1次元であり，巡回恒等式
\[
 [x,\partial_xW]+[y,\partial_yW]=0
\]
がその非零生成元を与える。
従ってpotential対角複体を片側のaugmentationでテンソルした複体は，
この極小解消である。
内部次数が下に有界な自由複体に対するgraded Nakayamaにより，
対角複体自体も完全となる。
巡回Hessianのtensor flip対称性が無捩れな自己双対性を与える
\cite[Proposition 5.1.9, Theorem 5.3.1, Proposition 8.1.3(ii)]{Ginzburg2007}。
potentialの重みが6であるため，内部次数のずれは6である。

普遍複体の各内部次数は有限射影 $T$-加群の有界複体であり，
augmentation項も補題\ref{lem:degree-one-elliptic-top-as-open}により有限射影である。
全閉点で完全であるから，これらは局所的に分裂完全である。
従って任意の係数変更後も完全である。
Hessianの自己双対性は係数環上の恒等式なので，
相対CY同型も任意の係数変更に両立する。
\end{proof}

\begin{lemma}\label{lem:degree-one-elliptic-universal-lower-pbw}
$S=T[c_{ij}:i+2j<6]$ とし，$A_S$ を
\eqref{eq:degree-one-elliptic-potential}で定める。
係数の活性filtration次数を0とすれば
\[
 \operatorname{gr}_{\mathrm{act}}A_S
       \simeq A^{\mathrm{top}}\otimes_T S.
\]
各filtration部分と各随伴次数部分は有限射影 $S$-加群であり，
この同型は任意の可換 $S$-代数への係数変更に両立する。
また $A_S$ は $S$ 上相対3-Calabi--Yauである。
$\deg c_{ij}=6-i-2j$，$\deg(\alpha,\beta,h)=0$ として全内部次数を付ければ，
相対CY同型の内部次数のずれは6である。
\end{lemma}
\begin{proof}
複素数値の最高次パラメータを固定する。
補題\ref{lem:degree-one-elliptic-top-relative-cy}で最高次potentialのCY性が
確立しているので，Etingof--GinzburgのProposition 7.3.1(ii)から，
全ての低次potentialに対するPBW性とCY性が従う\cite{EG2010}。
この適用は具体的な有限開集合上の各CY topに対するものである。
同論文のgeneric potentialの主張を有限開条件と読み替える必要はない。

普遍係数への移行は次のように行う。
活性次数の上限 $N$ を固定し，最高次代数の次数 $0,\ldots,N$ の部分を
局所的に自由化する。その基底を単語の線型結合で $A_S$ に持ち上げる。
最高次関係式による次数降下から，これらは $F_NA_S$ を生成する。
持上げた元の間に $S$-線型関係があれば，任意の複素最高次パラメータと
任意の複素低次係数へ特殊化する。
固定topに対するPBW性により全係数が零になり，
被約有限型の普遍係数環上で元の係数も零となる。
従って局所的に基底が得られ，随伴次数同型と有限射影性が従う。
filtrationの短完全列は射影加群の列として分裂するので，
任意の非被約係数への基底変更にも両立する。

相対CY性には\eqref{eq:degree-one-elliptic-potential}のpotential対角複体を使う。
活性filtrationを入れた随伴次数複体は，完全なtopの対角複体の係数変更である。
filtrationは下に有界で各元は有限次数を持つため，
最高次数のcycleを境界として持ち上げる次数降下により元の複体も完全である。
自己双対性とこの完全性は任意の係数変更に両立する。
\end{proof}

\begin{remark}\label{rem:degree-one-general-cyclic-lower}
同じ証明は，$-h\sum c_{ij}\operatorname{Sym}(x^iy^j)$ の代わりに
重み6未満の任意の巡回potentialを加えた場合にも適用できる。
上の12係数の表示は，古典Weierstrass方程式と中心商の規格化に合わせたものである。
\end{remark}

% !TeX root = ../nc-del-pezzo.tex
\subsection{23未知数の中心方程式と普遍的な中心超曲面}

以下では前項の最高次係数環を $T$ と書く。
すなわち，$\Delta=4\alpha^3+9\beta^2$ と
$D_{\mathrm{AS},1}=(1+\alpha h^4)(5+\alpha h^4)$ を可逆化した
$\mathbb C[\alpha,\beta,h]$ である。
低次係数を加えた環と可換多項式を
\[
 S=T[c_{ij}:i+2j<6],\qquad
 F=\frac12z^2+\frac13y^3+\alpha x^4y+\beta x^6
       +\sum_{i+2j<6}c_{ij}x^iy^j
\]
と置く。添字集合は12個の元を持ち，定数項 $c_{00}$ も含む。
$A$ は前項のpotential代数である。

\begin{theorem}\label{thm:degree-one-finite-central-family}
$\Delta(\alpha_0,\beta_0)\ne0$ とし，
$p=(\alpha_0,\beta_0,0)$ とする。
$p$ の主開近傍 $V\subset\operatorname{Spec}T$ を適切に選ぶと，
$S_V=\mathcal O(V)[c_{ij}:i+2j<6]$ 上で次が成り立つ。
\begin{enumerate}
 \item $A$ は有限表示で，各活性filtration部分は有限射影である。
       任意の可換係数変更に対してPBW同型と相対3-CY性が保たれる。
 \item 活性filtration次数6の中心元 $\Omega$ があり，
       $\Omega\bmod h=F$ である。
       $\Omega$ を構成する局所化の分母は $\alpha,\beta,h$ だけに依存し，
       全ての低次係数は任意に動かせる。
 \item $\Omega$ は任意の係数変更後も左右の非零因子である。
       $B=A/(\Omega)$ は普遍的に平坦であり，
       \[
        \operatorname{gr}_{\mathrm{act}}B
         \simeq\bigl(A^{\mathrm{top}}/(\Omega^{\mathrm{top}})\bigr)
                 \otimes_{\mathcal O(V)}S_V,
        \qquad
        H_{A^{\mathrm{top}}/(\Omega^{\mathrm{top}})}(t)
          =\frac{1-t^6}{(1-t)(1-t^2)(1-t^3)}.
       \]
       各随伴次数部分は有限射影である。
 \item $\deg c_{ij}=6-i-2j$ とすると，$\Omega$ は全内部次数6の斉次元に選べる。
\end{enumerate}
\end{theorem}
\begin{proof}
PBW性と相対CY性は前項による。中心元の構成には
定理\ref{thm:elliptic-finite-central-family}の有限行列法を使うが，
ここでは必要な次数，核の計算，規格化を明記する。

まず最高次代数の次数9以下だけで整列語 $x^iy^jz^k$ を基底にする。
非整列語の隣接交換による関係列は，$h=0$ で
$w-\operatorname{ord}(w)$ となる。
非整列語成分の行列式は $\det(I+hK_d)$ の形で，$h=0$ で1である。
有限個のこれらの行列式を可逆化すれば，各次数で整列語が生成系となり，
既に得たHilbert級数と局所自由性から基底となる。
低次係数を戻したときの行列は活性次数に関して三角であり，
その対角ブロックは最高次の場合と同じである。
従って同じ，最高次係数だけに依存する局所化で $F_9A$ の基底を得る。

$A/hA$ が可換で，$A$ に $h$-捩れがないことから，
\[
 D:F_6A\longrightarrow F_7A\oplus F_8A\oplus F_9A,
 \qquad v\longmapsto h^{-1}([v,x],[v,y],[v,z])
\]
は正則な行列で表せる。
ここでも $h^{-1}$ は係数環の局所化ではなく，$hA$ 内の一意的な割算である。
重み $(1,2,3)$ の次数 $0,1,\ldots,6$ の単項式数は
\[
 1,\ 1,\ 2,\ 3,\ 4,\ 5,\ 7
\]
なので，$F_6A$ の階数は23である。

$h=0$ の任意の複素特殊化において
\begin{equation}\label{eq:degree-one-low-poisson-center}
 \ker D=\mathbb C\cdot1\oplus\mathbb C\cdot F
\end{equation}
である。実際，$U^{\mathrm{ab}}=F-z^2/2$ と置き，
$v=g_0(x,y)+zg_1(x,y)+kz^2$ と書く。
ここで $k\in\mathbb C$ であり，重みからこの形で尽くされる。
$\{v,x\}=0$ の $z^0$ 係数は
$ (\partial_yU^{\mathrm{ab}})g_1=0$ である。
$\partial_yU^{\mathrm{ab}}$ の最高 $y$ 次数の項は $y^2$ なので非零であり，
$g_1=0$ を得る。
残りの $\{v,x\}=\{v,y\}=0$ は
\[
 \partial_yg_0=2k\partial_yU^{\mathrm{ab}},\qquad
 \partial_xg_0=2k\partial_xU^{\mathrm{ab}}
\]
を与える。従って $v=2kF+\text{constant}$ である。

低次係数を零とし，出力の次数から対応する生成元の重みを引くと，
$D$ は活性次数ごとのブロックに分かれる。
$p$ における階数は，次数 $1,\ldots,6$ でそれぞれ
$1,2,3,4,5,6$ である。
各ブロックからこの本数の独立な行を選び，
定数項と整列基底における $z^2$ 係数を読む二つの行を加える。
得られる23次正方行列 $M_1$ は $p$ で可逆である。
分母を払った $\det M_1$ の非零開集合へ縮める。
全低次係数を戻しても，同じ行選択による行列は活性次数に関して三角で，
対角ブロックは変わらない。その行列式に低次係数は現れない。

選んだ21本の交換子方程式を零，定数項を $c_{00}$，
$z^2$ 係数を $1/2$ とした23本の線型方程式の唯一の解を $\Omega$ とする。
任意の $h\ne0$ の複素特殊化で，
Etingof--GinzburgのTheorem 3.3.2(ii)は $F_6A$ 内の非定数中心元を与える。
これは既に確認したCY最高次potentialに適用する全lower項の定理であり，
genericパラメータだけの中心定理ではない。
従って $\dim\ker D\ge2$ である。
$M_1$ の可逆性から，定数項と $z^2$ 係数を読む
$\ker D\to\mathbb C^2$ は単射だから同型であり，
規格化した解は未選択の交換子方程式も満たす。
$h=0$ では\eqref{eq:degree-one-low-poisson-center}により解は $F$ となる。
普遍係数環は被約で有限型なので，全複素点で零となる残りの方程式は
係数環上でも恒等的に零である。
こうして中心性と $\Omega\bmod h=F$ を得る。

最高活性部分は低次係数に依存せず，$z^2$ 係数が $1/2$ なので非零である。
最高次代数の全複素ファイバーは整域だから，
次数別の乗法写像 $A^{\mathrm{top}}_{m-6}\to A^{\mathrm{top}}_m$ は
全ファイバーで単射である。
有限局所自由加群の小行列式の判定により，これは局所分裂単射で余核も有限局所自由である。
従って中心元は任意基底変更後も正則である。
活性最高項の次数降下により低次項を含む乗法もstrictな単射となり，
表示した随伴次数商と普遍的平坦性を得る。
最後に低次係数へ重み $6-i-2j$ を与えたスカラー作用を使う。
全方程式と規格化は同変であり，解の一意性から $\Omega$ は内部斉次次数6である。
\end{proof}

\begin{remark}
$D_{\mathrm{AS},1}\ne0$ はambient正則性の明示的な十分条件である。
中心元の規格化には，さらに次数9以下の有限行列の非消滅条件を使う。
両者を混同して，全中心構成の基底を
$\Delta D_{\mathrm{AS},1}\ne0$ だけで尽くしたとは主張しない。
また，全次数を同時に整列語基底で自由化する一つの主開集合や，
有限開集合全体での等式 $Z(A)=S_V[\Omega]$ も主張しない。
\end{remark}

\subsection{全次元への拡張と古典Poisson対の範囲}

\begin{corollary}\label{cor:degree-one-elliptic-general-dimension}
$n\ge2$ とし，$n-1$ 個の中心一次変数 $s_1,\ldots,s_{n-1}$ を加える。
各 $c_{ij}$ に $s$ の次数 $6-i-2j$ の任意の斉次多項式を代入し，
その有限個の係数を独立なパラメータとする。
このとき，有限型係数環 $K$ 上で，有限表示かつ普遍的に次数別平坦な中心商族を得る。
全幾何ファイバーは左右Noetherであり，
\[
 H_{B_\xi}(t)=\frac{1-t^6}{(1-t)^n(1-t^2)(1-t^3)},\qquad
 B/(h)\simeq(K/(h))[s,x,y,z]/(F).
\]
最初の交換子は $s_i$ を中心とするJacobian括弧
\[
 \{x,y\}=z,\qquad \{y,z\}=\partial_xF,
 \qquad \{z,x\}=\partial_yF
\]
である。
\end{corollary}
\begin{proof}
局所化が最高次係数だけに依存し，定理が任意の係数変更と両立するので，
全ての表示した代入が許される。
内部次数ごとの平坦性とHilbert級数は，活性随伴次数による計算から従う。
体ファイバーの最高次ambientはNoether代数の中央多項式拡大であり，
filtered持上げと中心商もNoetherである。
第一交換子はpotential関係式の $h$ の一次項である。
ここではまだ，普遍係数環上の強Noether性を体ファイバーのNoether性から推論していない。
\end{proof}

\begin{proposition}\label{prop:degree-one-elliptic-semiclassical-coverage}
次数1の可換Jacobian模型を
\[
 F=\frac12z^2+\frac13y^3+g_4(s,x)y+g_6(s,x)
\]
と正規化し，選んだ非零Poisson方向に適合する $s$ を中心座標とする。
$g_4(0,x)=\alpha x^4$，$g_6(0,x)=\beta x^6$ が
$4\alpha^3+9\beta^2\ne0$ を満たすとき，この模型とPoisson方向は，
必要なら $h$ を非零スカラー倍して，上の族の半古典極限として実現される。
この条件は，古典Poisson対のパラメータ空間の非空な稠密開条件である。
\end{proposition}
\begin{proof}
全ての最高次以外の項を $c_{ij}(s)$ に代入すればよい。
重み6の一般式にある $z$ の一次項は平方完成で，$y^2$ 項は
$y$ の重み2の多項式による平行移動で除ける。
これらの三角変換は活性体積形式を保ち，残る定数倍は $h$ に吸収できる。
曲線 $z^2/2+y^3/3+\alpha x^4y+\beta x^6=0$ の滑らかさは
$\Delta\ne0$ と同値である。
非空性は，全 $n\ge2$ で
\[
 F=\frac12z^2+\frac13y^3+x^6+\sum_{i=1}^{n-1}s_i^6
\]
から分かる。このaffine coneは原点以外で滑らかであり，
重み付き射影空間の $y$ 軸・$z$ 軸上の安定化群が非自明な点を避ける。
従って滑らかな次数1 del Pezzo多様体を与える。
判別式非零は係数チャートで非空な主開条件なので稠密である。
\end{proof}

% !TeX root = ../nc-del-pezzo.tex
\subsection{次数1の相対楕円曲線と次数1の線束}
\label{subsec:degree-one-relative-curve}

定理\ref{thm:degree-one-finite-central-family}の可換最高次点
$p=(\alpha_0,\beta_0,0)$ とその近傍 $V$ を固定する。
以下では必要に応じて $V$ をさらに主開近傍へ縮める。
縮小は全て最高次係数だけに依存する。
補題\ref{lem:degree-one-stephenson-normal-form}の有限忠実平坦被覆を
$V'\to V$ と書き，その上で
\[
 a=1+kh^2,\qquad b=\rho h^2,\qquad c=-2,
 \qquad k^2=-\alpha,\quad \rho^3=6\beta-4k^3
\]
と置く。$V'$ の係数環は $\mathbb C[k,\rho,h]$ の局所化であり，
特に正則整域である。$\rho$ は可逆である。
さらに $9+\alpha h^4$ を可逆化してよい。
これにより $a-c=3+kh^2$ も可逆となる。
$a$ と $a+c$ の可逆性は既にAS正則開条件に含まれている。

\begin{lemma}\label{lem:degree-one-resolved-elliptic-curve}
$V$ を適切に縮めると，$V'$ 上に滑らかで固有な幾何学的整な種数1曲線
$f:E\to V'$，自己同型 $\sigma$，切断 $P,Q$ があり，次を満たす。
\begin{enumerate}
 \item $h=0$ では $E$ は
       $z^2/2+y^3/3+\alpha x^4y+\beta x^6=0$
       という重み付き射影曲線であり，$\sigma$ は恒等写像である。
 \item $L=\mathcal O_E(P)$ は相対次数1の線束である。
       $\sigma Q=P$ であり，$E$ から $\mathbb P^1$ への次数2の射 $p_1$ に対して
       \[
        L\otimes\sigma^*L\simeq\mathcal O_E(P+Q)
          \simeq p_1^*\mathcal O_{\mathbb P^1}(1).
       \]
 \item $L_0=\mathcal O_E$，
       $L_m=L\otimes\sigma^*L\otimes\cdots\otimes\sigma^{(m-1)*}L$ と置く。
       $f_*\mathcal O_E=\mathcal O_{V'}$ であり，
       $f_*L_m$ は $m>0$ で階数 $m$ の有限局所自由加群で，
       $R^1f_*L_m=0$ である。これらは任意の基底変更に両立する。
\end{enumerate}
\end{lemma}
\begin{proof}
\textbf{曲線の方程式と二つのchart。}
$(w_0:z_0),(w_1:z_1)$ を $\mathbb P^1\times\mathbb P^1$ の座標とし，
$t=z_0/w_0$，$q=z_1/w_1$ と書く。
Stephensonのpoint curve\cite[\S2.1, (2.1.9)]{Stephenson1997}を，
符号を変えて次の式で表す。
\begin{equation}\label{eq:degree-one-point-curve}
 \begin{split}
 H(t,q)={}&a(a^2-c^2)(t^2+q^2)
       +ab(a+c)(t^2q+tq^2)\\
       &+(a^2c-b^3-c^3)tq-ab^2(t+q)-a^2b.
 \end{split}
\end{equation}
同じ文字 $H$ でこの式の二重次数 $(2,2)$ の斉次化を表す。
$\delta=w_0z_1-z_0w_1$ とすれば $H|_{h=0}=-3\delta^2$ であり，
$H+3\delta^2$ は $h^2$ で割り切れる。
$\mathbb P^1\times\mathbb P^1\times V'$ を理想 $(h,\delta)$ でblow upし，
$H=0$ のstrict transformを $E$ とする。

有限座標の $h$-chartでは $q=t+hw$ と書け，その方程式は
$G(t,w)=h^{-2}H(t,t+hw)=0$ である。
ここで $G$ は多項式であり，$h$ を可逆化していない。
特殊方程式は
\begin{equation}\label{eq:degree-one-resolved-special-curve}
 G|_{h=0}=-3w^2-6kt^2-2\rho t^3-\rho=0.
\end{equation}
無限遠では $u=1/t$，$v=1/q$，$\eta=(v-u)/h$ を使う。
斉次方程式をこのchartで $h^2$ で割ると，特殊方程式は
\[
 -3\eta^2-6ku^2-2\rho u-\rho u^4=0.
\]
重なりでは $\eta=-w/(tq)$ である。
blowupのもう一つのchartで $h=e\delta$ と書くと，
strict transformの方程式は $-3+e^2(\cdots)=0$ となる。
従って $h=0$ ファイバーには $e=0$ の点がなく，上の二つの $h$-chartで尽くされる。

$x=1$ のchartにおける変数変換
\[
 y=\rho t+k,\qquad z=\rho w
\]
は
\begin{equation}\label{eq:degree-one-curve-classical-identification}
 \frac6{\rho^2}
 \left(\frac{(\rho w)^2}{2}+\frac{(\rho t+k)^3}{3}
                   -k^2(\rho t+k)+\frac{\rho^3+4k^3}{6}\right)
   =3w^2+6kt^2+2\rho t^3+\rho
\end{equation}
を満たす。無限遠の唯一の点は $(u,\eta)=(0,0)$ で，
そこでの $u$ 微分は $-2\rho\ne0$ である。
従って特殊曲線は表示した滑らかなWeierstrass曲線であり，
その滑らかさは $\Delta\ne0$ から従う。
この同定で唯一の無限遠点は $\mathcal O(1)$ の切断 $x$ の零点である。

$E\to V'$ は射影的である。
上のJacobian検査により，$h=0$ の全ファイバーの近傍で滑らかな相対曲線になる。
非平滑箇所の固有像を除き，必要なら幾何学的連結性の開条件も課すと，
滑らかで固有な幾何学的整な種数1曲線を得る。
こうした悪い閉集合は $h=0$ 上に点を持たない。
$V'\to V$ は有限なので，その像も閉集合である。
その像を $V$ から除き，$p$ を含む主開近傍を取れば，
被覆全体の上で所要条件を満たす。
以下で用いる行列の階数条件や有限次数の生成条件についても，
この有限像による同じ縮小を行う。

\textbf{shiftと逆写像。}
$C=a^2-c^2$ と置き，次の成分を持つ行列 $M=(m_{ij})$ を考える。
\[
 \begin{aligned}
 m_{11}&=c w_0z_1+a z_0w_1,\\
 m_{12}&=(a w_0+b z_0)w_1,\\
 m_{21}&=abc w_0w_1+b^2c w_0z_1+aCz_0z_1,\\
 m_{22}&=cCz_0w_1+ab(a+c)z_0z_1+aCw_0z_1.
 \end{aligned}
\]
直接展開すると $\det M=cH$ である。
$h=0$ の有限chartでは第1行が $(-t,1)$ で，
無限遠では第2行が $(-3,0)$ である。
従ってrank零の閉集合の固有像を除くと，$E$ 全体で $M$ のrankは1となる。
そのkernelが第三の点 $(w_2:z_2)$ を正則に定める。
有限chartでは，$r=z_2/w_2$ と書いて
\begin{equation}\label{eq:degree-one-point-shift}
 r=-\frac{cq+at}{a+bt},\qquad
 (a+bt)^2H(q,r)=(a^2-bcq)H(t,q)
\end{equation}
を得る。右の式は分母を払った多項式恒等式である。
従って $(t,q)\mapsto(q,r)$ は $h\ne0$ で曲線を保つ。

$h=0$ では第二点と第三点が一致する。
よって $(p_2,p_3)$ による $(h,\delta)$ の逆像は可逆理想 $(h)$ であり，
blowupの普遍性から $E$ からblowupへの射を得る。
その像は $h\ne0$ でstrict transformに入り，
$E$ の $h$-平坦性によって全体でもそこに入る。
これを $\sigma:E\to E$ とする。
有限chartでの持上げは
\begin{equation}\label{eq:degree-one-lifted-shift}
 t'=q=t+hw,\qquad
 w'=\frac{w-h\{k(t+q)+\rho tq\}}{a+bt}
\end{equation}
である。従って $h=0$ で $(t',w')=(t,w)$ であり，
無限遠を含めて $\sigma_0=\operatorname{id}$ となる。
分母が消える点では先のrank1行列で射を定めており，
この有理式だけで正則性を主張しているわけではない。

二つの $\mathbb P^1$ を交換する射 $J$ は，$H$ とblowupの理想を保つ。
外側の変数 $t,r$ に関する式 $a(t+r)+cq+btr=0$ の対称性から，
$J\sigma J$ は $\sigma$ の逆写像である。
これは $h\ne0$ の稠密開集合で確認でき，
分離性と $h$-平坦性によって全体に延びる。
従って逆写像についても未処理のbasepointはない。

\textbf{次数1の線束。}
元の積の上で
\[
 Q=((0:1),(0:1)),\qquad
 P=((0:1),(b:c-a))
\]
を取る。これらはblowupへ持ち上がり，無限遠chartでは
\[
 Q:(u,\eta)=(0,0),\qquad
 P:(u,\eta)=\left(0,\frac{\rho h}{c-a}\right)
\]
となる。chartの方程式への代入から両者の像は $E$ に入り，
滑らかな相対曲線の切断なので，相対次数1の有効Cartier因子である。
$Q$ における行列の第2行は
$\bigl(a(a^2-c^2),ab(a+c)\bigr)$ であり，
そのkernelは $(b:c-a)$ である。従って $\sigma Q=P$ である。

$p_1$ を第一射影とする。無限遠chartのstrict transformを $u=0$ に制限すると
\begin{equation}\label{eq:degree-one-infinity-divisor}
 G_\infty(0,\eta)
   =a(a+c)\eta\bigl((a-c)\eta+\rho h\bigr).
\end{equation}
これが二つの切断を与え，$h=0$ では重複度2で合流する。
従ってCartier因子として $p_1^{-1}(\infty)=P+Q$ である。
同じことは，正則な全空間上で $h\ne0$ の因子を比較し，
双方に垂直因子 $h=0$ がないことからも従う。
よって $L=\mathcal O_E(P)$ に対し
$L\otimes\sigma^*L=\mathcal O_E(P+Q)=p_1^*\mathcal O(1)$ となる。
特殊ファイバーで $L$ は重み付き曲線の $\mathcal O(1)$ に一致する。

最後に，種数1曲線の正次数線束に対するSerre双対とRiemann--Rochにより，
全幾何学的ファイバーで $H^1(L_m)=0$，$h^0(L_m)=m$ である。
固有平坦な曲線上のcohomology and base changeから，
有限局所自由性と任意の非被約基底変更も含む両立性を得る。
\end{proof}

\subsection{次数1の最高次中心商の相対THCR表示}

\begin{theorem}\label{thm:degree-one-relative-thcr}
$V$ をさらに有限回縮めると，有限忠実平坦被覆 $V'\to V$ の上で
\[
 B^{\mathrm{top}}_{V'}
   :=\bigl(A^{\mathrm{top}}/(\Omega^{\mathrm{top}})\bigr)_{V'}
       \simeq D:=\bigoplus_{m\ge0} f_*L_m
\]
が成り立つ。右辺の乗法は
$s\cdot t=s\otimes\sigma^{m*}t$，$s\in f_*L_m$ で定める。
同型は全ての可換基底変更に両立する。
\end{theorem}
\begin{proof}
$L=\mathcal O_E(P)$ の切断を有理関数として表示し，
$x=1\in H^0(L)$ とする。
前補題の因子等式から
$Y=t\in H^0(L_2)$ と取れる。
この規約で $x^2=1$，$xY=q$，$Yx=t$ である。
従って
\[
 y=\rho Y+kx^2\in D_2,\qquad
 z=\rho\,\frac{xY-Yx}{h}\in D_3
\]
を定められる。後者の正則性を説明する。
$D_3$ は係数環上平坦で，基底変更に両立し，$h=0$ では $\sigma=\operatorname{id}$ なので
$D/(h)$ は可換である。従って $xY-Yx\in hD_3$ であり，
$h$-捩れがないことから割算が一意に存在する。
有限chartで $z=\rho w$ となり，これは単に $h^{-1}$ を導入した元ではない。

これらの切断がambientの関係式を満たすことを示す。
$h\ne0$ の上では $r=t\circ\sigma^2$，$s=t\circ\sigma^3$ と置くと，
\eqref{eq:degree-one-point-shift}から
\[
 a(t+r)+cq+btr=0.
\]
さらに
\begin{equation}\label{eq:degree-one-second-relation-certificate}
 (a+bt)(a+bq)\{a(tr+qs)+cts+b\}=-H(t,q)
\end{equation}
である。これらはそれぞれ $f_4$，$f_5$ の $D$ における零像を意味する。
補題\ref{lem:degree-one-stephenson-normal-form}と $z=h^{-1}[x,y]$ から，
$A^{\mathrm{top}}\to D$ を得る。
各 $D_m$ に $h$-捩れがないので，この写像の関係式は $h=0$ を含めて成立する。

$h=0$ ではこの写像はWeierstrass曲線のsection環への通常の写像である。
無限遠での $t,w$ の極の位数はそれぞれ2,3であり，
次数 $1,2,3,4$ の切断空間は，有理関数表示でそれぞれ
\[
 (1),\qquad (1,t),\qquad (1,t,w),\qquad(1,t,t^2,w)
\]
を基底に持つ。最後の基底は $x^4,x^2y,y^2,xz$ の像でも張られる。
従ってこれら四次数で全射である。
各次数の余核は有限係数加群なので，その支持の像を除く有限開条件で，
被覆全体の上でもこれら四次数の全射性を確保できる。
全次数への拡張には次数2の線束を使う。
$m\ge3$ とし，$M=\sigma^{m*}L_2$ と置く。
$M$ は相対的に大域生成され，評価列のkernelは
\[
 0\longrightarrow M^{-1}\otimes f^*\det(f_*M)
   \longrightarrow f^*f_*M\longrightarrow M\longrightarrow0
\]
である。これを $L_m$ でtensorすると，kernelのファイバー次数は $m-2>0$ である。
その $H^1$ が消えるので，任意の係数変更後にも
\begin{equation}\label{eq:degree-one-thcr-generation}
 D_m\otimes D_2\longrightarrow D_{m+2}
 \quad(m\ge3)
\end{equation}
は全射となる。従って次数 $1,\ldots,4$ の全射性から $A^{\mathrm{top}}\to D$ の全射性が従う。
この段階で初めて $\Omega^{\mathrm{top}}$ の像を $D$ の中心元として扱える。

一般点上でこの中心元の像は零である。
実際，\eqref{eq:degree-one-lifted-shift}より
\[
 t\circ\sigma^6=t+6hw\pmod {h^2}
\]
であり，特殊曲線上で $w$ は零関数ではないから，
$\sigma^6$ は一般点で恒等写像ではない。
一方，体上のTHCRに次数6の非零中心切断 $v$ があれば，
$s_1,s_2\in H^0(L_m)$ との可換性を割り算して
\[
 \sigma^{6*}(s_1/s_2)=s_1/s_2
\]
を得る。$m\ge3$ では $L_m$ は非常に豊富なので，
その切断の比は関数体を生成し，$\sigma^6=\operatorname{id}$ となる。
これは矛盾である。$D_6$ は整な係数環上有限射影なので，
一般点で零となる $\Omega^{\mathrm{top}}$ の像は全体で零である。

従って $B^{\mathrm{top}}_{V'}\twoheadrightarrow D$ を得る。
定理\ref{thm:degree-one-finite-central-family}によれば，左辺の各次数は有限局所自由で，
階数は次数0で1，正次数 $m$ で $m$ である。
右辺も同じ階数を持つので全射は同型である。
構成した写像と両辺の各次数が基底変更に両立するため，
同型は任意の非被約基底変更後にも保たれる。
\end{proof}

\subsection{強Noether性と積代数}

\begin{corollary}\label{cor:degree-one-strong-noetherian}
上で選んだ $V$ の上では，最高次ambientと中心商は強左右Noetherである。
同じことが，全低次係数を許す族と
系\ref{cor:degree-one-elliptic-general-dimension}の全次元の斉次族にも成り立つ。
すなわち，許される任意のnoetherian可換係数環への変更後に左右Noetherである。
\end{corollary}
\begin{proof}
まずnoetherianな $\mathcal O(V')$-代数 $R$ を取る。
$E_R$ は $R$ 上射影的であり，$L_R$ は相対次数1なので相対的に豊富である。
自己同型 $\sigma_R$ は全ファイバーの線束の次数を保ち，
相対数値同値群に恒等的に作用する。
従ってKeelerの $\sigma$-ample判定とNoether定理
\cite[Theorems 7.13, 7.14]{Keeler2003}を適用でき，$D_R$ は左右Noetherである。
次数1の線束そのものの大域生成は仮定していない。

最高次ambientは正次数6の普遍的な中央非零因子を持ち，その商が $D_R$ である。
中央元での随伴次数は $D_R[v]$ である。
固定した内部次数の元に対する中央元による割算は有限回で止まるので，
斉次イデアルの生成元を持ち上げる議論によりambientも左右Noetherとなる。
通常のイデアルには内部次数の最高項を用いる。

次にnoetherianな $\mathcal O(V)$-代数 $R$ を取り，
$R'=R\otimes_{\mathcal O(V)}\mathcal O(V')$ と置く。
$R'$ は $R$ 上有限忠実平坦で，特にnoetherianである。
右イデアルの上昇列を $R'$ へ変更すると上で証明したNoether性によって停止する。
平坦性と忠実性から元の列も停止する。左イデアルも同様である。
これで最高次代数についての強Noether性が $V$ へ降下する。

全低次係数を含むambientには普遍PBW filtrationを使う。
その随伴次数は強Noetherな最高次代数の係数変更なので，filtered持上げでNoether性が従う。
中心商もNoetherである。
全次元の族では中央多項式環 $R[s_1,\ldots,s_{n-1}]$ を係数環として同じ議論を行う。
\end{proof}

\begin{corollary}\label{cor:degree-one-product-noetherian}
上記の最高次族，低次係数を含む族，および全次元の斉次族について，
許される任意のnoetherian基底 $R$ の上で
\[
 A_R^{\mathrm{op}}\otimes_R A_R,\qquad
 B_R^{\mathrm{op}}\otimes_R B_R
\]
は左右Noetherである。
\end{corollary}
\begin{proof}
有限忠実平坦被覆上でまず $D^{\mathrm{op}}\otimes_R D$ を考える。
偶数Veroneseは
\[
 D^{(2)}=D(E,L_2,\sigma^2)
\]
であり，$L_2$ の相対次数は2である。
反対環も，自己同型 $\sigma^{-2}$ と次数2の線束
$(\sigma^{-2})^*L_2$ によるTHCRとして表示できる。
従って補題\ref{lem:elliptic-bisection-noetherian}により
$(D^{(2)})^{\mathrm{op}}\otimes_R D^{(2)}$ は左右Noetherである。
同補題は任意のnoetherian基底に対する二つの次数2線束の主張なので，
ここで非被約基底も許される。

\eqref{eq:degree-one-thcr-generation}と，左右を交換した同じ評価列の証明により，
\[
 D_mD_2=D_{m+2}=D_2D_m\qquad(m\ge3).
\]
従って $D$ は $D^{(2)}$ の左加群としても右加群としても有限生成である。
具体的には偶数部分には $1$，奇数部分には $D_1$ と $D_3$ の有限な係数生成元を取ればよい。
よって $D^{\mathrm{op}}\otimes_R D$ は
$(D^{(2)})^{\mathrm{op}}\otimes_R D^{(2)}$ の左右の有限生成加群である。
その右イデアルは後者の環の右部分加群なので上昇列条件を満たし，左も同様である。

最高次ambientの積には中央元
$\Omega^{\mathrm{top}}\otimes1$ と $1\otimes\Omega^{\mathrm{top}}$ を順に戻す。
各因子と各中心商の $R$-平坦性により，これらは途中の商でも普遍的な非零因子である。
二度の正次数中央拡大で積のNoether性が得られる。
さらに前系と同じ忠実平坦降下を行う。
低次項と全次元への拡張にはtensor積の活性filtrationを使う。
その随伴次数は，二組の中央変数を加えた係数環上の最高次ambientの積であり，
今証明した任意係数変更後のNoether性を持つ。
filtered持上げと二つの中心元による商を取って結論を得る。
\end{proof}

\subsection{次数1の相対tailsと有限window}

\begin{theorem}\label{thm:degree-one-relative-tails}
上で選んだ最高次開集合上の全次元の斉次族を，
noetherianな可換係数環 $\Lambda$ へ変更する。
$B=A/(\Omega)$，$\mathcal Q=\QGr B$，$r=n-1$，$l=n+5$ と置くと，次が成り立つ。
\begin{enumerate}
 \item $A$ は $\Lambda$ 上homologically smoothであり，
       $A^!_\Lambda\simeq A(n+5)[-n-2]$ である。
 \item 左右の相対局所コホモロジーは
       $R\Gamma_{\ge1}(B)\simeq B^\vee(n-1)[-n-1]$ である。
       第\ref{sec:finite-tails}節の飽和性，有向ACM，相対 $\chi$，
       全シフトの豊富性，compact生成，proper性，相対Serre公式が成立する。
       特にcompact圏のSerre関手は $S_\Lambda\simeq(-r)[n]$ である。
 \item $E=\bigoplus_{j=0}^{n+4}\mathcal O(j)$ はcompact生成元であり，
       $R\!\operatorname{End}(E)$ は $\Lambda$ 上perfectである。
\end{enumerate}
これらは許される任意のnoetherian係数変更に両立する。
さらに，全可換斉次式が滑らかなdel Pezzo多様体を与える点の有限開近傍では，
正周期 $r=n-1$ の列自体も相対的に豊富である。
\end{theorem}
\begin{proof}
活性代数の相対3-CY同型の内部shiftは6である。
$n-1$ 個の中央一次変数の左右の差によるKoszul解消を組み合わせると，
定理\ref{thm:elliptic-relative-tails}の第一段落と同じ随伴計算から
\[
 A^!_\Lambda\simeq A(6+n-1)[-3-(n-1)]
\]
を得る。差の列の正則性は中央係数環上の平坦性から従い，
任意の係数変更後も保たれる。

強Noether性，各内部次数の有限射影性，次数6の普遍的な中央正則元が既に得られている。
従って第\ref{sec:finite-tails}節の相対双対，局所コホモロジー，
smashing局所化およびSerre双対の証明を適用できる。
有限自由性を用いた箇所は有限射影性だけで十分であり，
中央元のchange of ringsで $l-6=n-1$ を得る。
torsionの比較に現れる $\operatorname{Tor}_1$ は長さ6の次数帯に支持を持ち，
十分離れた切断の間では遷移が零になる。
このため，非正則な基底でも同節の局所化と双対の議論が成り立つ。

有限windowのためにaugmentation解消を明記する。
最高次potentialの普遍的に完全な対角解消を片側のaugmentationでtensorすると，
\[
 \begin{aligned}
 0&\longrightarrow A^{\mathrm{top}}(-6)
   \longrightarrow A^{\mathrm{top}}(-5)\oplus A^{\mathrm{top}}(-4)
                           \oplus A^{\mathrm{top}}(-3)\\
  &\longrightarrow A^{\mathrm{top}}(-1)\oplus A^{\mathrm{top}}(-2)
                           \oplus A^{\mathrm{top}}(-3)
   \longrightarrow A^{\mathrm{top}}\longrightarrow\Lambda\longrightarrow0
 \end{aligned}
\]
を得る。$h\ne0$ では次数3の可縮対があってもよく，
この表示は $h=0$ と非被約係数も含めて完全である。
片側tensorの完全性には，対角解消が片側射影加群の分裂解消であることを使う。
実際，各次数が係数環上射影なのでその自由bimodule項は片側射影であり，
augmentation項 $A^{\mathrm{top}}$ 自体も片側自由である。

中央変数の補助重みを1，活性変数の補助重みを $(M,2M,3M)$，$M>1$ とする。
低次potentialの項 $c_{ij}(s)x^iy^j$ の補助次数は
\[
 M(i+2j)+(6-i-2j)<6M\qquad(i+2j<6)
\]
なので，随伴次数ambientは最高次代数と中央多項式環のtensor積である。
上の解消と中央変数のKoszul解消をtensorし，
命題\ref{prop:finite-window-proper}の最高記号を消去する手順でfilteredに持ち上げる。
全補助重みが正なので，固定した基底ベクトルの誤差補正は有限回で終わる。
普遍係数上で得た解消の随伴次数は任意係数変更後にも完全であり，
同じ次数降下で解消の完全性が保たれる。
両端は $A(-n-5)$ と $A$ が一つずつ，中間の内部shiftは $1,\ldots,n+4$ である。

$B$ とtensorしたときのTorは $\Lambda$ と $\Lambda(-6)$ だけであり，双方ともtorsionである。
従ってtailsで有限完全複体を得て，長さ $n+5$ の任意の連続したtwistから
両方向に全twistを生成できる。これが表示したcompact生成元を与える。
最後に命題\ref{prop:finite-periodic-generation}の有限次数検査と
全係数変更後のSerre消滅を適用すれば，滑らかな可換点の有限開近傍で周期 $r$ の豊富性を得る。
\end{proof}

\begin{remark}\label{rem:degree-one-relative-scope}
曲線表示を用いたのは活性最高次中心商であり，その構成には有限忠実平坦被覆を許した。
Noether性と積のNoether性は被覆から元の最高次開集合へ降下している。
全lower項を持つ高次元のtailsを曲線のcoherent圏と同一視するものではない。
また，有限位数の $\sigma$ を持つ特殊点を除いておらず，中心全体の普遍的な記述は含まない。
第\ref{sec:finite-smoothness}節では，ここで得た積のNoether性と相対tailsの性質を使い，
全斉次式が滑らかな可換模型を与える点のまわりで，
さらに有限回縮めてsmooth properな族を得る。
この結論は，初めの明示的AS正則領域の全点や，全lower係数でのsmoothnessを意味しない。
\end{remark}

% !TeX root = ../nc-del-pezzo.tex
\section{有限パラメータ族の滑らかな開近傍}
\label{sec:finite-smoothness}

本節では両次数の分離型族と，次数2の一般楕円型族について，
滑らかな可換点を含む有限型開集合上で対角のperfect性を証明する．
properなdg代数一般についてsmoothnessの開性を仮定するのではなく，
二重次数付き対角の有限部分解消と，既に得た相対Serre双対を使う．
対角への射の持上げに必要な消滅は，特殊ファイバーの
$HH^1(X)=0$ と $X\times X$ の有限ホモロジー次元から得る．

以下，$S$ は有限型可換 $\C$-代数とし，$B$ は次のいずれかの全次元の斉次族を
$S$ へ係数変更したものとする：第\ref{sec:finite-tails}節の次数1・2の分離型族，
または定理\ref{thm:elliptic-relative-tails}の次数2の楕円型族．
後者では，各古典楕円型最高次点について同節で構成した有限開近傍上で作業する．
lower係数は任意であり，その中央変数への斉次多項式代入も許す．
$\mathcal Q=\QGr B$ と書き，$S$ の被約性・正則性は仮定しない．
$r=n-1$，$l=n+5$（次数1），$l=n+3$（次数2）とする．
\[
 E=\bigoplus_{i=0}^{l-1}\mathcal O(i),\qquad
 C=R\!\End_{\mathcal Q}(E),\qquad
 D=C^{\mathrm{op}}\otimes_S^L C
\]
と置く．$C$ は $S$-perfectであり，$C^*=R\!\Hom_S(C,S)$ は
定理\ref{thm:finite-tails-relative-serre}または
定理\ref{thm:elliptic-relative-tails}のSerre関手を表す可逆bimoduleである．

\subsection{Serre双対を使った逆向きのHomの有限性}

\begin{lemma}\label{lem:serre-pseudoperfect-to-perfect}
$D$ を $S$-perfectなdg代数とし，
$U=R\!\Hom_S(D,S)$ が可逆 $D$-bimoduleであるとする．
$M$ を $S$ 上perfectな右 $D$-加群，$P$ をperfectな右 $D$-加群とする．
このとき $R\!\Hom_D(M,P)$ は $S$ 上perfectで，任意の係数変更と可換する．
ここでは $M$ が $D$ 上perfectであることを仮定しない．
\end{lemma}
\begin{proof}
$Q=P\otimes_D^L U^{-1}$ と置くと $Q$ はperfectである．
coinductionの随伴から
\[
 R\!\Hom_D(M,U)\simeq R\!\Hom_S(M,S)
\]
を得る．この式を $Q=D$ から有限coneと直和因子に拡張すると
\begin{equation}
 R\!\Hom_D(M,P)
 \simeq R\!\Hom_S\bigl(R\!\Hom_D(Q,M),S\bigr).
 \label{eq:serre-pseudoperfect-hom}
\end{equation}
$Q$ のperfect性と $M$ の $S$-perfect性により，右辺の内側のHomは
$S$-perfectであり，任意の係数変更に両立する．
その双対にも同じ性質がある．可逆bimoduleとその逆の係数変更も可逆であるため，
この同型自体が基底変更と両立する．
\end{proof}

この補題は上の $D=C^{\mathrm{op}}\otimes_S^LC$ に適用できる．
実際，$D$ の $S$-双対は $C^{\mathrm{op}}$ と $C$ のSerre bimoduleの外積であり，
両因子が可逆だから可逆である．

\subsection{二重次数付き積とcoherentな対角}

$R=B^{\mathrm{op}}\otimes_S B$ に，両因子の内部次数による
$\mathbb Z^2$-次数を与える．$\Gr_2 R$ から，いずれか一方の次数について
torsionな加群で生成される局所化部分圏を割ったAbel圏を $\mathcal Q_2$ と書く．
商関手を $\pi_2$，自由加群の像を
$\mathcal O(a,b)=\pi_2 R(a,b)$ と書く．
この構成は二つのirrelevant idealの商を順に取るものであり，
総次数だけによる通常の $\QGr R$ とは区別する．

\begin{lemma}\label{lem:finite-family-product-diagonal}
次の性質が成り立つ．
\begin{enumerate}
\item $R$ は左右noetherianであり，この性質は任意のnoetherian係数変更で保たれる．
\item $\mathcal Q_2$ の導来圏は，外積twist $\mathcal O(a,b)$ でcompact生成される．
特に
\[
 G=\bigoplus_{0\le i,j<l}\mathcal O(-i,j)
\]
はcompact生成元で，$R\!\End_{\mathcal Q_2}(G)\simeq D$ である．
\item 二重次数付き右 $R$-加群
\[
 \Delta_{a,b}=\begin{cases}B_{a+b}&a,b\ge0,\\0&\text{それ以外}\end{cases},
 \qquad \xi\cdot(u^{\mathrm{op}}\otimes v)=u\xi v
\]
は有限生成かつ $S$-flatである．
$\mathcal D=\pi_2\Delta$ のMorita像
\[
 M=R\!\Hom_{\mathcal Q_2}(G,\mathcal D)
\]
は標準の対角 $C$-$C$-bimoduleであり，$S$-perfectである．
これらのHom，対角の同定およびMorita模型は任意の許される係数変更と両立する．
\end{enumerate}
\end{lemma}
\begin{proof}
まずnoetherian性を確認する．分離型の場合，ambient $A$ の中央変数の補助次数を1，
活性変数の補助次数を次数2では $(6,8,12)$，次数1では $(12,20,30)$ とする．
命題\ref{prop:finite-window-proper}の証明により，
$\operatorname{gr}_{\mathrm{aux}}A$ は全生成元を持つ可換多項式環であり，
PBW基底によってこのfiltrationは $S$ 上分裂する．
従って積filtrationについて
\[
 \operatorname{gr}_{\mathrm{aux}}(A^{\mathrm{op}}\otimes_S A)
   =\operatorname{gr}_{\mathrm{aux}}A^{\mathrm{op}}
       \otimes_S\operatorname{gr}_{\mathrm{aux}}A
\]
も可換多項式環である．filtered--gradedのnoetherian判定により
$A^{\mathrm{op}}\otimes_S A$ は左右noetherianであり，その商 $R$ も同じである．
任意のnoetherian係数変更後に同じPBW基底と補助次数を使えるので，分離型について(1)を得る．
この議論では，この補助次数に関する $\Omega$ の最高記号の形は使っていない．
楕円型の場合は系\ref{cor:elliptic-product-noetherian}を適用する．
これは相対曲線の二重切断環と中央正則元から積代数のnoetherian性を証明する結果であり，
一般には成り立たない「二つのNoether代数のtensor積もNoether」という推論をしていない．

次に，定理\ref{thm:finite-tails-relative-serre}の証明で用い，
定理\ref{thm:elliptic-relative-tails}でも確認した
左右共通の局所化bimodule $L$ を使う．
$B$ は $S$-flatなので，$R$ は各因子上flatである．
従って各因子からのinductionは完全で，その右随伴であるrestrictionはinjectiveを保つ．
このため，一方の次数についてのtorsionの導来関手を，残りの次数と可換する作用を保持したまま，
その因子上の既存の計算で求められる．まず下に有界な複体でこの比較を行い，
各torsion関手の有限cohomological dimensionによって非有界導来圏へ延長する．
第一因子と第二因子のsmashing局所化は可換し，その合成は
$L^{\mathrm{op}}\otimes_S^L L$ との導来tensorである．
各局所化の核は対応するirrelevant torsionであり，有限cohomological dimensionを持つ．
したがって二つの商を順に取った導来圏は，この合成局所化の像である．
具体的に，局所化された自由twist間のHomは
\begin{equation}
 R\!\Hom_{\mathcal Q_2}
       (\mathcal O(a,b),\mathcal O(c,e))
   \simeq L^{\mathrm{op}}_{c-a}\otimes_S^L L_{e-b}.
 \label{eq:product-twist-hom}
\end{equation}
局所化は直和を保つので，全外積twistはcompactであり，
それらからのHomが全て零なら局所化された加群自身が零となる．
有限windowの完全複体を第一因子と第二因子でそれぞれ使うと，
任意の外積twistが $G$ のthick閉包に入る．
反対側のtwistの符号を逆に選んだため，
第一因子のHom複体と合成は $C^{\mathrm{op}}$ のものになり，第二因子は $C$ のものになる．
上のHom公式と合成の外積性から(2)が従う．

$B$ の斉次生成元の重みの最大値を $w$ とする．
重み付き単語を内部次数 $a$ の境界で二つに切ると，
境界をまたぐ生成元が存在する場合でもその両側の残りの重みは $w$ 未満である．
従って $\Delta$ は有限個の成分
\[
 \Delta_{a,b}=B_{a+b}\qquad(0\le a,b<w)
\]
で $R$-加群として生成される．これらの成分は有限 $S$-加群であるから有限生成性を得る．
各成分が有限射影なので $S$-flatでもある．

対角の局所化を明示する．$\Delta$ を全二重次数に延長した
$\widetilde\Delta_{a,b}=B_{a+b}$ を用いてよい．
実際，$\Delta\subset\widetilde\Delta_{\{a\ge0\}}\subset\widetilde\Delta$ の
二つの商はそれぞれ第二・第一方向でtorsionであり，二重商では消える．
一方の次数を固定すると，もう一方の次数については $B$ のtwistである．
従って一方を局所化した成分は $L_{a+b}$ である．
残りの因子による局所化は $L\otimes_B^LL\simeq L$ により既に恒等となる．
したがって
\[
 R\!\Hom_{\mathcal Q_2}
       (\mathcal O(-i,j),\mathcal D)\simeq L_{i-j}.
\]
これは右 $C^{\mathrm{op}}\otimes C$-加群としての標準対角の $(i,j)$ 成分，
すなわち $R\!\Hom_{\mathcal Q}(\mathcal O(j),\mathcal O(i))$ である．
左右の作用も $u\xi v$ の式から通常の合成に一致する．
定理\ref{thm:finite-tails-cohomology}とその楕円型族への拡張により，この有限個の成分は
$S$-perfectで，任意の係数変更に両立する．
外積twist間のHomにも同じ性質があるから，(3)と全ての基底変更の主張を得る．
\end{proof}

\subsection{対角の有限部分解消からsmoothnessを得る}

\begin{theorem}[有限族のsmoothnessの有限開近傍]
\label{thm:finite-family-smooth-open}
$s\in\Spec S$ を複素閉点とし，$B_s$ が滑らかな可換del Pezzo多様体
$X$ の切断環であるとする．
このとき $s$ のZariski開近傍 $V$ 上で $C$ の対角は
$C^{\mathrm{op}}\otimes_S^LC$ 上perfectである．
従って $D(\QGr B)^c$ は $V$ 上smooth properであり，
任意の許されたnoetherian基底変更後にもsmooth properである．
\end{theorem}
\begin{proof}
補題\ref{lem:finite-family-product-diagonal}のnoetherian性により，
$\Delta$ の二重次数付き有限生成自由加群による解消
\[
 \cdots\longrightarrow F_2\longrightarrow F_1
 \longrightarrow F_0\longrightarrow\Delta\longrightarrow0
\]
を取れる．$R$ と $\Delta$ は $S$-flatなので，短完全列を順に使うと
全てのsyzygyが $S$-flatである．従ってこの解消は任意の係数変更後も完全である．
$q>2n$ を一つ固定し，次数 $-q,\ldots,0$ に置いた有限複体
\[
 P_q=[\pi_2F_q\longrightarrow\cdots\longrightarrow\pi_2F_0]
\]
を取る．$K_q$ をその左端の核とすると，$\mathcal Q_2$ の導来圏に三角
\begin{equation}
 K_q[q]\longrightarrow P_q\xrightarrow{f}\mathcal D
      \longrightarrow K_q[q+1]
 \label{eq:finite-diagonal-truncation}
\end{equation}
を得る．$K_q$ はcoherentであり，$P_q$ は外積twistの有限複体だからcompactである．
以後，Morita同値の下で同じ記号 $P_q,f$ を使い，$\mathcal D$ の像を $M$ と書く．

特殊ファイバーでは，補題\ref{lem:weighted-free-locus}により
$\mathcal Q_{2,s}=\Qcoh(X\times X)$ であり，
$\mathcal D_s=\mathcal O_{\Delta_X}$ となる．
滑らかで固有な $2n$ 次元多様体 $X\times X$ では，coherent層間の
$\Ext^j$ は $j>2n$ で零である．
三角\eqref{eq:finite-diagonal-truncation}に
$R\!\Hom(\mathcal O_{\Delta_X},-)$ を適用すると，$q>2n$ により
\[
 \Hom(\mathcal D_s,P_{q,s})\xrightarrow{\sim}
       \Hom(\mathcal D_s,\mathcal D_s),\qquad
 \Ext^1(\mathcal D_s,P_{q,s})\xrightarrow{\sim}
       \Ext^1(\mathcal D_s,\mathcal D_s)
\]
となる．右側の第一の群は恒等射を含み，第二の群は
\[
 HH^1(X)=H^0(X,T_X)\oplus H^1(X,\mathcal O_X)=0
\]
である．ここではHKRとKodaira消滅，および次数1・2の全滑らか可換模型についての
$H^0(T_X)=0$ \cite[補題8.7]{PoissonDP}を適用した．
従って $f_s$ の右逆 $i_s:M_s\to P_{q,s}$ があり，
$\Ext^1_{D_s}(M_s,P_{q,s})=0$ である．

次が有限型開近傍へ移る際の要点である．
補題\ref{lem:serre-pseudoperfect-to-perfect}により
\[
 K=R\!\Hom_D(M,P_q)
\]
は $S$-perfectで，$K\otimes_S^L\C(s)$ は特殊ファイバーの同じHom複体に一致する．
従って $H^1(K\otimes_S^L\C(s))=0$．
$s$ の近傍で $K$ を有界有限自由複体で表す．
特殊ファイバーで非零な微分の小行列式を有限個反転し，
単元を持つ微分をcontractibleな二項直和因子として消去すると，
特殊ファイバーの微分が全て零の模型を取れる．
その次数1の項は $H^1(K_s)=0$ により零である．
従ってこの模型の次数0の元は全てcycleであり，
$i_s\in H^0(K_s)$ は近傍を縮めた後に $i\in H^0(K)$ へ持ち上がる．
この段階では単にcycleを係数ごとに持ち上げたのではなく，
次数1の項の消滅によって微分条件も満たしている．

合成 $g=fi:M\to M$ は特殊ファイバーで恒等である．
$M$ は $S$-perfectなので $\operatorname{Cone}(g)$ も $S$-perfectであり，
特殊ファイバーで零である．有限自由複体の支持を除くことにより，
さらに $s$ の開近傍を縮めればこのconeは零になる．
従って $g$ は $D$-加群として同型である．
$ig^{-1}$ は $f$ の右逆となり，$M$ はperfectな $D$-加群 $P_q$ の直和因子である．
これが対角のperfect性である．

ここで行った局所化は，有界有限自由複体の消去と支持の除去という有限個の操作だけである．
従って得られた $V$ は有限型基底のZariski開集合であり，形式近傍に限定されない．
対角の有限cone・直和因子による表示は導来基底変更で保たれ，
proper性とMorita模型の基底変更は既に証明済みなので最後の主張も従う．
\end{proof}

\begin{corollary}[有限型族が非可換del Pezzo条件を満たす開集合]
\label{cor:finite-family-ncdp-open}
定理\ref{thm:finite-family-smooth-open}の開近傍を，
系\ref{cor:finite-periodic-ampleness}または
定理\ref{thm:elliptic-relative-tails}の反標準豊富性の開近傍と交わらせる．
この開集合上の族は
$\mathfrak M^{\OO,\mathrm{hyp}}_{n,d}$ の族を定める．
特に，各幾何学的ファイバーのcoherent圏のホモロジー次元はちょうど $n$ で，
その有界導来圏はcompact圏と一致する．
\end{corollary}
\begin{proof}
任意の点 $t\in V$ に対し $k=\overline{\kappa(t)}$ を選び，
対応する幾何学的ファイバー上のcoherent対象 $F$ を取る．
命題\ref{prop:finite-full-shift-ample}とその楕円型族への拡張，
および有限cohomological dimensionにより，
$R\!\Hom(E,F)$ は $k$ 上の有界有限次元複体である．
smoothなdg代数では，基底上perfectな加群はperfectである：
有限projective対角分解をその加群とtensorすればよい．
従って $F$ はcompactであり，coherentな有界複体もcompactである．
逆にcompact対象はtwistの有限coneと直和因子から作られるので，coherentな有界複体である．
Serre公式から
\[
 \Ext^i(F,F')\cong
       \Ext^{n-i}(F',F(-r))^*
\]
を得る．右辺は $i>n$ で零であり，
$\Ext^n(\mathcal O,\mathcal O(-r))\cong k$ によりホモロジー次元は $n$ である．

相対proper性・基底変更・Serre公式・有向ACMは
第\ref{sec:finite-tails}節および定理\ref{thm:elliptic-relative-tails}，
反標準豊富性は上記の開集合の選択により成り立つ．
基本シフトとその $r$ 乗による周期も同じ表示から得られる．
分離型の幾何学的ファイバーのambient AS正則性は
系\ref{cor:finite-relative-augmentation-ext}で確認した．
楕円型では定理\ref{thm:elliptic-relative-tails}の相対CY式と
\cite[Lemma 1.2]{RRZ2014}を幾何学的体ファイバーに適用すれば，同じambient AS正則性を得る．
中心元の普遍的正則性は，各族の構成で証明済みである．
従って定義\ref{def:nc-dp-stack}および重み付き超曲面部分問題の全条件が成立する．
\end{proof}

\begin{remark}[この開近傍定理の範囲]
構成した族の全ての有限パラメータでsmoothnessを主張したわけではない．
本節は滑らかな可換点を含むZariski開集合を構成し，
それが定義した非可換多様体の条件を満たすことを示す．
また非正則基底上でcoherentな有界導来圏全体を相対perfect圏と同一視しない．
有限型の表示族を得たことから，全ての非可換点を被覆するatlasや
全モジュライの成分数が従うわけでもない．
次数2の楕円型族で用いたのは，片側の強Noether性だけでなく，
$B^{\mathrm{op}}\otimes_S B$ のnoetherian性である．
開集合は有限個の行列と有限自由複体から得られるが，その判別式の明示公式はここでは計算していない．
\end{remark}

\subsection{現在の表示族をatlasにする際の接空間の障害}

\begin{proposition}[一つの量子化パラメータによる表示の限界]
\label{prop:one-parameter-atlas-obstruction}
本稿の分離型族または楕円型族の量子化パラメータを $h$ と書く．
$h=0$ の滑らかな可換点 $s$ の像を $X$ とし，
$h$ に関する一次交換子が定めるPoisson構造を $\pi_s$ とする．
この表示族のKodaira--Spencer写像の像を，通常の可換変形の接空間で割ると，
その像は $\C\pi_s$ に含まれ，次元は高々1である．
従って，これらの表示族から
$\mathfrak M^{\OO}_{n,d}$ への射は $s$ で形式的にsmoothではない．
係数や座標枠を動かした同種の族を並べても，そのまま可換点を含むsmooth atlasにはならない．
\end{proposition}
\begin{proof}
活性生成元の関係式は $[u,v]=hJ_{uv}$ の形であり，中央変数の交換子は零である．
双対数上で $h=\epsilon a$，他のパラメータを
$\lambda=\lambda(s)+\epsilon\nu$ と置くと，
\[
 [u,v]=\epsilon a\,J_{uv}(\lambda(s)),\qquad \epsilon^2=0.
\]
従って，係数や座標枠の変化は $h=0$ の一次交換子に寄与せず，
Poisson成分は $a\pi_s$ である．
この成分は一次量子化比較で保たれるので
（命題\ref{prop:abel-poisson-formal-comparison}），
tailsのAbel変形に移っても同じ結論となる．

一方，定理\ref{thm:abel-formal-ring}と
定理\ref{thm:all-artin-are-nc-dp}により，
特殊ファイバーを固定した全変形の接空間を通常の可換変形で割った空間は
\[
 H^0(X,\wedge^2T_X),\qquad
 \dim H^0(X,\wedge^2T_X)
 =\begin{cases}
   n&d=1,\\
   \binom{n+1}{2}&d=2
  \end{cases}
\]
である．いずれも $n\ge2$ では1より大きい．
$HH^1(X)=0$ により，この不足を微小自己同値で補うこともできない．
smoothな射なら，指定した原像を通る任意の双対数変形が持ち上がるため，
この接空間への写像は全射でなければならない．上の次元評価はそれに反する．
\end{proof}

この障害は，個々の有限族が非可換多様体を与えることとは両立する．
また，様々な $\pi_s$ を選んで半古典的なPoisson方向を稠密に被覆できることと，
固定した $h=0$ の原像で全一次変形を持ち上げることも異なる．
atlas構成には，一つの共通因子 $h$ を通じて全交換子を作る表示を越えて，
独立な一次Poisson方向を同時に含む有限型模型を構成する課題が残る．



% !TeX root = ../nc-del-pezzo.tex
% 2026-11-20 JST / v0.8
\section{独立Poisson方向を含む有限型ambient族}
\label{sec:simultaneous-poisson}

本節では一つの共通因子 $h$ を使わず，次数1の全線型方向と
次数2のPl\"ucker錐の全座標を関係式の係数として用いる。
全ての可換点を含む一つの有限型開集合上で，通常の次数付きambient代数の
有限表示と普遍的平坦性を証明する。
本節ではambientと第一平方零近傍の中心商を扱う。
有限パラメータの同時中心商とその全Artin普遍性は，
続く第\ref{sec:relative-center-lift}--\ref{sec:simultaneous-formal-comparison}節で証明する。
係数基底上の強Noether性，相対smoothnessやatlasは別の問題である。

\subsection{次数2における座標を分離しない関係式}

$E$ を基底 $x_0,\ldots,x_n$ を持つ $(n+1)$ 次元複素ベクトル空間とする。
$\mathcal F=\Sym^4E$ は四次多項式 $U$ の係数空間であり，
\[
 \mathcal C=\Spec\mathbb C[p_{ij}\mid 0\le i<j\le n]/(q_{ijk\ell})
       \subset\bigwedge^2E^*
\]
をrankが高々2の交代形式のPl\"ucker錐とする。
$p_{ji}=-p_{ij}$，$p_{ii}=0$ と定め，普遍交代形式を $P$ と書く。
$\mathcal F\times\mathcal C$ は被約な有限型アフィンスキームである。
以下 $U$ の係数と $p_{ij}$ の内部次数は0とする。

$\operatorname{Sym}:\Sym E\to T(E)$ は完全対称化である。
$S_0=\mathcal O(\mathcal F\times\mathcal C)$ 上の代数を
\begin{equation}\label{eq:simultaneous-ambient-relations}
 A_{U,P}=
 S_0\langle x_0,\ldots,x_n,z\rangle\Big/
 \left(
 [x_i,x_j]-p_{ij}z,\quad
 [x_i,z]+\sum_jp_{ij}\operatorname{Sym}(\partial_{x_j}U)
 \right)
\end{equation}
で定める。生成元の重みは $\deg x_i=1$，$\deg z=2$ である。
関係式はそれぞれ内部次数2と3の斉次式で，係数は有限個の多項式である。
この表示は任意の可換係数環への基底変更と両立する。

\begin{lemma}[座標変換と非零方向での表示]
\label{lem:simultaneous-coordinate-free}
表示\eqref{eq:simultaneous-ambient-relations}は $\operatorname{GL}(E)$ と両立する。
複素点 $(U,P)$ において $P\ne0$ なら，$K=\ker(P:E\to E^*)$ は
$A_{U,P}$ の中央な次数1の元からなる部分空間であり，$\dim K=n-1$ である。
$K$ の基底 $s_1,\ldots,s_{n-1}$ と補空間の基底 $u,v$ を取り，
$h=P(u,v)\ne0$ と置けば，代数は
\begin{equation}\label{eq:simultaneous-active-presentation}
 \mathbb C[s_1,\ldots,s_{n-1}]\langle u,v,z\rangle/
 \bigl([u,v]-hz,\ [v,z]-h\partial_u\operatorname{Sym}U,
                 [z,u]-h\partial_v\operatorname{Sym}U\bigr)
\end{equation}
となる。ここでは中央変数 $s_i$ を係数として完全対称化し，
$\partial_u,\partial_v$ は巡回微分である。
$P=0$ のファイバーは $\mathbb C[x_0,\ldots,x_n,z]$ である。
\end{lemma}
\begin{proof}
$a,b\in E$ に対する関係式は，基底を使わず
\[
 [a,b]=P(a,b)z,\qquad
 [a,z]=-\operatorname{Sym}(\partial_{P(a,-)}U)
\]
と書ける。右辺の $\partial_{P(a,-)}$ は線型汎関数
$P(a,-)\in E^*$ に対応する可換多項式環の導分である。
完全対称化は線型座標変換と可換するから，この表示は自然である。
$a\in K$ なら両右辺が零となり，$a$ は全生成元と可換する。
非零Pl\"ucker形式のrankは2なので，残りの主張は基底変換による。
中央変数を含む単項式の完全対称化は，中央変数を外へ出した
活性変数 $u,v$ の完全対称化に一致する。
また，対称化した巡回語の巡回微分は，可換微分の完全対称化に等しい。
\end{proof}

\subsection{重根を許す最高次AS開集合}

以下では四次式に重根があってもよい。
二生成元の巡回potential
\begin{equation}\label{eq:simultaneous-cubic-potential}
 W_{h,f}=u^2v^2-uvuv+h^2\operatorname{Sym}(f),
 \qquad f\in\mathbb C[u,v]_4
\end{equation}
を考える。巡回語を巡回平均によって $\langle u,v\rangle^{\otimes4}$ に
持ち上げたものを $w_{h,f}$ と書く。

\begin{lemma}[全可換最高次点を含む正則開集合]
\label{lem:all-binary-quartic-as-open}
$(f,h)$ の有限型係数空間には，$h=0$ の全点を含む
$\operatorname{GL}_2$ と両立するZariski開集合 $\mathcal U_{\mathrm{CY}}$ があり，
その全複素点で $J(W_{h,f})$ は3次元Noetherな三次Calabi--Yau代数である。
対応する三生成元最高次代数
\[
 \mathbb C\langle u,v,z\rangle/
 ([u,v]-hz,\ [v,z]-h\operatorname{Sym}(f_u),
                 [z,u]-h\operatorname{Sym}(f_v))
\]
は $h=0$ も含めて3次元AS正則であり，Hilbert級数は
$((1-t)^2(1-t^2))^{-1}$ である。
\end{lemma}
\begin{proof}
Mori--Ueyama\cite[Theorem~3.4]{MoriUeyama2019}によれば，
$J(w)$ が3次元Calabi--Yauとなるための必要十分条件は，
二つの一次端点微分が線型独立であり，二階端点微分を並べた
$2\times2$ 行列 $M(w)$ の全成分の共通零点が
$\mathbb P^1\times\mathbb P^1$ に存在しないことである。
端点微分は巡回平均したtensorの最初，または最後の文字を除く微分であり，
通常の巡回微分とは非零定数倍だけ異なる。
線型従属性は有限行列の小行列式で定義される閉条件である。
後者の失敗も，相対的な $\mathbb P^1\times\mathbb P^1$ 内の閉集合の
properな射影像として閉集合を定める。
その和の補集合を $\mathcal U_{\mathrm{CY}}$ とする。

$h=0$ ではpotentialは $W_0=u^2v^2-uvuv$ であり，
二つの巡回微分は
\[
 uv^2+v^2u-2vuv,\qquad vu^2+u^2v-2uvu
\]
である。これらは線型独立である。
$M(w_0)$ を非零定数倍して二重線型化すると
\[
 \begin{pmatrix}
 v_1v_2&u_1v_2-2v_1u_2\\
 v_1u_2-2u_1v_2&u_1u_2
 \end{pmatrix}.
\]
対角成分がともに零なら，二つの射影点が有効であることから
$u_1=v_2=0$ または $v_1=u_2=0$ となる。
各場合に非対角成分の一つが非零なので，共通零点は存在しない。
従って $f$ の重根の有無にかかわらず，全ての $h=0$ が開集合に属する。
Noether性と標準Hilbert級数は同論文の三次AS正則代数の結果による。

$h\ne0$ では $z=h^{-1}[u,v]$ を消去すると，
三生成元代数は文字通り $J(W_{h,f})$ と同型になる。
potential複体内の $z$ とその関係の可縮対も消去できるので，
三生成元potentialはCY-potentialである。
これは補題\ref{lem:elliptic-top-relative-cy}と同じ消去である。
$h=0$ では三生成元表示は可換多項式環となる。
このファイバーで $z$ の消去は行わない。
最後に，活性基底を $g\in\operatorname{GL}_2$ で変えると
$h$ は $\det(g)h$ となり，$W_{h,f}$ は座標変換後の
potentialの $\det(g)^2$ 倍となる。
上の二条件はこの変換で保たれる。
\end{proof}

四次式の判別式を可逆化せず，前補題の十分条件を
全生成空間 $E$ 上の一つの行列式で与えられる。
$U\in\Sym^4E$ の対称四重線型化を
$\widetilde U:(E^*)^4\to\mathbb C$ とし，
$\widetilde U(\xi,\xi,\xi,\xi)=U(\xi)$ と規格化する。
$P^\flat(a)=P(a,-)$ と置く。
対称な双線型式
\begin{equation}\label{eq:simultaneous-quartic-operator}
 T_{U,P}:\Sym^2E\longrightarrow\Sym^2E,\qquad
 T_{U,P}(ab)=\widetilde U(P^\flat a,P^\flat b,-,-)
\end{equation}
は自然な線型作用素を定める。
右辺は二つの残った双対変数に関する二次式として $\Sym^2E$ の元と見る。

\begin{lemma}[明示的な行列式によるAS十分条件]
\label{lem:simultaneous-explicit-as-determinant}
\begin{equation}\label{eq:simultaneous-degree-two-open}
 g_2(U,P)=\det(1+4T_{U,P})
\end{equation}
は $\operatorname{GL}(E)$-不変な係数多項式で，
$g_2(U,0)=1$ である。
非零rank2形式 $P$ に対し $K=\ker P$ と置くと，
$T_{U,P}$ は $K\cdot E\subset\Sym^2E$ を零に送り，
その商 $\Sym^2(E/K)$ 上の作用素を $\overline T$ と書けば
\[
 g_2(U,P)=\det(1+4\overline T).
\]
$g_2(U,P)\ne0$ なら，商 $E/K$ の四次式を活性最高次項に持つ
三生成元potentialはCYである。
従って $g_2\ne0$ は，最高次四次式の重根も許す有限AS開条件である。
\end{lemma}
\begin{proof}
作用素はtensorの縮約だけで定義されるので，座標変換により共役される。
係数の多項式性と $g_2(U,0)=1$ も明らかである。
$P^\flat K=0$ だから $T_{U,P}(K\cdot E)=0$ であり，
$\Sym^2E$ のこの部分空間に適合する基底で $1+4T$ は
対角ブロックが $1$ と $1+4\overline T$ の三角行列になる。
商作用素は，$f=U\bmod K\in\Sym^4(E/K)$ と
$P$ が誘導する交代形式によって\eqref{eq:simultaneous-quartic-operator}を
構成したものに等しい。

従って $E=\langle u,v\rangle$ の場合を考えればよい。
$P(u,v)=h$ とし，$w$ を\eqref{eq:simultaneous-cubic-potential}の巡回平均とする。
まず二つの端点微分は，全ての $f,h$ で線型独立である。
仮に非零な線型結合が零なら，基底変換後に最初の文字 $u$ による
端点微分を零とできる。
このとき $w$ の全ての語は最初の文字が $v$ であり，
巡回不変性から全ての文字が $v$ でなければならない。
従って $w$ は $v^4$ の定数倍となり，完全対称tensorになる。
一方，$w_0$ の $uuvv$ と $uvuv$ の係数はそれぞれ $1/4$ と $-1/2$ なので，
$w_0$ は完全対称ではない。
完全対称tensorの部分空間は基底変換で保たれ，
$h^2\operatorname{Sym}(f)$ を加えてもその商での $w_0$ の非零類は変わらない。
これは矛盾である。

$M(w)$ の二重線型化を評価する二つの射影点を
$p=(u_1,v_1)$，$q=(u_2,v_2)$ とする。
完全対称部分の行列は対称なので，$4M(w)$ の非対角成分の差は
\[
 3(u_1v_2-v_1u_2)
\]
である。従って全成分の共通零点が存在すれば $p=q$ であり，
代表を $p=(r,s)$ に揃えられる。
この対角上では
\[
 4M(w)(p,p)=
 \begin{pmatrix}s^2&-rs\\-rs&r^2\end{pmatrix}
       +\frac{h^2}{3}\operatorname{Hess}_f(p).
\]
ここでHessianは通常の可換二階微分である。
$a=su-rv\in E$ と置くと $P^\flat a=hp$ であり，
\[
 T_{f,P}(a^2)=h^2\widetilde f(p,p,-,-)
             =\frac{h^2}{12}\operatorname{Hess}_f(p).
\]
従って上の対称行列は，$\Sym^2E$ の元として
$(1+4T_{f,P})(a^2)$ に一致する。
$p\ne0$ なので $a^2\ne0$ であり，$g_2\ne0$ の下では零にならない。
Mori--Ueyamaの二条件を満たすため，前補題の結論を得る。
\end{proof}

一般のbinary quarticを
\[
 f=a u^4+4b u^3v+6c u^2v^2+4d uv^3+e v^4
\]
と書くと，基底 $(u^2,uv,v^2)$ 上で
\[
 T=h^2\begin{pmatrix}
 c&-b&a\\2d&-2c&2b\\e&-d&c
 \end{pmatrix}.
\]
従って，標準的なbinary quartic不変量の規格化を
\[
 I=ae-4bd+3c^2,\qquad
 J=ace+2bcd-ad^2-b^2e-c^3
\]
とすると，直接の行列式計算から
\begin{equation}\label{eq:simultaneous-binary-invariant-open}
 g_2=1-16Ih^4+128Jh^6
\end{equation}
を得る。
例えば $f=\alpha(u^4+v^4)/4+\beta u^2v^2$ の場合，
基底 $(u^2,uv,v^2)$ に関する $1+4T$ は
\[
 \begin{pmatrix}
 1+\frac23\beta h^2&0&\alpha h^2\\
 0&1-\frac43\beta h^2&0\\
 \alpha h^2&0&1+\frac23\beta h^2
 \end{pmatrix}
\]
であり，十分条件は具体的に
\begin{equation}\label{eq:simultaneous-even-quartic-determinant}
 g_2=\left(1-\frac43\beta h^2\right)
       \left(\left(1+\frac23\beta h^2\right)^2-\alpha^2h^4\right)\ne0
\end{equation}
となる。これは十分条件であり，非零となることをAS正則性の必要条件とは主張しない。

\begin{theorem}[次数2の全Pl\"ucker方向の有限型平坦ambient族]
\label{thm:simultaneous-poisson-ambient}
\eqref{eq:simultaneous-degree-two-open}の明示的な多項式 $g_2$ に対し，
$S=S_0[g_2^{-1}]$ 上の\eqref{eq:simultaneous-ambient-relations}は，
各内部次数で有限局所自由な有限表示代数となる。
任意の可換 $S$-代数への基底変更に対して同じ表示と平坦性が成立し，
次数別の階数は
\[
 H_A(t)=\frac1{(1-t)^{n+1}(1-t^2)}
\]
の係数である。
特に開集合 $\Spec S$ は，四次式の係数空間の全零切断
$\mathcal F\times\{0\}$ を含む。
ここで四次式に対する滑らかさや一般性は仮定しない。
\end{theorem}
\begin{proof}
\emph{全複素ファイバーのHilbert級数。}
$P=0$ は補題\ref{lem:simultaneous-coordinate-free}から従う。
$P\ne0$ では同補題の活性表示を使う。
商 $E/\ker P$ の四次式が活性最高次項となり，
補題\ref{lem:simultaneous-explicit-as-determinant}によってそのpotentialはCYである。
全てのlower potentialについてEtingof--Ginzburgの
Proposition~7.3.1(ii)\cite{EG2010}を適用できる。
中央変数を係数にする段階を明記する。
最高次パラメータを固定し，まず各 $u^iv^j$，$i+j\le3$ の係数を
独立な中央変数 $c_{ij}$ とする。
最高次代数の各次数の基底を同じ次数の語の線型結合で持ち上げると，
最高次数を下げる帰納法により，これらは各有限filtration部分を生成する。
これらの間の多項式係数の関係は，$c_{ij}$ の全複素特殊化で
固定topのPBW定理により零となる。
被約な多項式係数環上で係数も零となるため，各filtration部分はこの基底を持つ。
従って任意の中央係数への基底変更にもPBW同型が保たれる。
これは補題\ref{lem:elliptic-universal-lower-pbw}の証明で用いた議論であり，
最高次式を偶四次式に限定する必要はない。

ここで $c_{ij}$ に $s_1,\ldots,s_{n-1}$ の次数 $4-i-j$ の多項式を代入する。
活性次数 $d$ の基底語は全内部次数でも $d$ であるから，
結果は次数付き $\mathbb C[s_1,\ldots,s_{n-1}]$-加群として，
最高次代数の次数別基底を持つ自由加群となる。
従ってHilbert級数は
\[
 \frac1{(1-t)^{n-1}}\cdot\frac1{(1-t)^2(1-t^2)}
\]
であり，$P=0$ と一致する。

\emph{普遍的平坦性。}
各 $A_m$ は有限個の重み付き語と，それらの次数以内の関係式による
有限表示 $S$-加群である。
上で全閉点のファイバー次元が同じであることを証明した。
$S$ は被約な有限型 $\mathbb C$-代数なので，その関係行列の階数を超える
小行列式は恒等的に零である。
所要階数の非零小行列式の開集合上で行列を消去すると，残る行列は零であり，
$A_m$ は有限自由となる。
これらの開集合は全閉点，従って $\Spec S$ を被覆する。
各 $A_m$ は有限局所自由である。
その基底変更は普遍関係式の基底変更に一致するので，
非被約環を含む任意の $S$-代数上で同じ階数と平坦性が保たれる。
この最後の主張では，非零rank2形式の基底を非被約環上で選ぶ必要はない。
\end{proof}

\subsection{次数1の全線型方向と明示的な正則開集合}

次数1については，楕円判別式を除いたAS開条件そのものが
方向の座標を使わずに記述できる。
まず第\ref{sec:elliptic-degree-one-finite}節で用いた
Weierstrass最高次項について，AS正則性の仮定を広げる。

\begin{lemma}[特異なWeierstrass最高次項も含むAS開集合]
\label{lem:degree-one-all-top-as-open}
\[
 T_1=\mathbb C[\alpha,\beta,h,D_1^{-1}],\qquad
 D_1=(1+\alpha h^4)(5+\alpha h^4)
\]
と置く。potential
\[
 \Phi=xyz-yxz-h\left(\frac12z^2+\frac13y^3+
                       \alpha x^4y+\beta x^6\right)
\]
の代数は，各内部次数で有限局所自由な $T_1$-代数であり，
相対的な3-Calabi--Yau代数である。
Hilbert級数は $((1-t)(1-t^2)(1-t^3))^{-1}$ で，
これらの主張は任意の係数変更に両立する。
全複素ファイバーはAS正則Noether整域である。
ここでは $\Delta=4\alpha^3+9\beta^2\ne0$ を仮定しない。
\end{lemma}
\begin{proof}
まず複素点で示す。$h=0$ では重み付き多項式環である。
$h\ne0$ かつ $(\alpha,\beta)\ne(0,0)$ なら，
$k^2=-\alpha$ の根の少なくとも一つに対し
$6\beta-4k^3\ne0$ である。
実際 $\alpha\ne0$ の場合に両方の符号で零なら
$\beta=k=0$ となり矛盾し，$\alpha=0$ の場合は $\beta\ne0$ である。
その根を取り，$\rho^3=6\beta-4k^3$ の非零根を選ぶ。
補題\ref{lem:degree-one-stephenson-normal-form}の計算は
$\rho\ne0$ だけを使うので，$\Delta=0$ の点でも同じ変数変換を行える。
得られるStephensonのパラメータは
\[
 a=1+kh^2,\quad b=\rho h^2,\quad c=-2,
 \qquad a(a+c)(c^2-ac-a^2)=-D_1\ne0.
\]
従ってStephenson\cite[Theorem~1.1]{Stephenson1997}のAS正則開集合に入り，
補題\ref{lem:degree-one-elliptic-top-relative-cy}のpotential複体も完全となる。
これは複素点ごとの根の選択であり，$\Delta=0$ の上で
補題\ref{lem:degree-one-stephenson-normal-form}の有限被覆全体を
可逆な座標変換に使っているわけではない。

残る $(\alpha,\beta)=(0,0)$ では
\[
 [x,y]=hz,\quad [y,z]=0,\quad [z,x]=hy^2.
\]
これは $\mathbb C[y,z][x;\delta]$，
$\delta(y)=hz$，$\delta(z)=-hy^2$ というOre拡大であり，
重み付きPBW基底を持つAS正則Noether整域である。
potentialのCY性は第\ref{sec:finite-auxiliary-quantization}節の分離型
$P(x)=0$，$Q(y)=y^3/3$ の場合にも含まれる。
よって全複素ファイバーで標準Hilbert級数とpotential複体の完全性を得た。

$T_1$ は被約有限型なので，一定階数の有限関係行列から
各内部次数の有限局所自由性が従う。
potential対角複体の各内部次数は有限射影加群の有界複体であり，
全閉点で完全だから局所的に分裂完全である。
Hessianの自己双対性と合わせ，
補題\ref{lem:degree-one-elliptic-top-relative-cy}と同じ証明により
相対CY性と任意基底変更との両立を得る。
\end{proof}

ここからは $E$ を基底 $x_1,\ldots,x_n$ を持つ $n$ 次元空間とし，
$A_4\in\Sym^4E$，$A_6\in\Sym^6E$，$D\in E^*$ を独立なパラメータとする。
$D_i=D(x_i)$ と置き，同じ記号 $D$ で多項式環の定数係数導分
$\sum_iD_i\partial_{x_i}$ を表す。
$A_4(D)$ は四次tensor $A_4$ と $D^{\otimes4}$ の縮約である。

\begin{theorem}[次数1の全線型方向の有限型平坦ambient族]
\label{thm:simultaneous-degree-one-ambient}
全係数空間
\[
 \mathcal F_1\times E^*
   =(\Sym^4E\times\Sym^6E)\times E^*
\]
上で，次の有限個の斉次関係式を考える。
\begin{equation}\label{eq:simultaneous-degree-one-relations}
 \begin{gathered}
 [x_i,x_j]=0,\qquad [x_i,y]=D_i z,\\
 [z,x_i]=D_i(y^2+A_4(x)),\\
 [y,z]=\operatorname{Sym}\bigl((D A_4(x))y\bigr)+D A_6(x),
 \qquad \deg(x_i,y,z)=(1,2,3).
 \end{gathered}
\end{equation}
ここで完全対称化は $x_i,y$ について行う。
明示的な主開集合
\begin{equation}\label{eq:simultaneous-degree-one-open}
 (1+A_4(D))(5+A_4(D))\ne0
\end{equation}
上では，この代数 $A_1$ は次数ごとに有限局所自由であり，
任意基底変更後も
\[
 H_{A_1}(t)=\frac1{(1-t)^n(1-t^2)(1-t^3)}
\]
を持つ。この開集合は $D=0$ の全係数点を含む。
関係式と開条件は $\operatorname{GL}(E)$ と両立する。
\end{theorem}
\begin{proof}
自然な導分とtensorの縮約を用いているので，座標変換との両立は直接従う。
$D=0$ のファイバーは可換多項式環である。
$D\ne0$ の複素点では，$\ker D$ の基底 $s_1,\ldots,s_{n-1}$ と
補空間の基底 $u$ を取り，$h=D(u)\ne0$ と置く。
表示関係式から全ての $s_i$ が中央となる。
残る関係式は，中央 $s_i$ を係数とする三生成元potential
\[
 uyz-yuz-h\left(\frac12z^2+
       \frac13y^3+\operatorname{Sym}(A_4(s,u)y)+A_6(s,u)\right)
\]
の関係式である。
完全対称化と巡回微分の両立は
補題\ref{lem:simultaneous-coordinate-free}と同様である。

活性重み $(1,2,3)$ に関する最高次式は
$y^3/3+\alpha u^4y+\beta u^6$ となる。
$E\to E/\ker D$ において $x_i$ の像は $(D_i/h)u$ だから
\[
 \alpha=h^{-4}A_4(D),\qquad \beta=h^{-6}A_6(D).
\]
従って補題\ref{lem:degree-one-all-top-as-open}の十分条件は
文字通り\eqref{eq:simultaneous-degree-one-open}である。
固定CY topへの全lower potentialのPBW定理と，その中央係数への拡張を用いると，
各複素ファイバーのHilbert級数は，最高次代数のHilbert級数を
$(1-t)^{n-1}$ で割ったものになる。
この中央係数への拡張は定理\ref{thm:simultaneous-poisson-ambient}の証明と同じであり，
補題\ref{lem:degree-one-elliptic-universal-lower-pbw}の証明も
最高次potentialがCYであることだけを使っている。

最後に，全係数基底は被約有限型であり，各内部次数は有限表示加群である。
全閉点でHilbert次元が一定なので，関係行列の一定階数から
有限局所自由性と任意の非被約基底変更への両立が従う。
\end{proof}

\subsection{両次数の一次中心商と方向の退化}

\begin{proposition}[全一次Poisson方向の中心商]
\label{prop:simultaneous-first-order-center}
定理\ref{thm:simultaneous-poisson-ambient}または
定理\ref{thm:simultaneous-degree-one-ambient}の基底環を $S$ と書き，
零切断のイデアルをそれぞれ $J=(p_{ij})S$ または $J=(D_i)S$ とする。
$S^{(1)}=S/J^2$，$A^{(1)}=A\otimes_S S^{(1)}$ と置く。
次の元は $A^{(1)}$ の中心元である。
\[
 \Omega^{(1)}=
 \begin{cases}
 \frac12z^2+\frac13y^3+\operatorname{Sym}(A_4(x)y)+A_6(x),&d=1,\\
 \frac12z^2+\operatorname{Sym}(U),&d=2.
 \end{cases}
\]
$B^{(1)}=A^{(1)}/(\Omega^{(1)})$ は次数ごとに有限局所自由である。
この商は任意基底変更後も正則中心商となる。
次数1の零切断における一次commutatorは，任意の $D\in E^*$ に対して
\[
 \{x_i,y\}=D_i z,\quad
 \{z,x_i\}=D_i(y^2+A_4),\quad
 \{y,z\}=(D A_4)y+D A_6,
 \qquad \{x_i,x_j\}=0
\]
であり，次数1の全一次Poisson方向を与える。
次数2では，任意の
$Q\in\bigwedge^2E^*$ に対して
\[
 \{x_i,x_j\}=Q_{ij}z,\qquad
 \{x_i,z\}=-\sum_jQ_{ij}\partial_{x_j}U
\]
であり，次数2の全一次Poisson方向を与える。
\end{proposition}
\begin{proof}
$A^{(1)}$ は平坦であり，$J/J^2$ を係数とする一次commutatorは
可換多項式環上の二重導分である。
表示関係式から得られる上の二重導分について
$F=(\Omega^{(1)})^{\mathrm{ab}}$ はCasimirである。
次数1では $U=y^3/3+A_4(x)y+A_6(x)$ と置けば
\[
 \{x_i,y\}=D_iF_z,\quad
 \{z,x_i\}=D_iF_y,\quad
 \{y,z\}=D(F)
\]
というJacobian型の式なので，$\{x_i,F\}=\{y,F\}=\{z,F\}=0$ を直接得る。
次数2でも $\{x_i,F\}$ の二つの項は相殺し，
$\{z,F\}=\sum_{i,j}Q_{ij}U_{x_i}U_{x_j}=0$ である。
従って $[\Omega^{(1)},a]=0$ が全生成元 $a$ について成立する。

零切断上の $F$ は $z$ に関してmonicな二次多項式の非零定数倍である。
その掛算の各内部次数は分裂単射で，余核は $z$ の指数が0または1の
単項式を基底に持つ。
この分裂を平方零厚化に持ち上げると，
有限射影加群の掛算行列は同じ分裂単射と可逆な平方零摂動に分解される。
従って $\Omega^{(1)}$ は普遍的に正則であり，余核も有限局所自由である。
最後に次数1の方向空間は $E^*$ そのものであり，
次数2のPl\"ucker関係は二次なので，頂点の接空間は
$\bigwedge^2E^*$ 全体である。
次数1の式は定数の重み付き座標変更と方程式の非零スカラー倍の下で，
第\ref{sec:equivariant-quantization}節の規格化と一致する。
$F=z^2-f_4$ という既存の規格化では $U=-f_4/2$，$Q=2P$ とすれば，
第\ref{sec:equivariant-quantization}節の初括弧と一致する。
\end{proof}

\begin{remark}[楕円型方向だけでは不十分であること]
滑らかな次数2模型でも，分岐四次超曲面の滑らかな点における接線上への
四次式の制限は重根を持つ。
従って楕円型最高次項だけの開集合を方向空間全体に要求することはできない。
次数1でも活性方向のWeierstrass判別式は
$\mathbb P^{n-1}$ 上の正次数斉次式であり，$n\ge2$ なら零点を持つ
（恒等的に零の場合は全方向が該当する）。
本節のambient構成で必要なのは，楕円曲線の滑らかさではなく
補題\ref{lem:all-binary-quartic-as-open}と
補題\ref{lem:degree-one-all-top-as-open}のCY条件である。
\end{remark}

\begin{remark}[ambientから中心商への論理]
本節は両次数の全独立Poisson方向を含む有限型ambient族を与える。
一方向ごとの中心元の存在だけから，方向空間の零切断を含む同時持上げは従わない。
この貼合せと有限代数化を後続の中心構成で証明する。
また，有限生成元と有限関係式を持つ任意の表示族で，
全次数の平坦性が有限開条件になるという一般原理も用いていない。
本節では明示的な有限行列式と具体的CY判定から全複素ファイバーの
Hilbert級数を先に固定したことが，普遍的平坦性の根拠である。
後続の全Artin普遍性を合わせても，有限基底上の相対tails，smoothnessと
大域的な全同値の商を確認するまでは，モジュライのatlasとは呼ばない。
\end{remark}

% !TeX root = ../nc-del-pezzo.tex
\section{中央多項式係数上の同じ ambient 代数における中心持上げ}
\label{sec:relative-center-lift}

ここでは，既に得た平坦な potential 代数の中心を持ち上げる．
別に選んだ星積の中心を，代数の同定なしに利用することはしない．
Etingof--Ginzburg の Proposition~8.1.1(i) は，形式 potential の中心持上げに
孤立特異点条件を課さず，その証明も Rees 変数を中央係数とする相対形式を用いる
\cite[\S8.1--8.2]{EG2010}．
以下では中央変数を任意個許し，同じ代数への同定，内部次数，
および有限パラメータへの移行を明記する．
使用する unimodularity 定理は複素数上の滑らかな affine 多様体についての
Dolgushev の定理であり，その未確認の相対版を仮定しない．

\subsection{形式中心と内部次数}

内部次数零の係数環と，正次数の中央変数を
\[
 R_0=\mathbb C[c_1,\ldots,c_N,\rho(c)^{-1}],\qquad
 R=R_0[s_1,\ldots,s_m],\qquad \deg s_j=v_j>0
\]
とする．$\rho=1$ を許し，$N,m$ は零でもよい．
$h$ の内部次数は零とし，活性変数の重みを
\[
 (w_x,w_y,w_z)=(1,1,2)\quad\text{または}\quad(1,2,3),
 \qquad d=w_x+w_y+w_z=2w_z
\]
とする．$Q=R[x,y,z]$ と置き，
\begin{equation}\label{eq:relative-center-classical-potential}
 F=\frac12z^2+U(s,x,y)\in Q_d,\qquad
 \Phi=xyz-yxz-h\operatorname{Sym}(F)
\end{equation}
を考える．$\operatorname{Sym}$ は活性変数の完全対称化であり，
$R$ に関して線型に延長する．
potential は巡回語として読み，その関係式は
\begin{equation}\label{eq:relative-center-potential-relations}
 [x,y]=hz,\qquad
 [y,z]=h\partial_x\operatorname{Sym}(U),\qquad
 [z,x]=h\partial_y\operatorname{Sym}(U)
\end{equation}
である．右辺の微分は巡回微分であり，可換化すると通常の偏微分となる．

\begin{lemma}[同じ相対 potential 代数の規格化形式中心]
\label{lem:relative-formal-center-lift}
$\widehat A$ を\eqref{eq:relative-center-potential-relations}の
$h$ 進完備な代数とする．次を仮定する．
\begin{enumerate}
 \item $\widehat A$ は $h$ 進分離完備かつ $h$-torsion-free であり，
       与えられた生成元による還元は $\widehat A/h\widehat A=Q$ である．
 \item $R[[h]]$ に相対的な対角は perfect であり，連続な完備包絡代数
       $\widehat A^{\mathrm e}_{R[[h]]}$ に関して
       \[
        \underline{R\operatorname{Hom}}_{\widehat A^{\mathrm e}_{R[[h]]}}
          (\widehat A,\widehat A^{\mathrm e}_{R[[h]]})
          \simeq\widehat A(d)[-3]
       \]
       という無捩れな相対 Calabi--Yau 同型を持つ．
\end{enumerate}
このとき，同じ $\widehat A$ の中に内部次数 $d$ の中心元
\[
 \widehat\Omega\in Z(\widehat A)_d,\qquad
 \widehat\Omega\bmod h=F
\]
が存在する．完全対称化による形式 PBW 座標で，その $z^2$ 係数を
ちょうど $1/2$ とすることができる．
各 $h$ 係数は $s_1,\ldots,s_m$ の多項式であり，中央変数の有理関数への
係数拡大は行わない．$U$ の最高次部分の判別式や $F$ の滑らかさは仮定しない．
\end{lemma}

\begin{proof}
\emph{形式 PBW 座標と相対 formality．}
完全対称化による連続な $R[[h]]$-線型写像
\[
 \sigma:Q[[h]]\longrightarrow\widehat A
\]
は還元後に恒等写像となる．全射性は $h$ 係数を順次持ち上げることで従い，
単射性は $h$-torsion-free 性と $h$ 進分離性から従う．
従って $\sigma$ は内部次数を保つ同型である．
ここでは平坦性を中心の存在から導いておらず，仮定1を先に使用している．
この座標で乗法を $Q[[h]]$ へ移すと，$R$ に関して線型で
内部次数零の形式的結合変形 $\mu_h$ を得る．

活性三変数の Kontsevich 写像を $R$-線型に延長して用いる
\cite[\S4, \S6]{Kontsevich}．
そのグラフは $x,y,z$ のみを微分し，係数 $c,s$ を微分しない．
相対多重微分作用素から全相対 Hochschild cochain への包含は
多項式 HKR の擬同型である．さらにこの擬同型と formality は
内部次数零の部分に制限しても擬同型である．
実際，相対 bar 分解と多項式 Koszul 分解の比較およびホモトピーを
次数を保つものに取ると，その比較は各内部次数で成立する．
多重微分作用素上では，線型な重み付き $\mathbb G_m$ 作用に関する
Kontsevich の同変性により全ての構造写像が内部次数を保存する．
従って $h$ 進 Maurer--Cartan 対応を内部次数零の複体に適用できる．
単位を保つ正規化 cochain を使えば，
\begin{equation}\label{eq:relative-center-formality-identification}
 \pi_h\in h\bigwedge^2_Q\operatorname{Der}_R(Q)[[h]],\qquad
 [\pi_h,\pi_h]=0,\qquad\deg\pi_h=0
\end{equation}
と，内部次数を保ち $R$ を固定する gauge 同型
\[
 T=1+O(h):(Q[[h]],\star_{\pi_h})
       \xrightarrow{\ \sim\ }(Q[[h]],\mu_h)
       \xrightarrow[\sigma]{\ \sim\ }\widehat A
\]
を得る．以下の $T$ はこの合成を表す．
一階交換子の規約を $[a,b]_{\star}=h\{a,b\}+O(h^2)$ と取ると，
\eqref{eq:relative-center-potential-relations}により
\[
 \pi_h=h\pi_F+O(h^2),\qquad
 \pi_F=F_x\partial_y\wedge\partial_z
        +F_y\partial_z\wedge\partial_x
        +F_z\partial_x\wedge\partial_y .
\]
以上は形式的な同型類の対応だけでなく，中心元を元の代数へ戻す
具体的な存在同型 $T$ を与える．

\emph{相対 CY から複素数上の絶対 CY へ．}
$\ell=N+m$ と置き，完備包絡代数を
\[
 B=\widehat A\widehat\otimes_{\mathbb C[[h]]}\widehat A^{\mathrm{op}},
 \qquad
 B_R=\widehat A\widehat\otimes_{R[[h]]}\widehat A^{\mathrm{op}}
\]
とする．$B_R$ は $B$ を中央差
\[
 c_i\otimes1-1\otimes c_i\quad(1\le i\le N),\qquad
 s_j\otimes1-1\otimes s_j\quad(1\le j\le m)
\]
で割ったものである．これらは中央正則列をなす．
実際，$h$ による還元後には局所化した多項式環の独立な中央差であり，
途中の各商も $h$-torsion-free となることを，形式 PBW 座標と
$h$ 係数による帰納法で確認できる．
この正則列の Koszul 分解により $B_R$ は $B$ 上 perfect であり，
内部次数のずれを省略すれば
\[
 R\operatorname{Hom}_B(B_R,B)\simeq B_R[-\ell]
\]
である．Koszul 最上外積は列挙した中央差の順序で自明化する．
導来 change of rings と仮定2から
\[
 \begin{split}
 R\operatorname{Hom}_B(\widehat A,B)
 &\simeq R\operatorname{Hom}_{B_R}
       (\widehat A,R\operatorname{Hom}_B(B_R,B))\\
 &\simeq\widehat A[-\ell-3]
 \end{split}
\]
を得る．対角の perfect 性もこの合成から従う．
内部次数を戻したずれは $d+\sum_jv_j$ であり，
左右の作用に自己同型の捻りはない．

\emph{絶対 unimodularity と体積形式の重み零成分．}
$X=\Spec Q$ は滑らかな複素 affine 多様体であり，
標準体積形式
\[
 \omega=dc_1\wedge\cdots\wedge dc_N
       \wedge ds_1\wedge\cdots\wedge ds_m
       \wedge dx\wedge dy\wedge dz
\]
はその正準束を自明化する．
\eqref{eq:relative-center-formality-identification}の $\pi_h$ を
$X$ 上の絶対二重ベクトルとも読む．
全変数についての Kontsevich グラフ公式は，この $\pi_h$ 上では
相対グラフ公式と一致する．各辺の微分添字は二重ベクトルの成分から来るため，
$\partial_{c_i}$ または $\partial_{s_j}$ を含む項が一つも生じないからである．
従って $\pi_h$ は，先ほどの絶対 CY 代数と同型な星積の
絶対 Kontsevich 類でもある．

Dolgushev の Theorem~3 は，正準束が自明な滑らかな複素 affine 多様体の
一変数形式量子化について，無捩れな Van den Bergh 双対化 bimodule と
その形式 Poisson 類の unimodularity が同値であることを述べる
\cite[Theorem~3]{DolgushevModular2009}．
同論文の Definition~1 と Proposition~1 により，この条件は
可逆元 $f\in Q[[h]]^\times$ が存在して
\begin{equation}\label{eq:relative-center-invariant-volume}
 L_{\pi_h}(f\omega)=0
\end{equation}
となることと同値である．今の場合は直前で全ての仮定を満たしている．

$Q$ の内部次数は非負であり，$Q_0=R_0$ である．
$f$ の各 $h$ 係数を内部次数零へ射影したものを
$f_0\in R_0[[h]]$ とする．$f\bmod h$ が $Q$ の単元であることから
$f_0\bmod h\in R_0^\times$ であり，$f_0$ は可逆である．
$\pi_h$ は内部次数零なので，\eqref{eq:relative-center-invariant-volume}の
内部次数 $\deg\omega=d+\sum_jv_j$ 成分は
\[
 0=L_{\pi_h}(f_0\omega)=f_0L_{\pi_h}\omega
\]
となる．最後の等号では $\pi_h$ が $R$ に相対的であることを使った．
従って $L_{\pi_h}\omega=0$ であり，活性成分以外を持たないことから
\[
 L_{\pi_h}(dx\wedge dy\wedge dz)=0
 \quad\text{in }\Omega^\bullet_{Q/R}[[h]]
\]
を得る．この重み射影によって，未知の体積形式を
標準の相対体積形式へ取り替えられる．

\emph{相対 Poincar\'e ホモトピーと Casimir．}
相対一形式
\[
 \alpha_h=\iota_{\pi_h}(dx\wedge dy\wedge dz)
     \in h\Omega^1_{Q/R}[[h]]
\]
は閉じており，内部次数 $d$ を持つ．
活性三変数の多項式 de Rham 複体において，通常の活性次数
（$x,y,z,dx,dy,dz$ をそれぞれ次数1とする）を使う．
その正次数 $r$ の成分上で，Euler 場
$E=x\partial_x+y\partial_y+z\partial_z$ による
$r^{-1}\iota_E$ は Poincar\'e ホモトピーを与える．
一形式の各単項式では $r>0$ なので，必要な割算は正整数だけである．
このホモトピーは中央係数を固定し，重み付き内部次数も保存する．
従って
\[
 \psi_h\in hQ_d[[h]],\qquad d_{Q/R}\psi_h=\alpha_h
\]
を取れる．$\alpha_h/h\bmod h=d_{Q/R}F$ なので，
$\psi_h$ に $hR_d$ の元を加えることにより
$\psi_h/h\bmod h=F$ とできる．
活性次元が3であるため，$\pi_h$ と $d_{Q/R}\psi_h$ の双対表示から
$[\pi_h,\psi_h]=0$ となり，$h$-torsion-free 性により
$\psi_h/h$ も $\pi_h$ の Casimir である．

\emph{中心へ戻す操作と規格化．}
$\pi_h$ で twist した Kontsevich の cochain 写像は，
Poisson cohomology の次数零を星積の Hochschild cohomology の次数零，
すなわちその中心へ送る．その $h$ による還元は HKR の恒等写像である．
従って Casimir $\psi_h/h$ の像を同型 $T$ で戻すと，
同じ $\widehat A$ の内部次数 $d$ の中心元で，還元が $F$ となるものを得る．
完全対称化座標での $z^2$ 係数を $a(h)$ とすれば，
$2w_z=d$ と全ての $s_j$ の正次数性により
\[
 a(h)\in R_0[[h]],\qquad a(h)\equiv\frac12\pmod h .
\]
$2a(h)$ は中央単元なので，得た元をこれで割れば所要の
$\widehat\Omega$ となる．
全ての段階で係数環は $R_0[s]$ のままであり，$s$ の分母は導入していない．
\end{proof}

\begin{remark}[平坦性・CY 性の入力と判別式]
前補題は，任意の potential 代数の平坦性を中心持上げから証明するものではない．
本稿の適用では，定理\ref{thm:simultaneous-degree-one-ambient}と
定理\ref{thm:simultaneous-poisson-ambient}の平坦な同時 ambient 族を，
Poisson 方向を自明化した chart に引き戻す．
その表示は\eqref{eq:relative-center-potential-relations}であり，
次数別平坦性から仮定1を満たす．
仮定2は既存の相対 potential 対角複体から得られる．
形式段階だけなら，先に確立した平坦性の下でこの複体と各項を $h$ 進完備化し，
還元後の三変数 Koszul 分解から cycle を順次持ち上げても完全性が従う．
巡回 Hessian の自己双対性は係数環上の恒等式であるため，無捩れな CY 同型を与える．
これは補題\ref{lem:elliptic-top-relative-cy}，
\ref{lem:degree-one-elliptic-top-relative-cy}および
普遍低次項の補題\ref{lem:elliptic-universal-lower-pbw}，
\ref{lem:degree-one-elliptic-universal-lower-pbw}と同じ機構であり，
楕円判別式を使わない形式段階の論法である．
とくに活性最高次式に重根がある場合にも，上の入力が満たされる限り
中心持上げを適用できる．
\end{remark}

\subsection{有限中心方程式による代数化}

\begin{proposition}[規格化中心の有限局所化による代数化]
\label{prop:relative-finite-center-lift}
$g\in R_0[h]$ は $g(c,0)\in R_0^\times$ を満たすとし，
$S=R_0[h,g^{-1}]$ とする．
中央変数 $s_1,\ldots,s_m$ と活性変数 $x,y,z$ による
同じ関係式\eqref{eq:relative-center-potential-relations}の
次数付き $S$-代数を $A$ とする．
各 $A_q$ は有限射影 $S$-加群であり，$h$ 進完備化は
補題\ref{lem:relative-formal-center-lift}の仮定を満たすと仮定する．
このとき，$f\in S$，$f\equiv1\pmod h$ と
\[
 \Omega\in Z(A_f)_d,\qquad
 \Omega\bmod h=\operatorname{Sym}(F)
\]
が存在する．有限個の必要次数で完全対称化を座標に取ることができ，
その座標における $\Omega$ の $z^2$ 係数は $1/2$ となる．
中心性，次数，還元および規格化の恒等式は，任意の可換 $S_f$-代数への
基底変更後にも成立する．基底変更の平坦性や被約性は必要ない．
追加の局所化は係数 $c,h$ だけに依存し，中央変数 $s$ は全ての値を許す．
\end{proposition}

\begin{proof}
$s$ も正次数の代数生成元として扱い，基底環 $S$ に含めないことが重要である．
各内部次数の可換多項式空間は有限自由 $S$-加群であり，
完全対称化はそれから $A_q$ への写像を与える．
必要な有限個の次数
\[
 q=d,\quad d+w_x,\quad d+w_y,\quad d+w_z
\]
でこの写像は $h=0$ 上の同型である．
有限射影加群間の写像が同型でない locus は閉集合であり，$V(h)$ と交わらない．
従って $1\bmod h$ となる一つの元を可逆化すれば，これらを同時に同型にできる．
実際，その閉集合の定義イデアル $I$ に対して $I+(h)=S$ なので，
$I$ 内に $1\bmod h$ となる元がある．
この局所化を済ませ，なお係数環を $S$ と書く．
$z^2$ 係数を読む $S$-線型写像を
$\lambda:A_d\to S$ とする．

$A/hA$ は可換であり，各次数の射影性により $A$ は $h$-torsion-free だから，
\[
 \begin{gathered}
 M=A_d,\qquad N'=A_{d+w_x}\oplus A_{d+w_y}\oplus A_{d+w_z},\\
 D:M\longrightarrow N',\qquad
 D(a)=h^{-1}([a,x],[a,y],[a,z])
 \end{gathered}
\]
は正則な $S$-線型写像である．ここで $h^{-1}$ は $hA$ 内での一意な割算を表し，
係数環では $h$ を可逆化していない．
$F_{\mathrm{sym}}=\operatorname{Sym}(F)\in A_d$ と書く．
一階交換子が Jacobian 括弧であり $F$ がその Casimir であることから
$D(F_{\mathrm{sym}})\in hN'$ である．

求める元を $\Omega=F_{\mathrm{sym}}+h\eta$ と置くと，
中心性と規格化は有限線型方程式
\begin{equation}\label{eq:relative-center-finite-linear-system}
 D\eta=-h^{-1}D(F_{\mathrm{sym}}),\qquad\lambda(\eta)=0
\end{equation}
になる．$s_j$ は既に中央なので，これら以外の交換子方程式は不要である．
写像と右辺を
\[
 L=(D,\lambda):M\longrightarrow N'\oplus S,\qquad
 b=\bigl(-h^{-1}D(F_{\mathrm{sym}}),0\bigr)
\]
と置く．前補題の規格化形式中心から
$\widehat\eta=(\widehat\Omega-F_{\mathrm{sym}})/h$ を作ると，
これは\eqref{eq:relative-center-finite-linear-system}の $h$ 進形式解である．

有限 $S$-加群 $C=\operatorname{coker}L$ における $b$ の像を $\bar b$ とする．
任意の素イデアル $\mathfrak p\supset(h)$ に対し，形式解は
$S_{\mathfrak p}$ の $h$ 進完備化へ写る．
Noether 局所環 $S_{\mathfrak p}$ のこの完備化は忠実平坦なので，
$\bar b_{\mathfrak p}=0$ となる．
従って有限巡回加群 $S\bar b$ の台は $V(h)$ と交わらず，
\[
 \operatorname{Ann}_S(\bar b)+(h)=S
\]
である．よって消滅イデアルの中に $f_1\equiv1\pmod h$ を取れ，
$S_{f_1}$ 上では $b$ は $L$ の像に入る．
その有限線型解 $\eta$ により所要の $\Omega$ を得る．
最初の座標のための局所化と合わせ，分母を払えば，元の $S$ の
$f\equiv1\pmod h$ による一回の主開局所化にまとめられる．

得られた中心性と規格化は有限表示代数内の実際の恒等式である．
従って任意の基底変更後にも保たれる．
この証明は形式級数への $h=1$ の代入ではなく，
有限加群の一つの余核類を消すことで有限パラメータの元を得るものである．
\end{proof}

\begin{remark}[この命題が与える範囲]
上の結果は，全内部次数4または6の中心元を，元の ambient 族の
零切断全体を含む有限主開集合上で与える．
局所的に選んだ Poisson 方向を変えたときの貼合せ，中心商の普遍正則性，
Noether 性，相対 tails および変形の完全性は，この命題の結論には含めない．
また，退化した古典 potential の Poisson 中心全体が $F$ で生成されるとは
主張していない．証明が使うのは一つの規格化 Casimir の存在だけである．
\end{remark}

% !TeX root = ../nc-del-pezzo.tex
% 2026-11-20 JST / v0.8
\section{次数2の全方向中心商の有限代数化}
\label{sec:simultaneous-degree-two-center}

本節では次数2について，四次式の全係数と全Pl\"ucker方向を
同時に扱う中心元を構成する。
第\ref{sec:relative-center-lift}節の有限局所持上げを方向空間上で貼り合わせ，
Pl\"ucker錐の頂点まで降ろす。
係数を特殊化した後の中心の次元が増えても，普遍係数環上の中心の核を
先に計算すれば，この構成を全係数について行える。
得られる主開集合は可換零切断全体を含むが，その追加の非消滅多項式の
明示公式は計算していない。

$N=n+1\ge3$，$E=\C^N$，$\mathcal F=\Sym^4E=\Spec R$ とする。
$F=z^2/2+U$ はこの全係数上の次数2の可換方程式である。
射影四次超曲面 $U=0$ が滑らかである係数の開集合を
$\mathcal T\subset\mathcal F$ と書く。
第\ref{sec:simultaneous-poisson}節のPl\"ucker錐を $\mathcal C$，
その頂点のイデアルを $J$ と書き，$X=\mathcal F\times\mathcal C$ と置く。
まず $g_2\ne0$ の主開 $X_0\subset X$ 上で，同節の代数 $A=A_{U,P}$ を使う。

\subsection{方向空間上の古典中心と消滅}

$G=\Gr(2,E^*)$ 上の標準完全列を
\[
 0\longrightarrow K\longrightarrow E\otimes\mathcal O_G
        \longrightarrow Q\longrightarrow0,
 \qquad \rank Q=2
\]
とし，$L=\det Q^*=\mathcal O_G(-1)$ と置く。
方向解消
\[
 q:Y=\mathcal F\times\operatorname{Tot}(L)\longrightarrow X
\]
はproperであり，頂点の外で同型である。
実際，$Y$ は $\mathcal F\times G\times\mathcal C$ 内の
閉incidence部分で，$G$ への射影は線束 $L$ の全空間である。
以下，$K,L$ の $Y$ への引戻しも同じ記号で書く。
零切断を $Z=\mathcal F\times G$，$L$ のtautological sectionを $\tau$ とする。
局所的には $\tau=h\ell$ であり，普遍形式は $P=hP_0$ となる。

\begin{lemma}[全方向での次数4の古典中心]
\label{lem:simultaneous-degree-two-classical-kernel}
$\mathcal F\times G$ 上の普遍係数環で，
$h^{-1}[-,-]\bmod h$ により得られるPoisson括弧を考える。
その内部次数4の中心の層は
\[
 \operatorname{Sym}^4K\oplus\mathcal O F
\]
である。$U\in\mathcal T$ における全方向のファイバーについても同じ式が成立する。
$z^2$ の係数を取る写像は，第二の直和因子で非零である。
全四次式の特殊化に関して，kernelの形成が基底変更と可換するとは主張しない。
\end{lemma}
\begin{proof}
$K$ の基底を $s_1,\ldots,s_{N-2}$，活性商の基底を $u,v$ とし，
$P_0(u,v)=1$ と規格化する。
次数4の元は一意的に
\[
 w=f_4(s,u,v)+zf_2(s,u,v)+cz^2
\]
と書ける。ここで $f_i$ は通常の次数 $i$ の多項式で，
$c$ は内部次数0の係数である。
$\{w,u\}=\{w,v\}=0$ の $z$ に関する各係数を比較すると
\[
 f_2U_u=f_2U_v=0,\quad
 (f_2)_u=(f_2)_v=0,\quad
 (f_4)_u=2cU_u,\quad (f_4)_v=2cU_v
\]
を得る。
普遍係数の各Grassmann chartは整域であり，$U_u$ は零多項式ではない。
従ってまず普遍係数環上で $f_2=0$ が従う。
最後の二式より $f_4-2cU$ は中央変数 $s_i$ だけに依存する。
この計算はchartの任意の非空開への局所化でも保たれるので，中心の層を与える。

滑らかな射影四次式は非零の線型方向に沿って恒等的に定数にはならない。
実際，その変数を含まない斉次四次式は，対応する射影座標軸の点で特異になる。
従って少なくとも $U_u,U_v$ の一方は零でなく，整域性から $f_2=0$ である。
従って滑らかな四次式の各ファイバーでも同じ核の計算が成立する。
逆に $F$ と $s_i$ はPoisson中心に属するので，これで全てである。
\end{proof}

\begin{lemma}[中心の差を貼り合わせるための消滅]
\label{lem:simultaneous-degree-two-bwb}
全ての $m\ge0$ に対して
\[
 H^1(G,\operatorname{Sym}^4K\otimes\mathcal O_G(m))=0,
 \qquad H^0(G,\operatorname{Sym}^4K)=0
\]
である。従って，$X$ の任意の主開 $D(g)$ の逆像 $Y_g$ でも
\[
 H^1(Y_g,\operatorname{Sym}^4K)=0
\]
となる。
\end{lemma}
\begin{proof}
$N\ge4$ の場合，Borel--Weil--Bottの重み規約を
\cite[Theorem~2.3]{GuoLuSharpe2017}に合わせる。
同文献の標準商束の双対が本稿の $K$ であり，重みは
\[
 (m,m,4,0,\ldots,0)
\]
である。$\rho=(N-1,N-2,\ldots,0)$ を加えると，
$m=0,1$ では長さ2の置換で優越重みになり，$m=2,3$ では成分が重複し，
$m\ge4$ では既に優越重みである。
従って非零コホモロジーが現れ得る次数は順に2，なし，0である。
特に上記の二つの消滅を得る。
$N=3$ では $G\simeq\PP^2$，$K$ は定数直線のtensorを除いて
$\mathcal O_{\PP^2}(-1)$ であるから，
$\mathcal O_{\PP^2}(m-4)$ の通常のコホモロジーから同じ結論が従う。

$\operatorname{Tot}(L)\to G$ はaffineであり，構造層の押出しは
$\bigoplus_{m\ge0}\mathcal O_G(m)$ である。
$\mathcal F$ はaffineなので，上の消滅を係数環 $R$ とtensorして
$H^1(Y,\operatorname{Sym}^4K)=0$ を得る。
有限affine被覆の\v{C}ech複体を使えば，このコホモロジーは
直和およびfiltered colimitと可換する。
主開への制限は，その複体を $q^*g$ で局所化することに等しい。
従って $Y_g$ 上でも消滅は保たれる。
\end{proof}

\subsection{普遍係数上の有限中心層}

次数4の完全対称化写像を
\[
 V_4\otimes\mathcal O_{X_0}\longrightarrow A_4,
 \qquad
 V_4=\operatorname{Sym}^4E\oplus z\operatorname{Sym}^2E\oplus\C z^2
\]
と書く。ここで中央でない $z$ を含む語も全ての文字について対称化する。
この写像は零切断上で同型である。
両辺は同じ階数の有限局所自由加群なので，その同型 locus は
零切断を含む $\operatorname{GL}(E)$-不変な開集合である。
以下ではまずこの開集合へ制限し，$\lambda:A_4\to\mathcal O$
をこの座標における $z^2$ 係数とする。
整列基底での係数は用いない。

\begin{lemma}[規格化中心元の局所構造]
\label{lem:simultaneous-degree-two-central-torsor}
$Z=\mathcal F\times G$ の $\operatorname{GL}(E)$-不変な開近傍
$V\subset q^{-1}(X_0)$ が存在し，
その上の次数4中心の層 $\mathcal Z_4$ は有限局所自由で，
\begin{equation}\label{eq:simultaneous-degree-two-center-extension}
 0\longrightarrow\operatorname{Sym}^4K
 \longrightarrow\mathcal Z_4
 \xrightarrow{\lambda}\mathcal O_V\longrightarrow0
\end{equation}
を満たす。この列は任意の基底変更に対して完全なままである。
従って $\lambda(\Omega)=1/2$ を満たす局所中心元の集合は，
$\operatorname{Sym}^4K$ の加法的torsorである。
ここで基底変更後の $\mathcal Z_4$ は中心の部分加群であり，
基底変更後の交換子のkernel全体との一致は主張しない。
滑らかな係数の零切断 $Z_{\mathcal T}=\mathcal T\times G$ の近傍では，
以下の除した交換子 $\mathcal D$ のkernelの形成は任意基底変更と両立する。
この場合でも，零切断の通常の交換子は零となるため，その中心はより大きい。
\end{lemma}
\begin{proof}
$\widetilde A=q^*A$ と置く。
$h=0$ では可換になるから，各交換子は局所的に $h\widetilde A$ に属する。
各次数の平坦性により $h$-torsionはなく，局所的な割算は一意的である。
局所フレームの変更まで含めると，この割算は大域的には
\[
 \mathcal D:\widetilde A_4\longrightarrow
 \bigl(E^*\otimes\widetilde A_5\oplus\widetilde A_6\bigr)\otimes L^*
\]
を定める。すなわち，$\mathcal D$ に $\tau$ を掛けると
全ての $x_i$ と $z$ に対する交換子の組になる。
この $L^*$ の捩れは，フレーム $\ell'=a\ell$ の下で
$h'=a^{-1}h$，$h'^{-1}[-,-]=a h^{-1}[-,-]$ となることから分かる。
従って $\ker\mathcal D$ は通常の中心の層である。

$K$ の切断は関係式から厳密に中央である。
完全対称化による $\operatorname{Sym}^4K\subset\widetilde A_4$ は
部分束で，$\mathcal D$ と $\lambda$ により零に送られる。
商束 $\mathcal C_4=\widetilde A_4/\operatorname{Sym}^4K$ 上の写像
\[
 (\overline{\mathcal D},\overline\lambda):
 \mathcal C_4\longrightarrow
 \bigl(E^*\otimes\widetilde A_5\oplus\widetilde A_6\bigr)\otimes L^*
       \oplus\mathcal O
\]
を $j$ と書く。
補題\ref{lem:simultaneous-degree-two-classical-kernel}から，
滑らかな四次式の零方向点では $j$ はファイバー上単射である。
従ってある最大小行列式は零でなく，整スキーム $Y$ の一般点で単射となる。
$\mathcal C_4$ は局所自由であるから，$\ker j$ はtorsion部分加群として零である。
これにより，$j$ は普遍係数上の\emph{層の写像として}単射である。
ここでは全ファイバー上の単射性を要求していない。

各Grassmann chartでは $K$ と $Q$ の基底を選べる。
このとき $\widetilde A$ は中央多項式係数を持つ
\eqref{eq:simultaneous-active-presentation}そのものである。
命題\ref{prop:relative-finite-center-lift}を適用すると，
$h=0$ の全点の近傍に，次数4，$z^2$ 係数 $1/2$ の中心元が存在する。
この適用では，Grassmann chartと四次式の係数を内部次数0，
中央変数 $s_i$ を内部次数1として扱う。
最高次CY条件は $g_2\ne0$ により満たされる。

開集合を不変に選ぶ点も明記する。
上の単射 $j$ の余核における $(0,1/2)$ の類が生成する
coherent部分加群を取り，その支持を除く。
これはその類が\emph{茎で零になる} locusであり，
ファイバーだけで零になる条件とは異なる。
局所中心の存在から，この開集合は $Z$ を含む。
全ての写像は自然なので，得られる開集合 $V$ は
$\operatorname{GL}(E)$-不変である。
ここでは $\overline\lambda=1/2$，$\overline{\mathcal D}=0$ を満たす
商束の切断が一意的に存在し，局所的に $\widetilde A_4$ へ持ち上がる。
その二つの持上げの差は $\operatorname{Sym}^4K$ に属する。

得られた商束の切断を $e$ とすると，$j(e)=(0,1/2)$ である。
$j$ の単射性から $\ker\overline{\mathcal D}=\mathcal O e$ となり，
$\overline\lambda(e)=1/2$ によりこの直線は $\mathcal C_4$ の直和因子である。
従って $\mathcal Z_4$ はこの直線の逆像であり，
$\operatorname{Sym}^4K\oplus\mathcal O$ と局所同型で，
$\widetilde A_4$ の部分束でもある。
この分裂により，表示した短完全列自体は任意基底変更後も完全である。

さらに $Z_{\mathcal T}$ の近傍では，$j$ の最大小行列式により
$j$ を分裂単射とできる。
$\mathcal C_4=\mathcal O e\oplus\ker\overline\lambda$ と分けると，
$\operatorname{coker}j$ は
$\overline{\mathcal D}|_{\ker\overline\lambda}$ の余核に等しい。
従ってこの近傍では $\mathcal D$ の像も部分束であり，
$\ker\mathcal D$ の形成は任意基底変更と可換する。
退化した四次式の係数において，この最後の性質までは主張しない。
\end{proof}

\subsection{錐への降下と普遍的に正則な中心商}

\begin{theorem}[次数2の全方向の有限型中心商]
\label{thm:simultaneous-degree-two-central-family}
$X=\mathcal F\times\mathcal C$ 上に
$g(U,0)=1$ を満たす $\operatorname{GL}(E)$-不変な正則関数 $g$ が存在する。
その主開 $X_g$ は $g_2\ne0$ に含まれ，$S=\mathcal O(X_g)$ 上の
第\ref{sec:simultaneous-poisson}節の同じambient代数 $A_S$ は，
次の性質を持つ次数4の中心元 $\Omega$ を持つ。
\begin{enumerate}
\item 完全対称化座標で $\lambda(\Omega)=1/2$ であり，
      $\Omega\bmod J=F$ である。
      $\Omega$ は $\operatorname{GL}(E)$ と両立するように選べる。
\item さらに $g(U,-P)=g(U,P)$ および
      $\Omega\equiv z^2/2+\operatorname{Sym}U\pmod{J^2A_S}$ を要求できる。
\item 任意の可換 $S$-代数への基底変更後も $\Omega$ は左右の非零因子である。
      有限表示代数 $B=A_S/(\Omega)$ は次数ごとに有限局所自由で，
\[
 H_B(t)=\frac{1-t^4}{(1-t)^{n+1}(1-t^2)}
\]
      の係数を階数に持つ。
\end{enumerate}
特にこれは，四次式の全係数の零切断と，その全Pl\"ucker方向を
一つの有限型開近傍で扱う中心超曲面族である。
\end{theorem}
\begin{proof}
\emph{錐側の不変な主開集合。}
補題\ref{lem:simultaneous-degree-two-central-torsor}の $V$ の閉補集合を
properな $q$ で写し，その像を $X_0$ から除く。
得られる $W\subset X$ は全係数の零切断を含む不変な開集合であり，
$q^{-1}(W)\subset V$ となる。
$X$ はaffineなので，閉補集合のイデアル $I_0$ と零切断のイデアル $J$ は
$I_0+J=\mathcal O(X)$ を満たす。
従って $g\in I_0$，$g\equiv1\pmod J$ を選べる。
$\operatorname{GL}(E)$ は標数0で線型簡約であり，$I_0,J$ は不変だから，
Reynolds作用素を適用して $g$ を不変に選び直せる。
この操作は $g\in I_0$ と $g-1\in J$ をともに保つ。
座標環の作用は代数的であるため，ここで使う表現はrationalかつ局所有限である。
必要なら $g(U,P)g(U,-P)$ に替え，主開を符号反転でも保つ。

\emph{貼合せと零切断での規格化。}
$Y_g=q^{-1}(X_g)$ 上で
\eqref{eq:simultaneous-degree-two-center-extension}を考える。
補題\ref{lem:simultaneous-degree-two-bwb}の $H^1$ 消滅により，
定数切断 $1/2$ は大域中心切断 $\widetilde\Omega$ へ持ち上がる。
零切断上の $\widetilde\Omega-F$ は
補題\ref{lem:simultaneous-degree-two-classical-kernel}の普遍係数上の計算から，
$\operatorname{Sym}^4K$ の大域切断である。
補題\ref{lem:simultaneous-degree-two-bwb}の $H^0$ 消滅からこの差は零となる。

\emph{錐への降下と同変性。}
Pl\"ucker錐は正規であり，$\mathcal F$ はaffine空間である。
従って $X_g$ は正規整スキームである。
$q:Y_g\to X_g$ はproper birationalで，$Y_g$ は被約整スキームなので，
\cite[Tag~0AY8, Lemma~37.53.6]{StacksProject}により
$q_*\mathcal O_{Y_g}=\mathcal O_{X_g}$ である。
$A_4$ の有限局所自由性とprojection formulaから
\[
 q_*q^*A_4=A_4
\]
となり，$\widetilde\Omega$ は一意的に $\Omega\in(A_S)_4$ へ降下する。
その交換子の引戻しは零である。
同じprojection formulaを次数5と6に適用すれば，交換子自体も零となる。
規格化と零切断での等式も保たれる。

同変性のために，まず
\[
 \Gamma(Y_g,\mathcal Z_4)\xrightarrow{\lambda}
 \Gamma(Y_g,\mathcal O_{Y_g})
\]
の全射にrational表現の不変部分を取る。
代数的な同変準連接層の大域切断は局所有限なrational表現であり，
線型簡約性により不変部分を取る関手は完全である。
ここでは有限affine被覆の\v{C}ech複体と体上の平坦性により
\[
 \Gamma(\operatorname{GL}(E)\times Y_g,\operatorname{pr}_2^*\mathcal Z_4)
 \simeq\mathcal O(\operatorname{GL}(E))\otimes\Gamma(Y_g,\mathcal Z_4)
\]
となり，同変構造が通常のcoactionを与えることからrational性が従う。
従って不変な切断 $1/2$ の不変な持上げを取れる。
$q$ とprojection formulaは同変なので，降下した $\Omega$ も同変である。
この選択の一意性は主張しない。

\emph{第一平方零近傍との一致。}
生成元を固定して語の順序を反転すると，関係式から
$A_{U,P}^{\mathrm{op}}\simeq A_{U,-P}$ を得る。
この反転は完全対称化座標を固定し，中心を中心へ送る。
$\Omega$ と反転後の $\Omega(U,-P)$ の平均を取れば，
同変性，$\lambda=1/2$，古典極限を保ったまま，係数を $P$ に関して偶にできる。
$J/J^2$ 上では $P\mapsto-P$ が $-1$ として作用するから，
古典極限と偶性により主張(2)を得る。

\emph{正則性と全次数の平坦性。}
各複素ファイバーのambientは整域である。
$P=0$ では多項式環であり，$P\ne0$ では活性最高次代数が
Mori--Ueyama\cite[Corollary~3.8]{MoriUeyama2019}により整域である。
中央多項式変数の添加と，全lower項に対するPBW filtrationを使えば，
全内部次数のambientも整域となる。
完全対称化写像が同型で $\lambda(\Omega)=1/2$ であるため，
$\Omega$ は各ファイバーで非零である。
従って全ての内部次数 $m$ で，掛算
\[
 (A_S)_{m-4}\xrightarrow{\ \Omega\ }(A_S)_m
\]
は全閉点のファイバー上単射である。
これは有限局所自由加群の写像であるから，最大小行列式によって
局所的に分裂単射となり，余核も有限局所自由になる。
この議論は各 $m$ に対して既に選んだ同じ基底 $X_g$ 上で成立し，
無限個の開集合の共通部分を取る必要はない。
分裂は任意基底変更を保つので，$\Omega$ の普遍正則性，
$B$ の普遍的平坦性と所要のHilbert級数を得る。
\end{proof}

\begin{remark}[適用範囲と残る圏論的問題]
本節は最高次binary quarticの楕円判別式を反転していない。
従って滑らかな全四次式 $U$ の接線方向を含む，全ての方向を零切断で扱う。
中心元の存在と普遍正則性は，退化した四次式も含む全係数空間で成立する。
一方，特殊化後の中心全体が普遍中心の基底変更で得られるとは限らない。
例えば $U$ が二つの活性変数に全く依存しない場合，古典中心に
$z\operatorname{Sym}^2K$ も入り，
特殊ファイバーのkernelの次元は増える。
普遍係数上での層の単射性と，ファイバー上での単射性を区別したことが，
この中心の増加を許しながら全係数を扱える理由である。

また，得られたのは同じambient上の有限表示・普遍的に正則な中心商である。
この全方向族の強Noether性，積Noether性，相対tailsのsmooth proper性は
別の入力を要する。一方，全Abel Artin変形に対する形式的完全性は，
後の定理\ref{thm:simultaneous-central-all-artin}で証明する。
有限基底上の圏論的性質と大域的な射を確認するまでは，
本定理を全モジュライのatlasの構成とは呼ばない。
\end{remark}

% !TeX root = ../nc-del-pezzo.tex
% 2026-11-20 JST / v0.8
\section{次数1の同時中心元と有限型中心商}
\label{sec:simultaneous-degree-one-center}

定理\ref{thm:simultaneous-degree-one-ambient}のambient族について，
独立な全方向 $D$ を同時に含む中心元を構成する。
次数1では方向空間上の $H^1$ は自動的に消えないが，
その障害が現れる次数まで中心元のjetを明示し，残る高次を貼り合わせる。
この有限jetは完成した中心元の公式ではなく，
全高次の中心元の存在証明を開始する入力である。

\subsection{座標に依存しない四次までの中心jet}

$E=\langle x_1,\ldots,x_n\rangle$，$A=A_4(x)$，$B=A_6(x)$ と書き，
$D=\sum_iD_i\partial_{x_i}$ は $x$ だけに作用させる。
$A^{(j)}=D^jA$，$B^{(j)}=D^jB$ と置く。
以下の $\operatorname{Sym}$ は $x_i,y,z$ の完全対称化であり，
$x$ だけの多項式の積は既に可換なので対称化を省略する。
零切断のイデアルを $J=(D_1,\ldots,D_n)$ とする。

\begin{lemma}[四次の中心jet]
\label{lem:degree-one-fourth-center-jet}
次の全内部次数6の元を考える。
\begin{align}
 \Omega_{\le4}={}&\frac12z^2+\frac13y^3+\operatorname{Sym}(Ay)+B
   -\frac16\operatorname{Sym}\bigl(A^{(2)}y^2+B^{(2)}y\bigr)
   +\frac1{12}\bigl(A^{(1)}\bigr)^2\notag\\
 &+\operatorname{Sym}\!\left(
       \left(\frac{AA^{(4)}}{180}
                  +\frac{A^{(1)}A^{(3)}}{72}\right)y\right)\notag\\
 &+\frac{AB^{(4)}}{180}
       +\frac{A^{(1)}B^{(3)}}{120}
       -\frac{A^{(2)}B^{(2)}}{120}
       -\frac{A^{(3)}B^{(1)}}{180}.
 \label{eq:degree-one-fourth-center-jet}
\end{align}
これは $\operatorname{GL}(E)$ と両立し，完全対称化座標での
$z^2$ 係数は $1/2$ である。
全生成元 $a\in E\cup\{y,z\}$ に対し
\[
 [\Omega_{\le4},a]\in J^6A_1
\]
が成り立つ。
\end{lemma}
\begin{proof}
方向を分離して $D=h\partial_u$ とし，$\ker D$ の変数を中央係数へ含める。
このとき $A=\sum_{i=0}^4a_i u^i$，$B=\sum_{i=0}^6b_i u^i$ と書ける。
関係式は
\[
 yu=uy-hz,\quad
 zu=uz+h(y^2+A),\quad
 zy=yz-h\bigl(\operatorname{Sym}(A'y)+B'\bigr)
\]
である。ここだけ $'$ は $u$ 微分を表す。
各置換の交換子補正は $h$ を一つ増やすので，
$h^6$ 以上の項を捨てながら順序 $u<y<z$ に整列する操作は有限回で終わる。
同じ $h$ 次数の項では通常の隣接交換によって転倒数が下がり，
補正項では残りの許容 $h$ 次数が下がることが停止性の根拠である。

\eqref{eq:degree-one-fourth-center-jet}を $F+h^2C_2+h^4C_4$ と書き，
上の三置換で全交換子を展開すると，$h$ 次数0から5の各係数が零となる。
これは $\mathbb Q[a_0,\ldots,a_4,b_0,\ldots,b_6]$ 上の恒等式であり，
低次係数や最高次係数の特殊化を用いない。
三つの交換子と全係数の厳密な有限証明書を
\texttt{computations/verify\_degree\_one\_center\_jet.py}に収録する。
この計算が示すのは $h^6$ を法とした恒等式であって，
$\Omega_{\le4}$ 自体の厳密な中心性ではない。

全ての項が $D$ による導分と完全対称化で定義されるため，
線型座標変換と両立する。
従って任意の非零方向のchartで同じ恒等式が成り立つ。
ambientの次数別局所自由性と零切断上の可換性により，
完全対称化は零切断の任意の有限冪零厚化上で同型である。
方向の斉次次数5以下の各係数が全方向で零になることから，
元の普遍係数環で表示した $J^6$ への包含を得る。
$z^2$ を含む項は最初の項だけなので規格化も従う。
\end{proof}

\subsection{高次の貼合せ障害の消滅}

$G=\mathbb P(E^*)$ 上の完全列を
\[
 0\longrightarrow K\longrightarrow E\otimes\mathcal O_G
       \longrightarrow\mathcal O_G(1)\longrightarrow0
\]
とし，$Y=\operatorname{Tot}(\mathcal O_G(-1))\to E^*$ を零切断のblowupとする。
例外因子のイデアルを $I$ とすると，
$I^m/I^{m+1}\simeq\mathcal O_G(m)$ である。

\begin{lemma}\label{lem:degree-one-center-high-cohomology}
$n\ge2$，$m\ge5$ に対し
\[
 H^1(G,\operatorname{Sym}^6K\otimes\mathcal O_G(m))=0.
\]
この消滅は，四次式と六次式の任意のaffine係数空間との積でも保たれる。
\end{lemma}
\begin{proof}
商が線束で標数が0なので，偏極微分による完全列
\[
 0\longrightarrow\operatorname{Sym}^6K(m)
  \longrightarrow\operatorname{Sym}^6E\otimes\mathcal O_G(m)
  \xrightarrow{\delta}
       \operatorname{Sym}^5E\otimes\mathcal O_G(m+1)
  \longrightarrow0
\]
がある。$m\ge0$ では $H^1(G,\mathcal O_G(m))=0$ だから，
全切断上の $\delta$ の全射性を示せばよい。
二組の変数 $u_i,t_i$ を使って
\[
 \delta=\sum_i t_i\partial_{u_i},\qquad
 \varepsilon=\sum_i u_i\partial_{t_i}
\]
と書く。標的の元 $p$ の二重次数は $(5,m+1)$ であり，
$w=m-4\ge1$ と置けば $[\delta,\varepsilon]p=wp$ である。
直接の交換子計算から，
\[
 q=\sum_{j=0}^{5}
    \frac{(-1)^j\varepsilon^{j+1}\delta^jp}
    {(j+1)!\,w(w+1)\cdots(w+j)}
 \quad\text{は}\quad \delta q=p
\]
を満たす。実際，隣接する項が相殺し，末尾は $\delta^6p=0$ で消える。
従って全射である。affine係数環とのtensorは $\mathbb C$ 上完全であり，
同じ有限affine被覆の\v{C}ech複体を使えば最後の主張も従う。
\end{proof}

\begin{lemma}[判別式条件を除いた有限局所中心]
\label{lem:degree-one-all-top-finite-center}
補題\ref{lem:degree-one-all-top-as-open}の環
$T_1=\mathbb C[\alpha,\beta,h,((1+\alpha h^4)(5+\alpha h^4))^{-1}]$
を用い，
\[
 R=T_1[c_{ij}:i+2j<6],\qquad
 U=\frac13y^3+\alpha u^4y+\beta u^6
       +\sum_{i+2j<6}c_{ij}\operatorname{Sym}(u^iy^j)
\]
とする。任意の $p=(\alpha_0,\beta_0,0)\in\operatorname{Spec}T_1$ に対し，
$p$ の主開近傍 $V$ 上で，potential
$uyz-yuz-h(z^2/2+U)$ の同じ代数に中心元 $\Omega_V$ がある。
その活性filtration次数は6以下で，$\Omega_V\bmod h=z^2/2+U^{\mathrm{ab}}$ である。
局所化の分母は $\alpha,\beta,h$ だけに依存する。
$\deg c_{ij}=6-i-2j$ とすると $\Omega_V$ は全内部次数6である。
特に $\Delta=4\alpha^3+9\beta^2$ が零でもこの結論は成立する。
\end{lemma}
\begin{proof}
定理\ref{thm:degree-one-finite-central-family}の23未知数の構成を使う。
その証明で判別式条件が必要だった入力は最高次potentialのCY性であり，
今回は補題\ref{lem:degree-one-all-top-as-open}がこれを全ての $p$ で与える。
全lower項へのPBW性とCY性は
\cite[Proposition~7.3.1(ii)]{EG2010}から従い，
普遍係数環への移行は
補題\ref{lem:degree-one-elliptic-universal-lower-pbw}と同じ有限関係行列の議論による。

以下，中心行列の仮定も確認する。
次数9以下で整列語を基底にする局所化は，最高次の有限関係行列を使って選べる。
その分母は $h=0$ で非零であり，最高次係数だけに依存する。
低次係数を加えても活性次数に関する三角性により同じ局所化でよい。
従って
\[
 \mathcal D:F_6A\longrightarrow F_7A\oplus F_8A\oplus F_9A,
 \qquad v\longmapsto h^{-1}([v,u],[v,y],[v,z])
\]
は正則な有限行列である。
$h=0$ で $v=g_0(u,y)+zg_1(u,y)+kz^2$ と書けば，
$\partial_yU^{\mathrm{ab}}$ の最高 $y$ 次数の項が $y^2$ なので，
定理\ref{thm:degree-one-finite-central-family}の計算は任意の
$\alpha,\beta,c_{ij}$ に対し
\[
 \ker\mathcal D=\mathbb C\cdot1
              \oplus\mathbb C\cdot(z^2/2+U^{\mathrm{ab}})
\]
を与える。従ってlower項が零の $p$ では，活性次数1から6のブロックの階数は
$1,2,3,4,5,6$ である。
各ブロックから合計21本の独立な行を選び，定数項と整列基底の
$z^2$ 係数を読む二行を加えると，$p$ で可逆な23次行列を得る。
lower項を戻しても行列は三角で対角ブロックが変わらないので，
その可逆化は最高次係数だけの主開集合 $V$ で実現できる。

選んだ交換子成分を零，定数項を $c_{00}$，$z^2$ 係数を $1/2$ とする
唯一の解を $\Omega_V$ とする。
$h\ne0$ の各複素点で，\cite[Theorem~3.3.2(ii)]{EG2010}は
非定数中心元を $F_6A$ 内に与える。
この定理の仮定は重みの和 $1+2+3=6$ の斉次CY最高次potentialと
任意のlower potentialであり，generic性や楕円曲線の滑らかさではない。
従って $\dim\ker\mathcal D\ge2$ であり，23次行列の可逆性から
定数項と $z^2$ 係数を読む $\ker\mathcal D\to\mathbb C^2$ は同型となる。
解 $\Omega_V$ は未選択の交換子方程式も全て満たす。
$h=0$ では表示した古典核から解は $z^2/2+U^{\mathrm{ab}}$ となる。
普遍係数環は被約な有限型環なので，全複素点での恒等式は環上の恒等式である。
最後に低次係数の重みを含むスカラー作用と解の一意性から，全内部次数6を得る。
\end{proof}

\begin{lemma}[規格化中心jetの局所持上げ]
\label{lem:degree-one-normalized-center-jets}
$Y$ の例外因子に沿う形式完成上で，完全対称化座標の
$z^2$ 係数を $1/2$ に規格化した次数6の中心元は局所的に存在する。
局所座標 $h$ で $I=(h)$ としたとき，$h^{-1}\operatorname{ad}$ の核を使って
規格化中心jetを定める。
その $I^{r+1}$ を法とするjetから $I^r$ を法とするjetへの射は局所全射であり，
同じjetを持つ持上げの差は
\[
 \operatorname{Sym}^6K\otimes I^r/I^{r+1}
        \simeq\operatorname{Sym}^6K(r)
\]
である。
\end{lemma}
\begin{proof}
各方向chartで $D=h\partial_u$ とし，$K$ の基底を中央変数 $s$ とする。
$A_4(s,u)$ と $A_6(s,u)$ の最高活性係数をそれぞれ $\alpha,\beta$ とする。
補題\ref{lem:degree-one-all-top-finite-center}の低次係数に
$s$ の所要次数の斉次多項式を代入する。
その分母は $\alpha,\beta,h$ だけに依存するため，$s$ の多項式の可逆化は生じない。
得られた中心元は内部次数6で，$h=0$ では
$F=z^2/2+y^3/3+A_4(s,u)y+A_6(s,u)$ となる。
完全対称化座標の $z^2$ 係数を $\lambda$ とすれば
$\lambda=1/2\pmod h$ だから，形式単元 $2\lambda$ で割って所要の規格化を得る。
これは活性変数や $s$ に依存しない内部次数0の中央単元である。
以上は楕円判別式を用いない有限行列による証明である。
局所的な形式中心元の存在は，補題\ref{lem:relative-formal-center-lift}を
同じpotentialに適用しても得られる。

古典側の次数6のPoisson中心は
\[
 \operatorname{Sym}^6K\oplus\mathcal O\,F
\]
である。実際，中心候補を $g_0(s,u,y)+zg_1(s,u,y)+cz^2$ と書く。
$\{-,u\}$ の $z^0$ 係数から $(y^2+A_4)g_1=0$ を得るが，
この因子は $y$ に関してmonicなので任意の係数環上で $g_1=0$ である。
残りの式は $\partial_y g_0=2c\partial_yF$，
$\partial_u g_0=2c\partial_uF$ を与え，
$g_0-2c(F-z^2/2)$ は $s$ だけの次数6多項式となる。

局所的な規格化形式中心元を一つ $\Omega$ とする。
別の規格化中心jetとの差の最低 $h$ 係数は，上の古典核のうち
$z^2$ 係数が零の部分，すなわち $\operatorname{Sym}^6K$ に属する。
この多項式は量子代数でも厳密に中央なので，それを引いて $h$ で割る操作を反復できる。
従って任意の規格化中心jetは $\Omega$ のjetに中央な $s$ 多項式のjetを加えたものであり，
局所全射性と表示した差の層を得る。
ここで $[A_1,A_1]\subset hA_1$ なので，
$h^{-1}\operatorname{ad}$ は $h^r$ を法とする元にも良定義である。
別の例外因子座標への変更は $h$ の単元倍であり，その核を変えない。
\end{proof}

\subsection{形式中心の降下と有限線型方程式による代数化}

\begin{theorem}[次数1の全方向の有限型中心超曲面族]
\label{thm:simultaneous-degree-one-central-family}
定理\ref{thm:simultaneous-degree-one-ambient}の基底環を
\[
 S=\mathcal O(\mathcal F_1\times E^*)[g_1^{-1}],
 \qquad g_1=(1+A_4(D))(5+A_4(D)),\qquad J=(D_i)S
\]
とする。$g\equiv1\pmod J$ を満たす
$\operatorname{GL}(E)$-不変な $g\in S$ があり，$S_g$ 上の同じambient代数に
全内部次数6の中心元 $\Omega$ を構成できる。
$\Omega$ は $\operatorname{GL}(E)$ と両立し，
完全対称化座標の $z^2$ 係数は $1/2$，
\[
 \Omega-\Omega_{\le4}\in J^5(A_1)_6
\]
である。特に零切断上では $\Omega=F$ となる。

$\Omega$ は任意の基底変更後も左右の非零因子であり，
$B_1=A_1/(\Omega)$ は有限表示かつ次数別に有限局所自由である。
任意の非被約基底変更にも両立し，Hilbert級数は
\[
 H_{B_1}(t)=\frac{1-t^6}{(1-t)^n(1-t^2)(1-t^3)}
\]
である。各幾何学的ファイバーは左右Noetherである。
\end{theorem}
\begin{proof}
\emph{全高次の形式中心。}
係数空間 $\mathcal F_1$ との積を暗黙に取る。
補題\ref{lem:degree-one-fourth-center-jet}より，
$\Omega_{\le4}$ は $I^5$ を法とする大域的な規格化中心jetである。
交換子が $I^6$ に入ることを，$h^{-1}\operatorname{ad}$ の方程式に用いている。
$I^r$ を法とするjetが得られたとし，$r\ge5$ とする。
補題\ref{lem:degree-one-normalized-center-jets}により局所的には次のjetへ持ち上がり，
その貼合せ障害は
$H^1(G,\operatorname{Sym}^6K(r))$ に属する。
補題\ref{lem:degree-one-center-high-cohomology}でこれは零である。
従って帰納的に，互いに両立する全次数の中心jetを得る。

\emph{方向のblowupからの形式降下。}
$W_6=(\operatorname{Sym}(E\oplus\mathbb Cy\oplus\mathbb Cz))_6$
を可換重み付き多項式の次数6部分とする。
完全対称化 $W_6\otimes S\to(A_1)_6$ は零切断上で同型だから，
任意の冪零厚化，従って形式完成上でも同型である。
一方，$Y=\operatorname{Tot}(\mathcal O_G(-1))$ について
\[
 \mathcal O_Y/I^r=\bigoplus_{m=0}^{r-1}\mathcal O_G(m),\qquad
 H^0(G,\mathcal O_G(m))=\operatorname{Sym}^mE
\]
である。
従って上で得た大域jetは，文字通り $(A_1)_6/J^r(A_1)_6$ の元を与え，
逆極限は $J$-進完成 $\widehat{(A_1)_6}$ 内の中心元 $\widehat\Omega$ となる。
これは別のstar代数への移送ではなく，同じ有限表示代数の完成内の中心元である。

\emph{有限な方程式。}
必要なら，零切断全体を含む不変な主開集合へさらに縮め，
上の完全対称化を完成前にも同型にする。
この操作は可能である。非同型となる閉集合は不変かつ零切断と交わらないので，
その定義イデアルには $1\pmod J$ となる元がある。
$\operatorname{GL}(E)$ のReynolds作用素で平均すれば，同じ性質を持つ不変元を選べる。
$z^2$ 係数を読む写像を $\lambda:(A_1)_6\to S$ とし，
$N=J^5\ker\lambda$ と置く。
有限 $S$-加群の射
\[
 L:N\longrightarrow
       E^*\otimes(A_1)_7\oplus(A_1)_8\oplus(A_1)_9,
 \qquad \nu\longmapsto([\nu,x_i]_i,[\nu,y],[\nu,z])
\]
を使うと，求める補正は
$L\nu=-\operatorname{ad}(\Omega_{\le4})$ という有限線型方程式の解である。
この方程式は $\widehat\Omega-\Omega_{\le4}\in\widehat N$ によって完成上で解を持つ。

この有限余核を $M$，右辺の類を $m\in M$ とする。
Noether環上の有限加群の完成の完全性から，
$m\in\bigcap_{k\ge0}J^kM$ である。
この交叉を $N_\infty$ とすると，Artin--Rees補題により
$N_\infty=JN_\infty$ となる。
有限生成加群への行列式の議論から，ある $a\in J$ について
$(1-a)N_\infty=0$ を得る。
従って $1-a$ を可逆化すれば元の有限線型方程式が解ける。
この論法は，$S$ 全体の $J$-進完成を忠実平坦と仮定していない。

右辺は不変で $L$ は同変なので，$\operatorname{Ann}(m)$ は不変イデアルである。
$1-a$ をReynolds作用素で平均して，$g\equiv1\pmod J$ となる不変元を得る。
既に行った主開への縮小もこの $g$ に含める。
$g$ の可逆化後に得た解 $\nu$ 自体も平均すれば，
$\Omega=\Omega_{\le4}+\nu$ は全ての規格化と合同式を保つ同変な中心元となる。
係数環の元と解はいずれも有限な式であり，形式級数への有限値代入は行っていない。

\emph{正則性と商の平坦性。}
任意の複素ファイバーで，$D=0$ のambientは多項式環である。
$D\ne0$ では，$\ker D$ の中央多項式環と活性三生成元のPBW filtrationを用いる。
その最高次代数は補題\ref{lem:degree-one-all-top-as-open}のAS正則整域の
中央多項式拡大なので，ambientファイバーも左右Noether整域である。
完全対称化が同型で $\lambda(\Omega)=1/2$ だから，$\Omega$ は全ファイバーで非零である。
従って各次数の掛算
\[
 (A_1)_{m-6}\xrightarrow{\ \Omega\ }(A_1)_m
\]
は全ファイバーで単射となる。
有限局所自由加群の一定階数判定により，これは局所分裂単射で余核も有限局所自由である。
全ての次数でこの議論が成立するため，任意基底変更後の正則性と表示したHilbert級数を得る。
有限表示とファイバーのNoether性は中心商を取っても保たれる。
\end{proof}

\begin{remark}
本定理の追加の主開条件 $g$ は存在を示したもので，その明示的な最小行列式は計算していない。
$g_1\ne0$ というambientの明示条件だけで全中心構成が完成したとは主張しない。
また，この全方向の中心商族について強Noether性，積Noether性，
有限型基底全体での相対tailsのsmooth proper性はここでは証明していない。
同じ中心商族の全Abel Artin変形に対する形式的完全性は
定理\ref{thm:simultaneous-central-all-artin}で証明する。
それも有限型基底全体のtails族や，全モジュライのatlasの構成とは区別される。
\end{remark}

% !TeX root = ../nc-del-pezzo.tex
% 2026-11-20 JST / v0.9
\section{次数2の全最高次式とGorenstein曲線}
\label{sec:degree-two-gorenstein-curve}

次数2の相対切断環表示は，楕円判別式の非零性を仮定せずに構成できる。
本節では節点・尖点・可約曲線と非被約な二重 $\PP^1$ を同時に含める。
必要なのは滑らかさではなく，$\PP^1$ の有限平坦二重被覆，その双対化層と
成分ごとの線束の次数である。
係数環への任意のnoetherian基底変更を許し，積代数のNoether性も証明する。

\subsection{全二元四次式のpoint curveとその解消}

独立な係数を用いて
\[
 f(u,v)=a u^4+4b u^3v+6c u^2v^2+4d uv^3+e v^4,
 \qquad T=\Spec\C[a,b,c,d,e,h]
\]
と置く。活性最高次代数は
\[
 A^{\mathrm{top}}=
 \mathcal O_T\langle u,v,z\rangle/
 ([u,v]-hz,\ [v,z]-h\operatorname{Sym}f_u,
                 [z,u]-h\operatorname{Sym}f_v)
\]
であり，内部次数は $(1,1,2)$ である。
補題\ref{lem:all-binary-quartic-as-open}と
\ref{lem:simultaneous-explicit-as-determinant}のAS開集合上で考える。
以下の追加開条件も全て $h=0$ の全係数を含む。

\begin{lemma}\label{lem:degree-two-all-top-gorenstein-curve}
$T$ の $\operatorname{GL}_2$ と両立する有限型開集合 $V$ で
$\{h=0\}\subset V$ を満たすものが存在し，その上に射影的な平坦曲線族
$\pi:\mathcal E\to V$，自己同型 $\sigma$，および射
\[
 p_1:\mathcal E\longrightarrow\PP^1_V,
 \qquad L=p_1^*\mathcal O_{\PP^1_V}(1)
\]
があり，次を満たす。
\begin{enumerate}
 \item $p_1$ は次数2の有限平坦射である。
       各幾何ファイバーは連結Gorenstein曲線で，
       $H^0(\mathcal O)=\kappa(\bar s)$，算術種数1，双対化層は自明である。
 \item $h=0$ では $\mathcal E$ は
       $z^2/2+f(u,v)=0\subset\PP(1,1,2)$ と同型で，
       $\sigma=1$，$L=\mathcal O(1)$ である。
       この主張は $f=0$ も含む。
 \item $u,v$ は $\pi_*L$ の基底で，$L$ は相対的に豊富かつ大域生成される。
       $z\in H^0(\mathcal E,L\otimes\sigma^*L)$ があり，
       twisted積について $uv-vu=hz$ となる。
 \item 各幾何ファイバーの全既約成分上で，
       $\sigma^*L$ と $L$ の次数は等しい。
\end{enumerate}
開集合 $V$ は有限個のproperな悪条件の像を除いて得られ，
必要なら $h=0$ 上で1となる多項式による主開に選べる。
\end{lemma}
\begin{proof}
二点 $p=(r:s)$，$q=(R:S)$ に対し
\[
 \begin{split}
 A(p,q)&=a rR+b(rS+sR)+c sS,\\
 B(p,q)&=b rR+c(rS+sR)+d sS,\\
 C(p,q)&=c rR+d(rS+sR)+e sS
 \end{split}
\]
と置く。三次関係式の多重線型化を第三点に関する行列にすると，
非零定数倍を除いて
\begin{equation}\label{eq:all-quartic-point-matrix}
 M(p,q)=
 \begin{pmatrix}
 sS+4h^2A&rS-2sR+4h^2B\\
 sR-2rS+4h^2B&rR+4h^2C
 \end{pmatrix}
\end{equation}
となる。これは補題\ref{lem:simultaneous-explicit-as-determinant}の
$4M(w)$ と同じ行列である。
$\delta=rS-sR$，$H=\det M$ とすると
\begin{equation}\label{eq:all-quartic-blowup-equation}
 H=2\delta^2+h^2Q,
 \qquad Q(p,p)|_{h=0}=4f(p).
\end{equation}
例えば最後の等式は，$M|_{h=0,p=q}$ の余因子行列が
$\bigl(\begin{smallmatrix}r^2&rs\\rs&s^2\end{smallmatrix}\bigr)$
であることから直接従う。

$\PP^1_T\times_T\PP^1_T$ を $(h,\delta)$ でblow upし，
$H=0$ のstrict transformを $\mathcal E$ とする。
$h$-chartでは $\delta=hw$ と書く。
$r=R=1$，$t=s/r$，$q=S/R=t+hw$ の表示で，その式は
\[
 G(t,w)=h^{-2}H(t,t+hw)=0,
 \qquad G|_{h=0}=2w^2+4f(1,t).
\]
ここで $h^{-2}H(t,t+hw)$ は多項式である。
他方のchartで $h=k\delta$ と書くと，式は $2+k^2Q=0$ になる。
$h=0$ 上で $k=0$ の点は存在しないから，特殊ファイバーを全て
$h$-chartで見ている。逆座標のchartでは
$t'=1/t$，$q'=1/q$，$w'=-w/(tq)$ と貼り合わさり，
特殊ファイバーは上記の重み付き四次曲線となる。
その $\PP^1$ への射は，$f$ の重根や零性によらず有限平坦な二重被覆である。

この $h$-chartの周囲は $T$ 上滑らかな相対曲面であり，
$G|_{h=0}$ は各ファイバーで非零因子である。
相対Cartier因子の平坦性の判定により，$\mathcal E\to T$ は
$h=0$ の全点の近くで平坦かつ局所完全交叉となる。
非平坦locus等のproperな像を除けば，この性質は族全体で成り立つ。
第一射影 $p_1$ はproperで，特殊ファイバー上で有限である。
非準有限locusのproperな像を除くと $p_1$ は有限となる。
さらに相対平坦性の局所判定を，$\mathcal E$ と $\PP^1_T$ の
$T$-平坦性および特殊ファイバーの有限平坦性に適用すると，
$p_1$ は特殊ファイバーの近くで平坦である。
その非平坦locusと階数が2でないlocusの像も除き，
$p_1$ を次数2の有限平坦射にできる。
これらの像はいずれも $h=0$ と交わらない。

標数0の有限平坦二重被覆では，traceにより
\[
 (p_1)_*\mathcal O_{\mathcal E}
       =\mathcal O_{\PP^1_V}\oplus N^{-1}
\]
と分裂し，$N$ は線束である。平坦曲線族のEuler標数は一定で，
特殊ファイバーのEuler標数は0だから，$N$ の相対次数は2である。
従って局所的に $N=\mathcal O_{\PP^1_V}(2)\otimes\pi_0^*M_0$ と書ける。
ここで $\pi_0:\PP^1_V\to V$，$M_0$ は基底上の線束である。
被覆の代数構造は $N^{-2}\to\mathcal O$ によって決まり，
各幾何ファイバーは $w^2=g_4$ 型の二重被覆である。
随伴公式は
$\omega_{\mathcal E/V}=p_1^*(\omega_{\PP^1_V/V}\otimes N)$
を与えるので，各ファイバーの双対化層は自明である。
また $H^0(\mathcal O)=1$，$H^1(\mathcal O)=1$ であり，
$\pi_*\mathcal O=\mathcal O_V$ は任意基底変更と両立する。
このtrace分裂から $\pi_*L=\mathcal O_V^2$ とその任意基底変更も従う。
$L$ は有限射による $\mathcal O(1)$ の引戻しなので豊富であり，
$u,v$ によって大域生成される。

次にshiftを構成する。AS開条件により行列\eqref{eq:all-quartic-point-matrix}は
$H=0$ 上で階数1であり，その核が第三点を一意的に定める。
有限座標では $A=A(1,t;1,q)$ 等と略記し，
\[
 D=1+4h^2C(t,q),\qquad
 r_{\mathrm{new}}=\frac{2q-t-4h^2B(t,q)}{D}
\]
と書ける。直接展開により
\begin{equation}\label{eq:all-quartic-shift-identity}
 D^2H(q,r_{\mathrm{new}})=K(q)H(t,q),
\end{equation}
ただし
\[
 \begin{split}
 K(q)={}&1+8h^2(c+2dq+eq^2)\\
       &+16h^4\{c^2-bd+(cd-be)q+(d^2-ce)q^2\}
 \end{split}
\]
である。分母を払った多項式恒等式として読む。
解消上では
\begin{equation}\label{eq:all-quartic-regular-shift}
 t_{\mathrm{new}}=q=t+hw,\qquad
 w_{\mathrm{new}}=
 \frac{w-4h\{B(t,q)+qC(t,q)\}}{1+4h^2C(t,q)}.
\end{equation}
分母は $h=0$ で1である。
逆座標の同じ式も合わせれば，特殊ファイバー全体の近傍で正則な射となる。
座標を持たない記述では，核による第三点への射の下で
新しい $\delta$ は $h$ で割り切れ，blow upへの持上げを定める。
未定義locusのproperな像を除けば大域的な射となる。
$H(p,q)=H(q,p)$ であり，二点の交換を $\iota$ とすると
$\iota\sigma\iota$ が逆shiftである。
これらの恒等式は普遍係数の一般点で確認され，
$\mathcal E\to V$ の平坦性と整な係数環により全体へ延長される。
従って $\sigma\in\Aut_V(\mathcal E)$ であり，式から $\sigma|_{h=0}=1$ となる。
第二射影は $p_2=p_1\sigma$ である。
$z=\delta/h$ は $L\otimes p_2^*\mathcal O(1)=L\otimes\sigma^*L$ の切断で，
定義から $uv-vu=hz$ である。

最後に各幾何ファイバーの既約成分を考える。
被約整な場合，$L$ の次数は2である。
可約被約な場合は二つの $\PP^1$ からなり，各成分の $L$ の次数は1である。
非被約な場合は $\PP^1$ の二重構造で，その被約曲線上の次数は1である。
有限平坦二重被覆のtrace表示がこれらの全てを尽くす。
従って自己同型が成分を交換しても $L$ の成分ごとの次数は保存される。

ここまでの悪条件は座標変換と両立し，有限個の閉集合のproper像だけを除いた。
その閉集合のイデアルは $(h)$ と互いに素であるから，
$g\equiv1\pmod h$ を満たす $g$ による主開を選べる。
必要ならReynolds作用素を用いて座標変換と両立する主開に取る。
\end{proof}

\subsection{非被約ファイバーを含むコホモロジー}

\begin{lemma}\label{lem:degree-two-double-cover-cohomology}
前補題の任意の幾何ファイバーを $C$ とする。
線束 $M$ の次数が $C_{\mathrm{red}}$ の全既約成分上で正ならば
$H^1(C,M)=0$ である。
特に
\[
 L_m=L\otimes\sigma^*L\otimes\cdots\otimes(\sigma^{m-1})^*L
 \quad(m\ge1)
\]
に対して，$\pi_*L_m$ は階数 $2m$ の有限局所自由加群で，
高次直像は零，任意基底変更と両立する。
さらに $m\ge2$ なら
\[
 \pi_*L_m\otimes\pi_*L\longrightarrow\pi_*L_{m+1}
\]
は全射であり，この全射も任意基底変更後に保たれる。
\end{lemma}
\begin{proof}
双対化層が自明なので，Serre双対により $H^1(C,M)$ は
$H^0(C,M^{-1})$ の双対である。
$C$ が被約なら，$M^{-1}$ は全成分で負の次数を持つため切断は零である。
非被約な場合，trace表示の四次式は恒等的に零であり，
nilradicalは $C_{\mathrm{red}}=\PP^1$ 上の $\mathcal O(-2)$ である。
$M|_{\PP^1}=\mathcal O(r)$，$r>0$ とすると
\[
 0\longrightarrow\mathcal O_{\PP^1}(-r-2)
   \longrightarrow M^{-1}
   \longrightarrow\mathcal O_{\PP^1}(-r)\longrightarrow0
\]
があるので，この場合も $H^0(C,M^{-1})=0$ となる。
非被約性を無視して被約曲線の次数だけから消滅を結論していない。

$L_m$ は成分ごとに $mL$ と数値同値なので，上の消滅と
$\chi(L_m)=2m$ を適用する。
proper平坦な局所完全交叉族上の線束のコホモロジーと基底変更により，
その導来大域切断は次数0の有限局所自由加群となる。
従って非平坦な基底変更も含めて両立する。
評価列
\[
 0\longrightarrow (\sigma^m)^*L^{-1}
  \longrightarrow\mathcal O_{\mathcal E}^{\oplus2}
  \longrightarrow(\sigma^m)^*L\longrightarrow0
\]
に $L_m$ をtensorすると，核の線束は成分ごとに $(m-1)L$ と数値同値である。
$m\ge2$ ならその $R^1\pi_*$ は零であるから，全射とその基底変更を得る。
\end{proof}

\begin{theorem}\label{thm:degree-two-all-top-thcr}
前補題の開集合を，なお $h=0$ の全係数を含む有限主開 $V$ に制限する。
同じ $A^{\mathrm{top}}$ に次数4の中心元 $\Omega^{\mathrm{top}}$ を取り，
$\Omega^{\mathrm{top}}\bmod h=z^2/2+f$，完全対称座標の $z^2$ 係数を $1/2$ とできる。
この中心元は普遍的に左右正則であり，任意の可換係数変更に両立する同型
\[
 A^{\mathrm{top}}/(\Omega^{\mathrm{top}})
       \simeq D:=\bigoplus_{m\ge0}H^0(\mathcal E,L_m)
\]
が存在する。右辺の積は $s\cdot t=s\otimes(\sigma^m)^*t$ である。
\end{theorem}
\begin{proof}
命題\ref{prop:relative-finite-center-lift}を，中央な正次数変数を持たない
係数環 $\C[a,b,c,d,e]$ 上に適用する。
AS開上の普遍PBW性とpotential対角複体によりその仮定は成立する。
従って全零切断を含む有限主開上で規格化した中心元を得る。
各複素ファイバーのambientはAS正則整域であり，中心元は $z^2$ 係数が非零である。
各内部次数の掛算写像の階数判定により，その余核は有限局所自由となる。
従って中心元は任意の基底変更後も正則で，商のHilbert級数は
$(1-t^4)/((1-t)^2(1-t^2))$ である。

前補題の $u,v,z$ は $A^{\mathrm{top}}$ の全関係式を満たす。
$h\ne0$ では $z=h^{-1}[u,v]$ を消去すると行列
\eqref{eq:all-quartic-point-matrix}の二つの関係そのものであり，
各 $D_m$ の係数環上の平坦性から $h=0$ まで恒等式が延長される。
よって $A^{\mathrm{top}}\to D$ を得る。
次数1では既に全射であり，次数2の切断 $u^2,uv,v^2,z$ は $h=0$ で基底である。
この一つの次数の有限行列の行列式を可逆化すれば，次数2でも全射となる。
前補題の乗法全射から全次数で全射となるので，
無限個の開条件を課していない。

普遍係数の一般ファイバーは整な滑らかな種数1曲線である。
そのtwisted切断環に正次数4の非零中心切断があれば，
同次数の切断の比と可換性を比較して $\sigma^4=1$ が従う。
ここで $m\ge2$ なら $L_m$ は次数4以上で非常に豊富なので，
その同次数の切断比は曲線の函数体を生成する。
しかし式\eqref{eq:all-quartic-regular-shift}は
$t\circ\sigma^4=t+4hw\pmod {h^2}$ を与え，
一般の特殊曲線上で $w$ は零でない。
従って $\sigma^4$ は一般点で恒等写像でなく，
$\Omega^{\mathrm{top}}$ の $D$ における像は零である。
$D_4$ は整な普遍係数環上で局所自由なので，
この消滅は特異・可約・非被約ファイバーを含めて成立する。
中心全体の次元が特殊ファイバーで一定とは仮定していない。

以上より中心商から $D$ への全射を得る。
両辺の次数0の階数は1，正次数 $m$ の階数は $2m$ である。
従って各次数で有限局所自由加群の同型であり，
全次数の同型と任意基底変更との両立が従う。
\end{proof}

\subsection{強Noether性と二つの族の積}

\begin{lemma}\label{lem:degree-two-gorenstein-bisection}
上の $V$ から任意のnoetherian可換環 $R$ へ基底変更して得られる
二つの曲線とtwisted切断環を $(\mathcal E_i,L_i,\sigma_i,D_i)$ とする。
二つの係数変更は異なってもよい。このとき $D_i$ と $D_1\otimes_R D_2$ は
左右noetherianである。反対環を一方の因子に用いても同じである。
\end{lemma}
\begin{proof}
$L_i$ は相対的に豊富で，各ファイバーの全積分曲線上で
$\sigma_i^*L_i$ と数値同値である。
従って $L_{i,m}$ は相対的に $L_i^m$ と数値同値であり，
Keelerのample filter判定\cite[Theorem 1.3]{Keeler2003}から左右
$\sigma_i$-ampleとなる。実際，任意の線束 $H$ に対して
$H\otimes L_{i,m}$ は $m\gg0$ で豊富である。
同文献のTheorem 7.14により $D_i$ は左右noetherianである。
ここで曲線の滑らかさ・整性は仮定していない。

積について補題\ref{lem:elliptic-bisection-noetherian}の証明を拡張する。
その際に使う入力を明記する。
$X=\mathcal E_1\times_R\mathcal E_2$ はnoetherianかつ $R$-flatである。
$\mathcal L_{m,n}=L_{1,m}\boxtimes L_{2,n}$ は大域生成され，
相対K\"unneth公式により
\[
 (D_1)_m\otimes_R(D_2)_n=H^0(X,\mathcal L_{m,n}).
\]
$m,n>0$ ならその導来大域切断は次数0の有限射影加群である。
次数0の因子についても，$\pi_*\mathcal O=R$ とその基底変更から上式が成り立つ。

二重斉次右イデアル $I$ の各成分を評価して得るイデアル層を
$J_{m,n}\subset\mathcal O_X$ とする。
右側に掛ける線束の大域生成により $J_{m,n}$ は両添字について増加する。
$X$ のNoether性から，その和はある $J_{N,N}=J$ に等しい。
一つの評価完全列
\[
 0\longrightarrow K\longrightarrow I_{N,N}\otimes_R\mathcal O_X
   \longrightarrow J\otimes\mathcal L_{N,N}\longrightarrow0
\]
を取る。$H=L_1\boxtimes L_2$ は相対的に豊富であり，
$\tau=\sigma_1^N\times\sigma_2^N$ とすると
$\tau^*\mathcal L_{p,q}\otimes H^{-k}$ は $p,q\ge k$ で相対的にnefである。
これは二つの曲線上の成分ごとの次数がそれぞれ
$(p-k)\deg L_1$，$(q-k)\deg L_2$ であることによる。
非被約構造は数値同値の判定を変えない。
相対Fujita消滅\cite[Theorem 1.5, Erratum Theorem E2.5]{Keeler2003}から，
$p,q\gg0$ で $H^1(X,K\otimes\tau^*\mathcal L_{p,q})=0$ を得る。
さらに $X$ の $R$-平坦性と正次数の導来大域切断の有限射影性により，
$I_{N,N}$ が $R$-flatでなくてもprojection formulaが適用できる。
従って
\[
 I_{N,N}(D_1)_p(D_2)_q
 =H^0(X,J\otimes\mathcal L_{N+p,N+q})
 =I_{N+p,N+q}\qquad(p,q\gg0).
\]
残る有限幅の行と列は，各 $D_i$ 上の有限射影加群の部分加群として有限生成である。
これで全二重斉次右イデアルの有限生成性を得る。
総次数と辞書式の最高成分による帰納で一般の右イデアルも有限生成である。
反対環は $\sigma_i^{-1}$ と $(\sigma_i^{-1})^*L_i$ の切断環なので，
左側にも同じ証明が成立する。
\end{proof}

\begin{corollary}\label{cor:degree-two-all-top-strong-product}
全binary quarticの零切断を含む上の有限主開 $V$ 上で，
$A^{\mathrm{top}}$ とその規格化中心商は強左右Noetherである。
任意のlower potentialを加え，中央変数 $s_1,\ldots,s_r$ の多項式を
その係数に代入した普遍PBW代数を $A$ とする。
任意の次数付き中央商 $B$ に対し，許される全noetherian係数環 $R$ 上で
\[
 A_R,\quad B_R,\quad
 A_R^{\mathrm{op}}\otimes_R A_R,\quad
 B_R^{\mathrm{op}}\otimes_R B_R
\]
は左右noetherianである。
開条件は最高次の $f,h$ だけに依存し，lower係数には条件を課さない。
\end{corollary}
\begin{proof}
前定理と補題を用いる。
正次数4の普遍的に正則な中央元による拡大は，内部次数ごとには有限な
中央元進filtrationを持つので，中心商のNoether性からambientへ持ち上がる。
積では二つの中央元の正則性が次数別平坦性により保たれ，
同じ議論を二度適用する。
中央変数を加えた場合，片側には $R[s_1,\ldots,s_r]$，
積には二組の中央変数を持つnoetherian係数環を用いる。
活性次数による普遍PBW filtrationの随伴次数は，それぞれ最高次ambientと
その積である。増大filtrationの最高記号を消去してNoether性が持ち上がる。
中央商とその積は，対応するambientと積の商だから同じ性質を持つ。
\end{proof}

この系は方向を固定するchart上の結果である。
方向解消から元の同時族への降下では，proper射を忠実平坦射として扱うことはできない。
零方向の完成と非零方向の有限開被覆を用いる降下を，別に適用する必要がある。

% !TeX root = ../nc-del-pezzo.tex
% 2026-11-20 JST / v0.9
\section{次数1の特異な最高次曲線を含む相対切断環}
\label{sec:degree-one-integral-curve}

次数1の最高次式について，楕円判別式の非零条件を除く。
必要な曲線は滑らかとは限らないが，幾何学的に整なGorenstein種数1曲線として扱える。
とくに $\alpha=\beta=0$ の尖点曲線も同じ構成に入る。
本節の有限開条件は最高次係数と一つの方向パラメータ $h$ にだけ依存する。
全独立方向への移行では，方向空間からの降下をさらに行う必要がある。

\subsection{三次根を用いないpoint曲線}

補題\ref{lem:degree-one-all-top-as-open}の係数環を $T_1$ とする。
以下では必要に応じて，$h=0$ の全点を含む有限主開近傍へ制限する。
有限忠実平坦被覆
\[
 T'_1=T_1[k]/(k^2+\alpha),\qquad
 \gamma=6\beta-4k^3,\qquad a=1+kh^2,\quad c=-2
\]
を使う。この被覆の階数は2であり，$T'_1$ は
$\mathbb C[k,\gamma,h]$ の局所化である。
$\gamma$ の三次根を取らず，$\gamma$ も楕円判別式も可逆化しない。
$9+\alpha h^4$ を可逆化すれば $a-c$ が可逆であり，
$a(a+c)=-(1+\alpha h^4)$ より $a,a+c$ も可逆である。

$h\ne0$ で $z$ を消去し，$y=Y+kx^2$ と変数変換すると，巡回potentialは
\[
 W=a x^2Y^2+\frac c2xYxY
       +h^2\left(\frac13Y^3+\frac\gamma6x^6\right)
\]
となる。従って二つの関係式は
\begin{equation}\label{eq:degree-one-unscaled-relations}
 \begin{aligned}
 f_4&=a(x^2Y+Yx^2)+cxYx+h^2Y^2,\\
 f_5&=a(Y^2x+xY^2)+cYxY+\gamma h^2x^5.
 \end{aligned}
\end{equation}
これは補題\ref{lem:degree-one-stephenson-normal-form}の計算から
$Y$ の尺度変更だけを取り除いたもので，$\gamma=0$ でも可逆な変数変換である。

$t=z_0/w_0$，$q=z_1/w_1$ を
$\mathbb P^1\times\mathbb P^1$ の有限座標とし，
$C=a^2-c^2$ と置く。次の式とその二重次数 $(2,2)$ の斉次化を $H$ と書く：
\begin{equation}\label{eq:degree-one-integral-point-curve}
 \begin{split}
 H(t,q)={}&aC(t^2+q^2)+ah^2(a+c)(t^2q+tq^2)\\
 &+(a^2c-\gamma h^6-c^3)tq
     -a\gamma h^4(t+q)-a^2\gamma h^2.
 \end{split}
\end{equation}
$\gamma=\rho^3\ne0$ では，これは既存のStephenson曲線に
$t/\rho,q/\rho$ を代入して $\rho^2$ 倍した式である。
この由来は式の発見に用いるが，以下の恒等式は全 $\gamma$ で成り立つ多項式恒等式である。

\begin{lemma}\label{lem:degree-one-integral-curve}
$\Spec T_1$ の $h=0$ 全体を含む有限主開 $V$ を選ぶと，
その全逆像 $V'$ 上に，平坦で射影的な曲線族 $f:E\to V'$，
自己同型 $\sigma$，相対smooth locus内の切断 $P,Q$ が存在する。
全幾何学的ファイバーは整なGorenstein曲線で算術種数1を持ち，次を満たす。
\begin{enumerate}
 \item $h=0$ の曲線は
 $z^2/2+y^3/3+\alpha x^4y+\beta x^6=0$ であり，$\sigma_0=\operatorname{id}$ である。
 \item $L=\mathcal O_E(P)$ は相対次数1であり，
 $\sigma Q=P$，$L_2:=L\otimes\sigma^*L\simeq p_1^*\mathcal O_{\mathbb P^1}(1)$ である。
 $L_2$ は相対的に大域生成される。
 \item $L_m=L\otimes\sigma^*L\otimes\cdots\otimes\sigma^{(m-1)*}L$ と置くと，
 $f_*\mathcal O_E=\mathcal O_{V'}$，
 $f_*L_m$ は $m>0$ で階数 $m$ の有限局所自由加群，$R^1f_*L_m=0$ である。
 これらは任意の基底変更に両立する。
\end{enumerate}
\end{lemma}
\begin{proof}
\textbf{曲線と特殊ファイバー。}
$\delta=w_0z_1-z_0w_1$ と置くと，
$H|_{h=0}=-3\delta^2$ であり，$H+3\delta^2$ は $h^2$ で割り切れる。
積の射影直線を $(h,\delta)$ に沿ってblow upし，$H=0$ のstrict transformを取る。
$h$-chartの有限座標では $q=t+hw$ と書け，その方程式は
$G=h^{-2}H(t,t+hw)=0$ となる。無限遠では
$u=1/t,v=1/q,\eta=(v-u)/h$ を用いる。
両式は多項式であり，特殊方程式はそれぞれ
\begin{equation}\label{eq:degree-one-integral-special-charts}
 \begin{aligned}
 G_0&=-3w^2-6kt^2-2t^3-\gamma,\\
 G_{\infty,0}&=-3\eta^2-6ku^2-2u-\gamma u^4.
 \end{aligned}
\end{equation}
もう一方のblowup chartで $h=e\delta$ とすると，strict transformの式は
$-3+e^2(\cdots)=0$ である。
従って $h=0$ 上には $e=0$ の点がなく，二つの表示した $h$-chartで尽くされる。

$y=t+k,z=w,\alpha=-k^2,\beta=(\gamma+4k^3)/6$ により
\[
 6\left(\frac{w^2}{2}+\frac{(t+k)^3}{3}
             -k^2(t+k)+\frac{\gamma+4k^3}{6}\right)=-G_0.
\]
$G_0$ は $w$ について二次，$t$ について最高次係数が非零の三次なので，
任意の代数閉体ファイバーで既約である。
無限遠点は $(u,\eta)=(0,0)$ だけで，$u$ 微分は常に $-2$ である。
従って全ての特殊曲線は幾何学的に整であり，無限遠点は滑らかである。
二重点を持つ場合も尖点の場合も，重み付きWeierstrass曲線のGorenstein性と
算術種数1は保たれる。

strict transform上で $h$ は非零因子である。
二つの特殊chartは $w^2$ または $\eta^2$ の係数が $-3$ である相対超曲面なので，
平坦性の局所判定により，$h=0$ の全逆像の近傍で $f$ は平坦かつ相対Gorensteinとなる。
その外の閉集合の固有像を除く。
平坦で固有かつ有限表示な射の幾何学的被約ファイバーの開性
\cite[Tag 0C0E]{StacksProject}を適用し，全幾何学的ファイバーを被約にする。
算術種数は平坦射影族で一定なので1である。
整性は後で次数1の豊富な線束を使って示す。
この操作で特異なファイバーを除いてはいない。

\textbf{shift。}
次の行列を使う：
\[
 \begin{aligned}
 m_{11}&=cw_0z_1+az_0w_1,&
 m_{12}&=(aw_0+h^2z_0)w_1,\\
 m_{21}&=a\gamma h^2cw_0w_1+\gamma h^4cw_0z_1+aCz_0z_1,\\
 m_{22}&=cCz_0w_1+ah^2(a+c)z_0z_1+aCw_0z_1.
 \end{aligned}
\]
$\det M=cH$ である。$h=0$ の有限chartでは第1行が $(-t,1)$，
無限遠では第2行が $(-3,0)$ なので，rank零箇所の固有像を除くと
$M$ のkernelは全曲線上で第三点を定める。
その有限表示は
\begin{equation}\label{eq:degree-one-integral-shift}
 r=-\frac{cq+at}{a+h^2t},\qquad
 (a+h^2t)^2H(q,r)=(a^2-h^2cq)H(t,q).
\end{equation}
$h=0$ 上で第二点と第三点が一致するので，
blowupの中心理想の逆像は $(h)$ となる。
blowupの普遍性によりshiftはstrict transformへ持ち上がり，有限chartでは
\begin{equation}\label{eq:degree-one-integral-lifted-shift}
 t'=t+hw=q,\qquad
 w'=\frac{w-h\{k(t+q)+tq\}}{a+h^2t}
\end{equation}
となる。$h=0$ で恒等写像である。
射の正則性はrank1の行列で示しているため，表示分母が消える点も含まれる。
二つの射影直線の交換を $J$ と書くと，
$a(t+r)+cq+h^2tr=0$ の外側の変数の対称性から $J\sigma J=\sigma^{-1}$ を得る。
$h\ne0$ 上の恒等式は，$h$-平坦性と分離性により全体で成り立つ。

\textbf{線束とコホモロジー。}
元の積の上で
\[
 Q=((0:1),(0:1)),\qquad
 P=((0:1),(h^2:c-a))
\]
を取る。無限遠chartの持上げは
$Q:(u,\eta)=(0,0)$，$P:(u,\eta)=(0,h/(c-a))$ である。
両者は $h=0$ 上で相対smooth locusに入り，その性質の開性により全体でもそうできる。
従って $P,Q$ は相対次数1の有効Cartier因子である。
行列のkernelを $Q$ で計算すると $\sigma Q=P$ である。
無限遠chartの厳密な因子分解
\begin{equation}\label{eq:degree-one-integral-infinity}
 G_\infty(0,\eta)=a(a+c)\eta\bigl((a-c)\eta+h\bigr)
\end{equation}
により $p_1^{-1}(\infty)=P+Q$ がCartier因子として成り立つ。
とくに $h=0$ で二切断が合流する点でも，$G_\infty$ の $u$ 微分は $-2$ という単元である。
従って局所的に $G_\infty(u,\eta)=G_\infty(0,\eta)+uK$，$K$ は単元と書け，
曲線上の $u$ は $\eta((a-c)\eta+h)$ の単元倍となる。
これは交点での重複度も含む等式を与える。
従って $L\otimes\sigma^*L=p_1^*\mathcal O(1)$ であり，後者は大域生成される。

$L$ は全ての $h=0$ ファイバーで次数1の線束なので豊富である。
相対豊富性の開性\cite[Proposition 2.5]{Keeler2003}により，
全零切断を保って縮めれば $L$ は相対的に豊富となる。
有効Cartier切断で定義したため，全幾何学的ファイバーの $L$ の全次数は1である。
被約な射影曲線上では，豊富な線束の次数は各既約成分で正の整数であり，
全次数はそれらの和である。従って各ファイバーの既約成分は一つしかない。
既に被約性を確保したので，全幾何学的ファイバーは整である。

整なGorenstein種数1曲線上では，双対化線束は自明である。
実際，その次数は0で，Serre双対により非零切断を持つため，零因子を持てない。
正次数の線束 $N$ には $H^1(N)=H^0(N^{-1})^*=0$ が成り立ち，
Riemann--Rochから $h^0(N)=\deg N$ を得る
\cite[Tag 0BS6]{StacksProject}。
これを全ファイバーの $L_m$ に適用し，proper flatな族のcohomology and base changeを使う。
これで非被約基底を含む主張3が従う。

最後に，各縮小を元の $\Spec T_1$ 上で行えることを確認する。
$V'\to V$ を有限被覆の全逆像として保ち，除く閉集合の有限像を基底から除く。
全ての悪い集合は $h=0$ の全点と交わらない。
残る開集合の補集合のイデアルを $I$ とすれば $I+(h)=(1)$ なので，
$g\in I$，$g\equiv1\pmod h$ を選んだ主開で全構成を行える。
\end{proof}

\subsection{同じ最高次中心商の切断環表示}

\begin{theorem}\label{thm:degree-one-all-top-thcr}
$\Spec T_1$ の $h=0$ 全体を含む有限主開 $V$ 上で，
同じ最高次ambient $A^{\mathrm{top}}$ に次数6の普遍的な中央非零因子
$\Omega^{\mathrm{top}}$ を選べる。
完全対称化座標の $z^2$ 係数は $1/2$，$h=0$ での値は古典の最高次式である。
補題\ref{lem:degree-one-integral-curve}の階数2被覆 $V'\to V$ 上で
\[
 \bigl(A^{\mathrm{top}}/(\Omega^{\mathrm{top}})\bigr)_{V'}
       \simeq D(E,L,\sigma):=\bigoplus_{m\ge0}f_*L_m
\]
が成り立つ。同型は任意の可換基底変更に両立する。
\end{theorem}
\begin{proof}
最高次中心は補題\ref{lem:degree-one-all-top-finite-center}でlower項を零にして得る。
全零切断を一つの有限主開で扱うには，次数6と交換子の有限行列の余核を用いる。
各零切断点で規格化中心が存在するので，その余核類の支持は零切断に交わらない。
支持のイデアルから $g\equiv1\pmod h$ を選べば，有限線型方程式の解を全体で選べる。
同時に有限個の完全対称化座標と $z^2$ 係数の単元性を確保する。
これは命題\ref{prop:relative-finite-center-lift}と同じ有限余核の操作である。
各ファイバーのambientは整域であり，規格化により中心は非零である。
各内部次数の乗法行列は全ファイバーで単射なので，局所的に分裂した単射となる。
従って中心は普遍的に正則で，商の正次数 $m$ の階数は $m$ である。

$D$ の切断を有理関数で表示し，
\[
 x=1\in D_1,\quad Y=t\in D_2,\quad
 y=Y+kx^2,\quad z=\frac{xY-Yx}{h}\in D_3
\]
と置く。$D/(h)$ は可換であり $D_3$ は $h$-平坦なので，最後の割算は正則で一意である。
有限chartで $z=w$ となる。
式\eqref{eq:degree-one-integral-shift}と，$s=t\circ\sigma^3$ に対する恒等式
\[
 (a+h^2t)(a+h^2q)\{a(tr+qs)+cts+\gamma h^2\}=-H(t,q)
\]
から，式\eqref{eq:degree-one-unscaled-relations}の二関係式が成立する。
従って $h\ne0$ 上でambientからの写像があり，$D_m$ の $h$-平坦性により
元の三生成元の全関係式は $h=0$ を含めて成立する。

特殊曲線では無限遠の $t,w$ の極の位数は2,3であり，
次数1から4の切断はそれぞれ
$(1)$，$(1,t)$，$(1,t,w)$，$(1,t,t^2,w)$ を基底に持つ。
このことは尖点の場合も同じである。
有限な四つの余核を消して，これらの次数の全射性を確保する。
$m\ge3$ では，$\sigma^{m*}L_2$ の二生成元の評価列を $L_m$ でtensorすると，
そのkernelのファイバー次数は $m-2>0$ となる。
前補題の $H^1$ 消滅と基底変更により
\[
 D_mD_2=D_{m+2}\qquad(m\ge3)
\]
が成り立ち，全次数で $A^{\mathrm{top}}\to D$ は全射となる。

一般点では $\Omega^{\mathrm{top}}$ の像は零である。
式\eqref{eq:degree-one-integral-lifted-shift}により
$t\circ\sigma^6=t+6hw\pmod {h^2}$ なので，$\sigma^6$ は一般点で恒等写像でない。
一方，非零の次数6中心切断があれば，非常に豊富な高次数の切断の比との可換性から
$\sigma^6$ は関数体に恒等に作用する。
一般点の曲線は滑らかで整だから矛盾となる。
$D_6$ は整な普遍係数環上で捩れを持たないので，中心の像は全体で零である。
得られた中心商から $D$ への全射は，各次数で同じ階数の有限局所自由加群の間の写像であり同型となる。
全ての構成が基底変更に両立するため，最後の主張も従う。
\end{proof}

\subsection{強Noether性，積，および全低次項}

\begin{lemma}\label{lem:integral-genus-one-bisection}
補題\ref{lem:elliptic-bisection-noetherian}の滑らかさの仮定は，
平坦なGorenstein曲線族で全幾何学的ファイバーが整かつ算術種数1である，
という仮定に弱められる。ただし次数2の線束の相対大域生成を仮定する。
従ってその二つのTHCRのtensor積は，任意のnoetherian基底変更後も左右Noetherである。
\end{lemma}
\begin{proof}
正次数線束の $H^1$ 消滅と次数別有限射影性は前補題と同じRiemann--Rochと双対で従う。
自己同型は各整な曲線ファイバーの次数を保つので，相対数値同値群には恒等的に作用する。
従ってKeelerのTheorems 7.13--7.14\cite{Keeler2003}は特異な曲線族にも適用できる。
既存補題の二重評価イデアル $J_{m,n}$ は，大域生成により同じく増加し，
noetherianな積のスキーム上で有限段階で安定する。
二曲線からの引戻しの外積は，各ファイバーの次数を調べると同じ相対nef性を持つ。
Keelerの相対Fujita消滅（同論文Theorem 1.5と訂正Theorem E2.5）は
滑らかさを仮定しないので，安定した一つの評価列のkernelの高象限での $H^1$ を消す。
積のスキームは基底上平坦であり，正次数の外積線束の導来切断は有限射影で次数0に集中する。
従って評価イデアル成分自体の平坦性を仮定せずprojection formulaを使える。
残る有限幅の行と列の議論も，片側THCRのNoether性と有限射影性だけを使う。
これらは既存補題の証明の全入力であり，滑らかさを使う箇所は残らない。
\end{proof}

\begin{theorem}\label{thm:degree-one-all-top-noether}
ある有限主開 $V\subset\Spec T_1$ は $h=0$ の全点を含み，次を満たす。
最高次ambient，全12個のlower係数を独立に加えたambient，および
$n-1$ 個の中央一次変数にlower係数を斉次多項式で代入した全 $n\ge2$ のambientについて，
任意の許されるnoetherian可換係数環 $R$ 上で
\[
 A_R,\qquad A_R^{\mathrm{op}}\otimes_R A_R
\]
は左右Noetherである。従って同じambientの任意の中心商 $B_R$ と
$B_R^{\mathrm{op}}\otimes_R B_R$ も左右Noetherである。
特に，選んだ規格化中心商にもこれらの性質が成り立つ。
$V$ の分母は $\alpha,\beta,h$ だけに依存する。
\end{theorem}
\begin{proof}
まず階数2の有限忠実平坦被覆の上で，前定理の $D=D(E,L,\sigma)$ を使う。
$L$ の相対次数は1で豊富であり，自己同型は曲線ファイバーの数値同値を保つ。
Keelerの定理から，任意のnoetherian係数変更後に $D$ は左右Noetherである。
偶数Veronese $D^{(2)}=D(E,L_2,\sigma^2)$ には前補題を適用できる。
左右の評価列から
$D_mD_2=D_{m+2}=D_2D_m$ ($m\ge3$) となるため，
$D$ は $D^{(2)}$ の左右有限生成加群である。
従って $D^{\mathrm{op}}\otimes_R D$ も
$(D^{(2)})^{\mathrm{op}}\otimes_R D^{(2)}$ の左右有限生成加群であり，左右Noetherである。

正次数6の普遍的な中央非零因子 $\Omega^{\mathrm{top}}$ を戻す。
その随伴次数は $D[v]$ であり，内部次数が非負なので斉次イデアルの生成元の持上げは有限回で止まる。
積には二つの中心元を順に戻す。各代数と中心商の係数環上の平坦性により，
途中の商でも二つの元は正則である。
これで最高次ambientとその積のNoether性を得る。
有限忠実平坦被覆への基底変更はイデアルの上昇列の停止を検出するので，結論は $V$ へ降下する。

全lower項の活性filtrationの随伴次数は最高次ambientである。
そのPBW同型は任意の中央係数変更に両立するから，filtered持上げによりNoether性を得る。
積にも両因子の活性filtrationを用いれば同じ結論となる。
全次元への移行では中央多項式環 $R[s_1,\ldots,s_{n-1}]$，
積では二組の中央変数を持つ多項式環を係数環として同じ議論を行う。
中心商はNoether環の商であり，その積もambientの積の商なので結論が従う。
\end{proof}

\begin{corollary}\label{cor:degree-one-all-top-tails}
本節の有限開集合上で，規格化した普遍正則な斉次中心元を用いた全次元の中心商について，
定理\ref{thm:degree-one-relative-tails}の相対CY性，局所コホモロジー，
飽和性，有向ACM，相対 $\chi$，全シフトの豊富性，compact生成とproper性，
および $S_R\simeq(-(n-1))[n]$ が成り立つ。
compact生成窓は $\bigoplus_{j=0}^{n+4}\mathcal O(j)$ である。
全斉次式が滑らかな可換del Pezzo模型を与える零方向点の有限開近傍では，
正周期 $n-1$ の列が豊富となり，tailsのcompact圏は相対的にsmooth properとなる。
\end{corollary}
\begin{proof}
相対potential対角複体は補題\ref{lem:degree-one-all-top-as-open}により
楕円判別式によらず普遍的に完全であり，中央差のKoszul複体を加えたCY shiftは $n+5$ である。
全lower項のfiltered持上げと次数6の普遍的に正則な中心商も既に得ている。
これらと前定理の強Noether性が，定理\ref{thm:degree-one-relative-tails}の証明で用いた入力の全てである。
従ってその局所コホモロジー，Serre双対，有限windowの証明を適用できる。
さらに積のNoether性が得られたので，第\ref{sec:finite-smoothness}節の
対角の有限部分解消と $HH^1(X)=0$ による持上げを適用できる。
最後の縮小は全斉次係数にも依存してよく，最高次曲線自体の滑らかさは要求しない。
\end{proof}

\begin{remark}
整な種数1曲線を用いたのは活性な最高次中心商の表示だけである。
全lower項を持つ高次元のtailsが曲線のcoherent圏になるという主張ではない。
また本節だけから，全独立方向を持つ同時族の基底でNoether性が降下するとは結論していない。
明示的なAS開条件の全域での結論でもなく，零切断を含む追加の有限主開を選んでいる。
\end{remark}

% !TeX root = ../nc-del-pezzo.tex
% 2026-11-20 JST / v0.9
\section{完成と頂点外の被覆による全方向族のNoether性}
\label{sec:noether-descent}

方向解消は頂点上で平坦ではないため，その上のNoether性を
忠実平坦降下として直接用いることはできない。
本節では，次数別に有限な代数に対するKrull交叉の議論と，
係数環の完成および頂点外の主開集合からなる一つの忠実平坦被覆を使う。
これにより，退化した活性最高次式を含む相対曲線表示を，
全独立Poisson方向の同時族へ移す。

\subsection{次数別有限性とJacobson根基}

\begin{lemma}\label{lem:graded-jacobson-noether}
$R$ を可換Noether環，$J\subset\operatorname{Jac}(R)$ をイデアルとする。
$C=\bigoplus_{d\ge0}C_d$ を $R$ が中央に作用する非負次数付き代数とし，
各 $C_d$ は有限 $R$-加群であると仮定する。
$C/JC$ が右Noetherなら $C$ も右Noetherである。
左側についても同じ主張が成立する。
$C_d$ の $R$-平坦性や $R$ の完備性はこの補題の仮定に含めない。
\end{lemma}
\begin{proof}
$J=(a_1,\ldots,a_s)$ と書く。
$J$-進随伴次数環には中央変数を持つ多項式環からの全射
\[
 (C/JC)[T_1,\ldots,T_s]\longrightarrow
 \operatorname{gr}_J C
   =\bigoplus_{r\ge0}J^rC/J^{r+1}C,
 \qquad T_i\longmapsto a_i\bmod J^2C
\]
がある。従って中央変数についてのHilbert基底定理により
$\operatorname{gr}_J C$ は右Noetherである。
ここには $J$-進次数と内部次数の二つの次数があり，
$\deg(T_i)=(1,0)$ である。

内部次数に関して斉次な右イデアル $I\subset C$ を取る。
誘導された $J$-進filtrationから
\[
 \operatorname{gr}_J I
 =\bigoplus_{r\ge0}
       (I\cap J^rC)/(I\cap J^{r+1}C)
 \subset\operatorname{gr}_J C
\]
を得る。これは二重斉次な右イデアルなので，有限個の二重斉次生成元を選べる。
その持上げを $f_1,\ldots,f_t\in I$ とし，
内部次数 $\deg f_i=d_i$，$J$-進次数 $r_i$ を記す。
$H=\sum_i f_iC\subset I$ と置く。
初項を順に消去することで，任意の $x\in I_d$ と任意の $r\ge0$ に対して
\[
 x\in H_d+J^rC_d
\]
となる。実際，各段階で内部次数 $d-d_i$ の係数だけを用いれば，
誤差の内部次数を $d$ に保ちながら $J$-進次数を一つずつ上げられる。
各 $r$ について必要なのは有限回の消去であり，無限和を $C$ の元とみなす必要はない。

$C_d/H_d$ は有限 $R$-加群である。
$J\subset\operatorname{Jac}(R)$ なのでKrull交叉
\cite[Tag~00IQ, Lemma~10.51.5(2)]{StacksProject}から
\[
 \bigcap_{r\ge0}J^r(C_d/H_d)=0
\]
であり，$x\in H_d$ が従う。従って $I=H$ である。
これで全ての斉次右イデアルが有限生成となった。

一般の右イデアルについては，内部次数の最高項が作る斉次右イデアルを取り，
その有限個の斉次生成元を元の右イデアルへ持ち上げる。
最高項の消去で内部次数が真に下がるので，今度は有限回で厳密に停止する。
従って一般の右イデアルも有限生成である。
反対代数に適用すれば左側も従う。
\end{proof}

\begin{corollary}\label{cor:simultaneous-local-noether}
両次数の同時有限中心族 $A$，$B=A/(\Omega)$ を，
定理\ref{thm:simultaneous-degree-two-central-family}または
定理\ref{thm:simultaneous-degree-one-central-family}の基底から
可換Noether環 $R$ へ変更する。
方向座標の生成するイデアル $J_R$ が $\operatorname{Jac}(R)$ に含まれるなら，
\[
 A_R,\qquad B_R,\qquad
 A_R^{\mathrm{op}}\otimes_R A_R,\qquad
 B_R^{\mathrm{op}}\otimes_R B_R
\]
は左右Noetherである。
特に，可換零方向点へ閉点を送る任意のNoether局所基底に成立し，
Artin環や完備局所環に限定されない。
\end{corollary}
\begin{proof}
$A_R/J_RA_R$ は重み付き可換多項式環であり，
$B_R/J_RB_R$ はその可換超曲面商である。
いずれも $R/J_R$ 上有限型なのでNoetherである。
両代数の各内部次数は有限 $R$-加群であるから前補題を適用できる。
積代数には二つの内部次数の和を用いる。
各総次数には有限個のtensor成分しかなく，係数環上有限である。
$J_R$ を法とする積代数も有限型可換代数なので，同じ補題が適用できる。
\end{proof}

\subsection{完成と主開集合の忠実平坦被覆}

\begin{lemma}\label{lem:completion-punctured-flat-cover}
$R$ を可換Noether環，$J=(a_1,\ldots,a_s)$ を任意のイデアルとし，
$\widehat R=\varprojlim_r R/J^r$ と置く。
このとき
\begin{equation}\label{eq:completion-punctured-flat-cover}
 R\longrightarrow R'=
     \widehat R\times\prod_{i=1}^s R[a_i^{-1}]
\end{equation}
は忠実平坦であり，$R'$ はNoether環である。
また $J\widehat R\subset\operatorname{Jac}(\widehat R)$ である。
零環となる因子は除いてよい。
\end{lemma}
\begin{proof}
$R\to\widehat R$ の平坦性，$\widehat R$ のNoether性と
$\widehat R/J\widehat R=R/J$ は
\cite[Tags~00MB, 031C, 05GH]{StacksProject}による。
局所化も平坦であり，有限積は加群として有限直和だから，
表示した射は平坦である。$\Spec R$ のうち $V(J)$ の点は
$\Spec(\widehat R/J\widehat R)=\Spec(R/J)$ を通じて
$\Spec\widehat R$ の像に入る。それ以外の点は，いずれかの
$\Spec R[a_i^{-1}]$ の像に入る。
従ってスペクトル上全射であり，
\cite[Tag~00HQ, Lemma~10.39.16]{StacksProject}から忠実平坦である。
各因子はNoetherで，因子数は有限だから $R'$ もNoetherである。
最後に $u\in J\widehat R$ と $v\in\widehat R$ に対し，
$\sum_{r\ge0}(uv)^r$ は $J\widehat R$-進位相で収束して
$1-uv$ の逆元になる。これは $J\widehat R$ がJacobson根基に入ることを意味する。
\end{proof}

\begin{proposition}\label{prop:noether-completion-punctured-gluing}
$R,J=(a_1,\ldots,a_s)$ と $C=\bigoplus_{d\ge0}C_d$ は
補題\ref{lem:graded-jacobson-noether}と同じとするが，
$J\subset\operatorname{Jac}(R)$ は仮定しない。
次の二条件を仮定する。
\begin{enumerate}
 \item $C/JC$ は左右Noetherである。
 \item 各 $C[a_i^{-1}]$ は左右Noetherである。
\end{enumerate}
このとき $C$ は左右Noetherである。
\end{proposition}
\begin{proof}
$\widehat C=C\otimes_R\widehat R=\bigoplus_d(C_d\otimes_R\widehat R)$
を考える。これは係数環を完成した通常の次数付き代数であり，
環 $C$ 全体の $J$-進完成を意味しない。
各次数は有限 $\widehat R$-加群で，
\[
 \widehat C/J\widehat C\simeq C/JC
\]
である。前補題のJacobson根基条件と
補題\ref{lem:graded-jacobson-noether}により $\widehat C$ は左右Noetherとなる。
従って
\[
 C\otimes_R R'\simeq
   \widehat C\times\prod_i C[a_i^{-1}]
\]
は左右Noetherである。
右イデアルの上昇列を $R'$ へtensorすると，$R'$ の $R$-平坦性により
右イデアルの上昇列のままであり，有限段階で停止する。
その後の各商が $R'$ へのtensorで零になることと忠実性から，
元の上昇列も同じ段階で停止する。左側も同様である。
\end{proof}

この命題では，完成だけを忠実平坦と仮定していない。
零切断から離れた主開集合を有限個加えることが，忠実性を保証している。
また，それぞれの局所環上のNoether性から，一様な有限開近傍を
推測する議論も用いていない。

\begin{remark}\label{rem:jacobson-localization-cover}
完成の代わりに乗法的集合 $1+J$ による局所化
$R^\circ=R[(1+J)^{-1}]$ を用いることもできる。
実際，$j\in J$，$r\in R$，$s,t\in1+J$ なら
\[
 1-\frac js\frac rt=\frac{st-jr}{st}
\]
の分子と分母はいずれも $1+J$ に属するので，
$JR^\circ\subset\operatorname{Jac}(R^\circ)$ である。
$R^\circ$ はNoetherな平坦局所化で，$R^\circ/JR^\circ=R/J$ だから，
$R\to R^\circ\times\prod_iR[a_i^{-1}]$ もスペクトル上全射で忠実平坦である。
従って補題\ref{lem:graded-jacobson-noether}による同じ降下が成立する。
ここでは $1+J$ の全元を可逆化しており，この局所化が有限表示であるとは主張しない。
一つの元の可逆化に置き換えたのでもなく，降下にはNoether性と忠実平坦性で十分である。
\end{remark}

\subsection{方向解消上の曲線表示から同時族へ}

\begin{theorem}\label{thm:simultaneous-strong-product-noether}
両次数 $d=1,2$，全 $n\ge2$ について，同時有限中心族の基底を
全可換係数の零切断を含む $\operatorname{GL}(E)$-不変な有限主開集合へ
さらに縮めることができ，その上の同じambient $A$ と中心商 $B=A/(\Omega)$ は
強左右Noetherである。さらに，任意の可換Noether係数環 $R$ への変更後に
\[
 A_R^{\mathrm{op}}\otimes_R A_R,\qquad
 B_R^{\mathrm{op}}\otimes_R B_R
\]
も左右Noetherとなる。
この主開集合は，活性最高次式の楕円判別式が零になる方向も含む。
\end{theorem}
\begin{proof}
\emph{方向解消上の入力。}
中心構成後のaffine基底を $X_0$，零方向のイデアルを $J$ と書く。
次数1では $\mathcal F_1\times E^*$，次数2では
$\mathcal F_2\times\mathcal C$ の既に選んだ主開である。
方向解消を $q:Y_0\to X_0$ とすると，次数1では
$\operatorname{Tot}\mathcal O_{\PP(E^*)}(-1)$，次数2では
$\operatorname{Tot}\mathcal O_{\Gr(2,E^*)}(-1)$ との係数空間の積からの基底変更である。
$q$ はproperであり，$X_0\setminus V(J)$ 上で同型である。

方向束と中央部分束を自明化した各chartでは，
引戻したambientは，一つの活性方向 $h$ と中央一次変数 $n-1$ 個を持つ
既存のPBW族である。
次数1には定理\ref{thm:degree-one-all-top-noether}，
次数2には系\ref{cor:degree-two-all-top-strong-product}を適用する。
これらは全活性最高次係数の $h=0$ を含む有限開近傍で，
任意lower係数，任意個数の中央多項式変数と，積代数のNoether性を扱う。
従って $q^{-1}V(J)$ 全体を含む開集合 $W\subset Y_0$ を取り，
その上でambientと積代数が強左右Noetherになるようにできる。
これは局所的な性質であり，座標変更は代数の同型を与える。
必要なら $W$ の全ての $\operatorname{GL}(E)$-移動の和を取り，$W$ を不変にする。

\emph{基底の有限主開の選択。}
proper性により $q(Y_0\setminus W)$ は閉で，$V(J)$ と交わらない。
その閉集合を除いた不変開集合を $U\subset X_0$ と置く。
$X_0$ はaffineであるから，$X_0\setminus U$ の定義イデアル $I_0$ と
$J$ は $I_0+J=\mathcal O(X_0)$ を満たす。
従って $g\in I_0$，$g\equiv1\pmod J$ を取れる。
二つのイデアルは不変なので，標数0のReynolds作用素を用いれば
$g$ を不変にできる。$X_g\subset U$ を新しい基底とする。
この操作により $q^{-1}(X_g)\subset W$ であり，
全係数の零切断を保ったまま，全方向のchart入力を同時に適用できる。

\emph{任意Noether基底上の頂点外。}
$S=\mathcal O(X_g)$ とし，任意の可換Noether $S$-代数 $R$ を取る。
方向座標の像を $a_1,\ldots,a_s$，その生成するイデアルを $J_R$ と書く。
$D(a_i)$ 上では方向解消が同型である。
従って $A_{R[a_i^{-1}]}$ とその積代数は，方向束の自明化による
有限なaffine被覆上で，上記の強Noether定理の基底変更として表される。
各affine被覆を有限個の主開集合で細分し，その係数環の有限積へ忠実平坦に変更すれば，
Noether性は右・左イデアルの上昇列を用いて降下する。
従って $A_{R[a_i^{-1}]}$ と
$A_{R[a_i^{-1}]}^{\mathrm{op}}\otimes_{R[a_i^{-1}]}A_{R[a_i^{-1}]}$
はともに左右Noetherである。

\emph{頂点を含む基底全体。}
$A_R/J_RA_R$ は $R/J_R$ 上の重み付き多項式環である。
各内部次数は有限局所自由であり，積代数の各総次数も有限局所自由である。
従って命題\ref{prop:noether-completion-punctured-gluing}を
$A_R$ と $A_R^{\mathrm{op}}\otimes_R A_R$ にそれぞれ適用できる。
後者の $J_R$ による商も有限型可換代数である。
これで両者の左右Noether性を得る。
$B_R$ は $A_R$ の商であり，
$B_R^{\mathrm{op}}\otimes_R B_R$ は $A_R^{\mathrm{op}}\otimes_R A_R$ の商なので，
同じ結論が成立する。
$R$ は任意だったから，これは主張した強Noether性と積の性質である。
\end{proof}

\begin{remark}\label{rem:noether-descent-scope}
本定理は方向解消を忠実平坦とみなしていない。
proper性は有限開集合の選択に，忠実平坦性は
\eqref{eq:completion-punctured-flat-cover}による別の係数変更に用いた。
また，特殊ファイバーの中心全体が基底変更と可換することも仮定していない。
中心元の既存の有限構成と，退化曲線を含む最高次代数のNoether性で十分である。
追加の主開多項式 $g$ は存在を得たものであり，最小の明示公式は求めていない。
相対tailsやsmooth proper性は，本定理を入力として別に確認する。
\end{remark}

% !TeX root = ../nc-del-pezzo.tex
% 2026-11-20 JST / v0.10
\section{全方向族の相対双対とsmooth properな有限開近傍}
\label{sec:simultaneous-relative-geometry}

本節では定理\ref{thm:simultaneous-strong-product-noether}を使い，
同じ全方向中心商の有限基底上の圏論的性質を証明する。
方向を分離する座標は頂点で選べないため，
相対CY同型は方向解消からの平坦降下ではなく，
次数別に有限なHom複体と，全ファイバーのCY同型から構成する。
最後に第\ref{sec:finite-smoothness}節の有限部分対角の議論を適用する。

\subsection{次数付き加群のファイバー別projective性}

\begin{lemma}\label{lem:graded-fiberwise-projective}
$S$ を可換Noether環，$T=\bigoplus_{j\ge0}T_j$ を
$T_0=S$，各 $T_j$ が有限射影 $S$-加群である左右Noetherな次数付き代数とする。
$M$ を有限生成次数付き右 $T$-加群とし，$S$ 上平坦と仮定する。
全ての幾何学的ファイバーについて
$\operatorname{pd}_{T_s}M_s\le a$ なら，
$M$ は長さ $a$ 以下の有限生成次数付きprojective解消を持つ。
この解消は任意の係数変更後も完全である。
\end{lemma}
\begin{proof}
Noether性により，有限生成次数付き自由加群による部分解消を
$a+2$ 段より先まで取れる。
$M$ と自由項が $S$ 上平坦なので，各短完全列から全syzygyが $S$ 上平坦であり，
解消の各部分は任意の係数変更に対して完全のままである。
各syzygyの内部次数部分は有限 $S$-加群なので，有限射影でもある。
$a$ 番目のsyzygyを $K$ とする。
その全ファイバーは次数付きprojectiveであるから，$K$ 自身のprojective性を示せばよい。

有限自由表示を含む解消
\[
 \cdots\longrightarrow F_2\longrightarrow F_1
       \longrightarrow F_0\longrightarrow K\longrightarrow0
\]
を取り，$N=F_0$ または $N=K$ とする。
次数零のHomを取った複体の最初の三項
\[
 C^0=\Hom_T(F_0,N)_0\longrightarrow
 C^1=\Hom_T(F_1,N)_0\longrightarrow
 C^2=\Hom_T(F_2,N)_0
\]
は有限射影 $S$-加群である。
各ファイバーでは $\Ext^1_{T_s}(K_s,N_s)_0=0$ だから，
この三項複体は中央で完全である。
ある点での二つの微分の階数を $u,v$ とすると
$u+v=\rank C^1$ である。
対応する二つの非零小行列式を可逆化し，複体の合成が零であることを使って
単位行列の部分複体を消去すると，中央の項全体が消える。
従って $\ker(C^0\to C^1)$ は局所的に有限自由で，任意基底変更に可換する。
ここで単に「ファイバーのExt消滅からHomが可換する」と仮定せず，
次数零の有限三項複体に還元した。

以上より $\Hom_T(K,F_0)_0$ と $\Hom_T(K,K)_0$ は
有限射影 $S$-加群であり，任意基底変更と可換する。
写像
\[
 \Hom_T(K,F_0)_0\longrightarrow\Hom_T(K,K)_0
\]
は全ファイバーで全射である。
その余核は有限 $S$-加群だから零であり，恒等射を持ち上げると
$F_0\to K$ は次数付きで分裂する。
従って $K$ は有限生成次数付きprojectiveである。
元の解消をここで切ればよく，既に確認した $S$-平坦性が任意基底変更後の完全性を与える。
\end{proof}

\subsection{同じambientの無捩れな相対CY同型}

以下，$S$ は全方向中心族に対して選んだ，被約な有限型普遍係数環とする。
必要な縮小は全可換係数の零切断を保って行う。
記号を
\[
 N=n+2,\qquad
 (l,e)=\begin{cases}(n+5,6)&d=1,\\(n+3,4)&d=2,\end{cases}
 \qquad r=l-e=n-1
\]
と固定する。

\begin{theorem}\label{thm:simultaneous-relative-cy}
零切断全体を含む有限主開集合へさらに制限すると，
同時ambient $A$ は $S$ 上homologically smoothであり，
次数付きbimoduleの同型
\begin{equation}\label{eq:simultaneous-relative-cy}
 A^!_S=\underline{R\Hom}_{A^e_S}(A,A^e_S)
          \simeq A(l)[-N]
\end{equation}
を持つ。
この同型とhomological smoothnessは任意の可換係数変更で保たれる。
従って非被約な係数変更後にもNakayama捻りは現れない。
\end{theorem}
\begin{proof}
\emph{ファイバーの入力。}
零方向では $A$ は $N$ 変数の重み付き可換多項式環なので，
Koszul対角解消から $N$-CY，内部shift $l$ である。
非零方向では中央な一次変数 $n-1$ 個と活性三変数を分ける。
次数2の最高次代数には補題\ref{lem:all-binary-quartic-as-open}，
次数1には補題\ref{lem:degree-one-all-top-as-open}を適用する。
どちらも楕円判別式の非零性を仮定しない。
全lower項を含むpotential複体の活性filtrationの随伴次数は
最高次CY代数のpotential複体である。
その完全性と巡回Hessianの自己双対性から，
中央多項式環上の相対3-CY同型を得る。
これは\cite[Proposition~7.3.1(ii)]{EG2010}のfiltered論法であり，
補題\ref{lem:elliptic-universal-lower-pbw}および
次数1の対応するPBW証明と同じく，中央係数へ適用できる。
次に中央差 $s_i\otimes1-1\otimes s_i$ のKoszul解消を組み合わせると，
長さは $3+n-1=N$，内部shiftは $e+n-1=l$ になる。
これで全幾何学的ファイバーの無捩れなCY同型を得た。

\emph{有限な相対対角解消。}
定理\ref{thm:simultaneous-strong-product-noether}により $A^e_S$ はNoetherである。
総内部次数に関する各部分は有限射影 $S$-加群であり，
対角 $A$ は有限生成かつ $S$-平坦である。
補題\ref{lem:graded-fiberwise-projective}を $T=A^e_S$，$M=A$ に適用すると，
長さ $N$ の有限生成次数付きprojective対角解消 $P^\bullet$ を得る。
従ってambientは相対的にhomologically smoothである。

この解消を双対化した複体の各内部次数は，有界な有限射影 $S$-複体である。
全ファイバーのコホモロジーが解消次数 $N$ だけに集中するため，
有限行列の消去から
\[
 \underline{R\Hom}_{A^e_S}(A,A^e_S)=H[-N],
 \qquad H=\underline{\Ext}^{N}_{A^e_S}(A,A^e_S)
\]
となり，各 $H_j$ は有限射影で任意基底変更に可換する。
そのファイバーは $A_s(l)_j$ である。
特に $H_j=0$ ($j<-l$) で，$L=H_{-l}$ は階数1の射影 $S$-加群である。
左作用による自然な写像
\[
 A(l)\otimes_S L\longrightarrow H
\]
は全ファイバーで同型なので，内部次数ごとの有限射影性から同型である。

\emph{左右の作用と基底線束。}
任意の $a\in A_j$ に対する，$L$ 上の左作用と右作用の差は
$H_{j-l}$ を値に取る $S$-線型写像である。
各複素ファイバーのCY同型では最低次数の生成元は中央なので，この差は全複素点で零である。
$S$ は被約な有限型 $\C$-代数，$H_{j-l}$ は射影なので，差そのものが零となる。
従って上の同型はbimoduleの同型である。
被約性を使ったのは普遍係数環上でこの恒等式を確認する段階だけである。
確認後の恒等式は任意の非被約基底にも移る。

零方向のイデアルを $J$ とする。
Koszul体積形式は $L/JL\simeq S/J$ の自明化を与える。
その生成元を $L$ の切断へ持ち上げ，生成しない点の閉集合を除けば，
零切断全体を含む有限主開で $L\simeq S$ となる。
同変性が必要なら，古典体積に対応する1次元characterで $L$ をtwistし，
Reynolds作用素によりこの持上げを同変に選ぶ。
その零点集合は不変なので，前節と同様に $g\equiv1\pmod J$ の
不変な主開を選べる。これが式\eqref{eq:simultaneous-relative-cy}である。
有限projective対角解消，その双対，同型は任意の係数変更で保たれる。
\end{proof}

\subsection{有限compact生成窓と相対tails}

\begin{lemma}\label{lem:simultaneous-finite-window}
全零切断を含む有限開近傍において，局所的にambientのaugmentationは
有限生成次数付き自由解消を持ち，その両端は $A$ と $A(-l)$ 一つずつ，
他の項の内部shiftは $1,\ldots,l-1$ に入る。
この解消は任意の係数変更で完全である。
従って
\[
 E=\bigoplus_{j=0}^{l-1}\mathcal O(j)
\]
は $D(\QGr B)$ のcompact生成元となる。
\end{lemma}
\begin{proof}
全体ファイバーのaugmentationは大域次元 $N$ 以下の解消を持つ。
補題\ref{lem:graded-fiberwise-projective}を $T=A$，$M=S$ に適用し，
長さ $N$ の有限projective解消を取る。
有限生成次数付きprojective右 $A$-加群 $P$ は，
$P/PA_+$ という有限生成次数付きprojective $S$-加群から誘導される。
実際，次数ごとに $S$-線型な持上げを取り $(P/PA_+)\otimes_S A\to P$ を作ると，
次数付きNakayamaの帰納により全射である。
$P$ のprojective性で分裂した核は $A_+$ を法として零であり，
再び次数付きNakayamaにより零となるので，同型である。
従って係数空間の有限な局所自明化上で，解消の各項を次数付き自由にできる。

零方向の複素点を一つ固定する。
そこでのambientは重み付き多項式環であり，augmentationの最小解消はKoszul解消である。
既に選んだ有限自由解消をこの最小解消へ縮める操作は，
同じ内部shiftを持つ二項間の非零なスカラー成分を有限回消去することからなる。
これらの成分は係数環 $A_0=S$ の元なので，
有限個の非零小行列式を可逆化すれば，同じ消去を点の有限開近傍で実行できる。
得られた解消の自由項は特殊ファイバーのKoszul解消と同じshiftを持つ。
その両端は $A$ と $A(-l)$ で，真部分集合の重みの和は $1$ から $l-1$ の間にある。
元の解消の $S$-平坦syzygyと分裂消去から，任意基底変更後の完全性も保たれる。
零切断は準コンパクトなので有限個のこのような開集合で覆える。
必要ならその不変開集合内に，零切断を含む一つの有限主開を選ぶ。

この解消を $B=A/(\Omega)$ とtensorしたときのTorは
$S$ と $S(-e)$ だけであり，ともにirrelevant torsionである。
従ってtailsでは両端に $\mathcal O$ と $\mathcal O(-l)$ を一つずつ持つ有限完全複体となる。
そのシフトにより，長さ $l$ の連続窓から両方向に全twistを生成できる。
前定理の相対CY同型と中央元のchange of ringsを用いる
定理\ref{thm:finite-relative-local-cohomology}の証明により，
tails局所化はsmashingで全twistがcompact生成する。
従って表示した $E$ がcompact生成元である。
これらの生成操作は係数基底で局所的に行え，有限な主開被覆で検査できる。
\end{proof}

\begin{theorem}\label{thm:simultaneous-relative-tails}
上で選んだ全零切断の有限主開近傍から，任意の可換Noether環 $R$ へ係数変更する。
その上の同じ中心商 $B_R$ について，次が成り立つ。
\begin{enumerate}
 \item $(\QGr B_R,\pi B_R)$ は $R$-flatなAbel圏とflatな構造対象を与え，
       許された基底変更と両立する。
 \item 左右の相対局所コホモロジーは
       \[
        R\Gamma_{(B_R)_+}(B_R)\simeq B_R^\vee(r)[-n-1]
       \]
       であり，torsionの導来関手は有限cohomological dimensionを持つsmashing関手である。
       切断環は飽和し，有向ACM，相対 $\chi$，全シフトの相対豊富性が成立する。
 \item 全ての整数 $t$ について
       \[
        R\Hom(\mathcal O,\mathcal O(t))\simeq
        \begin{cases}
         (B_R)_t,&t\ge0,\\
         \Hom_R((B_R)_{-t-r},R)[-n],&t\le-r,\\
         0,&-r<t<0.
        \end{cases}
       \]
       $E=\bigoplus_{j=0}^{l-1}\mathcal O(j)$ はcompact生成元であり，
       $R\End(E)$ は $R$ 上perfectである。
 \item compact圏は相対properで，その相対Serre関手は
       \[
          S_R\simeq(-(n-1))[n]
       \]
       である。以上の構成と同型は任意の許される基底変更に両立する。
\end{enumerate}
\end{theorem}
\begin{proof}
Noether性と積のNoether性は定理\ref{thm:simultaneous-strong-product-noether}，
次数別有限射影性と普遍正則性は同時中心商の構成による。
Abel平坦性には，命題\ref{prop:finite-pbw-abel-family}の
作用付き加群・Gabriel商・injectiveのcoflat性の証明を適用する。
その証明で必要なのは係数平坦性と有限生成Noether性であり，分離型の関係式ではない。

定理\ref{thm:simultaneous-relative-cy}により相対有限対角解消と無捩れな双対性がある。
各 $A/A_{\ge q}$，$B/B_{\ge q}$ は係数環上有限射影なので，
補題\ref{lem:relative-proper-module-duality}の証明が適用できる。
中央正則元の二項解消とchange of ringsから
\[
 R\Hom_{B_R}(M,B_R)=D_RM(r)[-n-1]
\]
を係数上perfectな $M$ に対して得る。
$M=B_R/(B_R)_{\ge q}$ としてcolimitを取れば局所コホモロジーの表示を得る。
ambientと中心商のtorsion比較における $\operatorname{Tor}_1$ は
長さ $e$ の内部次数帯にだけ支持を持ち，十分離れた切断間で遷移が零になる。
従って定理\ref{thm:finite-relative-local-cohomology}のsmashingの証明も適用できる。

局所化三角から(3)のHom公式が従う。
有限型自由解消にこの局所コホモロジーを適用する
命題\ref{prop:finite-full-shift-ample}の証明から，
coherent対象のSerre消滅，$\chi$ と全シフトの豊富性を得る。
有限窓は補題\ref{lem:simultaneous-finite-window}による。
左右共通の局所化bimoduleとその相対双対を使う
定理\ref{thm:finite-tails-relative-serre}の証明は，
上記のCY同型・局所コホモロジー・compact生成だけを用いるので，
ここでも無捩れな相対Serre公式を与える。
有限窓のHom公式により $R\End(E)$ は有限個の係数perfect複体から作られ，proper性も得る。
これらは固定した普遍係数上の構成の基底変更であり，非正則・非被約な $R$ を許す。
\end{proof}

\subsection{滑らかな全可換模型を含むnc del Pezzo族}

\begin{theorem}\label{thm:simultaneous-smooth-ncdp-open}
両次数 $d=1,2$ と全 $n\ge2$ について，滑らかな可換正規形の係数空間を
$\mathcal F_d^{\mathrm{sm}}$ とする。
全方向有限中心族には，$\mathcal F_d^{\mathrm{sm}}$ の零切断全体を含む
有限型の不変開集合 $V$ があり，その上で
\[
 (\QGr B_V,\mathcal O)\in\mathfrak M^{\mathcal O,\mathrm{hyp}}_{n,d}(V)
\]
となる。
特にcompact圏は任意の許されるNoether基底変更後もsmooth properで，
反標準周期 $r=n-1$ の列は相対的に豊富である。
全幾何学的ファイバーのcoherent圏はホモロジー次元 $n$ を持ち，
その有界導来圏はcompact圏と一致する。
\end{theorem}
\begin{proof}
前定理により，$C=R\End(E)$ は係数上perfectで，その双対は可逆Serre bimoduleである。
定理\ref{thm:simultaneous-strong-product-noether}により
$B^{\mathrm{op}}\otimes_S B$ は任意Noether基底変更後もNoetherである。
従って補題\ref{lem:finite-family-product-diagonal}の証明に必要な入力が全て揃う。
二つのirrelevant torsionを順に割った積圏は外積twistでcompact生成され，
有限窓のMorita代数は $C^{\mathrm{op}}\otimes_S^LC$ である。
二重次数付き対角 $\Delta_{a,b}=B_{a+b}$ は有限生成かつ $S$-平坦で，
そのMorita像は標準の対角bimoduleとなる。

任意の滑らかな可換零方向点 $s$ を固定する。
$\Delta$ の有限生成自由解消を $q>2n$ で切り，得られるcompact対象を $P_q$ とする。
特殊ファイバーは $X\times X$ 上の対角であり，
そのホモロジー次元は $2n$，また $HH^1(X)=0$ である。
従って対角への写像 $P_{q,s}\to\mathcal O_{\Delta_X}$ は右逆を持ち，
この右逆の持上げに関わる $\Ext^1$ は零である。
相対Serre双対から得られる
補題\ref{lem:serre-pseudoperfect-to-perfect}によって，
対角から $P_q$ へのHom複体は係数上perfectで基底変更に可換する。
定理\ref{thm:finite-family-smooth-open}の有限行列の持上げと
可逆性の開性を適用すると，$s$ の有限開近傍で右逆を持ち上げられる。
対角はperfect対象の直和因子となるのでsmoothである。
これは対角が既にperfectだと仮定してからその開性を用いる議論ではない。

命題\ref{prop:finite-periodic-generation}と
系\ref{cor:finite-periodic-ampleness}を用いて，
同じ近傍を正周期 $r$ の豊富性が成立する開集合と交わらせる。
全可換係数の滑らかな零切断は準コンパクトなので，有限個の近傍の和で覆える。
必要ならその群による移動の和を取り，不変な有限型開集合 $V$ を得る。
構成した対角の有限cone・直和因子による表示と豊富性は任意基底変更に保たれる。

coherent対象が幾何学的ファイバーでcompactとなること，
Serre公式からホモロジー次元が $n$ となることは，
系\ref{cor:finite-family-ncdp-open}の証明をそのまま適用する。
非正則な係数基底上の全有界coherent導来圏までcompactと同一視する主張ではない。
その他の飽和性・ACM・周期・Hilbert関数と構造対象の平坦性は前定理と中心商の構成による。
\end{proof}

\subsection{低次数Hochschildコホモロジーの開近傍での剛性}

\begin{proposition}\label{prop:simultaneous-hochschild-rigid-open}
定理\ref{thm:simultaneous-smooth-ncdp-open}の開集合 $V$ は，
滑らかな可換模型の零切断全体を保ってさらに縮小し，次の性質を満たすように取れる。
有限窓のdg代数を $C$ とし，相対Hochschild複体を
\[
 \mathcal K=R\!\mathcal Hom_{C^{\mathrm{op}}\otimes^L_{\mathcal O_V}C}(C,C)
\]
と置くと，$\mathcal K$ は $\mathcal O_V$ 上perfectであり，
任意の許された基底変更に両立する。
任意のnoetherian affine基底変更 $\Spec R\to V$ に対して
\begin{equation}\label{eq:simultaneous-low-hochschild}
 HH_R^i(C_R)=0\quad(i<0),\qquad
 R\xrightarrow{\sim}HH_R^0(C_R),\qquad HH_R^1(C_R)=0
\end{equation}
が成り立つ。中央の同型は恒等bimodule写像が与える標準写像である。
非affine基底では，同じ主張を相対Hochschild複体のコホモロジー層について読む。
また，$V$ は座標変換に関して不変に選べる。
\end{proposition}
\begin{proof}
$C$ のsmoothnessは対角 $C$ が $C^{\mathrm{op}}\otimes^LC$ 上perfectであることを意味する。
この有限coneと直和因子による表示に $R\!\mathcal Hom(-,C)$ を施し，
$C$ の係数上のperfect性を使うと，$\mathcal K$ も係数上perfectである。
同じ表示から，$\mathcal K$ とその単位写像
$\eta:\mathcal O_V\to\mathcal K$ は任意の導来基底変更に両立する。

単位写像のconeを $\mathcal F=\operatorname{Cone}(\eta)$ と置く。
滑らかな可換零方向点 $s$ のファイバーでは，Morita不変性により
$\mathcal K_s\simeq R\!\Hom_{X_s\times X_s}
(\mathcal O_{\Delta_{X_s}},\mathcal O_{\Delta_{X_s}})$ である。
従って負次数コホモロジーは零，零次数は $H^0(X_s,\mathcal O_{X_s})=\C(s)$ であり，
一次はHKR，Kodaira消滅および \cite[補題8.7]{PoissonDP} により
$H^0(X_s,T_{X_s})\oplus H^1(X_s,\mathcal O_{X_s})=0$ となる。
ゆえに $H^i(\mathcal F_s)=0$ ($i<2$) である。

perfect複体のファイバーのコホモロジー次元の上半連続性
\cite[Tag 0BDI]{StacksProject}から
\[
 W=\{v\in V\mid H^i(\mathcal F\otimes^L\kappa(v))=0
                 \text{ for all }i<2\}
\]
は開集合である。局所的に関係する次数は有限個だけであり，
$W$ は滑らかな可換零切断全体を含む。
さらに，任意の $v\in W$ の近傍で有限自由模型の微分から
可縮な二項直和因子を順に消去すると，$\mathcal F$ は
次数 $2$ 以上だけに項を持つ有界有限自由複体で表せる
\cite[Tag 0BCD]{StacksProject}。
従ってこの次数の下界は任意の非平坦な基底変更後も保たれる。
$R\to\mathcal K_R\to\mathcal F_R$ のコホモロジー完全列から
式\eqref{eq:simultaneous-low-hochschild}を得る。
$\mathcal K$ と $\eta$ はMorita同値で保たれるので，$W$ の定義は座標変換に不変である。
有限型基底の開部分 $W$ を新たな $V$ と取ればよい。
\end{proof}

\begin{corollary}\label{cor:noncommutative-diagonal-artin-rigidity}
前命題の開集合 $V$ の任意の複素幾何学的点 $s$ を取る。
$(R,\mathfrak m)$ を剰余体 $\C$ の局所Artin環とし，$A_R$ を $C_s$ の
smooth properな $R$-線形dg持上げとする。
$P,Q\in\Perf(A_R^{\mathrm e})$ の還元がともに標準対角 $C_s$ に同型であれば，
還元上の任意の同型 $P_0\to Q_0$ は $P\simeq Q$ へ持ち上がる。
特に，perfectな可逆bimoduleで表されたMorita自己同値が特殊ファイバーで恒等なら，
それは持上げ上でも恒等Morita自己同値に同型である。
\end{corollary}
\begin{proof}
前命題は $HH^1(C_s)=0$ を与えるので，
補題\ref{lem:diagonal-bimodule-rigidity}を $A_R$ に適用する。
最後の主張では $Q=A_R$ と取る。
\end{proof}

\begin{remark}
ここで得たのは，構成した有限族の実際の非可換点でも成立する
低次数Hochschildコホモロジーの消滅と，存在するperfect bimodule間の剛性である。
一般のAbel変形に対するHochschild比較やdg模型の存在を，この命題だけから結論しない。
また，自己同値群の有限型表現可能性や全モジュライのatlasを与えるものでもない。
\end{remark}


\begin{remark}\label{rem:simultaneous-geometric-scope}
本節により，全独立Poisson方向を含む同じ有限中心商は，
滑らかな可換零切断の有限型開近傍で実際の非可換del Pezzo族となる。
定理\ref{thm:simultaneous-central-all-artin}と合わせれば，
各可換点の横断slice上で全framed Abel Artin変形を実現する有限型nc族を得る。
全 $\mathcal O$ 保存同値の有限表示による認識は第\ref{sec:finite-full-equivalences}節で行う。
この族による全モジュライの被覆，有限型atlas，
および可換locusに接しない成分の不存在は別の主張である。
また，ここで選んだ追加の開集合は有限行列による存在証明であり，
最小の分母やsmoothnessの判別式は明示していない。
\end{remark}

% !TeX root = ../nc-del-pezzo.tex
% 2026-11-20 JST / v0.8
\section{同時中心商による全Artin変形の有限代数的表示}
\label{sec:simultaneous-formal-comparison}

本節では，直前の二つの有限型中心商族そのものを全Abel Artin変形と比較する。
比較のために，有限型基底全体での強Noether性や相対tailsのsmoothnessを
先に仮定する必要はない。Artin基底上の平坦な線型変形の理論と，
既に求めた全形式環を使う。

滑らかな可換次数 $d\in\{1,2\}$ のdel Pezzo多様体 $X$ を固定する。
その標準方程式を表す係数点を $f$ とし，係数空間の中で
座標変換群の軌道の接空間に相補的なaffine線型切断を取る。
$f$ の開近傍へ制限したものを $T$ と書く。
これは滑らかな次元 $b_{n,d}$ の有限型スキームであり，
その可換族のKodaira--Spencer写像は同型である。
実際，標準方程式の係数表示から通常のdel Pezzoスタックへの射は滑らかで，
$H^0(X,T_X)=0$ であるから，その微分の核が座標軌道の接空間に等しい
\cite[定理7.3，定理8.8]{PoissonDP}
（定理\ref{thm:abel-formal-ring}の証明も参照）。
有限な自己同型群を割らず，特殊ファイバーの同一視を固定する。

方向スキームを
\[
 C_1=E^*\quad(\dim E=n),\qquad
 C_2=\operatorname{Cone}(\Gr(2,E^*))\quad(\dim E=n+1)
\]
と書く。定理\ref{thm:simultaneous-degree-one-central-family}または
定理\ref{thm:simultaneous-degree-two-central-family}の族を
$T\times C_d$ に引き戻し，得られる零切断の有限型開近傍を $V$ とする。
その同じambientと中心商を $A_V$，$B_V=A_V/(\Omega)$ と書く。
追加の開条件は零方向で可逆なので，$(f,0)$ における完備局所環は
\begin{equation}\label{eq:simultaneous-family-formal-ring}
 C_X=\widehat{\mathcal O}_{V,(f,0)}\simeq
 \begin{cases}
  \mathbb C[[t_1,\ldots,t_{b_{n,1}},D_1,\ldots,D_n]],&d=1,\\[2pt]
  \mathbb C[[t_1,\ldots,t_{b_{n,2}},p_{ij}]]/I_{\mathrm{Pl}},&d=2.
 \end{cases}
\end{equation}
ここで $I_{\mathrm{Pl}}$ は全Pl\"ucker二次式のイデアルである。

\begin{lemma}[Artin中心商の平坦Abel変形]
\label{lem:simultaneous-artin-abel-family}
任意の局所Artin $\mathbb C$-代数 $R$ と，$(f,0)$ を特殊点とする射
$a:\Spec R\to V$ に対し，
\[
 (\operatorname{QGr}(B_a),\pi B_a),\qquad B_a=a^*B_V,
\]
は $X$ のframedな構造対象付き平坦Abel変形を与える。
この構成は任意のArtin基底変更と両立する。
\end{lemma}
\begin{proof}
中心商の普遍的な次数別有限局所自由性から，各 $(B_a)_q$ は有限自由であり，
その還元は $X$ の重み付き切断環 $B_0$ の次数 $q$ 部分である。
従って
\[
 b_{ij}=(B_a)_{j-i}\quad(i\le j),\qquad b_{ij}=0\quad(i>j)
\]
は，$X$ の切断 $\mathbb Z$-代数の平坦線型変形である。
De Deken--Lowen \cite[Proposition~4.4, p.~495]{DDL2011}により
$\operatorname{Qmod}(b)$ は $\operatorname{Qcoh}(X)$ の平坦Abel変形となる。
同文献Proposition~4.6は有限生成性とNoether性をArtin変形へ持ち上げる。
従ってtails topologyによる $\operatorname{Qmod}(b)$ は
通常のtorsion商 $\operatorname{QGr}(B_a)$ に等しい
\cite[\S3.1, \S3.3]{DDL2011}。

構造対象 $\pi B_a$ は $R$-平坦である。
実際，$B_a$ の次数別平坦性はgraded加群としての $R$-平坦性を意味し，
その係数tensor関手に完全な商関手 $\pi$ を合成しても完全性は保たれる。
還元後の構造対象は $\mathcal O_X$ である。
基底変更との両立は上の線型変形の自然性，または
補題\ref{lem:tails-arbitrary-base-change}から従う。
\end{proof}

\begin{theorem}[同じ有限型中心商の全Artin普遍性]
\label{thm:simultaneous-central-all-artin}
上で選んだ可換sliceと中心元に対し，自然変換
\[
 Q^{\mathrm{fin}}_X:
 \operatorname{Hom}_{\mathrm{cont},\mathrm{loc}}(C_X,-)
 \longrightarrow\mathrm{Def}^{\mathcal O}_X,
 \qquad a\longmapsto[(\operatorname{QGr}B_a,\pi B_a)]
\]
は全ての局所Artin環上で全単射である。
従って，任意のframedな構造対象付き平坦Abel Artin変形は，
一つの固定した有限表示中心商族 $B_V$ のArtin基底変更として実現される。
自然同型類を射とするframed群oidについても，左辺を離散群oidと見て同値となる。
各変形は定理\ref{thm:all-artin-are-nc-dp}のnc del Pezzo条件を満たす。
\end{theorem}
\begin{proof}
前補題により，表示した自然変換が定まる。
その接写像を，命題\ref{prop:quantization-tangent-isomorphism}の完全列
\[
 0\longrightarrow H^1(X,T_X)
  \longrightarrow T\mathrm{Def}^{\mathcal O}_X
  \xrightarrow{\beta}H^0(X,\wedge^2T_X)\longrightarrow0
\]
によって調べる。
零方向では $\Omega=F$ であり，ambientも可換多項式環である。
従って可換方向での接写像は $T$ のKodaira--Spencer同型である。

Poisson方向では，ambientの定義関係の一次交換子が
方向パラメータ $D_i$ または $p_{ij}$ に対応するJacobian括弧を与える。
この対応は，標準方程式の規格化に由来する非零定数を座標へ吸収して，
$H^0(X,\wedge^2T_X)$ の恒等写像とできる。
中心元は第一平方零近傍で $\operatorname{Sym}F$ に等しく，
その中心商の一次交換子も同じ括弧である。
これがtails上の $\beta$ と一致する点は，次の局所化で確認できる。
一次変形では二重交換子が零であり，非零斉次被覆元の冪は左右Ore条件を満たす。
次数2の $D_+(x_i)$，次数1のそれらと $D_+(yz)$ 上で局所化し，
次数零を取れば，商則による同じPoisson括弧が得られる。
これらの局所化は次数1の可逆元を持つため強次数付きであり，
tailsの対応する局所化と一致する。
これは命題\ref{prop:quantization-tangent-isomorphism}のOre被覆の議論を
同じ一次交換子に適用したものであり，有限パラメータ上のOre性は要求していない。
従って接写像は完全列の核と商で同型を誘導し，全体として同型である。

定理\ref{thm:abel-formal-ring}により，標的の全変形関手は $R_X$ で
prorepresentされる。
その環と\eqref{eq:simultaneous-family-formal-ring}の $C_X$ は同じ表示を持つので，
一つの同型 $R_X\simeq C_X$ を選べる。
Yoneda対応によって $Q^{\mathrm{fin}}_X$ をこの環の連続局所自己写像と見れば，
直前の接写像の同型性は余接写像の同型性を意味する。
補題\ref{lem:formal-endomorphism-inverse}からこの自己写像は同型である。
これで全Artin環上の全単射を得る。
次数2ではPl\"ucker関係を持つ\emph{全形式環}を比較しており，
二次障害を接空間だけの議論で消してはいない。

最後に，補題\ref{lem:abel-infinitesimal-aut-rigidity}と
定理\ref{thm:quantization-framed-groupoid}によりframed自己同型群は自明である。
従って同型類の自然な全単射は群oidの同値へ一意に持ち上がる。
nc del Pezzo条件は，全Abel Artin変形について既に証明した
定理\ref{thm:all-artin-are-nc-dp}を適用すれば従う。
\end{proof}

\begin{remark}[有限型族の形式的完全性とatlas]
本定理の有限型族は，その完備化で全変形を表すという意味で
一方向族の接空間の不足を解消する。
中心元の選択や可換sliceの変更は，形式パラメータの自己同型を変え得る。
グラフ量子化の既存の対応と厳密に同じ座標を与えることは主張しない。

また，本定理は有限型基底 $V$ 全体のtailsがsmooth properな族をなすことや，
全O保存同値によるモジュライのatlasを与えることとは異なる。
全方向族の強Noether性・積Noether性，相対tailsと相対Serre式，
有限開近傍へのsmoothness，および大域的な射の認識を引き続き確認する必要がある。
ここで証明したのは，同じ有限表示代数族のArtin基底変更に関する完全性である。
\end{remark}


\section{大域的なモジュライへ進むための問題}
\label{sec:questions}
主な残問は、一般の非可換点での有限表示と変形の局所完全性、
許容な表示の全 $\OO$ 保存同値による商、ならびに全モジュライの代数性である。
各可換点の形式環が整であることから、可換点における形式分枝は一つと分かる。
可換locusに接しない非可換成分が存在しないことまでは証明していない。
第\ref{sec:finite-auxiliary-quantization}節の分離型族では，
第\ref{sec:finite-tails}節によりtailsの飽和性，proper性，相対Serre公式まで得られた。
滑らかな可換点の開近傍では反標準列も豊富になる。
第\ref{sec:elliptic-relative}節と第\ref{sec:elliptic-degree-one-finite}節により，両次数の一般楕円型族でも
係数基底上の強Noether性，相対tailsと積のNoether性が得られる。
両方の族について，第\ref{sec:finite-smoothness}節は滑らかな可換点を含む
有限開近傍上のsmoothnessと非可換del Pezzo条件を証明した。
次の課題は，この開集合を具体的な有限行列で記述し，
より大きい表示族の局所完全性と全 $\OO$ 保存同値による商を調べることである。
特に命題\ref{prop:one-parameter-atlas-obstruction}の接空間の障害により，
共通因子 $h$ による表示を越えて独立なPoisson方向を同時に含める必要がある。
第\ref{sec:simultaneous-poisson}--\ref{sec:simultaneous-formal-comparison}節では，
両次数で独立Poisson方向を含む有限型中心商族を構成し，
その同じ族の全Artin普遍性まで証明した。
全方向族の強・積Noether性，相対tailsとSerre公式，
滑らかな可換零切断の有限開近傍でのsmooth proper性も証明した。
活性最高次式の特異・可約・非被約な曲線を扱い，
完成と頂点外を合わせた忠実平坦被覆を使うことが必要だった。
次の核心は，この有限型nc族について一般点の全 $\mathcal O$ 保存同値を認識し，
商と被覆の範囲を確定することである。
追加の中心分母・Noether性の開集合・smoothnessの有限行列を明示し，
どの境界を残せるかも調べる。
半古典極限で稠密なPoisson対を実現することと，全非可換対象を被覆することは異なる。
有限パラメータの非可換対象から古典Poisson多様体を自然に復元する射も構成していない。

三次三様体の残余圏による幾何化を次数1・2へ移すには、
例外部分と残余圏の相対構成、同じ古典多様体との比較を新たに証明する必要がある。
接空間や形式環の一致だけでKuznetsov成分の同値を結論しない。
公表に向けては、本稿の形式局所定理と条件を明示した復元定理を一つの論理単位とし、
未解決の大域分類をその定理の仮定に紛れ込ませない。

\begingroup
\small
\bibliographystyle{alpha}
\bibliography{../references}
\endgroup
\end{document}
